% =====================================================================
%  Lezione zero 1 — Lavorare in squadra: la cultura del feedback
%  Fonte: trascrizione della prima "lezione zero" (GD0), prima giornata di
%  un workshop di due giorni tenuto prima dell'inizio del corso di Game
%  Design, dedicato alla figura professionale e al lavoro in team.
%  Pacchetti richiesti nel preambolo del libro:
%    \usepackage[italian]{babel}, \usepackage[T1]{fontenc},
%    \usepackage[utf8]{inputenc} (solo pdfLaTeX), \usepackage{booktabs}
%
%  Collocazione consigliata nel libro: le lezioni zero precedono il corso,
%  quindi vanno in una parte separata PRIMA di GD1, con numerazione propria
%  (0.1, 0.2, ...) così che i capitoli del corso conservino i numeri 1, 2, 3.
%  Nel file principale del libro, ad esempio:
%
%    \part{Lezioni zero: il game designer nell'industria}
%    \renewcommand{\thechapter}{0.\arabic{chapter}}
%    \setcounter{chapter}{0}
%    % =====================================================================
%  Lezione zero 1 — Lavorare in squadra: la cultura del feedback
%  Fonte: trascrizione della prima "lezione zero" (GD0), prima giornata di
%  un workshop di due giorni tenuto prima dell'inizio del corso di Game
%  Design, dedicato alla figura professionale e al lavoro in team.
%  Pacchetti richiesti nel preambolo del libro:
%    \usepackage[italian]{babel}, \usepackage[T1]{fontenc},
%    \usepackage[utf8]{inputenc} (solo pdfLaTeX), \usepackage{booktabs}
%
%  Collocazione consigliata nel libro: le lezioni zero precedono il corso,
%  quindi vanno in una parte separata PRIMA di GD1, con numerazione propria
%  (0.1, 0.2, ...) così che i capitoli del corso conservino i numeri 1, 2, 3.
%  Nel file principale del libro, ad esempio:
%
%    \part{Lezioni zero: il game designer nell'industria}
%    \renewcommand{\thechapter}{0.\arabic{chapter}}
%    \setcounter{chapter}{0}
%    % =====================================================================
%  Lezione zero 1 — Lavorare in squadra: la cultura del feedback
%  Fonte: trascrizione della prima "lezione zero" (GD0), prima giornata di
%  un workshop di due giorni tenuto prima dell'inizio del corso di Game
%  Design, dedicato alla figura professionale e al lavoro in team.
%  Pacchetti richiesti nel preambolo del libro:
%    \usepackage[italian]{babel}, \usepackage[T1]{fontenc},
%    \usepackage[utf8]{inputenc} (solo pdfLaTeX), \usepackage{booktabs}
%
%  Collocazione consigliata nel libro: le lezioni zero precedono il corso,
%  quindi vanno in una parte separata PRIMA di GD1, con numerazione propria
%  (0.1, 0.2, ...) così che i capitoli del corso conservino i numeri 1, 2, 3.
%  Nel file principale del libro, ad esempio:
%
%    \part{Lezioni zero: il game designer nell'industria}
%    \renewcommand{\thechapter}{0.\arabic{chapter}}
%    \setcounter{chapter}{0}
%    % =====================================================================
%  Lezione zero 1 — Lavorare in squadra: la cultura del feedback
%  Fonte: trascrizione della prima "lezione zero" (GD0), prima giornata di
%  un workshop di due giorni tenuto prima dell'inizio del corso di Game
%  Design, dedicato alla figura professionale e al lavoro in team.
%  Pacchetti richiesti nel preambolo del libro:
%    \usepackage[italian]{babel}, \usepackage[T1]{fontenc},
%    \usepackage[utf8]{inputenc} (solo pdfLaTeX), \usepackage{booktabs}
%
%  Collocazione consigliata nel libro: le lezioni zero precedono il corso,
%  quindi vanno in una parte separata PRIMA di GD1, con numerazione propria
%  (0.1, 0.2, ...) così che i capitoli del corso conservino i numeri 1, 2, 3.
%  Nel file principale del libro, ad esempio:
%
%    \part{Lezioni zero: il game designer nell'industria}
%    \renewcommand{\thechapter}{0.\arabic{chapter}}
%    \setcounter{chapter}{0}
%    \input{Seminari/GD0_1_cultura_del_feedback}
%    % ... altre lezioni zero ...
%    \part{Il corso di Game Design}
%    \renewcommand{\thechapter}{\arabic{chapter}}
%    \setcounter{chapter}{0}
%    \input{GD1/GD1_introduzione_al_game_design}
%    % ...
%
%  (Con hyperref aggiungere anche \renewcommand{\theHchapter}{z.\arabic{chapter}}
%  nella parte zero, per evitare ancore duplicate.)
%  Tutte le etichette di questo capitolo hanno il prefisso "z1:" per non
%  collidere con quelle dei capitoli del corso.
% =====================================================================

\chapter[La cultura del feedback]{Lavorare in squadra: la cultura del feedback}
\label{cap:z1-feedback}

Questo capitolo apre le \emph{lezioni zero}: incontri che precedono l'inizio
vero e proprio del corso di Game Design e che non riguardano una tecnica di
design in senso stretto, ma \textbf{la figura professionale} di chi lavora
nell'industria dei videogiochi. Prima di imparare a progettare un gioco
conviene infatti capire in che tipo di ambiente lo si progetterà: un ambiente
fatto di team, di gerarchie, di scadenze e, soprattutto, di persone che devono
riuscire a lavorare insieme.

La lezione è la prima giornata di un \textbf{workshop di due giorni} (circa
dieci ore in totale, pause pranzo comprese) tenuto da Stefano
Cozzi\footnote{Il cognome è riportato come compare nella trascrizione
automatica, che in più punti storpia i nomi propri; andrebbe verificato.
Nella stessa trascrizione lo studio per cui lavora appare come
``Stormine'', ``Stormline'' e ``Stormind'': si tratta di Stormind Games,
studio con sede a Catania, coerentemente con quanto detto in lezione.}, un
professionista con circa quindici anni di esperienza nel settore, tra
Blizzard Entertainment, Riot Games e, oggi, Stormind Games. Il workshop è
nato da un confronto con i coordinatori dell'accademia\footnote{Nella
trascrizione l'accademia è resa in modi diversi (``Hive'', ``a dive'',
``D.I.V.''); i coordinatori citati sono Enrico e Andrea, insieme al
direttore Diego Ricchiuti. Nel resto del capitolo si usa semplicemente
``l'accademia'', come negli altri capitoli.}: serviva un momento dedicato a
\emph{come ci si comporta in un ambiente professionale}, cosa che riguarda sia
il lavoro futuro sia i progetti di gruppo svolti durante il percorso di
studi, dove ci si deve già abituare a comportarsi ``come se si fosse in
un'azienda''.

La prima giornata ha un unico grande tema, sezionato in più parti:
\textbf{la cultura del feedback}. La seconda giornata (che sarà oggetto della
prossima lezione zero) sarà più pratica: \emph{time management},
\emph{brainstorming} e il ``guidare da dentro'' un team anche senza esserne
il leader (sezione~\ref{sec:z1-domani}).

Il percorso del capitolo segue quello della lezione:
\begin{enumerate}
  \item il formato del workshop e il percorso professionale del relatore
        (sezione~\ref{sec:z1-relatore});
  \item perché le aziende tengono tanto al lavoro di squadra e al
        \emph{fit} culturale (sezione~\ref{sec:z1-squadra});
  \item che cos'è il feedback e da che cosa dipende il modo di darlo
        (sezione~\ref{sec:z1-cose-feedback});
  \item le cinque fasi di vita di un team e il valore del
        \emph{post-mortem} (sezioni~\ref{sec:z1-fasi-team}
        e~\ref{sec:z1-qa-fasi});
  \item il confronto con i colleghi e alcune domande sul lavoro nei grandi
        studi (sezioni~\ref{sec:z1-rivalita} e~\ref{sec:z1-industria});
  \item il circolo virtuoso tra feedback, miglioramento e fiducia, e i suoi
        effetti sulla \emph{retention} e sulla cultura aziendale
        (sezioni~\ref{sec:z1-triangolo} e~\ref{sec:z1-cultura});
  \item il ``gioco del traditore'' e il principio del
        \textbf{default to trust} (sezioni~\ref{sec:z1-among-us}
        e~\ref{sec:z1-default-trust});
  \item le traiettorie di carriera e la valutazione delle performance
        (sezione~\ref{sec:z1-traiettorie});
  \item come si dà, come si riceve e quando si dà un feedback
        (sezioni~\ref{sec:z1-dare}, \ref{sec:z1-ricevere}
        e~\ref{sec:z1-quando});
  \item il feedback come regalo, le domande finali e le anticipazioni
        della seconda giornata (sezioni~\ref{sec:z1-regalo},
        \ref{sec:z1-qa-finale} e~\ref{sec:z1-domani}).
\end{enumerate}

Una premessa di metodo: il relatore ha alternato continuamente teoria,
domande alla classe e racconti della propria esperienza. I racconti non sono
digressioni: sono gli esempi concreti con cui i principi vengono resi
tangibili, e per questo sono riportati con cura.

% ---------------------------------------------------------------------
\section{Il workshop e il relatore}
\label{sec:z1-relatore}

\subsection{Formato e regole di partecipazione}

Il relatore ha cercato di ``riassumere quindici anni di esperienza in dieci
ore''. Ha materiale per circa il doppio delle ore disponibili: le cose
strettamente necessarie vengono trattate per prime, il resto è tenuto ``in
riserva'' come \emph{nice to have}. Le regole del workshop:
\begin{itemize}
  \item le domande si fanno \textbf{alla fine di ogni slide}, così da chiudere
        prima il discorso; non c'è un'unica sessione finale di domande e
        risposte (anche se ci sarà), perché il relatore preferisce fermarsi
        spesso e non fare ``uno spiegone di dieci ore'';
  \item le domande non devono per forza riguardare la slide appena vista: se
        una slide dà lo spunto per altro, se ne può parlare; al massimo sarà
        il relatore a rispondere in breve o a rimandare a una slide
        successiva;
  \item si può chiedere anche di esperienze passate, opinioni e sensazioni;
        se per NDA o per serietà professionale non può rispondere, lo dirà
        esplicitamente (``non lasciate che sia un \emph{blocker}'');
  \item su Zoom si usa l'alzata di mano o la chat, che il relatore tiene
        aperta sul secondo monitor;
  \item non ci sono pause ufficiali: la lezione è registrata e ognuno è
        libero di assentarsi quando serve, cercando però di non disturbare il
        flusso della lezione;
  \item le slide verranno fornite integralmente tramite la segreteria
        dell'accademia (Vanessa), insieme alle registrazioni: le slide senza
        il parlato servono poco, ma sono un buon promemoria.
\end{itemize}

\subsection{Il percorso professionale}

Il relatore si presenta (scherzosamente) come la ``guida escursionistica''
che accompagna la classe nella ``foresta incantata'' della cultura del
feedback, perché per molti è un territorio sconosciuto. Il suo percorso:
\begin{itemize}
  \item in Italia ha lavorato per aziende che fornivano servizi a terzi, tra
        cui il lavoro su \emph{Warhammer Online}, e per un canale televisivo
        sul digitale terrestre dedicato ai videogiochi\footnote{Nella
        trascrizione: ``Digital Game Network Channel''. Il nome esatto del
        canale andrebbe verificato.};
  \item a \textbf{Blizzard Entertainment} (circa sei anni, a Parigi) ha fatto
        il \emph{community manager} per l'Italia (\emph{StarCraft~II},
        \emph{World of Warcraft} da \emph{Mists of Pandaria} in poi,
        \emph{Hearthstone}, \emph{Diablo~III}, le espansioni di
        \emph{StarCraft}, il lancio di Battle.net), per poi diventare il
        manager dei community manager europei;
  \item a \textbf{Riot Games} (circa cinque anni, a Dublino) è stato
        brevemente community manager italiano e poi \emph{publishing
        manager} di un team di circa trenta persone per l'Europa. Il
        \emph{publishing} non è solo community: è l'unione di
        \textbf{community, marketing e PR};
  \item a \textbf{PG Esports} si è occupato di eventi dal vivo e online
        (ad esempio la Milano Games Week e tornei di \emph{Valorant} e
        \emph{League of Legends}), restando in contatto con Riot e Blizzard;
  \item da circa un anno è tornato nello sviluppo, a \textbf{Stormind
        Games}, dove oggi si occupa di risorse umane e organizzazione del
        lavoro; lo studio verrà presentato più ampiamente al \emph{Career
        Day} dell'accademia insieme alla sua HR manager;
  \item fa inoltre consulenze aziendali anche al di fuori dei videogiochi e
        insegna in master universitari (IULM di Milano).
\end{itemize}

Il suo ruolo ricorrente, sia a Blizzard sia a Riot, è stato quello di
\emph{Mr.~Wolf} (``per chi ha visto \emph{Pulp Fiction}''): quello che
risolve i problemi. Veniva inserito nei team più in difficoltà, ad esempio
dal punto di vista delle performance, per sistemarli e passare al successivo.
Era presente in alcuni dei momenti più critici, come il lancio di
\emph{Diablo~III} o il lancio di espansioni di \emph{World of
Warcraft}\footnote{La trascrizione riporta ``il lancio con i Gary's'' e
``World of Trenor'': si tratta verosimilmente di \emph{Mists of Pandaria} e
di \emph{Warlords of Draenor}.}, a fare \emph{damage control}.

La scoperta fondamentale di questi anni è che i problemi da risolvere
\textbf{quasi mai riguardavano il prodotto} (non era certo lui a sistemare i
server di \emph{Diablo~III}), ma i team che c'erano dietro: persone che
litigavano, non collaboravano, non riuscivano a stare nella stessa stanza
senza urlarsi contro. E nel 90\% dei casi tutto questo si riduce a
\textbf{problemi di comunicazione}: è raro che perfino le incompatibilità
caratteriali non siano risolvibili parlando.

\subsection{Struttura non vuol dire burocrazia}

A Stormind il relatore sta accompagnando un team passato da 30 a 95 persone
in meno di un anno, attraverso tutti i cambiamenti di cui parla il workshop.
Darsi una struttura viene spesso percepito male, come ``più processi, più
procedure'', ma bisogna distinguere:
\begin{itemize}
  \item un processo come ``timbrare il cartellino'' può essere davvero
        invasivo e noioso;
  \item definire \emph{come funziona un documento di game design, chi lo
        approva, chi lo visiona e chi no} è invece una struttura necessaria.
\end{itemize}
Uno studio di 30 persone nella stessa sede può farne a meno: quando c'è un
dubbio ci si alza e si va alla scrivania del collega. Uno studio di 95
persone sparse per l'Italia (e in parte per l'Europa) non può più farlo:
servono una struttura diversa, un modo di lavorare diverso e una
pianificazione del futuro diversa.

% ---------------------------------------------------------------------
\section{Perché il lavoro è sempre lavoro di squadra}
\label{sec:z1-squadra}

\subsection{Non si va in ufficio per fare amicizia}

Il primo messaggio è netto: \textbf{non si va in ufficio per fare amici} o
per trovare compagni di vita. Si va in ufficio perché c'è un lavoro da fare, e
lo scopo è riuscire a farlo \emph{anche con chi sta antipatico}. Con le
persone simpatiche è facile: diventeranno amici, e non sarà mai un problema.
È con chi sta antipatico che le cose si complicano, ed è lì che emergono i
\textbf{veri professionisti}: quelli che riescono a lavorare con chiunque
senza che l'antipatia pesi sul modo in cui lavorano e collaborano, mettendo il
prodotto e il team davanti al proprio ego e alle proprie preferenze personali.

La frase che fa da filo conduttore a tutto il workshop (e che il relatore
promette di riproporre ``fino alla nausea'') è:

\begin{quote}
  \emph{La capacità di un gruppo di persone di fare cose straordinarie
  dipende dal modo in cui queste persone operano insieme.}
\end{quote}

Le aziende raramente pensano che l'innesto di una singola persona possa
cambiare il destino di un progetto. Questo è tanto più vero quanto più grande
è il team, ma vale anche per i piccoli studi, anche per un gruppo di cinque
persone. Come negli sport di squadra: l'individuo conta fino a un certo punto,
e anche se si ha Dončić in squadra, se si gioca uno contro cinque si perde
comunque\footnote{Il riferimento è alle finali NBA in corso nei giorni del
workshop tra Dallas Mavericks e Boston Celtics (giugno 2024), il che
permette di datare approssimativamente la lezione.}.

\subsection{Lo \emph{useless genius} e il costo delle persone tossiche}

Molti dipartimenti HR usano la definizione di \textbf{useless
genius}\footnote{Il concetto è più noto in letteratura manageriale come
\emph{brilliant jerk}; in lezione è usata l'espressione \emph{useless
genius}.}: si può essere bravissimi, a volte i migliori nel proprio lavoro,
ma se non ci si sa comportare, relazionare e operare in un gruppo si è
fondamentalmente inutili ai fini del gruppo. Nel migliore dei casi inutili,
nel peggiore \emph{dannosi}.

In passato le aziende tendevano a ``chiudere un occhio'': la persona è
fondamentale, la teniamo. Oggi la domanda è diventata: \emph{quanto costa
all'azienda, oltre al suo stipendio, tenere una persona che crea un ambiente
tossico e peggiora chi le sta intorno?} Ci sono costi sommersi che le aziende
cominciano a vedere, e la tendenza è liberarsi di queste persone il prima
possibile. Il ragionamento è semplice:
\begin{itemize}
  \item le \textbf{competenze si possono trasmettere} o sostituire: si trova
        qualcuno che fa lo stesso lavoro magari nel doppio del tempo, oppure
        più persone che se lo dividono;
  \item il \textbf{carattere delle persone} invece cambia molto raramente.
\end{itemize}

\subsection{Il fit culturale}

Quando qualcuno \textbf{non è un \emph{fit} culturale} per un'azienda è
meglio separarsi presto. Ma non essere adatti a \emph{quell'}azienda non
significa non essere adatti a lavorare: è perfettamente legittimo che una
persona inadatta a un'azienda sia adattissima a un'altra.

Il relatore l'ha vissuto in prima persona. Dopo sei anni a Blizzard viene
contattato da Riot e affronta un processo di selezione lunghissimo (``quasi
undici colloqui''), al termine del quale viene scartato. Scrive allora una
mail che gli ha ``cambiato la traiettoria di carriera'': chiede
semplicemente perché non è stato ritenuto adatto. La risposta: in nessun
colloquio si era esposto abbastanza da far capire \emph{perché} volesse
lasciare Blizzard, e quindi pensavano che in realtà stesse bene dov'era. Il
motivo era che, per educazione professionale, non gli piace parlare male del
proprio posto di lavoro; portata all'eccesso, però, questa riservatezza era
diventata un ``è tutto bellissimo, ma vorrei provare Riot'', senza un perché.
Lo invitano comunque a un incontro dal vivo (``dal vivo è diverso''), e alla
fine viene assunto.

Una volta arrivato, i sei mesi di prova (\emph{probation}) sono serviti tutti
per smettere di ragionare e lavorare ``come a Blizzard''. La differenza
culturale era enorme, tanto che in Riot esisteva un gruppo di supporto
interno tra ex dipendenti Blizzard: un collega gliel'ha descritto come
``uscire da una relazione tossica e doversi riabituare a una relazione
normale''\footnote{Il relatore precisa che parla della Blizzard di allora,
che non esiste più: oggi è parte di Microsoft e nessuno dei suoi ex colleghi
vi lavora ancora, anche a causa dei licenziamenti di massa.}.
\begin{description}
  \item[Blizzard] era \emph{developer first}: il focus era chiaramente sugli
        sviluppatori, prima di tutti gli altri. Dopo sei anni c'erano ancora
        meeting a cui non poteva partecipare perché non abbastanza
        \emph{senior}, informazioni che non poteva sapere, piani che non
        poteva vedere.
  \item[Riot] era \emph{player first}: prima i giocatori, e i dipendenti come
        mezzo per trasferire il feedback dei giocatori dentro il gioco. Il
        secondo giorno lo portano in una stanza e gli mostrano il piano dei
        cinque anni successivi. Lui, sorpreso, pensa a un errore (``sono qui
        da oggi, non vorrei mettervi nei guai''); la risposta è: ``sei un
        dipendente Riot o no? Allora non ci sono livelli di accesso: devi
        sapere queste cose per fare meglio il tuo lavoro''.
\end{description}

Da questa esperienza ha imparato molto sulla \textbf{trasparenza}: condividere
le informazioni aiuta il lavoro. Non sempre, però: a volte non si condivide
qualcosa non per segretezza, ma perché non è ancora pronta per essere
presentata, esattamente come non si mostra un gioco quando è ancora solo
teoria (``te lo faccio vedere quando ho qualcosa di tangibile'').

La lezione per gli studenti è doppia: le aziende valuteranno il \emph{fit}
culturale sempre di più, ma capire se si è un \emph{fit} richiede anche
\textbf{lavoro da parte del candidato}. Bisogna avere la presenza di spirito
di dire: ``conosco i valori di questa azienda, non mando nemmeno il CV,
perché so che durerei un mese e mezzo''.

% ---------------------------------------------------------------------
\section{Che cos'è il feedback}
\label{sec:z1-cose-feedback}

\subsection{Le prime risposte della classe}

La prima domanda alla classe è: \emph{quando sentite la parola
``feedback'', a cosa pensate?} Le risposte raccolte, a voce e in chat:
\begin{itemize}
  \item un \textbf{riscontro} su ciò che si fa e su come lo si fa, sul lavoro
        o sul modo in cui ci si interfaccia con il gruppo; può essere
        lavorativo oppure comunicativo e comportamentale. Chi ha risposto lo
        vede come qualcosa che, se l'azienda non dà, va attivamente
        richiesto, perché è importante sapere come sta andando per
        migliorarsi;
  \item né positivo né negativo, ma \textbf{costruttivo}: in entrambi i casi
        serve a migliorarsi;
  \item uno \emph{scambio di opinioni}, un \emph{riscontro produttivo};
  \item un feedback ``di reazione'', non verbale: si fa una cosa e si
        osserva come cambia l'ambiente o il comportamento delle persone
        intorno (anche perché a voce la gente spesso si tiene le cose
        negative);
  \item per un designer, anche il \textbf{linguaggio del corpo}: osservando
        un giocatore giocare si può dedurre l'emozione che prova (frustrazione,
        incomprensione).
\end{itemize}

Sulla forma (scritto o orale?) le opinioni divergono: c'è chi lo preferisce
\emph{mille volte} scritto; chi pensa che dipenda dalle dimensioni
dell'azienda (orale nelle piccole, scritto in quelle con team di centinaia di
persone); chi pensa che dipenda dall'importanza del feedback (scritto se
tecnico o pieno di dati, da tenere sotto gli occhi; orale se è un consiglio
veloce).

Il relatore usa due esempi per mostrare che \emph{dipende}:
\begin{itemize}
  \item un feedback sulla programmazione ha senso darlo scritto, magari
        commentando direttamente il codice, più che ``andando a prendere un
        caffè'';
  \item ma se il feedback è ``ieri, durante il meeting, ho percepito male il
        fatto che mi interrompessi di continuo'', ce lo si aspetta
        \textbf{orale}.
\end{itemize}

\subsection{Da che cosa dipende il modo di dare feedback}

Lo \emph{spoiler} della giornata è che \textbf{non esiste una risposta univoca}
su come dare feedback. Il workshop non insegna ``il'' modo, ma fornisce
\emph{tutte le opzioni}, perché il modo giusto cambia in base a molti fattori:
\begin{itemize}
  \item a chi lo si dà: un collega, un superiore, il proprio capo, un proprio
        \emph{report} (un sottoposto);
  \item se è personale o professionale;
  \item se è la prima volta o la terza volta che si dà lo stesso feedback;
  \item se riguarda un progetto importante o secondario.
\end{itemize}

Manca però l'elemento fondamentale, che la classe individua subito:
\textbf{l'altro}. A parità di tutti gli altri fattori, il \emph{carattere
della persona che riceve il feedback} dovrebbe far cambiare il modo in cui lo
si dà. Per questo, prima della lezione, agli studenti era stato chiesto di
fare un test di personalità: non un test di psicologia avanzata, ma uno
strumento per rendersi conto di quanta varietà ci sia anche solo in una
classe di 53 persone (il che non vuol dire trovare 52 modi diversi di dare
feedback, ma lavorarci un po' sì).

Il relatore porta se stesso come esempio. Preferisce:
\begin{itemize}
  \item i feedback \textbf{positivi scritti}, perché è a disagio a ricevere
        complimenti in pubblico;
  \item i feedback sul proprio lavoro \textbf{diretti}: ``vai subito al
        punto e non girarmi intorno''.
\end{itemize}
È molto orgoglioso del suo lavoro e, quando qualcuno lo critica, la sua
primissima reazione istintiva è ``come ti permetti''. Negli anni ha imparato
a farla durare un secondo, ma chi gli dà feedback di persona dice di
vederla ancora negli occhi. Per mitigarla, dice esplicitamente a chi gli dà
feedback: \emph{``ignora il modo in cui reagisco nei primi cinque secondi,
vai avanti, poi mi passa''}. È considerato da tutti molto aperto al feedback
e non ha problemi a dire ``hai ragione, facciamo diversamente'': ma ci sono
reazioni istintive e caratteriali che è difficile controllare al cento per
cento.

\subsection{Con chi si lavora meglio}

Guardando le persone con cui si lavora abitualmente (anche nei progetti
dell'accademia), si noterà che si lavora meglio con due profili:
\begin{enumerate}
  \item chi è \textbf{esattamente come noi}, ma qui è un 50 e 50: alcuni
        caratteri vanno d'accordo con i propri simili, altri non vogliono
        nessuno con lo stesso carattere nel gruppo (e si potrebbe rientrare
        in questa categoria);
  \item chi è \textbf{affine}: ha qualcosa in comune con noi, ma non tutto,
        e proprio per questo può dare punti di vista differenti.
\end{enumerate}
Il lavoro della giornata è capire come valorizzare il secondo caso: ottenere
feedback da persone con caratteri, visioni e background diversi, per
completare una visione del proprio lavoro che è necessariamente limitata. È
come esplorare una mappa con la nebbia di guerra: un'altra persona, al di là
del muro, vede cose che noi non vediamo.

\subsection{Il feedback come somma di piccoli miglioramenti}

Per il relatore il feedback è questo: \textbf{applicare tanti piccoli
miglioramenti periodici}, che si sommano nel tempo. Migliorare dello 0,1\%
ogni giorno porta a fine anno a un miglioramento enorme, e questo vale per i
prodotti, per i giochi, per le persone e per i team.

\paragraph{L'esempio del ciclismo britannico.} La squadra di ciclismo
britannica, fino a una quindicina d'anni fa, non aveva vinto quasi nulla.
Cambiato il team manager, tutti si aspettavano grandi cambiamenti sportivi
(convocazioni, strategie di gara). Arrivarono invece \emph{minuscoli
cambiamenti non sportivi}: un altro chef in trasferta, materassi migliori nei
camper per riposare meglio, un sellino più comodo a fine tappa. Cambiamenti
marginali che da soli non avrebbero cambiato niente; eppure, dopo due o tre
anni di questo lavoro quotidiano, la squadra ha cominciato a vincere a
ripetizione, Olimpiadi comprese, cambiando il modo in cui il ciclismo è
visto in Gran Bretagna\footnote{È la nota strategia delle \emph{marginal
gains} associata a Dave Brailsford (British Cycling e Team Sky); il nome non
è citato in lezione.}. Certo, si è investito anche sulle persone giuste, ma
tutto è partito dai piccoli miglioramenti periodici.

Per gli studenti vale lo stesso: chiarire un conflitto con una persona con
cui si ha ancora poco a che fare, o una situazione che ha lasciato
``l'amaro in bocca'', sembra inutile oggi, ma darà i suoi frutti tra tre, sei,
nove mesi. Sono \emph{investimenti sul futuro}.

\paragraph{Peggiorare ogni giorno.} Esiste anche il percorso inverso:
lasciar correre, ``buttare sotto il tappeto'' le cose perché ``tanto è
inutile parlarne''. Queste cose tornano sempre fuori dopo otto mesi, come il
\emph{libro dei rancori} dei nani di Warhammer: ``ti ricordi quando, un anno
e mezzo fa, all'inizio del progetto, mi hai detto quella cosa? Mi sono
offeso''. A quel punto è quasi impossibile risolvere, perché nessuno ricorda
più con precisione cosa sia successo (``ma quale cena?'', ``forse era in
pizzeria'', ``no, eravamo da un'altra parte''). \textbf{Tenersi le cose nel
cassetto per tirarle fuori quando servono è sempre sconsigliato}: non servono
ed è impossibile usarle in modo costruttivo.

% ---------------------------------------------------------------------
\section{Le cinque fasi di vita di un team}
\label{sec:z1-fasi-team}

\subsection{Perché conoscerle}

Quasi tutti gli studenti lavorano già in gruppo (al primo anno su un gioco da
tavolo, negli anni successivi su un videogioco) e si saranno accorti che un
team attraversa delle fasi: il modo di lavorare insieme cambia nel tempo.
Spesso ci si fa prendere di sorpresa: ``eravamo così affiatati, cosa è
successo?''.

Un esempio da una consulenza fuori dai videogiochi: un'azienda farmaceutica
gli affida un team commerciale (rappresentanti che presentano medicinali a
ospedali, medici e privati). Il problema: ``prima andava tutto bene; ora che
siamo in diciotto dovremmo lavorare meno e meglio, e invece va peggio''. Va
peggio perché le persone nuove sono appena arrivate, mentre il management si
aspettava già di vedere i risultati dell'investimento.

Fino a pochi anni fa le fasi di un team venivano disegnate come una
\textbf{linea dritta verso l'alto}: le fasi erano le stesse, ma viste come un
miglioramento costante. Non si vedevano le parti \emph{fisiologiche} del
lavoro in team, cioè anche i peggioramenti. Accettarle è il primo passo per
accettare i fallimenti: nessun percorso personale o di team è lineare.

Il relatore ammette che il suo modo di lavorare è nato per il 70\% dal
proprio impegno a imparare, fare bene e rispettare le persone, e per il 30\%
dall'aver lavorato con ``incapaci colossali'' che gli hanno insegnato come
\emph{non} fare le cose. A Blizzard un manager gli fece capire che, pur non
sapendo ancora che tipo di \emph{people manager} sarebbe diventato, sapeva
benissimo quale tipo di manager \emph{non} voleva essere.

\subsection{Forming, storming, norming, performing, adjourning}

Le fasi sono cinque, un po' come le fasi del lutto. Molte aziende ne
considerano solo quattro (``e sono quelle che solitamente chiudono''); le
aziende al passo con i tempi ne considerano cinque\footnote{Il modello è
quello di Bruce Tuckman (1965), a cui la quinta fase, \emph{adjourning}, fu
aggiunta nel 1977. Il nome di Tuckman non è citato in lezione; la
trascrizione riporta la quinta fase come ``journing''.}.

\begin{description}
  \item[Forming] Si viene messi insieme e si comincia a lavorare. Non si sa
        nulla degli altri; spesso non si conoscono nemmeno i nomi, ed è
        normale.
  \item[Storming] Si cominciano a ``dare le testate come le capre in
        montagna per determinare il territorio'': questa cosa la faccio io,
        questa la vorrei fare così, chi decide? Si stabilisce chi fa cosa e
        come si divide il lavoro. Nel grafico è una \textbf{fase di
        peggioramento}: che il team litighi o discuta civilmente non cambia
        nulla, il team deve investire parte di tempo ed energie per risolvere
        queste questioni. È normale e \textbf{non è evitabile}.
  \item[Norming] Si comincia a capire come lavorano gli altri, cosa sanno
        fare e come far funzionare il gruppo; iniziano i miglioramenti. È una
        forma di \emph{accettazione} degli altri, che deriva dall'aver capito
        cosa devo fare io e cosa devono fare loro.
  \item[Performing] C'è uno sforzo \textbf{proattivo}: ``questa cosa
        toccherebbe a me, ma so che Alessio la fa molto meglio, quindi gli
        chiedo aiuto''. Si sfruttano i punti di forza degli altri, perché li
        si conosce e perché si sa che passare il testimone non avrà effetti
        negativi (il capo non dirà ``perché non l'hai fatta tu?'').
  \item[Adjourning] Il progetto finisce, si chiude o viene riorganizzato, il
        team viene sciolto. Quello che andrebbe fatto è analizzare
        \textbf{cosa ha funzionato, cosa no e cosa migliorare} per il
        progetto successivo.
\end{description}

\subsection{La fase che tutti saltano}

``Spoiler incredibile'': l'\emph{adjourning} è la fase che molte aziende
saltano, \textbf{perché c'è sempre il progetto dopo}.

A PG Esports, per la Milano Games Week (come per qualsiasi evento dal vivo a
cadenza annuale: Lucca Comics, Etna Comics, Comicon\dots), finita l'edizione
di un anno si cominciava dal giorno dopo a lavorare a quella successiva,
perché un anno sembra tanto ma non basta mai. Nei videogiochi è lo stesso:
finito un gioco, a parte le patch di mantenimento e del \emph{day one}, si
passa subito al successivo.

Il cambiamento introdotto dal relatore (in parte già praticato, ma da lui
strutturato) è stato: finita la Games Week, \textbf{tutti in ferie per una
settimana}; poi ci si ritrova per definire cosa ha funzionato e cosa no, con
un \textbf{post-mortem} (o \emph{debrief}); solo dopo si inizia a lavorare
all'edizione successiva. Senza questa fase manca la parte di
\emph{learning}, cioè di ciò che si è imparato dal progetto.

Il consiglio vale per gli studenti: quando si finisce il gioco da tavolo, e
poi il gioco digitale, bisogna trovare il tempo per il post-mortem,
altrimenti si butta via \textbf{il 90\% dell'utilità del progetto}. Una
struttura semplice per farlo è descritta nella sezione~\ref{sec:z1-qa-finale}.

% ---------------------------------------------------------------------
\section{Domande sulle fasi del team}
\label{sec:z1-qa-fasi}

\subsection{Lo storming che ritorna}

\paragraph{Capita che lo storming si ritagli spazio anche nelle altre fasi?}
Lo \emph{storming} è una fase precisa, che però molti manager, coordinatori e
team cercano di \textbf{saltare}: ``se io sono bravo come manager e loro sono
bravi come team, passeremo da forming a norming senza problemi''. È un errore,
perché senza una fase in cui si discute davvero non si sviluppa una
\emph{skill} fondamentale: \textbf{saper discutere}. Tutti conoscono persone
non abituate a farlo, per cui ogni minimo disaccordo diventa un affronto
personale (``avevamo deciso la pizza, perché andiamo al sushi? Perché mi fate
questo affronto?''): sono persone che non sono abituate ad avere conflitti
costruttivi.

Se si pensa di essere in \emph{norming} e continuano a emergere conflitti, il
punto è che \textbf{non si era davvero in norming}: lo storming non ``si
trascina'' nelle fasi successive, semplicemente non se n'è mai usciti.
Allo stesso modo, a chi chiede se ogni mini-fase di storming durante la vita
del gruppo sia frutto di uno storming fatto male o solo di carattere, la
risposta è già contenuta qui.

Un partecipante osserva che spesso, ancor prima di finire un progetto, si è
già messi su quello successivo: è il modo migliore per essere in costante
\emph{damage control}. Era l'esperienza del relatore a Blizzard: veniva
spostato fisicamente da un team all'altro, e finito il lancio di
\emph{Diablo~III} era già su \emph{World of Warcraft}, senza vivere il
``dopo''. Si è dovuto ritagliare da solo lo spazio per un debrief personale,
ed è lì che ha deciso che tipo di manager non voleva essere. Il consiglio:
\textbf{se non lo fa il team, fatelo voi}.

\subsection{Dove si passa la maggior parte del tempo}

\emph{In quale fase si passa più tempo nella propria carriera?} Gli studenti
rispondono \emph{storming} e \emph{norming}. Il relatore, in quindici anni,
calcola di aver passato in \emph{performing} circa \textbf{dieci mesi}. Non è
pessimismo: arrivare a performing deve essere sempre l'obiettivo, sia da
manager sia da membro del team; ma è difficile arrivarci, ed è normale che non
sempre ci si arrivi.

Lo storming può durare a lungo, ma con team molto esperti dura pochissimo,
perché c'è chi lo sa gestire. In realtà \textbf{la maggior parte del tempo si
passa in forming}, perché ogni equilibrio raggiunto vale per \emph{quelle}
persone, \emph{quel} progetto, \emph{quelle} tempistiche. Basta che cambi un
qualsiasi fattore (una deadline spostata, una persona nuova, qualcuno che si
dimette, un progetto cancellato) e si ricomincia, perché i bilanciamenti
precedenti non valgono più.

Esempio: un team di cinque persone affiatatissimo, in pieno performing. La
dirigenza, entusiasta, ``raddoppia le risorse'' e lo porta a dieci. O si
lascia tutto com'è (e i cinque nuovi non fanno niente), o bisogna ridiscutere
chi fa cosa e come i nuovi possono aiutare: che piaccia o no, si ricomincia
da forming.

Quindi la bravura di un manager, ma anche di un membro del team,
\textbf{non sta nel passare tanto tempo in performing, ma nell'accelerare il
passaggio da una fase all'altra}, e nel fare bene la fase di
\emph{adjourning}.

\subsection{Analizzare anche i successi}

C'è la tendenza a fare post-mortem solo dei progetti andati male: ``questa
cosa non ha funzionato, vediamo perché''. Va benissimo. Ma quando un progetto
va bene si dà per scontato che sia andato bene ``perché era buono''. Bisogna
invece chiedersi \emph{perché} è andato bene: era il prodotto? La
tempistica? Il modo in cui abbiamo fatto una cosa piuttosto che un'altra?
\textbf{Se non si capisce davvero il motivo del proprio successo, si
tenderanno a replicare le cose sbagliate.}

\paragraph{Blockbuster.} Quando arrivò la notizia che Netflix avrebbe
lanciato un servizio di streaming, Blockbuster (la catena di noleggio di
videocassette e DVD) aveva tutto il tempo per preparare un'offerta
concorrente. In una riunione del management raccontata in un libro (di cui il
relatore non ricorda il titolo), la direzione sembrava orientata a investire
nello streaming, finché un dirigente fece notare che il 12\% del fatturato
derivava dalle \emph{multe per riconsegna tardiva}, impossibili con lo
streaming, mentre secondo le stime lo streaming avrebbe prodotto il 10\% del
fatturato. ``Non possiamo rinunciare a questo 2\%'': niente streaming. Oggi
non c'è più un Blockbuster ``manco a cercarlo col lanternino''. Non avevano
capito qual era il vero motivo del loro successo, cioè il loro modello di
business\footnote{Le percentuali sono quelle riferite in lezione e non sono
state verificate.}.

\paragraph{The Founder.} Nel film sulla nascita di McDonald's c'è una scena in
cui il futuro socio del protagonista gli dice: ``tu pensi di essere nella
\emph{food industry}, e invece sei nella \emph{real estate industry}''. Non
dovrebbe investire sugli hamburger, ma sui terreni su cui costruisce i
ristoranti. Anche per un team vale lo stesso: bisogna capire che cosa
l'ha reso davvero di successo, oppure no.

\subsection{L'autoanalisi personale}

\paragraph{E quando si cambia azienda?} Uno studente chiede se il debrief sia
possibile anche quando si lascia un'azienda. Non si può tornare in ufficio a
discutere della propria esperienza, ma lo si può fare \textbf{per se stessi}.
Il relatore l'ha fatto, anche se all'epoca in modo poco consapevole, ed è
così che è passato dal \emph{community management} al \emph{people
management}. Analizzando cosa aveva fatto bene a Blizzard si era reso conto
che:
\begin{itemize}
  \item aveva costruito un ottimo rapporto di fiducia con i giocatori (con
        alcuni proprietari di fansite parla ancora oggi);
  \item pur senza la seniority, era diventato un punto di riferimento per il
        team, visto come imparziale ed equilibrato anche su temi che non lo
        riguardavano direttamente;
  \item sapeva \textbf{risolvere problemi}: durante il lancio di
        \emph{Diablo~III} era il più tranquillo di tutti, non perché
        sminuisse il problema, ma perché sapeva che si sarebbe risolto con i
        suoi tempi; il lavoro non era spegnere l'incendio ma ``far sfollare
        la gente'' e aspettare che tutto tornasse in ordine.
\end{itemize}
E aveva capito cosa sapeva fare meno bene, ad esempio \textbf{gestire il
proprio manager}. Avere un manager non significa che il manager gestisca te:
il manager gestisce il tuo lavoro, ma tu devi saper gestire i tuoi
\emph{stakeholder}. Spesso lo stakeholder è il proprio manager; da manager,
gli stakeholder saranno CEO, CFO e così via. Gestirli significa conoscerli,
sapere come portare loro le informazioni, come far presente che qualcosa non
va. Fa parte del lavoro, anche se a lui non piace, e ci ha lavorato.

Nel mese e mezzo di pausa prima di entrare in Riot si è concentrato su
questo: ``cosa porto come mio punto di forza?''. L'autoanalisi è culminata
nel chiedere espressamente alla persona che l'aveva assunto: ``prima mi avete
detto di no, ora mi avete preso: perché?''.

\paragraph{Come si capisce se un progetto è andato bene?} Dipende
dall'\textbf{obiettivo}. Se l'obiettivo era vendere milioni di copie e il
gioco le ha vendute, bisogna capire perché. Se, più ragionevolmente,
l'obiettivo era un gioco da tavolo che permettesse di mettere in pratica le
conoscenze del corso, va valutato per quello: ha permesso di farlo? Se sì,
perché? Se no, perché no? L'obiettivo va fissato \textbf{all'inizio del
progetto, non alla fine}, proprio per poterlo valutare dopo.

\subsection{Quando il team cambia di continuo}

\paragraph{Tutti questi cambiamenti non sono dannosi per il progetto?}
Assolutamente sì, ma quasi mai sono volontari. I progetti andati male sono
spesso progetti grandi.

Il relatore racconta di \emph{Cyberpunk 2077}. Conosceva persone in CD
Projekt Red (``il mondo è piccolissimo, le cose si sanno''). Alla Gamescom va
con Diego Ricchiuti, direttore dell'accademia, a vedere la presentazione del
primo gameplay, quella ``in cui regalavano il giubbottino da samurai''.
Finito il video, tra gli applausi, si guardano e capiscono entrambi che
\textbf{il gioco non è pronto}, a soli quattro mesi dall'uscita prevista:
Diego aveva notato il combattimento e la guida. Gli amici di CD Projekt Red
(uno aveva lavorato con lui a Blizzard per sei anni) gli confermano tra le
righe che ``ce la faremo, ma ci sono ancora cose da sistemare''. Il gioco fu
poi rimandato più volte\footnote{In lezione il relatore dice ``due, forse
tre volte''; l'uscita di \emph{Cyberpunk 2077} è stata effettivamente
rinviata tre volte prima del dicembre 2020.}.

Per deformazione professionale, chi ha sviluppato l'occhio critico non
riesce più a godersi un trailer: in un trailer di \emph{The Day Before} Diego
gli aveva fatto notare che l'effetto della fiammata di un fucile era stato
chiaramente aggiunto in postproduzione e non faceva parte del gameplay. ``È
il triste epilogo di questa storia'': non riescono più a godersi un trailer.

Il punto è che più il team è grande, più queste variabili vengono imposte, e
più è facile che il fallimento \textbf{non dipenda dal fatto che il prodotto
sia brutto} (\emph{Cyberpunk} è un bellissimo gioco sotto molti punti di
vista, e il relatore lo sta rigiocando ora che è stato corretto ed espanso).
Nei grandi gruppi qualcuno arriva, qualcuno si dimette, cambia il budget,
arriva il rinvio; nel caso di \emph{Anthem}, il team ha fatto così tante volte
forming--storming--forming--storming che ne è uscito ``un prodotto che
dovevamo espellere: l'abbiamo espulso, buttato fuori dal nostro organismo pur
di liberarcene''.

Tra i cambiamenti c'è anche la crescita verticale: se il capo viene promosso
perché lavora benissimo, bisogna scegliere un nuovo capo, ed è sempre un
caos (chi ci resta male, chi pensava di meritarlo, chi voleva il suo amico,
chi non sopporta il prescelto). Per questo il relatore è \textbf{contrario
alle promozioni durante i progetti}: si rischia di rovinare tutto il lavoro di
storming.

\subsection{Obiettivi chiari nei piccoli gruppi}

Nei gruppi piccoli, come quelli degli studenti, il problema deriva più
spesso dal \textbf{non avere chiaro che cosa si vuole ottenere}: un prodotto
puramente educativo e formativo, o un prodotto richiesto dall'accademia che,
se viene bene, si potrebbe anche vendere? Un gioco commerciale o un gioco
artistico? Va chiarito subito, perché ci saranno sempre momenti in cui
scegliere tra A e B, e a guidare la scelta sarà la domanda: \emph{quale delle
due opzioni ci porta più vicino all'obiettivo?} Se l'obiettivo non è chiaro,
non si sa cosa scegliere.

Questo porta molte aziende a fare \textbf{scelte di design che in realtà sono
scelte commerciali}. Le microtransazioni sono scelte commerciali: nessuno ha
mai pensato ``sai cosa piacerebbe tanto ai giocatori? Pagare le bustine''.
Prima le bustine si guadagnavano giocando; ora la tendenza è ``potresti
giocare, ma se paghi è meglio''. Si pensi alle polemiche su Ubisoft per
l'opzione di ``saltare il gioco'' pagando, o ai famosi \emph{boost} di livello
di \emph{World of Warcraft}: ``devo pagarti per saltare il gioco?''.

Una studentessa (Silvia) cita \emph{Final Fantasy XIV}, che a un certo punto fu
``resettato'' perché non funzionava: un esempio positivo. Il relatore aggiunge
quello negativo di \emph{Age of Conan}, che non è mai riuscito a fare quel
reset iniziale (``è una schifezza, ve lo ripresentiamo quando è pronto''). Ci
vuole consapevolezza, ma anche i mezzi: un conto è Square Enix, che può
permetterselo, un altro è un'azienda che non può chiudere tutto per un anno
continuando a pagare gli stipendi.

Riassumendo: nei piccoli gruppi i problemi sono più di
\textbf{comunicazione e chiarezza} (obiettivi non chiari o che si spostano);
nelle grandi aziende di più di \textbf{persone che vanno e vengono, obiettivi
e tempistiche che cambiano}. In accademia le scadenze sono fisse (il gioco da
tavolo si consegna a fine anno scolastico, a giugno, e giugno resta); in
azienda può succedere che una feature prevista per la \emph{vertical slice}
di luglio venga anticipata a maggio per portarla alla Gamescom, e un sistema
perfetto per consegnare a giugno può non esserlo più per consegnare a maggio.

\subsection{L'ambiente determina la qualità del prodotto}

\paragraph{È vero che chi passa per Electronic Arts finisce per fare brutti
giochi?} Uno studente cita BioWare, che faceva ottimi giochi finché era
autonoma. Il problema non è dei developer, che non sono ``improvvisamente
diventati incapaci'': è evidente che il problema è altrove, e che
\textbf{l'ambiente in cui si fa un gioco determina molto spesso la qualità del
prodotto}.

Un esempio locale positivo: Tommaso Bonanni, conosciuto dal relatore insieme a
Diego e già docente dell'accademia, ha fondato a Roma Caracal Games Studio,
applicando fin dall'inizio questi concetti di cultura aziendale. L'azienda è
rimasta piccola, ma ha molto chiaro cosa vuole fare e come. Quando si
predispone fin dall'inizio l'azienda e il team a lavorare in un certo modo è
fattibilissimo; è molto più difficile quando si è EA, si acquisisce
un'azienda e le si dice ``da oggi si fa così''.

Il relatore ha vissuto la transizione da Blizzard ad Activision Blizzard (uno
dei motivi per cui se n'è andato, volontariamente). Ha visto la cultura
interna cambiare per gradi:
\begin{enumerate}
  \item ``facciamo giochi fighissimi, perché i giochi devono essere
        fighissimi'';
  \item ``facciamo giochi fighi così vendono di più'';
  \item ``vendiamo di più facendo giochi fighi'' (sembra uguale, ma è molto
        diverso);
  \item ``facciamo più soldi''.
\end{enumerate}
Lo si capiva dalle scelte di design, osservando quale fosse il peso più
grande sulla bilancia. Lo stesso varrà per gli studenti, nei loro giochi e nei
colloqui di lavoro: capire subito che cosa vuole l'altra parte.

\paragraph{Leader e visione.} Nei gruppi dell'accademia, dal secondo anno, c'è
una persona che assume il ruolo di \emph{lead} o \emph{scrum master}. È
diverso dall'essere un \emph{game director} (``il gioco si fa come dico io,
aiutatemi a farlo con la mia visione''), come avviene nei grandi studi: da lì
nascono i nomi nei titoli, come \emph{Sid Meier's Civilization}, che indicano
la persona che ha fatto quel gioco, oppure un puro branding, come \emph{Tom
Clancy's Rainbow Six}.

\paragraph{Inclusion come jolly.} A margine, uno studente osserva che in molte
aziende ``la risposta giusta è sempre \emph{inclusion}'', usata come un
jolly. È comunissimo, conferma il relatore: \emph{diversity and inclusion}
viene spesso usata per dire ``così includiamo più clienti'', ma non è quello
il suo punto (``altro discorso'').

% ---------------------------------------------------------------------
\section{Confrontarsi con chi è più bravo}
\label{sec:z1-rivalita}

Arrivato a Riot dopo sei anni a Blizzard, il relatore pensava di avere un buon
bagaglio di esperienza. Si è trovato invece ``come un fante buttato in mezzo
ai Navy SEAL'': persone che avevano superato 11, 12, 13 colloqui, con un
livello molto alto.

Individua subito un collega francese che, nella sua testa, diventa ``la
versione migliore di me'': stesso ruolo (community manager), esperienza
simile (arrivava da Ubisoft, come quasi tutti i francesi di Riot), ma molto
più spigliato, sempre con la battuta pronta, ``l'anima della festa''. Il
relatore invece è introverso: un workshop come questo per lui è
\emph{emotivamente drenante}, non ama parlare in pubblico pur amando
condividere le proprie esperienze, e non riesce a sostenere conversazioni
come ``cosa hai fatto questo weekend?''. Lavorando raramente insieme, aveva
sviluppato una rivalità ``da allenatore di Pokémon'' (lui era Ash, il
relatore l'antagonista).

Dopo sei mesi, a un evento di \emph{team building} aziendale (a Tenerife), il
suo capo, molto attento a queste dinamiche, li mette apposta nello stesso
gruppo. A fine serata, complice l'open bar, il relatore gli confessa di averlo
mal sopportato dal primo giorno per invidia della sua estroversione. Il
collega ride e risponde che pochi giorni prima aveva detto al capo la stessa
cosa al contrario: invidiava a Stefano il suo essere sempre organizzato, sul
pezzo, presentarsi nella sala riunioni giusta (mentre lui finiva in quella
sbagliata), arrivare con gli appunti, con i dati a supporto e le analisi di
mercato, mentre lui spesso non aveva nemmeno letto la mail del meeting.
Un'invidia in due direzioni, che li ha portati a lavorare molto insieme e a
migliorarsi a vicenda.

Le lezioni:
\begin{itemize}
  \item non bisogna demoralizzarsi vedendo qualcuno che fa lo stesso lavoro
        ed è più avanti sotto alcuni aspetti: spesso c'è un bilanciamento. I
        ``punti caratteristica'' sono limitati: se li ha messi in carisma,
        non li avrà messi in saggezza;
  \item conviene \textbf{chiedere a queste persone come mai} sono così brave
        in qualcosa, come si fa con un progetto andato bene: magari non lo si
        farà mai bene quanto loro, ma si comincerà ad avere il
        \emph{mindset} giusto; e viceversa, loro saranno spinti a chiedere a
        noi;
  \item spesso \textbf{le proprie doti le scopriamo dagli altri}. Il relatore
        non si era mai accorto di essere ``quello che resta calmo durante le
        emergenze'' (il lancio di \emph{Diablo~III}, quello di
        \emph{Draenor}, i problemi dal vivo alla Games Week) finché non
        gliel'hanno fatto notare tutti: per lui era la normalità. A chi pensa
        ``non sono capace, forse non è il mio percorso'', può darsi che
        semplicemente nessuno gli abbia ancora detto ``questa cosa la sai
        fare davvero bene, concentrati su questo''.
\end{itemize}

% ---------------------------------------------------------------------
\section{Domande sul lavoro nell'industria}
\label{sec:z1-industria}

Diverse domande degli studenti hanno riguardato le condizioni di lavoro nei
grandi studi internazionali. Le risposte sono raccolte qui per argomento.

\subsection{Presenza, remoto e trasferimenti}

\paragraph{Grandi studi: quasi solo in presenza.} Sia a Blizzard sia a Riot si
lavorava \textbf{solo in presenza}, senza smart working: sei anni a Parigi e
cinque a Dublino. Si parlava solo inglese: anche nei meeting tra soli
italiani si tendeva a comunicare e a scrivere in inglese, per correttezza
verso chi si fosse aggiunto e per lasciare una traccia comprensibile a tutti.
Il francese serviva per fare la spesa e sopravvivere, non in ufficio.
L'ambiente era estremamente eterogeneo (francesi, polacchi, spagnoli e
molti altri): durante il colloquio dal vivo a Dublino, mentre si preoccupava
di essersi vestito in modo adeguato, vide passare dietro la parete a vetro un
collega svedese con i dreadlock, a piedi scalzi, con una pianta sotto un
braccio e un aquilone sotto l'altro, e capì di essere ``a posto così''.
Un'esposizione a molte culture e modi di pensare che il relatore considera
bellissima.

Anche oggi, tranne alcuni studi particolari, la maggior parte del lavoro si
svolge di persona: amici a Paradox lavorano a Stoccolma, i pochi rimasti a
Blizzard si sono dovuti trasferire a Irvine. Per posizioni \emph{junior} o
\emph{entry level} i grandi studi raramente concedono il lavoro da casa; per
profili senior può esserci qualche eccezione (ad esempio nove mesi da remoto
e tre in sede).

\paragraph{Relocation.} Spesso tocca spostarsi, ma il relatore non ha mai
pagato un trasloco di tasca propria: queste aziende offrono la
\emph{relocation} (aiuto a trovare casa, ad aprire il conto in banca e
supporto pratico), perché è un investimento sulla persona.

\paragraph{Colloqui mentre si lavora altrove.} Per i colloqui in presenza
prendeva ferie, senza dire dove andava. Il consiglio generale: per fare
colloqui \textbf{prendete ferie o permessi}, non datevi malati, ``che non è un
bel modo di comportarsi''.

\paragraph{Remoto ``doppio''.} Un collega lavora per l'azienda coreana che
gestisce \emph{PUBG}, con ufficio ad Amsterdam, ma lavora da Milano: un
``doppio remoto''. Si trova, ma in genere un ambiente anglofono con sedi in
Europa chiederà di trasferirsi in Europa.

\paragraph{Famiglia e trasferimenti.} Il relatore ha avuto la fortuna di
sposare una persona che fa lo stesso lavoro (conosciuta lavorando su
\emph{Warhammer Online}, poi con lui a Blizzard e a Riot, nella
localizzazione e nel doppiaggio, ad esempio di \emph{World of Warcraft} in
italiano): avendo competenze rare, spesso venivano assunti entrambi, al
massimo a pochi mesi di distanza. Per i developer la soluzione frequente è
che la famiglia segua la persona assunta, perché lo stipendio basta per
tutti. Ma non sempre c'è equilibrio, e a volte uno dei due deve sacrificarsi:
sua moglie ha lasciato Blizzard per seguirlo a Dublino (per poi essere
assunta da Riot); lui ha lasciato Riot per tornare in Italia quando lei ha
trovato lavoro qui. Ha dovuto rinunciare più volte a lavorare in Islanda per
CCP (\emph{EVE Online}), perché non era una scelta solo sua. Si fanno
compromessi, ma si riesce a far quadrare.

\paragraph{Ci si sposta per ogni progetto?} No: si sono spostati
esclusivamente per cambiare azienda. Cambiando progetto nella stessa azienda
è raro, quasi mai (``non l'ho mai visto succedere''), che venga chiesto di
trasferirsi. In Francia stavano benissimo (vivevano fuori Parigi; ``i
problemi con i francesi sono in realtà problemi con i parigini''), parlavano
francese ed erano vicini alla doppia cittadinanza: se ne sono andati perché
non vedevano un futuro a Blizzard, per il modo in cui l'azienda trattava i
colleghi.

\paragraph{Il mercato oggi.} Non esiste un \emph{silver bullet} per trovare
lavoro domani dove si vuole. Con tutti i \emph{layoff} recenti ci sono
migliaia di esperti del settore che cercano lavoro: serve pazienza, ma è
fattibile.

\paragraph{Studi indie e remoto.} Negli studi più piccoli il remoto è più
diffuso: a Stormind circa 60 persone su 95 lavorano da remoto, ma lo studio
si è strutturato fin dall'inizio per farlo. Serve consapevolezza. Bending
Spoons (che sviluppa app per mobile, non giochi, ed è l'azienda dell'app
Immuni), pur essendo grandissima e con sede a Milano, lavora da remoto e
lascia l'ufficio a chi vuole: lo smart working è posizionato quasi come un
\emph{benefit}, ma permette anche di avere un ufficio da 300 persone invece
che da 700. Le sedi fisiche costano decine di milioni all'anno (Blizzard ha
chiuso le sedi di Amsterdam e di Parigi). Tuttavia, per lavori molto
verticali e specifici gli studi preferiscono avere le persone dal vivo; i
ruoli più spesso remoti sono quelli che non fanno sviluppo (customer service,
PR\dots).

Anche gli artisti lavorano spesso in presenza: l'accademia può dare solo una
\emph{simulazione} di quanto stretto sarà il contatto tra programmazione,
design e art. ``Sarà ancora di più.''

\paragraph{La comunicazione da remoto.} È più difficile. Pur essendo
estremamente introverso, il relatore \textbf{obbliga tutti a usare la webcam}
nei meeting in cui conta: non in un aggiornamento con novanta persone
collegate, ma certamente in un colloquio a due sulla performance dell'ultimo
anno. La componente fisica è una parte enorme della
comunicazione\footnote{In lezione si parla di ``quasi l'ottanta per cento'';
il dato va preso come indicativo, poiché le percentuali di questo tipo
circolano in forme diverse e sono spesso citate fuori contesto.}: anche se si
è timidi, accendere la webcam è l'approssimazione più vicina all'essere dal
vivo. Negli ambienti dove si lavorava senza webcam la situazione era
``esponenzialmente peggiore''.

\subsection{Differenze culturali nel team}

\paragraph{Gestire sensibilità diverse.} A Riot i team erano divisi per
\emph{lingua}, non per nazione: il team di lingua spagnola includeva anche un
messicano, quello di lingua russa includeva russi e ucraini, proprio nel
periodo della guerra in Crimea. Tra i colleghi non c'era alcun problema (erano
amici), ma ci si è posti la domanda: \emph{come comunichiamo con una community
che sarà divisa su questo tema?} Nei giochi online la divisione era evidente:
in \emph{Rocket League} c'era chi, con un compagno di squadra con la bandiera
avversaria, si faceva otto autogol pur di farlo perdere. Allo stesso modo,
con molti americani mandati a Dublino a coprire posizioni scoperte, c'erano le
elezioni (Clinton--Trump, Biden--Trump) e varie guerre: serviva sempre una
certa sensibilità.

Imparare a navigare queste situazioni lo ha reso un professionista migliore:
gli ha insegnato a fermarsi e chiedersi ``io la percepisco così; gli altri
dieci, con background, storie e lingue diverse, come la penseranno?'', a
ragionare più \emph{long term} e più \emph{big picture}, con più empatia.

\paragraph{Stili di comunicazione.} Un partecipante che ha lavorato nel Regno
Unito chiede dello stile comunicativo italiano, a volte frainteso. Uno dei più
grandi complimenti ricevuti dal relatore è stato ``non sembri italiano''
(``gesticoli poco per essere italiano'', ``parli bene inglese per essere
italiano''). Ha imparato ad \textbf{adattare la comunicazione alla persona}:
con spagnoli e francesi la battuta sulla nazione funziona, perché c'è un
``odio--amore'' condiviso; con russi e giapponesi no. Doveva adattare anche
l'umorismo. Un collega metà russo e metà messicano diceva di aver preso ``il
lato peggiore di entrambe le nazionalità'', ed è stato un ottimo maestro per
imparare a leggere culture diverse.

Non bisogna generalizzare, ma gli stereotipi possono essere un punto di
partenza ragionevole, \emph{sapendo che troverete sempre le eccezioni}. E
vanno letti correttamente: gli americani tendono a essere più estroversi e ad
avere sempre qualcosa da dire non perché ``sono fatti così'', ma perché
vengono educati così (all'università tra i corsi c'è il \emph{debate}, saper
argomentare in pubblico); un collega giapponese può essere più introverso per
ragioni culturali. Nel dubbio si parte da lì, poi si conosce la persona.

\subsection{Negoziare, discriminazione e stipendi}

\paragraph{La famiglia come leva.} Avere un partner riduce le opzioni, ma le
aziende lo sanno: uno spostamento è per due. Si può usare come leva per
negoziare (un compenso diverso, un supporto per il trasloco): se un'azienda ha
deciso di volere quella persona, 2, 3, 4 o 10.000 euro in più all'anno, per
un'azienda che fattura centinaia di milioni o miliardi, non cambiano nulla.

\paragraph{Essere sposati è un motivo per non essere assunti?} No. In alcuni
paesi è proprio illegale discriminare su questo: in Francia era illegale
chiedere CV con data di nascita, foto, sesso o religione, e chiedere se si
vogliono figli. Purtroppo, dai racconti dei suoi studenti di master (in
settori come \emph{beauty and wellness} e \emph{food and wine}), in piccole
realtà italiane queste domande (``hai 25 anni, hai intenzione di
sposarti?'') vengono ancora fatte, e servono a filtrare i candidati. Nelle
grandi aziende ci sono sistemi per evitarlo; anzi, molte danno \emph{più}
benefit a chi ha famiglia (``porta pure la famiglia, è tutto pagato''),
perché le persone competenti hanno spesso una certa età e quindi,
ragionevolmente, una famiglia.

\paragraph{Gli stipendi.} Dipendono molto da dove si vive. Lo stipendio a Riot
non era paragonabile a quello di Blizzard perché in Francia era tassato al
52\% e in Irlanda al 17\%\footnote{Aliquote come riportate in lezione,
riferite alla situazione personale del relatore.}: nessuno stipendio
competitivo può compensare il 35\% di differenza di tasse. Negli Stati Uniti,
con pochissima previdenza pubblica, gli stipendi sono molto più alti perché
poi ci si paga tutto (``se hai un raffreddore paghi tu''). Bisogna valutare il
potere d'acquisto di quei soldi nel paese in cui si andrà a vivere.

\subsection{Freelance e \emph{work for hire}}

All'inizio della sessione pomeridiana uno studente che lavora nell'audio
chiede quali siano le condizioni per chi lavora come \emph{freelance}.
\begin{itemize}
  \item Dipende molto dal progetto: spesso il contratto è a progetto, per
        portare a termine una parte o un intero progetto.
  \item Non tutte le aziende usano freelance: Riot e Blizzard lo facevano
        pochissimo e preferivano assumere (magari trattando poi la persona
        ``come un freelance'', lasciandola a casa a fine progetto).
  \item Spesso i lavori offerti come freelance sono in realtà
        \textbf{sostitutivi di un full time} (``freelance, ma sei ore al
        giorno'').
  \item Si può fare, ma bisogna essere molto chiari nel definire l'impegno
        (quante ore, come vengono tracciate): non tutti specificano cosa
        significa davvero ``freelance''. Lavorando così per due o tre aziende
        si può arrivare a un tempo pieno.
  \item È più comune il \emph{work for hire}: aziende che forniscono personale
        ad altre aziende; si è assunti da X ma si lavora ``embeddati'' presso
        Y (una sorta di subappalto). Anche lo studio del relatore lo fa.
  \item Il puro freelance di quattro mesi esiste, ma di solito per quei
        quattro mesi diventa un lavoro a tempo pieno; è raro ``quattro mesi,
        due ore al giorno''. Nelle grandi aziende è improbabile lavorare nello
        stesso mese per Ubisoft, Riot e Blizzard tre ore ciascuna; è più
        realistico lavorare per più startup indie italiane, che non possono
        permettersi un'assunzione ma possono pagare due ore al giorno.
  \item Molto passa dai contatti: trovare le aziende giuste e rendere esplicite
        le condizioni.
\end{itemize}

% ---------------------------------------------------------------------
\section{Il circolo virtuoso: feedback, miglioramento, fiducia}
\label{sec:z1-triangolo}

\subsection{Il triangolo}

L'obiettivo a cui il workshop vuole arrivare è un \textbf{triangolo che si
autoalimenta}:
\begin{enumerate}
  \item si dà un \textbf{feedback} con l'obiettivo di migliorare le
        performance (inclusi elementi caratteriali);
  \item la persona \textbf{migliora} e vede di averne beneficiato;
  \item questo genera \textbf{fiducia}, che rende più facile dare e ricevere
        il feedback successivo.
\end{enumerate}
Togliendo uno qualsiasi dei tre elementi, il circolo non è altrettanto forte.

In qualsiasi fase si parta, il rispetto per i colleghi è inizialmente dettato
dall'\emph{educazione}, non dalla fiducia. La fiducia arriva quando i colleghi
cominciano a dire cose giuste e sensate e a \textbf{mantenere la parola}.
Dare e ricevere feedback è il modo in cui, in un certo senso,
``inganniamo il cervello'' per creare una \emph{case history} di fiducia: si
ha motivo di fidarsi di più, di ascoltare di più, di essere più ricettivi. E
si diventa sempre più bravi ad accettare il feedback e ad agire di
conseguenza.

\subsection{Accettare un feedback non vuol dire fare come dice l'altro}

Dieci anni fa il relatore non era ostile ai feedback, ma molto
\emph{difensivo}: ``sì, sì, va bene, ora ci guardo''; oppure, a chi gli diceva
``avresti dovuto farla così'', rispondeva spiegando perché l'aveva fatta in
quel modo. Oggi ha imparato a dire ``ok, grazie, ora ci ragiono'': non deve
difendere la propria posizione da chi viene solo a dire ``hai pensato a
questa cosa?''. Prima lo viveva come un attacco (``mi viene a dire che non so
fare il mio lavoro''), ora come un ``ottimo punto di vista, lo prendo in
considerazione''.

Con una postilla fondamentale: \textbf{accettare il feedback non significa
fare come dice l'altro}. ``Mi hai detto che preferisci questo mob verde. Ho
pensato anche a farlo verde? Sì. Penso comunque che giallo sia meglio? Sì, lo
faccio giallo.'' Il relatore accetta \emph{sempre} il feedback, valido o no:
garantisce di tenerlo a mente quando lavora. Che poi faccia come gli viene
detto è un altro discorso.

\subsection{Fiducia e scadenze}

La fiducia funziona allo stesso modo con gli impegni. Se si chiede a qualcuno
una cosa per venerdì e la consegna venerdì, la volta dopo gliela si chiederà
per venerdì. Se invece arriva il giovedì della settimana dopo, la volta dopo
gli si darà una \textbf{deadline anticipata}, come con chi arriva sempre in
ritardo agli appuntamenti (``se ti dico alle 7 arrivi alle 7 e mezza; la
prossima volta ti dico alle 6 e mezza'').

Al lavoro questo è dannosissimo, perché crea \textbf{deadline fittizie} che
non sono mai quelle vere, e nascono discussioni del tipo ``perché devo farlo
in due giorni se si fa in cinque?'' (``sì, se tu lo facessi in cinque''). La
conversazione giusta è un'altra: \emph{``che problemi hai avuto nel
consegnare venerdì?'' ``Questo, questo e questo.'' ``Cosa posso fare per
evitare che succeda?''}

Il relatore non vuole mai sentirsi dire ``non sono riuscito a farlo''.
Preferisce di gran lunga sentirsi dire, tre giorni prima, \textbf{``potrei non
riuscire a farlo, per questi motivi''}. Chi ha la visione d'insieme può dare
più budget, più persone, spostare la deadline, cambiare lo scopo del progetto:
può fare qualsiasi cosa per aiutare. Ma se lo si scopre il venerdì, cinque
minuti prima della deadline, non può fare più niente. Il messaggio è:
\textbf{``aiutami a poterti aiutare''}. Non penserà mai male di chi dice ``non
ce la faccio per questi motivi'', perché non sta dicendo ``non ho voglia'': sta
indicando i problemi, e lavorare su quei problemi è compito del manager. Non
può invece fidarsi di chi fino al giorno prima dice ``tranquillo, ce la
faccio'' e il giorno dopo ``non ce l'ho fatta''.

Questo è fondamentale per chi lavora in programmazione, design o art, perché
ci saranno sempre deadline: bisogna sempre \textbf{lavorare sulle aspettative},
dire al proprio manager (o al proprio report) cosa è lecito aspettarsi. Non
bisogna fare come i pinguini di \emph{Madagascar} (``sorridete e salutate, va
tutto bene''): se c'è un problema, va detto, altrimenti si peggiora solo la
propria situazione.

Da italiani si tende a non voler mai alzare la mano per chiedere aiuto
(``no, no, ce la facciamo, lo faccio io''). A volte non ce la si può fare, non
perché non si è capaci, ma perché vari fattori lo impediscono. Accettarlo è
importante: \textbf{è molto meglio avvisare e chiedere una mano che dire
``non ce l'ho fatta''}.

\subsection{Domande sul dare e ricevere feedback}

\paragraph{``Ne parliamo a fare, se poi fai come vuoi?''} Uno studente (Cristian)
riceve feedback che, dopo averli analizzati, non ritiene congrui; quando non
li segue gli viene detto ``ne parliamo a fare?''. C'è un problema di
\textbf{aspettative}: l'idea che accettare un feedback voglia dire fare come
dice l'altro. Accettarlo significa prendere in considerazione ciò che è stato
detto. Chi decide diversamente ha però un \textbf{onere}: argomentare perché
si fa comunque a modo proprio (``mi hai detto verde, la facciamo gialla
perché\dots''). Poi bisogna solo aspettare i risultati per vedere chi aveva
ragione.

La decisione spetta sempre \textbf{a chi ha la responsabilità di quella
decisione}. Se il game director deve scegliere tra A e B, anche se il
feedback dice B, la scelta è sua perché la responsabilità è sua. Non sarebbe
corretto essere ritenuti responsabili di una scelta delegata a chi dà
feedback. La risposta può essere semplicemente: ``l'ho tenuto in
considerazione, ma visto che ne sono responsabile io e le conseguenze sono
mie, ho scelto di fare X''.

\paragraph{E se l'altro fa ripicca?} Un partecipante chiede di chi, vedendo
il proprio feedback non accolto, si vendica ignorando a sua volta feedback
buoni. Per il relatore è un problema caratteriale: chi si comporta così ``è
già sulla strada per uscire''. Non c'è spazio per l'ego in team di questo
tipo. Il fenomeno nasce dall'idea che \emph{il feedback negativo faccia fare
brutta figura}, che il fatto che si potesse fare meglio tolga qualcosa alle
proprie capacità.

Non è così: \textbf{il feedback non è un giudizio}. Si può passare un anno a
Riot ricevendo solo feedback negativi e avere una valutazione annuale
stellare, perché ricevere molti feedback negativi significa essere una persona
approcciabile, che agisce sul feedback e invita gli altri a darlo. Il relatore
ha visto valutazioni molto peggiori per chi riceveva solo feedback positivi,
perché la gente aveva paura di dargliene di negativi: ma solo ``pacche sulle
spalle'' non fanno migliorare. Il problema va affrontato parlando con la
persona; a volte le conversazioni iniziano con \textbf{``ti do un feedback sul
feedback''}, cioè sul modo in cui l'altro ha dato il suo.

Esempi dalla sua esperienza:
\begin{itemize}
  \item lui chiede spesso di \textbf{non ricevere feedback senza esempi
        tangibili}, perché fa fatica ad astrarre: ``sei stato maleducato'' non
        significa nulla, ``ieri al meeting mi hai risposto male su questa
        cosa'' sì;
  \item ad altri serve l'opposto: ha ricevuto il feedback ``quando devi dirmi
        che ho sbagliato, puoi girarci un po' intorno e darmi del contesto?'',
        perché il suo stile diretto (``questa cosa non ha funzionato, dovremmo
        fare così'') feriva. Ha adattato il suo modo di dare feedback: con
        alcuni è diretto, con altri gira intorno. L'importante è saperlo.
\end{itemize}

\subsection{Non portare la soluzione}

Un consiglio centrale: \textbf{quando si dà un feedback, non si dà una
soluzione}. Non si dice ``io farei così'', ma ``secondo me questa cosa si
potrebbe fare diversamente, per questo e questo motivo: come potremmo farla
diversamente?''. È come chi smette di fumare perché vuole e chi smette perché
glielo dicono moglie e figli: se non si è convinti, spesso non si smette.

Quando il relatore diceva ``c'è questo problema, io farei così'', il meeting
successivo iniziava inevitabilmente con ``ci ho provato come hai detto tu, con
me non ha funzionato''. Sempre. Perché la soluzione era sua. La
responsabilità della soluzione deve essere invece \textbf{congiunta}: ogni
feedback deve creare un piccolo team che ha come obiettivo risolvere quel
problema. Questo toglie l'antagonismo (non ``ti dico che non sai fare il tuo
lavoro'', ma ``ti voglio aiutare a farlo meglio, lavoriamoci insieme'') e
sposta il focus: il problema non sei tu, è ciò che sta accadendo.

\paragraph{E se sono i piani alti a tenere all'oscuro dei problemi?} È una
questione di cultura aziendale. Il relatore fa sempre presente al proprio capo
il concetto di \textbf{service level agreement} (che verrà ripreso nella
seconda giornata): ``se mi dai questi input, ti garantisco questo output in
questi tempi''. Se il capo chiede delle cose senza dire budget, tempi e
perché, non può aspettarsi lo stesso risultato. È una comunicazione a due
vie: si dice al manager cosa serve per fare il lavoro, e lui deve fornirlo o
accettare che il risultato non sarà quello voluto.

\subsection{Amici, parenti e parole proibite}

\paragraph{``Non si lavora con amici e parenti''?} Per sua definizione
personale il relatore ha pochissimi amici (``stanno sulle dita di una mano'')
e una soglia molto alta per considerare amico un collega. Allo stesso tempo,
in quindici anni e con centinaia di colleghi, ha avuto solo \textbf{due}
esperienze in cui non è riuscito a lavorare con qualcuno: non serve essere
amici per far funzionare le cose. Con gli amici, anzi, gli viene più facile:
se dice a un amico che una cosa ``fa schifo'', l'amico non pensa ``come si
permette'', ma ``se perfino Stefano, che è mio amico, me lo dice, allora è
proprio così''. Diventa un rafforzativo, grazie alla fiducia costruita prima.
Il problema del lavorare con amici e parenti nasce dal \emph{non voler} dire
all'amico che sta sbagliando; se invece la confidenza permette di dirlo, si
può saltare la ``confezione'' del messaggio e passare direttamente al
contenuto.

\paragraph{Le frasi peggiori.} Alla domanda se esistano parole peggiori di
``ottimo lavoro'', il relatore indica la sua classifica: ``ottimo lavoro''
(detto in modo vuoto)\footnote{Il passaggio della trascrizione è
frammentario: si intende verosimilmente il complimento generico e di
circostanza, che non dice nulla di utile.}; poi la frase che lo ``fa uscire di
testa'', \textbf{``abbiamo sempre fatto così''} (``e infatti avete chiuso'',
come Blockbuster); infine \textbf{``siamo una famiglia''}, riferito al posto di
lavoro. Un'azienda non è mai una famiglia, perché non ci si può aspettare lo
stesso impegno da chi la possiede e da chi ci lavora: non c'è un
corrispettivo proporzionale tra impegno e ritorno. Per quanto inclusiva,
flessibile e razionale, non sarà mai una famiglia.

\subsection{Ego, muri e il peggior modo di scegliere un leader}

\paragraph{Chi mette davanti l'ego e alza muri.} Esistono, ma di solito durano
pochissimo in azienda. Non si tratta di avere un ego forte, ma di \textbf{non
saper ricevere feedback}. Spesso sono persone bravissime nel loro lavoro che
poi si chiedono perché tutti intorno diventino più bravi di loro: perché gli
altri accettano feedback e hanno più fonti da cui attingere, mentre loro si
sono rintanati nella loro torre.

Molti manager tendono ad \textbf{assumere persone un filo peggiori di loro},
per non fare brutta figura. Il relatore assume solo persone più brave di lui,
``perché non è una gara'': ha lavorato con molti community manager più bravi
di lui come community manager, ma lui era più bravo a fare il manager. Non
viene valutato su quanto sa fare il community manager, ma su quanto il suo
team sa farlo. Un partecipante chiede se la paura di avere qualcuno di
migliore esista anche in ambito accademico. Assolutamente sì: al relatore è
capitato di non essere preso come docente in un corso di formazione sui
videogiochi perché il responsabile aveva percepito che ne sapesse più di
lui\footnote{La trascrizione dice ``che non sapevo più di lui'', ma il senso
del passaggio richiede il contrario.}. ``Ma perché te ne importa? Tu sei quello che ha portato chi
sa fare il suo lavoro: ben venga.'' O si ha questa cultura, o si ha paura per
il proprio posto.

\paragraph{Promuovere il migliore.} Quando bisogna scegliere un leader per un
team (programmazione, design, art), molto spesso si sceglie \textbf{il migliore
tecnico}: il miglior programmatore, il miglior designer. È spesso la scelta
sbagliata: è come prendere un idraulico per fare l'elettricista. Fare il
\emph{people manager} non è come fare il designer o il programmatore; si può
essere un ottimo programmatore e un pessimo manager, e viceversa. Le aziende
l'hanno capito e lo fanno meno di prima. A Blizzard c'era un grave problema di
\emph{middle management}: molte persone in ruoli importanti erano state
promosse cercando le skill sbagliate, erano frustrate dal non saper fare il
proprio lavoro e schiacciavano chi stava sotto per paura di fare brutta
figura.

Un partecipante cita una teoria simile: una persona viene promossa finché non
raggiunge \emph{il proprio livello di incompetenza}, dove si blocca\footnote{È
il \emph{principio di Peter} (Laurence J.\ Peter, 1969), citato in lezione
come ``la teoria di Peter''.}. Un altro cita Michael Scott di \emph{The
Office}, su cui il relatore tornerà (sezione~\ref{sec:z1-quando}).

% ---------------------------------------------------------------------
\section{Perché la cultura conta: retention, costi e fiducia}
\label{sec:z1-cultura}

\subsection{L'ambiente di lavoro e la retention}

Perché tutto questo è così importante? Perché l'ambiente di lavoro conta
moltissimo per la \textbf{retention}, cioè per trattenere le persone e farle
lavorare in modo efficiente. Se contasse solo lo stipendio, molti non
lascerebbero mai Riot o Blizzard; e più si va avanti con l'età, più si impara
a dare valore anche al resto.

Inoltre, secondo gli studi citati nelle slide, chi vive il proprio lavoro in
modo negativo è \textbf{meno efficiente nel problem solving}: l'energia
mentale va verso ``come faccio a non farmi dare la colpa'' o ``come faccio a
fare bella figura'', non verso la soluzione del problema. In ambienti negativi,
dove ci si deve preoccupare della propria ``sopravvivenza'', si lavora
peggio.

\subsection{Il costo di chi se ne va}

Ogni volta che qualcuno si dimette o viene licenziato, il costo di trovare
una nuova persona, fare i colloqui, formarla e farle attraversare forming,
storming, norming\dots\ è circa \textbf{un terzo del suo stipendio}. Per questo
le aziende tendono a fare più \emph{retention} che \emph{acquisition}, almeno
con le figure considerate utili alla produzione.

Allo stesso tempo si vedono sempre più spesso \emph{layoff} e \emph{hiring
wave}: si licenziano 3.000 persone e sei mesi dopo se ne riassumono 2.500 o
3.500, ma tutte più junior. Il relatore ne ha scritto in un articolo su
Eurogamer.it, in cui spiega anche alcuni ``trucchi del mestiere'': perché i
licenziamenti avvengono sempre in certi periodi e sempre ``a gruppi'' (prima
Bethesda, poi Blizzard, poi un'altra\dots).

\paragraph{L'offerta di Riot.} A Riot esisteva una pratica (di cui il relatore
non ricorda il nome) per cui, firmato il contratto, si aveva \textbf{un anno
di tempo per dire ``non mi piace, non me la sento, me ne vado''} e ricevere
subito come buonuscita un anno di stipendio. Avevano stimato che fosse molto
meglio pagare subito un anno a qualcuno e lasciarlo andare che tenerlo due,
tre, quattro anni come \emph{underperformer}: la responsabilità della scelta
di assumere era dell'azienda\footnote{La pratica, nota come ``The Offer'', è
descritta in fonti pubbliche con condizioni diverse (un periodo più breve e
una somma pari a una frazione dello stipendio annuo). Si riportano le
condizioni come raccontate in lezione; andrebbero verificate.}. L'hanno usata
pochissimi, ma il solo fatto che esistesse faceva lavorare meglio anche lui da
manager, perché sapeva che chi non si trovava bene poteva andarsene con un
anno di tempo per trovare un altro lavoro.

\paragraph{Il ricircolo.} C'è un forte ricircolo nell'industria: molti ex
colleghi passano da Ubisoft a Paradox, da Paradox a Funcom, da Funcom a
Blizzard. Il relatore ha scelto di uscirne lavorando solo per aziende
\textbf{indipendenti}: Stormind non è un indie di cinque persone, ma è
indipendente, e non deve accontentare azionisti. Quando vede un'azienda con
utili miliardari licenziare migliaia di persone si chiede ``dov'è la linea''
(\emph{when enough is enough}): con azionisti che vogliono un ritorno, se
quello è l'obiettivo stanno facendo la cosa giusta, ma lui non condivide
l'obiettivo\footnote{In lezione si parla di ``6,5 miliardi di utile'' e di
``3.000 persone licenziate''; le cifre sono riportate come dette e non sono
state verificate.}.

\paragraph{L'open space vuoto.} A Riot, durante un'ondata di licenziamenti,
del suo team di trenta persone è rimasto solo lui: ``teniamo te come
coordinatore, poi ricostruirai il team''. Ha rifiutato di ricostruire un team
cancellato ``con un colpo di spugna'' e ha deciso che non avrebbe più lavorato
per una grande azienda. Nessuno aveva pensato di spostarlo, e ha lavorato per
tre mesi in un open space di cento metri quadri con la sua scrivania e altre
ventinove vuote. Su una lavagna aveva scritto i nomi dei suoi ragazzi e li
cancellava man mano che, anche grazie ai suoi \emph{referral}, trovavano
lavoro. Cancellato l'ultimo nome, si è dimesso: per lui il lavoro era finito.

\paragraph{Team che si dissolvono.} Uno studente cita un articolo su Larian
Studios, il cui management è contrario a ricostruire ogni volta il team a fine
progetto. È lo stesso problema: si pensa ai ``giochi Blizzard'' o ai ``giochi
BioWare'', ma le persone che li hanno fatti non ci sono più; è rimasto solo il
brand, e non ha senso avere le stesse aspettative da un gioco fatto da persone
diverse. E anche chi sopravvive ai layoff vive pensando che il prossimo
potrebbe essere lui: alcune persone, licenziate nell'ultima ondata a
Blizzard, gli hanno detto ``erano due anni che lo sapevo''. Vivevano a Cork,
``un paesotto irlandese dove o vai per lavorare a Blizzard o non ci vai'',
avevano comprato casa: non ci si può aspettare chissà che morale.

\subsection{La cultura è forte quando si lavora per gli altri}

\begin{quote}
  \emph{La cultura è forte quando le persone lavorano per gli altri e a favore
  degli altri, e non per se stesse.}
\end{quote}

Se si dà feedback per far migliorare l'altro e non per imporre il proprio
modo, se si aiuta per il bene del progetto e non per fare bella figura, la
cultura si rafforza. È dai \textbf{piccoli gesti} che emerge la cultura
aziendale: ad esempio, dal non dover chiedere ``perché mi stai aiutando?''.

\paragraph{L'incognita.} \textbf{L'incognita è spesso dannosa quanto le
cattive intenzioni.} In un esperimento (``mi pare a Stoccolma, cinque o sei
anni fa'') si distribuivano per strada banconote da 500 euro, e la maggior
parte delle persone le rifiutava chiedendosi ``dov'è la fregatura?''
\footnote{L'esperimento è riportato in modo vago in lezione e non è stato
possibile identificarlo con certezza.}. Lo stesso accadeva al relatore,
arrivato a PG e a Stormind da Riot: a ``vedo che sei in difficoltà, posso
aiutarti?'' la risposta era sempre ``perché?''. Alla fine rispondeva: ``perché
quando sarò in difficoltà io vorrei che tu mi aiutassi''.

Quando si comincia a lavorare per gli altri e si smette di pensare ``così avrò
la promozione, così avrò l'aumento'', promozioni, aumenti e rispetto arrivano
da soli. A Blizzard aveva fallito tre volte i colloqui per essere promosso a
manager (inteso come coordinatore del team); quando ha deciso di smettere di
pensarci e di fare solo il bene del proprio lavoro, la volta dopo è stato
promosso. Funziona così: chi aiuta i colleghi e fa andare bene il progetto
viene notato (``questa persona fa migliorare cinque persone, non solo se
stessa: vale di più'').

\subsection{Il lavoro è valutato come opera d'insieme}

Il relatore riconosce che buona parte del suo successo è venuto dal fatto che
il suo team lavorava bene, non dalla sua bravura personale. Anche i progetti
degli studenti, presentati a Diego, a Enrico, ad Andrea, verranno valutati
\textbf{come opera d'insieme}. Nessuno dice ``il \emph{combat system} di questo
gioco l'ha fatto X, che bravo'': si dice ``BioWare ha fatto un bel gioco'',
``Blizzard ha fatto un bel gioco''.

Lo ha replicato anche lato HR. In molte aziende ci sono \textbf{bonus
personali} legati a obiettivi (se il gioco vende X, se la feature è pronta in
tempo). Il relatore li ha dimezzati: metà restano personali, metà diventano
\textbf{bonus di team}, che tutti ricevono se il team raggiunge gli obiettivi.
Un incentivo concreto a collaborare: non solo si dice che il lavoro di gruppo è
importante, lo si dimostra dandogli più valore economico. Aveva proposto il
100\% di bonus di team, ma gli è stato concesso il 50 e 50; ha funzionato
comunque, anche se sarebbe stato ancora più potente. Ovviamente, alla
domanda ``perché mi aiuti?'' si può anche rispondere ``perché voglio il
bonus'': un aiuto dato ``per il motivo sbagliato'' era un compromesso che era
disposto ad accettare.

\subsection{Il coraggio di chiedere aiuto}

Un partecipante osserva che aiutarsi dovrebbe essere la norma. Il modello
diffuso, invece, è ``devi essere forte e indipendente, non chiedere mai una
mano'', e vale in tutte le direzioni (dal maschilismo alla ``donna forte e
indipendente che non ha bisogno di nessuno''): \textbf{non ammettere di avere
bisogno di aiuto è sempre uno dei motivi per cui le aziende vanno male}, che si
tratti di uomini, donne o persone non binarie. Il problema è non avere la
sicurezza di poter alzare la mano.

Il relatore ha cercato di togliere ogni conseguenza dall'aver alzato la mano:
se si \emph{depenalizza} qualcosa, le persone non sono più disincentivate a
farla. Ha detto espressamente che non avrebbe mai valutato nessuno per il
numero di volte in cui chiedeva aiuto o per il numero di feedback negativi
ricevuti, e questo ha funzionato.

\subsection{Perché si dà un feedback}

\emph{Perché avete dato il vostro ultimo feedback? Che cosa volevate ottenere?}
Le risposte della classe:
\begin{itemize}
  \item per \textbf{riconoscere il valore} di un collega, anche verso il
        manager, perché qualcuno che non aveva visibilità vedesse quanto
        quella persona fosse efficace;
  \item per \textbf{spingere una persona a superarsi};
  \item per aiutare un collega a rendere il prodotto all'altezza di ciò che
        voleva rappresentare;
  \item per \textbf{genuinità}: c'era qualcosa che poteva essere fatto meglio.
\end{itemize}
Il relatore aggiunge che la forma cambia col tempo: all'inizio era ``devo
darti un feedback, hai dieci minuti?''; ora è un messaggio informale su
Discord, perché ha costruito un rapporto di fiducia e sa che non verrà letto
come ``l'HR manager mi scrive, sono nei guai?'', ma come ``Stefano è così, mi
dà feedback, resterà tra noi e lo risolveremo insieme''.

Un altro motivo è \textbf{risolvere incompatibilità caratteriali}: le
conversazioni più difficili, del tipo ``io e te non stiamo funzionando, cosa
possiamo fare diversamente?''. Sono importanti perché, se due responsabili non
vanno d'accordo, il team se ne accorge, comincia a raccontarsi storie (``si
odiano, ci saranno problemi'') e, per difendere il proprio manager, comincia a
giudicare male il lavoro dell'altro gruppo. È la dinamica da fermare:
\textbf{valutare gli altri solo per difendere la propria fazione}.

Dare e ricevere feedback sono \textbf{due muscoli da allenare} (il relatore è
sempre stato più bravo a darlo che a riceverlo), e si impara col tempo. È come
fare colloqui: fare colloqui è una skill a parte, che non indica quanto si è
bravi nel proprio lavoro. Essere scartati non significa non saper fare quel
lavoro: spesso significa aver fatto male il colloquio o che è stato preferito
qualcun altro.

\paragraph{Feedback pubblici e privati.} Uno studente racconta di aver dato un
feedback negativo a un superiore, senza un motivo particolare, per dimostrare
a sé e ai colleghi che non si viene licenziati per questo. Il relatore coglie
l'occasione per precisare che \textbf{i feedback negativi si danno
individualmente}: lui dà moltissimi feedback di cui il resto dell'azienda non
sa nulla. Reagirebbe molto male se, in un meeting, qualcuno alzasse la mano per
dire ``comunque questa cosa non l'hai fatta bene'': quello non è feedback, è
discussione, accettabile solo se è un brainstorming e se l'opinione è stata
chiesta. Alcune aziende hanno piattaforme pubbliche di feedback, ma spesso
servono solo ad applaudire; lo stesso accade nelle aziende che
\emph{obbligano} a lasciare un certo numero di feedback a settimana, che
diventano ``lui è bravo, lei è competente, anche questa settimana l'abbiamo
fatto''. Bisognerebbe chiedere a chi le adotta cosa pensa di ottenere (forse
trasparenza), ma difficilmente una persona ``attaccata'' pubblicamente reagisce
bene.

% ---------------------------------------------------------------------
\section{Among Us: la caccia al traditore}
\label{sec:z1-among-us}

Prima della pausa pranzo il relatore mostra tre film di fantascienza di
qualche decennio fa e chiede che cosa li accomuni\footnote{I titoli non sono
leggibili nella trascrizione: viene citato \emph{Event Horizon} da uno
studente, e il film ``in mezzo'' è indicato come ``l'originale'',
verosimilmente \emph{La cosa} (\emph{The Thing}).}. In ognuno c'è un gruppo di
persone con un obiettivo e la presenza, reale o percepita, di qualcuno che
rema contro: \textbf{l'energia del gruppo viene spesa per trovare il
responsabile}.

\paragraph{L'esperimento degli alieni.} Nei master in presenza il relatore
propone un esercizio. Distribuisce segretamente cartellini plastificati con
scritto ``essere umano'' o ``alieno''; le carte alieno contengono istruzioni di
comportamento con cui gli alieni, che cercano di fingersi umani, si
tradirebbero (dare sempre ragione a chi ha parlato prima, finire ogni frase
con una domanda, parlare in un certo modo). Lo scopo: in dieci minuti formare
gruppi, e \textbf{vince il gruppo più numeroso senza alieni}; basta un alieno
per perdere. Si formano sempre gruppi piccoli (due, cinque, sei persone),
perché la tendenza è ``mi fido al massimo di voi due''.

Ma vincono i gruppi più grandi, per un motivo semplice: il relatore
\textbf{non distribuisce nessuna carta alieno}. Mostra a tutti entrambe le
carte, ma consegna solo carte ``umano''. Vuole simulare ciò che accade in
ufficio: la tendenza a percepire sempre qualcuno che rema contro (i manager,
l'altro team, quel collega), quando in realtà non c'è nessuno, ``è tutto nella
nostra testa''. Quando chiede le motivazioni, emergono sospetti
\emph{fabbricati}: ``faceva sempre domande'', ``parlava solo all'infinito''
(``no, parla così e basta''). Ci si fabbrica da soli i motivi per non fidarsi
degli altri, ed è comunissimo.

\paragraph{Among Us.} L'autosabotaggio è, per tornare ad \emph{Among Us},
\textbf{concentrarsi sul cercare cadaveri e traditori invece di fare le
task}: non si vince perché non si sono completate le task, mentre si cercava il
colpevole. Succede tantissimo negli uffici, ed è il motivo per cui l'80\% dei
problemi che il relatore risolveva erano di comunicazione. Entrando nei team
sentiva immancabilmente: ``facciamo il possibile, ma questa lotta continua con
l'altro gruppo è estenuante''. ``Perché c'è questo problema?'' ``Non lo so,
c'era già quando sono arrivato.'' Un paradosso: c'è un problema, nessuno sa
dire quale sia, ma c'è sempre stato e quindi lo si accetta.

\paragraph{L'esperimento delle scimmie.} È l'esempio, ``finto o reale'', delle
scimmie in gabbia\footnote{È il noto racconto delle ``cinque scimmie'', spesso
presentato come esperimento scientifico ma probabilmente apocrifo (lo stesso
relatore lo definisce ``finto barra reale''); nella trascrizione il numero
delle scimmie varia.}. In una gabbia con alcune scimmie pende una banana
elettrificata: chi la tocca fa prendere la scossa a tutte. Le scimmie imparano
e, quando una si avvicina alla banana, le altre la picchiano. Poi si
sostituisce una scimmia alla volta con una nuova: la nuova va verso la banana,
le altre la picchiano, e smette. Si continua finché nella gabbia ci sono solo
scimmie nuove, poi si toglie la corrente. Nessuna scimmia tocca la banana:
nessuna ha mai preso la scossa, ma tutte hanno imparato che chi tocca la
banana prende le botte.

Nei team è lo stesso: il motivo per cui si ``odia'' l'altro team si
trasforma, ma un vero motivo non c'è. Ne derivano la sfiducia, il non alzare
mai la mano per chiedere aiuto, il non dire mai ``bel lavoro'' agli altri
perché non sia mai che qualcuno ti veda fare un complimento ``a quelli''.
Risolvere queste dinamiche è lungo e tedioso, e sempre pieno di ``non mi
ricordo'', ``io non c'ero'' e ``abbiamo sempre fatto così'', la cosa più
frustrante per chi cerca di risolverle.

% ---------------------------------------------------------------------
\section{Il principio del default to trust}
\label{sec:z1-default-trust}

\subsection{La definizione}

Il relatore indica questa come la slide più importante della presentazione.
\textbf{Default to trust} è uno dei valori aziendali di Riot Games, che lui ha
tradotto, adattato e ``esportato'':

\begin{quote}
  \emph{Approccio ogni situazione di conflitto partendo dal presupposto che una
  persona abbia agito in un certo modo perché, con le informazioni a sua
  disposizione in quel momento, ha pensato che fosse la soluzione migliore, e
  non perché sia un'incapace o volesse danneggiare il progetto
  intenzionalmente.}
\end{quote}

In altre parole: \textbf{di default, mi fido}.

\subsection{L'esempio del ritardo al daily stand-up}

A Riot il relatore aveva un \emph{daily stand-up} dal vivo alle 15 con il suo
team di venti persone. Una ragazza arrivava regolarmente alle 15:15, per una,
due, tre settimane. Alla quarta ci ha parlato. \emph{Come iniziereste la
conversazione?} Le proposte della classe:
\begin{itemize}
  \item sottolineare l'importanza di esserci dall'inizio, perché potrebbe
        perdersi qualcosa che le interessa, e per rispetto dei colleghi
        presenti;
  \item ``devi arrivare in orario'', spiegando perché è importante;
  \item \textbf{chiedere perché}.
\end{itemize}

Il relatore osserva, anche dalle sue consulenze fuori dai videogiochi, che
\textbf{più si sale con la seniority}, più la risposta è ``devi smettere di
arrivare in ritardo'', spesso con un \emph{or else} finale (``altrimenti
succede quest'altra cosa''). Più si scende con l'età, e quindi con meno
abitudini consolidate, più la prima domanda diventa ``dimmi perché arrivi in
ritardo''. Già spiegare perché è importante è ``miglia avanti'' rispetto al
semplice ``smetti''.

Partire dal perché gli ha permesso di scoprire la causa: il meeting precedente
della ragazza, in un'altra zona del campus a un quarto d'ora di distanza,
finiva sempre tardi, e lei doveva attraversare il campus di corsa. Non aveva
la confidenza per dire all'altro manager, che insisteva perché si restasse
fino alla fine, che doveva andare; sacrificava quindi il meeting del relatore,
sapendo che era il più comprensivo dei due.

Il relatore si è offerto di \textbf{mediare}: ha detto all'altro manager che
o anticipava l'inizio del suo meeting, così da finire comunque in tempo,
oppure lasciava uscire la ragazza, perché quel meeting gli serviva. C'è anche
una questione di \textbf{economia di scala}, spesso sottovalutata: se venti
persone aspettano qualcuno per dieci minuti, si sono sprecati dieci minuti per
venti. Per lo stesso motivo il workshop è iniziato alle 10:05: iniziare alle
10:10 avrebbe significato far aspettare cinquanta persone cinque minuti in
più. A Riot c'era una scrivania vuota con un pupazzo gonfiabile, \textbf{Mr.\
Late}: sommando tutti i minuti di ritardo a meeting, call e conversazioni, in
un anno veniva fuori praticamente \emph{una persona a tempo pieno}. Serviva a
ricordare che anche cinque minuti contano, e che bisogna rispettare gli altri
arrivando in orario o almeno avvisando.

\paragraph{Gli effetti.} La conversazione era avvenuta nel primissimo periodo a
Riot, quando il team non aveva ancora motivo di fidarsi di lui. Risolvendo il
problema ha eliminato anche tutto lo strato di micro-dinamiche spiacevoli che
un manager impara a vedere (due colleghi che si guardano con la risatina
quando lei entra in ritardo). La ragazza ha cominciato a entrare tranquilla,
con una certa leggerezza: non le faceva piacere arrivare in ritardo (``ci sono
anche persone che lo fanno apposta, e quello è un problema da affrontare
diversamente''). Soprattutto, con i colleghi non parlava di lui come del capo
che minacciava, ma come di chi le aveva risolto un problema: era diventato
un \textbf{alleato}. La settimana dopo altre tre persone sono andate da lui
con un problema da risolvere. Ripetuto cinque, sei, sette, dodici volte, in un
paio di mesi poteva dire quasi qualsiasi cosa nei meeting sapendo che
richieste e critiche sarebbero state percepite come ``me lo chiede perché vuole
darmi una mano'', non per fare bella figura o per presentare i successi del
team ai suoi superiori.

\subsection{Condividere i successi, schermare i fallimenti}

Un altro elemento fondamentale (ripreso nella seconda giornata): il manager
\textbf{condivide i successi e protegge il team dai fallimenti}. ``Se va male
qualcosa è colpa mia; se va bene qualcosa, voglio che ne parliate voi.'' Il
lavoro del manager non dovrebbe essere valutato su ciò che presenta del lavoro
del team: tutti sanno comunque chi è il manager. Tanto vale lasciar presentare
il team, che potrà dire con orgoglio di aver presentato a Mike
Morhaime\footnote{Cofondatore ed ex presidente di Blizzard; nella
trascrizione ``Mike Morayme''.} la propria campagna europea per
\emph{StarCraft~II}. Ci sono manager che dicono ``come mi prendo le colpe, mi
prendo i meriti'', ma i meriti li hanno già: è inutile ostentarli. Con i
fallimenti invece il manager è la \textbf{rete di sicurezza} del team. Il
relatore ne ha fatto una questione di principio: \textbf{non scaricare mai la
colpa su qualcun altro}.

\paragraph{Colpevole o causa.} Un esempio che il relatore attribuisce a Pirate
Software, creatore di contenuti ed ex dipendente Blizzard\footnote{Il relatore
aggiunge che si tratterebbe del figlio dello sviluppatore a cui si sarebbe
ispirato il personaggio del giocatore ``senza vita'' nell'episodio di
\emph{South Park} dedicato a \emph{World of Warcraft}. L'informazione è
riportata come detta in lezione e non è verificata.}. Contattando il servizio
clienti Blizzard per motivi personali, aveva scoperto che i colleghi gli
avevano scritto nei ticket informazioni confidenziali che non avrebbero dovuto
condividere in quel modo, nemmeno con un collega. Non ha segnalato le persone
da punire: ha detto ``c'è un problema, evidentemente le policy non sono
chiare''. Quando gli hanno chiesto i nomi, ha rifiutato: ``se vi dico chi è,
punite quelli, e gli altri cinquecento continueranno a farlo; non avremo
cambiato la causa alla fonte''.

\begin{quote}
  \emph{È inutile trovare il colpevole. Voglio trovare la causa. Perché il
  colpevole cambia.}
\end{quote}

Il colpevole a volte è vittima della situazione: per questo la definizione
dice \emph{``con le informazioni a sua disposizione in quel momento''}. Se il
relatore avesse rimproverato la ragazza in ritardo, non solo non avrebbe
risolto il problema, ma avrebbe creato un'\textbf{ingiustizia}: lei, che
attraversava il campus di corsa per una colpa non sua e di cui non poteva
parlare, sarebbe stata anche rimproverata.

Nessuno fa un brutto lavoro apposta: nessuno è contento quando sbaglia o
perde. Perché al lavoro diamo per scontato che gli altri sbaglino
intenzionalmente? Se qualcuno ha fatto una cosa male, è perché in quel momento
credeva fosse la scelta migliore. La domanda giusta non è ``perché hai fatto
questa cosa che è stata un fallimento?'', ma:

\begin{quote}
  \emph{Che cosa ti ha fatto pensare che questa fosse la scelta giusta?}
\end{quote}

Magari erano limiti di budget, di tempo, di scopo: si agisce su quelli, non
sulla persona. Se si cambia la persona ma le si danno gli stessi fattori, è
molto probabile che arrivi alle stesse conclusioni, e ci si troverà a dover
``disciplinare'' anche lei.

\paragraph{La reputazione costruita sui rumor.} Durante il giro di
presentazione in un ufficio (non in nessuna delle aziende citate), al relatore
è stato detto: ``questo è il dipartimento di design; se ti serve qualcosa
chiedi a lei, non a lui, perché lui consegna sempre in ritardo''. Oltre a
trovarlo poco professionale, una settimana dopo ha chiesto: ``cosa ti ha fatto
questa persona?'' ``A me niente, però ho sentito che\dots''. Si torna alle
scimmie: magari quella persona non consegna in ritardo da mesi, ma ormai lo
dicono tutti. È un circolo vizioso da fermare subito.

\paragraph{Niente escalation inutili.} Lo stesso vale per il feedback: quando
qualcuno va dal relatore dicendo ``ho un problema con questa persona'', lui
chiede: \textbf{``perché ne stai parlando a me e non a lei?''}. ``Perché sei il
mio capo'' non basta. Passare da lui significa dare per scontato che l'altra
persona non abbia interesse a risolvere, mentre magari basta parlarle. È la
strada più facile (``lo dico a Stefano, ci pensa Stefano''), e proprio per
questo bisogna mettere dei paletti: il manager interviene quando serve, ma le
cose che si possono gestire da soli vanno gestite da soli. (Quando invece il
confronto diretto fallisce, l'escalation è legittima: si veda la
sezione~\ref{sec:z1-qa-finale}.)

% ---------------------------------------------------------------------
\section{Traiettorie di carriera e valutazione delle performance}
\label{sec:z1-traiettorie}

\subsection{Quale traiettoria vorreste?}

Il relatore mostra un grafico con quattro curve colorate, contrassegnate
anche da lettere (A, B, C, D) per eventuali persone daltoniche (``a me è
capitato che qualcuno scoprisse di esserlo durante una lezione''). Le curve
rappresentano la performance lavorativa nel tempo. Due domande: \emph{quale
vorreste essere? E quale pensate realisticamente di essere?} Un partecipante
aspira a B e pensa di essere C; quasi tutti scelgono B.

Perché non A, che sale più rapidamente? Un partecipante risponde che una
crescita così repentina può essere più dannosa che utile, come una ``rush'' di
studio molto faticosa che poi non si riesce a sostenere. Il relatore propone
una lettura suggerita da un amico psicologo, che spiega la forma del grafico:
di B, C e D \textbf{vediamo dove finiscono}; di A no. Potrebbe continuare a
salire oltre la scala del grafico, o crollare sotto D. Non saperlo ci spinge a
preferire B. (A chi obietta che anche le altre proseguono oltre il bordo, il
relatore precisa: il bordo destro è ``oggi'', e di A non sappiamo come sia
arrivata fin lì.)

\subsection{La carriera non è una linea retta}

Tutti immaginano la carriera come una linea: si parte da X e, attraverso
stage e posizioni, si arriva a Y. In realtà carriera e performance avranno un
andamento irregolare, con alti e bassi. Come non si accetta che un team possa
temporaneamente \emph{underperformare}, tanto meno lo si accetta da se stessi.
Bisogna accettare che ci saranno momenti bassi.

Per questo il relatore \textbf{detesta le valutazioni annuali delle
performance}: prendono un punto casuale della traiettoria, che se si è
fortunati è il picco più alto e se si è sfortunati il più basso, e lo usano
per valutare un intero anno. Insegna ai manager a fare più punti di controllo
e a guardare altre cose:
\begin{enumerate}
  \item \textbf{quanto tempo passa tra un picco e l'altro}, cioè quanto
        \emph{cooldown} serve a una persona. È il motivo per cui, dopo la
        Games Week, tutti staccavano una o due settimane. Il relatore stesso,
        che si impegna moltissimo e tende al burnout, non potrebbe rifare lo
        stesso workshop il lunedì e il martedì successivi con la stessa
        qualità;
  \item \textbf{la differenza tra la miglior performance e la peggiore}: che
        cosa ha portato a lavorare così bene, e così male? Con chi si
        lavorava, su cosa ci si concentrava, quanto tempo e budget c'erano?
        L'obiettivo è capire come \emph{replicare} le condizioni che fanno
        lavorare bene, invece del solito ``hai lavorato male, è colpa tua''.
        Il relatore ha avuto molti manager per cui quando lavorava male era
        colpa sua e quando lavorava bene ``che bravo'', ma sapeva che raramente
        era l'unico fattore della propria performance. È la fallacia del ``se
        sai fare 100, perché non fai sempre 100?'', come nella scena di
        \emph{Schindler's List} sull'operaio che produce una vite in pochi
        secondi: non sempre un ritmo è sostenibile;
  \item \textbf{il fallimento di oggi confrontato con quello di sei mesi fa}.
        Il relatore non è contento della presentazione di oggi (l'avrebbe
        migliorata per altri due anni), ma rispetto a quelle di cinque anni fa
        vede un grande miglioramento di chiarezza: il suo peggior fallimento di
        oggi equivale al suo primo successo di inizio carriera.
\end{enumerate}
Quello che conta è la \textbf{traiettoria} dall'inizio dell'anno a oggi, non
i singoli punti. E se si riesce a capire insieme quando e perché si ha
successo e quando e perché si fallisce, si può lavorare per replicare solo le
situazioni che fanno lavorare bene.

\paragraph{Lo stesso vale per i giochi.} È il motivo principale per cui i
giochi delle grandi case vanno bene o male. Quando si riescono a replicare le
condizioni che hanno prodotto \emph{Warcraft~II}, esce \emph{Warcraft~III};
quando non ci si riesce, esce ``il nuovo \emph{Dragon Age}'': la gente non è
arrabbiata perché il gioco non è bello, ma perché non è come gli altri
\emph{Dragon Age}; non ne replica il \emph{feeling}, anche se potrebbe
essere un gioco favoloso. L'eccezione positiva è il passaggio da \emph{Dune} a
\emph{Dune~II}: il primo era di fatto un'avventura grafica, il secondo ha
generato il genere degli strategici in tempo reale. Gli sviluppatori avevano
capito che il successo non stava nel tipo di gioco, ma nell'ambientazione che
attirava le persone: un bel gioco in quell'ambientazione sarebbe piaciuto
comunque. Il relatore comprò \emph{Dune~II} aspettandosi un seguito
dell'avventura e, pur sorpreso, lo finì.

\subsection{Lying, hiding, faking}

Quando si è valutati su un singolo punto nel tempo, e da quel punto dipendono
bonus, stipendio e il rischio di licenziamento, si comincia a fare
\textbf{lying, hiding e faking}:
\begin{description}
  \item[Lying] dire che è sempre andato tutto bene, che si finirà sempre in
        orario, che non c'è nessun problema;
  \item[Hiding] nascondere: ``non era un mio task'', ``non c'ero'', ``non ho
        visto niente'';
  \item[Faking] fingere, anche in cose apparentemente innocue come ``sì, sì,
        ho capito tutto'' quando si ha solo paura di dire ``me lo rispieghi?
        Non ho capito niente''.
\end{description}

\subsection{``Non ho capito'' e la conferma delle istruzioni}

Dire ``non ho capito, ripetimelo'' è ancora visto come una figura da sciocco.
Il relatore non si arrabbierà mai se qualcuno gli chiede di rispiegare
un'intera slide: le cose ovvie per lui possono non esserlo per altri. Più si
conosce qualcosa, \textbf{più è difficile spiegarla}, perché si danno per
scontate molte cose. Ne parlava proprio con Diego, che gli diceva che fare
l'insegnante diventa sempre più difficile man mano che si diventa esperti. È
come spiegare \emph{EVE Online}, uno dei giochi più complessi in circolazione:
chi lo conosce bene dà per scontato che ``ovviamente non puoi fare questa
cosa'', mentre per un nuovo giocatore non è affatto ovvio. A scuola chi dice
``non ho capito'' viene visto come lo ``scemo''; nessuno si chiede ``non è che
l'ho spiegato male?''.

Il racconto più ``traumatico'' dei ragazzi che fanno stage nelle grandi
aziende è sempre lo stesso: ``Cinzia, devi fare questa cosa.'' Cinzia la fa.
``Non te l'avevo chiesta così.'' Il problema è che nessuno ha
\textbf{confermato che Cinzia avesse capito}. Il consiglio, valido sempre:
\begin{itemize}
  \item quando si \emph{ricevono} istruzioni, ripeterle: ``ho capito bene?
        Vuoi che faccia questa cosa in questo modo?'';
  \item quando si \emph{danno} istruzioni, chiedere a chi le riceve di
        ripeterle, non perché lo si ritenga incapace, ma per capire che cosa
        ha capito dalla propria spiegazione (``quindi vuoi un quadrato?''
        ``No, ho detto un triangolo\dots'' ``Però mi hai detto una forma con
        quattro angoli.'' ``Hai ragione: voglio un quadrato'').
\end{itemize}
Frasi come ``non te lo ripeto due volte'' o ``stammi dietro'' portano solo a
fare le cose come si pensa che vadano fatte e a sentirsi dire a fine progetto
``non la volevo così'': che si tratti di una feature, di un design o di un
modello 3D. (Un partecipante osserva che lo stesso accade insegnando le lingue:
chi insegna dà per scontato, ad esempio, che si sappia dell'esistenza del
genere neutro in tedesco.)

\paragraph{Come convincere il team a chiedere?} Il relatore mette sempre come
premessa che \textbf{l'obiettivo non è trovare la persona più intelligente
nella stanza, ma fare un buon lavoro}. Che Gianni capisca alla prima e un altro
alla quinta non gli interessa, se alla fine entrambi lavorano bene: ripetere
cinque volte è comunque meno lavoro che far rifare le cose sette volte. È un
investimento a lungo termine, e non dissimula alcun fastidio: davvero non gli
dà fastidio. Ma ripetere le stesse parole non aiuta: ``non ho sentito'' si
risolve ripetendo; ``non ho capito'' si risolve cambiando spiegazione. Meglio
ancora: ``dimmi che cosa hai capito tu, e vediamo cosa non torna''. Se
intravedesse nel team uno sguardo del tipo ``è la quinta volta che lo
chiede'', interverrebbe: spesso chi giudica è proprio chi dice ``ho capito'' e
poi sbaglia. Preferisce di gran lunga chi chiede cinque volte.

\subsection{Il feedback più utile che abbiate ricevuto}

\emph{Qual è il feedback più utile che abbiate mai ricevuto?} Il relatore ne ha
almeno dieci che gli hanno cambiato il modo di lavorare. Le risposte:
\begin{itemize}
  \item un semplice \textbf{``grazie''} per un lavoro fatto, più utile di un
        ``questa cosa fa schifo'';
  \item un feedback sulla propria \textbf{impulsività} nel rispondere: fermarsi
        un attimo e capire di cosa si sta parlando prima di intervenire;
  \item \textbf{essere se stessi}, anche al lavoro, e non scusarsi
        continuamente per ogni cosa.
\end{itemize}
I feedback che toccano di più sono quelli \textbf{veri}. Il ``grazie''
significa ``ho visto quello che hai fatto e l'ho apprezzato''. Quelli negativi
che si ricordano sono quelli che hanno fatto realizzare ``questa cosa che
pensavo di saper fare, non la so fare'', oppure ``la so fare, ma questa volta
non mi è venuta'' (non è che se una cosa una volta non riesce non la si sa più
fare: può capitare di sbagliare una porta). Arrivano sempre da un
\textbf{punto cieco}: qualcosa che non si sapeva fosse percepito dagli altri.
Conta più ciò che le persone ti dicono di ciò che pensi che pensino di te; e
ciò che il relatore pensa che gli altri pensino di lui è sempre terribile,
mentre il loro giudizio si rivela molto più positivo del previsto.

\paragraph{Il denominatore comune siete voi.} Si potrebbe pensare che questi
concetti non siano applicabili da chi arriva in un'azienda senza poterne
cambiare la cultura. Ma c'è sempre un denominatore comune a cui applicarli:
\emph{se stessi}. La cultura aziendale può non spingere a dare feedback, può
vedere male chi chiede aiuto; ma il relatore ha visto che quando lui lo fa, le
persone intorno tendono a farlo, e un gruppo che lo fa lavora meglio. Perfino
i CEO più ostici a un certo punto gli chiedevano: ``perché il tuo team ora
performa meglio degli altri?''. La differenza è il default to trust: una
pratica che col tempo fa operare il team \emph{come un team} e non come un
insieme di individui. Un videogioco è un progetto corale: c'è chi fa giochi da
solo, ma non si fa da soli \emph{Red Dead Redemption~2}; e se il gruppo non è
affiatato e allineato non esce \emph{Red Dead Redemption~2}, esce
\emph{Anthem} (``oggi la mia fissazione è \emph{Anthem}, domani sarà
un'altra'').

% ---------------------------------------------------------------------
\section{Come dare un feedback}
\label{sec:z1-dare}

\subsection{Concentrarsi sull'azione e sui suoi effetti}

Il primo principio: \textbf{il feedback non riguarda mai la persona}. Non si
dice ``tu hai fatto questa cosa che non va bene'', ma ``questa cosa non va
bene: come facciamo a non farla più?''. Che sia un comportamento, un'azione,
una decisione o una strategia, ci si concentra su quello, e si
\textbf{esplicita l'effetto} che ha. Nell'esempio del ritardo l'effetto era
``quindici minuti di attesa moltiplicati per venti persone''. In un progetto
può essere una catena:
\begin{quote}
  se non rispetti le deadline, l'animazione che viene dopo resta ferma; se
  loro sono in ritardo, la \emph{vertical slice} è in ritardo; se la vertical
  slice non è pronta per la Gamescom, non andiamo alla Gamescom e non troviamo
  un publisher; se non troviamo un publisher, chiudiamo.
\end{quote}
Più la conseguenza è grave, più dovrebbe sembrare piccolo il gesto di dire
``datemi una mano'': è il \emph{male minore}. Poi, a cose fatte, ci si dà
feedback: cosa ha funzionato, cosa no, cosa fare di più, di meno, cosa smettere
di fare.

\subsection{Il pericolo degli ambienti solo positivi}

Gli ambienti dove ci si fa sempre grandi complimenti e nessuno parla dei
problemi creano difficoltà vere alle aziende. Il problema non è mai la
mancanza di complimenti, ma un \textbf{eccesso di ottimismo}: si comincia a
pensare che le cose andranno come si spera, perché tutti intorno danno solo
feedback positivi.

A Blizzard i feedback determinavano gli \textbf{aumenti annuali}: ognuno
doveva indicare cinque persone che avrebbero dato feedback su di lui, e
ovviamente tutti indicavano gli amici, o colleghi di progetti andati bene. Un
sistema nato con le migliori intenzioni diventava uno \emph{scambio di voti}
(``se mi metti follow ti metto follow''), senza alcun miglioramento.

Al contrario, \textbf{le conversazioni più brutte sono state le più utili}:
affrontare i licenziamenti a Riot, o la cancellazione di \emph{Project Titan}
a Blizzard, il progetto che doveva essere il nuovo MMO dopo \emph{World of
Warcraft} e da cui poi è nato \emph{Overwatch}\footnote{Nella trascrizione si
legge anche ``doveva essere Warcraft 2'', probabilmente un lapsus o un errore
di trascrizione.}. Andando a vedere perché lo si cancellava, c'erano solo
motivazioni razionali e pragmatiche, di business, chiare a tutti, non un ``Mike
Morhaime ci odia''. Sono anche le conversazioni in cui ha lavorato su se
stesso: è molto \textbf{avverso ai conflitti}, e questo a volte gli ha
procurato conseguenze molto più gravi di quanto sarebbe costato un piccolo
disagio. Dieci anni fa faceva fatica ad accettare questo difetto; ha dovuto
subire le conseguenze di conflitti evitati e poi cresciuti.

\subsection{Che cosa fare}

Le linee guida della slide, che il relatore legge rapidamente perché saranno
disponibili nel materiale:
\begin{description}
  \item[Attenzione al linguaggio del corpo] dare un feedback a braccia
        incrociate o guardando altrove è diverso dal guardare l'altro negli
        occhi.
  \item[Sii specifico] con esempi tangibili.
  \item[Esplicita gli effetti] come visto sopra.
  \item[Discuti il prima possibile] non aspettare un mese per parlare di
        qualcosa successo ieri, altrimenti nessuno ricorderà più bene cosa sia
        accaduto.
  \item[Concentrati sull'azione e sul problema, non sulla persona] il nemico è
        l'azione, non la persona che si ha davanti. Non bisogna far cambiare
        l'altro: bisogna evitare che il problema si ripresenti.
\end{description}

\paragraph{Gli insulti sono rivolti al ruolo.} Un esempio solo apparentemente
scollegato. Il relatore ha lavorato su \emph{League of Legends}, ``non
esattamente noto per avere una community da orsetti del cuore'', in varie
\emph{task force} sulla tossicità. A chi smetteva di giocare veniva inviato un
sondaggio sul perché; il 33\% indicava la tossicità della community. Dopo varie
campagne di sensibilizzazione e cambiamenti si è scesi al 30\%: sembra solo
un 3\% in meno, ma valeva circa \textbf{120 milioni di revenue in più
all'anno}, perché quei giocatori restavano esposti a skin e contenuti e
continuavano a comprare. Anche qui, piccoli guadagni che si sommano.

Lavorando con degli psicologi, aveva notato che gli insulti \textbf{non sono
quasi mai rivolti alla persona, ma al ruolo}: ``il mio jungler fa schifo'',
``il mio DPS non è capace'', ``il nostro portiere è una pippa''. Chiamare
qualcuno per ruolo lo \emph{depersonalizza} e rende più facile insultarlo. Lo
stesso nell'esempio del giro dell'ufficio: non ``non chiedere a Giulio'', ma
``non chiedere a quel grafico, consegna sempre in ritardo''. Concentrarsi
sull'azione, senza attaccare la persona, fa sì che il feedback non sia (e non
venga percepito come) un attacco personale.

\paragraph{Considera il carattere.} Qui tornano utili i test di personalità.
Servono più agli studenti che al relatore (che sa già che usciranno 53 profili
diversi, senza un pattern): bisogna conoscere il proprio profilo e
\textbf{condividerlo con chi lavora con noi}, chiedendo anche agli altri di
fare lo stesso. In ufficio il relatore ha usato una versione più avanzata, a
pagamento, con un report di circa trenta pagine: il \textbf{profilo
DISC}\footnote{Nella trascrizione ``disk''. Si tratta del modello DISC
(Dominance, Influence, Steadiness, Conscientiousness).}. Ha mappato tutta
l'azienda e ha fatto stampare il profilo di ciascuno sulla sua scrivania, così
che tutti sapessero con chi stavano parlando. Capire come ragionano gli altri
(``è introverso, non gliene parlo davanti a tutti'') è una delle cose che gli
sono tornate più utili: \emph{adattare ciò che si vuole dire a chi si ha
davanti}.

\paragraph{Assicurati di non disturbare.} Sotto il profilo di ognuno c'era un
piccolo \textbf{semaforo}: verde, ``parlami pure, interrompimi''; giallo,
``interrompimi solo se è urgente''; rosso, ``mi sto concentrando, non mi
parlare''.

\paragraph{Lascia finire il lavoro.} Durante gli eventi dal vivo il relatore si
era reso conto di quanto fosse grave dare feedback mentre qualcuno sta ancora
lavorando. Se una cosa sta andando male, dirlo in quel momento non aiuta: tra
la pressione di finire e quella del feedback, peggiorano entrambe le cose, e
anche il rapporto. Meglio lasciarla finire, anche male, e parlarne dopo.

\subsection{Che cosa non fare}

\begin{description}
  \item[Non basarti su supposizioni] ``mi hanno detto che ieri in un meeting
        hai fatto\dots'': o c'eri e l'hai visto, o non esiste.
  \item[Non giudicare la persona.]
  \item[Non fornire feedback vago.]
  \item[Non accettare e non fornire feedback di terze parti] Se qualcuno vi
        dicesse ``mi hanno detto che hai fatto un brutto lavoro'', il primo
        pensiero non sarebbe ``cosa ho fatto male?'', ma \emph{``chi te l'ha
        detto? Chi è la vipera che ho in seno? Chi è l'impostore?''}. Siamo
        fatti così. Ogni conversazione che inizia con ``mi hanno detto che'' è
        inutile: l'altro non ascolterà, penserà solo a chi dare la caccia, e
        spesso penserà alla persona sbagliata, creando antipatie tra persone
        che non c'entravano nulla. Quando al relatore dicono ``mi hanno detto
        che lui ha fatto\dots'', risponde ``vai a parlarne con lui''; se gli
        dicono ``mi hanno detto che tu hai fatto\dots'', risponde ``mandami la
        persona che te l'ha detto, altrimenti non ne parliamo''. Non per
        cattiveria: se chiedesse ``che cosa ho detto di preciso a Marco?'' e la
        risposta fosse ``non lo so, non c'ero'', di cosa si starebbe parlando?
        \textbf{Parlate con la fonte}, altrimenti fate come se quel feedback non
        esistesse.
  \item[Non dare consigli non richiesti] come visto, non portare la soluzione:
        la soluzione deve venire da chi riceve il feedback.
  \item[``Non mi piace'' non è un feedback] ``Penso che questo design non
        funzioni perché, visto che il nostro obiettivo è\dots'' ha argomenti;
        ``non mi piace il combat system'' è un'opinione. Anche ``io l'avrei
        fatto diversamente'' non è un feedback.
  \item[Accettare non vuol dire fare come proponi] deve essere chiaro a chi dà
        il feedback: l'aspettativa è essere ascoltati e presi in
        considerazione, ma se la responsabilità è dell'altro, deciderà
        l'altro.
  \item[Non offenderti] se il feedback non viene seguito. Restarci male
        significa pensare di saperne più dell'altro, cioè \emph{venire meno al
        default to trust}. ``Decidi tu'' significa ``mi fido di te, perché sei
        lo specialista, e mi fido che mi ascolterai e farai la cosa migliore
        per tutti'', che potrebbe anche non essere ciò che ho detto io. Se
        voglio che l'altro faccia sempre come dico io, mi sto fidando solo di
        me stesso.
\end{description}

\subsection{Il primo feedback e il debriefing}

La cosa più difficile è dare il \textbf{primo} feedback: è un muscolo da
allenare. Chi non l'ha mai fatto (soprattutto con questa struttura) dovrebbe
provarci, anche su minuzie (``abbiamo concordato di usare WhatsApp nel gruppo:
non potremmo usare Telegram?''), anche solo per vedere come va. Altre
indicazioni:
\begin{itemize}
  \item feedback \textbf{solo uno a uno}, mai davanti ad altri, soprattutto se
        negativo;
  \item \textbf{prepararsi} prima;
  \item fare il \textbf{debriefing del feedback}: dopo, fermarsi a pensare a
        com'è andato. Il relatore, quando media un feedback tra due persone,
        si sente spesso rispondere ``benissimo'' da una e ``malissimo''
        dall'altra: la percezione è molto diversa. Non si fa problemi a tornare
        dalla stessa persona: ``ieri ti ho dato un feedback: com'è andata? Sei
        soddisfatto? Come poteva andare meglio?''. Con le persone con cui si ha
        meno a che fare, questo aiuta molto la volta successiva.
\end{itemize}

% ---------------------------------------------------------------------
\section{Come ricevere un feedback}
\label{sec:z1-ricevere}

Per il relatore è la parte più difficile. Si riassume in quattro passi:
\begin{enumerate}
  \item \textbf{Ascolta}: si parte sempre ascoltando.
  \item \textbf{Chiedi dettagli}, se serve, ma per capire meglio, non per
        difenderti.
  \item \textbf{Accetta, a prescindere}: non si sta decidendo chi ha ragione;
        si sta ascoltando un'opinione. Non c'è nulla da cui difendersi.
  \item \textbf{Agisci}: fatti un piano. Può essere banale (``mi hanno detto
        che interrompo sempre: cerco di non interrompere''), oppure ``devo
        migliorare questa cosa: come?''. Pensa anche a \emph{chi} potrebbe
        aiutarti e a \emph{che supporto} ti serve: il piano non deve essere
        per forza ``devo risolvere tutto da solo''.
\end{enumerate}
Una domanda che molti sottovalutano: \emph{``sono junior in un'azienda: sarò
mica il primo ad avere questo problema?''}. Probabilmente qualcun altro l'ha
già avuto: basta chiedere ``tu come hai fatto?''.

\subsection{Il feedback positivo}

Tutto questo vale per i feedback costruttivi negativi, a cui si pensa di
default; ma esistono anche i \textbf{feedback positivi}, e vanno usati. Quando
si riceve un ``grazie'' o un ``bravo'', la domanda è: \emph{come faccio a
ricreare le stesse condizioni?}

Il relatore ha lavorato anche in aziende dove il feedback era solo negativo.
Il rischio? Si immagini l'impiegato ideale: vuole sempre migliorare, è sempre
pronto a dire ``magari sbaglio io''. Se nessuno gli dice cosa sta facendo
\emph{bene}, spinto dal desiderio di migliorare rischia di cambiare cose che
erano già buone, e di farle peggio. Il rinforzo positivo \textbf{non è solo una
gentilezza}: è un atto pratico, che evita che le persone peggiorino ciò che
sanno già fare bene.

\subsection{Il feedback accidentale}

Esiste infine il \textbf{feedback accidentale}, molto frequente negli uffici:
\begin{itemize}
  \item il ``chiedi a lei, non a quel grafico'' del giro dell'ufficio è un
        feedback negativo su una persona, dato accidentalmente, né richiesto
        né voluto;
  \item a Riot li chiamavano \textbf{drive-by}, come nelle scene dei film in
        cui passa un'auto, si abbassa il finestrino, si spara e si riparte. In
        molte aziende con una ``cultura del feedback'' in cui tutti devono
        dire la loro, si riceve una valanga di commenti del tipo ``comunque il
        sito a me non piace, ciao''. Non si sa perché, non si sa cosa fare di
        diverso, non si sa che rischi ci siano: \emph{a che cosa è servita
        questa conversazione?} A niente.
\end{itemize}

% ---------------------------------------------------------------------
\section{Quando, quanto e con chi}
\label{sec:z1-quando}

\subsection{Frequentemente e recentemente}

\emph{Quando avete ricevuto il vostro ultimo feedback?} Per gli studenti,
qualche giorno fa (nel corso) o qualche mese fa (al lavoro, dopo una
consegna). Come aveva osservato Silvia, il feedback andrebbe dato
spesso.

Il relatore usa l'esempio di una nave da crociera (``la Costa Concordia'') che
deve andare da A a B. La rotta ideale (azzurra) è quasi impossibile: nessun
progetto va esattamente come previsto. C'è la rotta ``panoramica'' (verde), che
gira intorno all'isola d'Elba; ci sono poi rotte (arancione e rosa) che a un
certo punto si allontanano molto e poi correggono bruscamente: qualcuno si è
accorto che si stava sbagliando, ``ma non è quello che avevo chiesto'',
``non è dove dovevamo andare''. Sono i famosi \textbf{sanity check} o
\emph{checkpoint}.

L'ideale è avere \textbf{tanti piccoli check per fare tante micro-correzioni}:
correggere la rotta di un paio di gradi è molto più rapido che farlo dopo
essersi allontanati. Bastano trenta secondi di ``ehi, vado bene?''. Spesso non
lo si fa per paura di sembrare insicuri o poco autonomi, ma per il relatore è
l'opposto: quando qualcuno gli chiede ``vado bene così?'' e lui risponde
``sì'', si alleggerisce anche del proprio carico mentale. Un capo che assegna
un compito e lo vede solo alla fine vive sempre col dubbio (``avranno capito?
Lo staranno facendo giusto? Arriverà in tempo?''); con un check di dieci
secondi al giorno (``tutto ok, Stefano'') è tranquillo, e la persona va avanti
sapendo che quanto fatto finora va bene. Non serve fare ogni giorno il giro del
dipartimento: anche ogni uno o due giorni basta a evitare la rotta verde.

Per il feedback è lo stesso: \textbf{va dato subito}. Dire oggi ``ieri hai
fatto questa cosa'' è molto diverso da dire ``negli ultimi sei mesi hai fatto
questa cosa''. Il relatore si è sentito dire più volte, alla valutazione di
dicembre, ``da maggio hai fatto questa cosa che non andava''. La sua risposta:
``e aspetti dicembre per dirmelo? Mi hai lasciato sbagliare per sei mesi?''. È
una mancanza di rispetto, e la conversazione non può più essere costruttiva.
Per questo non ama le valutazioni annuali: i manager crescono abituandosi a
prendere appunti delle cose negative per tirarle fuori alla valutazione
(``a marzo hai fatto questa cosa'' ``sì, però poi ho smesso'' ``sì, però a
marzo l'hai fatta'').

In sintesi:
\begin{description}
  \item[Frequentemente] riguarda l'intervallo tra un check e l'altro.
  \item[Recentemente] significa non lasciar passare un anno prima di dare il
        feedback, altrimenti non serve più a niente.
\end{description}

\subsection{Quanto dura un feedback}

Un feedback dura \textbf{da trenta secondi a quindici minuti}; raramente di
più. Chi non ne ha mai dato uno dovrebbe sforzarsi di darne almeno uno, magari
positivo per cominciare. Ogni feedback dato o ricevuto è un
\textbf{deposito nel proprio ``fondo fiduciario''} (\emph{trust fund}), che
costruisce la reputazione futura. Il relatore si è costruito la reputazione di
chi dà feedback il più possibile oggettivi e analitici: giusti o sbagliati che
siano, costringono a pensare. Non impone mai il suo punto di vista (``si fa
come dico io''), ma chiede ``tu come faresti? Come posso aiutarti?''. Il
feedback, per lui, serve a \emph{far crescere qualcun altro}.

\subsection{I manager peggiori: ``a che numero sto pensando?''}

I peggiori manager che il relatore abbia avuto erano quelli che volevano che si
facesse come volevano loro, anche quando chiedevano ``tu come la faresti?''.
In realtà la domanda era ``\emph{a che numero sto pensando?}'': se dicevi 5,
``no, dovremmo fare 9''; se dicevi 9, ``no, 5''. Lo stress costante di dover
indovinare cosa vuole il capo nasce dalla mancanza di fiducia (``se mi chiedi
come fare, fammelo fare'').

Lo stesso gli succede con i CEO che lo chiamano per ``sistemare la cultura
aziendale'' e poi dicono ``non sono sicuro, io farei diversamente''. La sua
risposta: ``lo so che faresti diversamente, per questo mi hai chiamato''. Se si
chiama un esperto, bisogna lasciargli fare il suo lavoro: non si va dal
meccanico a cambiare le gomme per poi dirgli ``non dovresti girare un po' di
più quel bullone?''. Chiedere un parere a un esperto significa dargli anche il
potere di scegliere. Il problema è il \emph{micromanagement}, o in generale il
non accettare punti di vista diversi dal proprio (magari perché si è il
fondatore e si pensa di essere l'unico a saper fare le cose).

\subsection{Dare feedback a un superiore egocentrico}

\paragraph{Come si dà un feedback negativo a un superiore con un grande ego?}
Il trucco è \textbf{non puntare sull'ego, ma sull'azione e sui suoi effetti}.
Non ``tu fai così e io mi sento così'': anzi, si alimenta il suo ego
rendendolo il \emph{main character} dell'avventura, l'eroe che risolve il
problema (``tu, grande manager, aiutaci a risolvere questo problema che
affligge noi che facciamo l'operatività''). Non è l'ideale (``io cambierei
manager''), ma pochissimi possono scegliersi i colleghi: il relatore stesso ha
quasi sempre ricevuto team preassemblati (``ecco le chiavi, questa è la tua
macchina''). Bisogna giocare con le carte che si hanno in mano.

\paragraph{Un esempio pratico.} Con un manager che fa micromanagement e mette
sempre ``il becco'' nel proprio lavoro, si può dire: ``se ti fidi di me per
questo lavoro, con le tue capacità potresti dedicarti a questa cosa più
strategica e di lungo termine: pensa quanto sarebbe più utile per l'azienda''.
In sostanza: ``guarda cosa potresti raggiungere se mi lasciassi fare il mio
lavoro''. Non come un ``tu fai il tuo e io faccio il mio'', ma come un
``ti sollevo da questa parte del carico, che so che ritieni importante: lascia
che ci pensi io, così puoi dedicarti ad altro'', che magari gli dà anche più
visibilità.

Si gioca su ciò che quella persona pensa di saper fare, e che probabilmente sa
anche fare. Anche Michael Scott di \emph{The Office}, per quanto preso in giro,
è raccontato come un venditore bravissimo, che ha chiuso i contratti più
importanti per l'azienda, e rispetto ai manager che gli succedono è forse
``uno dei meno peggiori''. Tutti, anche i manager peggiori, hanno punti di
forza: \textbf{giocate su quelli}.

% ---------------------------------------------------------------------
\section{Il feedback come regalo}
\label{sec:z1-regalo}

Per chiudere, il relatore risponde alla domanda d'apertura, che aveva lasciato
apposta in sospeso: che cos'è il feedback per lui?

\begin{quote}
  \emph{Il feedback è un regalo che mi viene fatto da qualcuno che so che vuole
  il mio interesse.}
\end{quote}

Ha costruito un team di persone che si fidano di lui e che non hanno motivo di
sabotarlo o di ferirlo: per questo interpreta anche le cose peggiori con il
default to trust. È un regalo, e cerca di prenderlo in considerazione il più
possibile.

Ci è arrivato partendo da se stesso. Oltre a essere avverso ai conflitti, si
definisce ``avverso alla comunicazione umana'': paradossalmente, una persona
introversa e timida finita a fare comunicazione, marketing e ora HR (``succede
spesso''). Prima di ordinare una pizza al telefono ha dieci minuti di
preparazione mentale, e se l'ordine delle domande cambia va in confusione. Per
lui un feedback di dieci minuti sono spesso \textbf{tre ore di preparazione}:
costruisce un \textbf{albero dei dialoghi}, come in un'avventura grafica
personale, con tutte le opzioni. Poi spesso finisce con ``hai fatto questa
cosa, andrebbe migliorata'' ``hai ragione, grazie'', e basta.

Così ha pensato: chissà quanti si fanno gli stessi problemi con lui, e passano
ore a prepararsi per dirgli ``questa cosa l'hai fatta male'', sapendo che
riceveranno comunque quell'occhiata fugace del ``come ti permetti''. Eppure
vengono a dirglielo. Non può far altro che rispondere: ``grazie, hai ragione,
cercherò di fare meglio''. È così che ha metabolizzato il ricevere feedback:
capendo che non c'è nessuna cattiva intenzione e che, se non proprio felice
(``ci ho messo una pietra sopra''), può almeno essere \emph{grato} che ci sia
ancora qualcuno che ha così a cuore il suo interesse da dargli feedback.

\paragraph{Il silenzio fa più paura.} Non bisogna spaventarsi dei feedback
negativi che emergono dai sondaggi aziendali: bisogna spaventarsi del
\textbf{silenzio}, quando nessuno dice più nulla. Significa che nessuno si
fida più che si sappiano risolvere i problemi.

\paragraph{La sindrome dell'impostore.} Il suo motivatore personale, all'inizio,
era soprattutto ``non voglio fallire'': per questo cercava di migliorare, e
migliorare è diventato accettare feedback, perché più feedback riceve, più
migliora. Il feedback è un \textbf{acceleratore}, e lui cerca di incentivare
gli altri a darglielo. Resta però sempre la vocina del ``e se si accorgono che
non sei capace?''. Con la \emph{impostor syndrome} convive da anni e
continuerà a conviverci; se ce l'avete, conviveteci anche voi. Avrete sempre
la sensazione che ``se ne accorgeranno'': forse sì, ma sapete fare anche
altre cose. Non concentratevi solo su ciò che non sapete fare. Gli errori
capitano: la perfezione non è di questo mondo.

\paragraph{Esercitarsi in accademia.} Per fare pratica con il feedback,
l'ambiente migliore è quello dei colleghi di corso: persone nella stessa
situazione, con esperienza simile e lo stesso obiettivo, in un contesto in cui
non si può essere licenziati e l'unico effetto a lungo termine è ``ho parlato
con questa persona, ora le cose vanno meglio''. E bisogna avere il coraggio di
parlare anche con i docenti (``dire a Diego: secondo me questa lezione potevi
farla in questo modo''; ``so che è una persona che lo accetta''), senza farsi
bloccare dalla gerarchia. Nell'azienda giusta (come quella in cui lavora ora) il
relatore va sempre dal CEO a dare feedback, sapendo che non verrà licenziato:
il CEO gli ha detto ``ti ho preso per darmi feedback''.

La frase di chiusura:

\begin{quote}
  \emph{Per superare le vostre difficoltà, tutto quello che vi viene chiesto è
  il coraggio di chiedere aiuto.}
\end{quote}

Chiedere aiuto non è una debolezza: è un segno di consapevolezza di sé
(\emph{self-awareness}) e di forza personale. Se tutti lo facessero, i
progetti andrebbero molto meglio; spesso invece ci si fa frenare dal ``non
voglio fare brutta figura'', ``non voglio passare per quello scemo'', e sono
proprio queste paure a rallentarci e a impedirci di raggiungere gli obiettivi.

% ---------------------------------------------------------------------
\section{Domande e risposte finali}
\label{sec:z1-qa-finale}

L'ultima mezz'ora è dedicata a domande libere.

\paragraph{E se una persona peggiora per uno o due anni di fila?} Se nessuno
gliel'ha mai detto, c'è un grosso problema, non suo ma dell'azienda. Se invece
c'è consapevolezza, il percorso è: ``cosa vuoi fare per risolvere?''; ``ho
provato questo, questo e questo''; ``visto che non ha funzionato, vuoi un mio
consiglio? Proviamo con quello?''; ``che supporto ti serve?''. Se, con tutto il
supporto, il tempo e le soluzioni concordate per mesi, il miglioramento non
arriva, per il relatore la soluzione è: \textbf{``dimmi un prezzo per andare
via''}. Preferisce pagare una buonuscita (\emph{severance package}) che
avviare la strada dei richiami formali sulla performance fino al
licenziamento per giusta causa. A un certo punto andare avanti non è un favore
per nessuno.

\paragraph{È vero che nei videogiochi si viene sfruttati? Meglio un piccolo
studio italiano?} Per il relatore si è quasi più a rischio a Blizzard che in un
piccolo studio di Roma o Torino. Nel piccolo studio si sa dove ci si va a
cacciare: non ci sarà il team building a Santo Domingo, ma ci si concentra su
un gioco alla volta e, se va bene, si fa il successivo. A Blizzard si può
lavorare su \emph{Overwatch}, che vende milioni, e perdere comunque il posto,
in modo totalmente scollegato dal proprio lavoro. Gli stipendi più alti
all'estero compensano anche questo rischio, che spesso non dipende da come va
il gioco (anzi, ``il gioco non ha fatto come volevamo'' è spesso una scusa per
licenziamenti che non avevano motivo).

\paragraph{Cambiare settore a trent'anni, da web developer full remote?} Lo
stipendio indicato dallo studente non è fuori dalla portata di vari studi
italiani, anzi è superabile. Non va esclusa la possibilità di lavorare da
remoto per aziende piccole ma estere. Alcuni amici del relatore, pagati con
stipendi italiani, si sono trasferiti a Bali o in paesi con un costo della vita
più basso; il collega di \emph{PUBG} percepisce uno stipendio olandese vivendo
a Milano, che per quanto cara non costa come i Paesi Bassi.

\paragraph{Come si diventa manager?} Dipende dall'azienda: alcune promuovono i
migliori (processo tendenzialmente sbagliato), altre propongono un percorso di
\emph{people management}. Diventare manager significa \textbf{abbandonare ciò
che si sa fare per imparare a fare altro}. Il relatore ha dovuto accettare che
non sarebbe più stato ``il miglior community manager italiano del mondo''
(``eravamo in pochi, quindi era facile''): dedicandosi alle persone non poteva
più fare il suo lavoro bene come prima. È stata una negoziazione con se
stesso: fare il manager significa non essere più l'esperto di qualcosa, ma
saper gestire e far crescere persone. Significa anche \textbf{cedere il
controllo della propria carriera}: se il team decidesse di non fare più nulla,
il manager ne pagherebbe le conseguenze. Il suo lavoro diventa scegliere le
persone giuste e farle lavorare bene. Lui lo esplicita sempre: sa di affidare
al team parte della propria carriera, e il team sa che lui combatterà per
loro, anche prendendosi le colpe.

\paragraph{L'inglese.} \textbf{Va assolutamente imparato}, anche lavorando in
Italia. Non saperlo preclude circa il 95\% del mercato (tutto quello fuori
dall'Italia) e soprattutto il \emph{networking} con persone che potrebbero dare
occasioni future: l'inglese è spesso l'unica lingua in comune con chi lavora
all'estero ed è oggi la lingua universale dei videogiochi.
\begin{itemize}
  \item \emph{Che livello serve?} Non importa B2 o C2: bisogna saper sostenere
        una conversazione, guardare almeno Netflix con sottotitoli in inglese,
        comunicare senza dodici pause e senza ``come si dice\dots''. Il
        certificato non verrà mai chiesto: verrà chiesto se si sa parlare
        inglese. Il relatore ha sentito persone con un C2 che non sapevano
        parlarlo.
  \item \emph{Come allenarlo senza occasioni di conversazione?} Netflix, Prime
        Video, Now TV, Udemy, Coursera, YouTube, Duolingo, siti di scambio
        linguistico con madrelingua, Discord: non sono mai esistiti così tanti
        mezzi gratuiti. Il relatore ha imparato l'inglese giocando (con una
        pronuncia sbagliata, perché leggeva senza sentire, su cui ha poi
        dovuto lavorare); guardare un film in inglese con i sottotitoli in
        italiano insegna anche la traduzione.
  \item Oltre all'ascolto passivo è fondamentale \textbf{parlare}: in una
        flotta di \emph{EVE Online} con 150 persone quasi tutte madrelingua
        inglese si è esposti a decine di accenti e modi di dire (mille modi di
        dire ``come stai'' oltre a \emph{how are you}); lavorando all'estero si
        impara a distinguere gli accenti (Manchester, Londra) e ad apprezzarli
        anche nei media.
\end{itemize}

\paragraph{Quanto deve sapere un junior?} Dipende da quanto è tecnico il
ruolo. Un game designer junior che non sa qualcosa può essere affiancato da un
senior; un programmatore che non sa programmare, in quel momento, non è utile
all'azienda.

\paragraph{Annunci ``junior con cinque anni di esperienza''.} ``Cercasi
ventenne con trent'anni di esperienza.'' Non è l'annuncio per un junior: è un
errore dell'azienda, che in realtà cerca un \emph{mid}. Riceverà profili
junior, si stupirà che non abbiano cinque anni di esperienza, rifarà l'annuncio
e avrà perso sei mesi. Il consiglio: tenere d'occhio quella pagina, perché
dopo qualche mese l'annuncio verrà corretto.

\paragraph{Una struttura per il post-mortem.} A chi chiede un formato di base,
il relatore anticipa quanto vedrà nella seconda giornata a proposito di
brainstorming: una lavagna \textbf{divisa in quattro aree}, su cui ognuno
attacca tutti i post-it che vuole:
\begin{center}
  \begin{tabular}{@{}p{0.44\textwidth}p{0.44\textwidth}@{}}
    \toprule
    \textbf{Smettere di fare} & \textbf{Iniziare a fare} \\
    ciò che non ha funzionato & ciò che non abbiamo fatto e andava fatto \\
    \midrule
    \textbf{Fare di meno} & \textbf{Fare di più} \\
    ciò che ha fatto perdere tempo rispetto all'impatto ottenuto &
    ciò che ha funzionato ma, fatto con troppa leggerezza, non ha dato i
    frutti sperati \\
    \bottomrule
  \end{tabular}
\end{center}
Con i post-it ognuno ha voce e tutto risulta catalogato; per il report basta
una foto della lavagna (o uno screenshot, se la lavagna è digitale, come Miro).

\paragraph{Come sbloccare la paralisi del feedback verso il management?} Uno
studente ha visto manager dire ``non abbiate paura, la mia porta è sempre
aperta'' senza che cambiasse nulla. C'è una dissonanza tra ciò che dicono e
ciò che fanno: è ``la mamma che dice vieni qua che non ti faccio niente''.
Qualcosa deve aver fatto pensare alle persone che non è sicuro. Prima va
sistemato quello; poi si può provare con i \textbf{feedback anonimi}.
L'obiettivo del relatore resta il feedback aperto e firmato, ma a Stormind ha
iniziato con feedback anonimi verso il management per verificare se
l'anonimato fosse un fattore. Non lo era (evidentemente non c'erano problemi a
parlare con il CEO), e la volta successiva intende abbandonarli. Ma è un buon
\emph{check}: se senza anonimato non arriva nulla e con l'anonimato arrivano
cento feedback, il problema è chiaro.

\paragraph{Quando è giusto rivolgersi a un superiore?} La strada giusta resta
parlare direttamente con l'interessato. Ma se si prova e non si trova una
soluzione, la cosa più ovvia è andare dal proprio manager insieme: ``ci
aiuti?''. Quando il relatore parlava di mediare i dialoghi, era proprio questo:
fare da moderatore, aiutare a sbloccare i punti fermi o, se non si riesce a
scegliere tra A e B, scegliere lui. L'escalation, quando necessaria, va fatta.

\paragraph{Il tempo libero veniva rispettato?} No, e se l'è dovuto imporre
lui. Il lavoro gli piaceva così tanto che continuare a casa non gli sembrava
lavorare. Gli si è ritorto contro con il lancio di \emph{Diablo~III}, fatto in
orario americano (le due o l'una di notte in Europa), con l'ufficio chiuso
perché a nessuno era stato chiesto di restare. Il lancio va malissimo, i
giocatori protestano, e lui, sveglio, risponde sui forum fino alle tre, quattro,
cinque del mattino: ``non posso risolvere il problema, ma resto qui ad
aggiornarvi con quello che riesco a sapere dal team di sviluppo''. Il giorno
dopo trova sulla scrivania un \textbf{richiamo formale dell'HR}: aver lavorato
fuori orario senza autorizzazione, condotta non professionale. Per lui è stato
un interruttore: ``ora me ne vado''. Due settimane dopo faceva colloqui con
Riot. Ancora oggi è contrario a quel richiamo, ma gli ha fatto capire come
ragiona un'azienda: ``hai un contratto 9--18, non mi interessa se l'hai fatto
con le migliori intenzioni''. Eppure per anni ha goduto di un'ottima
reputazione in quelle community come ``quello che è rimasto con noi fino alle
cinque del mattino'', e anche l'azienda ne ha beneficiato.

Ci sono poi aziende in cui, soprattutto per i developer, viene richiesto il
famigerato \textbf{crunch}: lavoro straordinario continuativo, senza pause.
Dipende molto dall'azienda; sta venendo (``grazie al cielo'') deprecato, perché
molte aziende capiscono che non è sano ed è visto male anche dal pubblico, ma
esiste ancora.

\paragraph{E se il tuo capo demotiva qualcuno del tuo team?} Qui ``diventa
personale'': il team del relatore si fida di lui perché sa che in questi casi
``diventa una bestia''. ``Te la puoi prendere con me quanto vuoi, ma il mio
team non lo tocchi.'' Lo gestisce direttamente con l'altro manager, per quanto
brutta possa diventare la conversazione. Lo sforzo più grande è \textbf{non
andarne poi a parlare con il team} (``ragazzi, ho fatto questa cosa''), perché
minerebbe la fiducia verso l'alto: non vuole un rapporto ottimo fino a lui e
pessimo oltre. L'obiettivo è che l'altro manager ``si faccia perdonare''
cambiando il proprio comportamento, e che il team percepisca da solo il
cambiamento.

\paragraph{Chiusura della giornata.} Il relatore chiede un feedback collettivo
in chat (``se domani ci dai \emph{more of the same} siamo contenti, oppure
ditemi cosa limare''), si scusa per i ripetuti crash della connessione e
ringrazia per la proattività nelle domande, che gli servono anche come
\emph{check}: se nessuno chiede nulla su un argomento, significa che può
andare avanti. Racconta di aver passato una settimana a decidere che cosa
mettere nelle slide e di averle poi preparate in poco tempo (rifacendone
alcune fino alle tre di notte). Ricorda infine di usare il test di
personalità per sé e per il proprio gruppo, facendolo fare a chi non l'ha
ancora fatto: tornerà molto utile per cooperare meglio.

% ---------------------------------------------------------------------
\section{Anticipazioni sulla seconda giornata}
\label{sec:z1-domani}

La seconda giornata del workshop sarà più pratica: ``oggi vi ho spiegato
perché, domani vi insegno come lo fanno queste aziende''. Gli argomenti
previsti:
\begin{description}
  \item[Time management] come organizzarsi una giornata professionale in
        azienda: come strutturare una \emph{task list}, cosa fare e cosa non
        fare, come e cosa prioritizzare.
  \item[Brainstorming] i gruppi dovranno fare moltissimi meeting (come facciamo
        il gioco, come lo chiamiamo, che scopo gli diamo). Se non si sanno fare
        bene, i brainstorming sono \textbf{dannosi ed eterni}: un brainstorming
        dovrebbe durare \textbf{al massimo due ore}. Verranno mostrati metodi
        che accorciano la durata, massimizzano l'effetto e, soprattutto,
        \textbf{fanno partecipare tutti}: non c'è niente di peggio di un
        brainstorming in cui parlano in due per il 90\% del tempo. Il classico
        ``ragazzi, non esistono idee sbagliate, sparate'' è il modo peggiore di
        farlo; non è colpa di nessuno, ma si può fare diversamente.
  \item[Guidare da dentro] come essere leader nel proprio team anche senza
        esserne il leader: sviluppare responsabilità personale e verso gli
        altri. Non essere il manager di qualcuno non significa non esserne
        responsabili: la persona accanto è responsabilità di tutti, perché si
        è un team e non un insieme di individui. ``Se qualcuno cade, tutti lo
        facciamo rialzare: non c'è l'addetto al rialzamento.'' È la versione
        \emph{advanced} della giornata: come comportarsi da manager senza
        averne il titolo.
\end{description}
Torneranno inoltre il \emph{service level agreement} e il tema della
condivisione dei successi e della protezione del team dai fallimenti.

% ---------------------------------------------------------------------
\section*{Punti chiave del capitolo}
\addcontentsline{toc}{section}{Punti chiave del capitolo}

\begin{itemize}
  \item ``La capacità di un gruppo di persone di fare cose straordinarie
        dipende dal modo in cui queste persone operano insieme.'' Il 90\% dei
        problemi dei team sono \textbf{problemi di comunicazione}.
  \item Le aziende si liberano sempre più presto dello \textbf{useless
        genius}: le competenze si trasmettono, il carattere no. Il
        \textbf{fit culturale} va valutato anche dal candidato.
  \item Non esiste un unico modo di dare feedback: dipende dal contesto e
        soprattutto da \textbf{chi lo riceve}. Il feedback è una somma di
        piccoli miglioramenti periodici; tenersi le cose nel cassetto non
        serve mai.
  \item I team attraversano \textbf{forming, storming, norming, performing e
        adjourning}. Lo storming non si salta; si passa la maggior parte del
        tempo in forming; la bravura sta nell'accelerare i passaggi e nel
        fare bene il \textbf{post-mortem}, anche dei successi.
  \item Gli obiettivi di un progetto vanno fissati \textbf{all'inizio}: solo
        così si può valutare se è andato bene e scegliere tra A e B.
  \item Il circolo virtuoso è \textbf{feedback $\rightarrow$ miglioramento
        $\rightarrow$ fiducia}. Accettare un feedback non significa fare come
        dice l'altro: decide chi ha la responsabilità, ma deve argomentare.
  \item ``Aiutami a poterti aiutare'': meglio dire in anticipo ``potrei non
        farcela'' che ``non ce l'ho fatta''. Chiedere aiuto non è una
        debolezza.
  \item Quando si dà feedback \textbf{non si porta la soluzione}: la si
        costruisce insieme a chi lo riceve.
  \item La cultura è forte quando si lavora \textbf{per gli altri}; il lavoro
        è valutato come opera d'insieme (anche con bonus di team).
  \item Come in \emph{Among Us}, i team sprecano energie a cercare il
        traditore invece di fare le task. \textbf{Default to trust}: chi ha
        sbagliato ha agito con le informazioni che aveva. Cercare
        \textbf{la causa, non il colpevole}: ``che cosa ti ha fatto pensare
        che fosse la scelta giusta?''.
  \item Carriere e performance non sono lineari: vanno valutate le
        \textbf{traiettorie}, non un punto dell'anno. Le valutazioni puntuali
        generano \emph{lying, hiding, faking}.
  \item Confermare sempre le istruzioni facendosele ripetere; ``non ho
        capito'' non è una colpa.
  \item Dare feedback: sull'\textbf{azione} e sui suoi \textbf{effetti}, specifico,
        presto, uno a uno, adattato al carattere, senza terze parti, senza
        ``non mi piace'', senza offendersi se non viene seguito; poi fare il
        debriefing.
  \item Ricevere feedback: \textbf{ascolta, chiedi dettagli, accetta, agisci}.
        Il feedback positivo serve a non peggiorare ciò che si fa bene; il
        feedback accidentale e il \emph{drive-by} non servono a nulla.
  \item Feedback \textbf{frequenti e recenti}: tanti piccoli \emph{sanity
        check} costano meno di una grande correzione di rotta.
  \item Il feedback è \textbf{un regalo} di chi ha a cuore il nostro
        interesse; fa paura il silenzio, non le critiche.
\end{itemize}

%    % ... altre lezioni zero ...
%    \part{Il corso di Game Design}
%    \renewcommand{\thechapter}{\arabic{chapter}}
%    \setcounter{chapter}{0}
%    % =====================================================================
%  Capitolo 1 — Introduzione al corso e al mestiere del game designer
%  Fonte: trascrizione della lezione GD1 (Game Design, primo anno)
%  Pacchetti richiesti nel preambolo del libro:
%    \usepackage[italian]{babel}, \usepackage[T1]{fontenc},
%    \usepackage[utf8]{inputenc} (solo pdfLaTeX), \usepackage{booktabs}
% =====================================================================

\chapter[Il mestiere del game designer]{Il mestiere del game designer: introduzione al corso}
\label{cap:introduzione-game-design}

Questo primo capitolo fa da soglia. Prima di affrontare metodi, framework e
documenti, è necessario capire \emph{in che modo} si impara il game design,
\emph{con chi} lo si fa, \emph{con quali strumenti} e, soprattutto,
\emph{che cosa fa concretamente} un game designer nel suo lavoro quotidiano.
Il capitolo segue l'impostazione della lezione inaugurale del corso: si parte
dal contesto formativo (il corso come simulazione di lavoro), si passa al
metodo di apprendimento (l'\emph{inquiry-based learning}), si elencano
letture e risorse di riferimento, e si arriva infine a una prima definizione
operativa del ruolo del designer, dei limiti entro cui lavora e delle sue
possibili specializzazioni.

Una premessa terminologica: il game design è una disciplina giovane e quasi
tutta la sua letteratura è scritta in inglese. Per questo motivo molti termini
tecnici verranno lasciati in lingua originale (\emph{pipeline},
\emph{pitching}, \emph{tuning}, \emph{feature team}, \dots): sono i termini che
si useranno nel lavoro, ed è importante imparare a riconoscerli e a usarli
correttamente fin da subito.

% ---------------------------------------------------------------------
\section{Il contesto: un corso come simulazione di lavoro}
\label{sec:simulazione-lavoro}

\subsection{Il principio fondamentale}

Il primo concetto da fissare è che il corso \textbf{si presenta come una
simulazione del lavoro}. Dal momento in cui si entra in aula si viene trattati
come \emph{potenziali professionisti dell'industria}: non soltanto studenti,
ma lavoratori effettivi di uno studio (l'accademia), con una gerarchia, dei
referenti, dei compiti e delle scadenze.

Ne derivano alcune conseguenze pratiche:
\begin{itemize}
  \item il docente non assegna un ``compitino'' da risolvere per accumulare
        punti o presenze: ogni \emph{task} è pensato per far dedicare tempo
        all'apprendimento di qualcosa di concreto;
  \item non consegnare un task non danneggia il docente (che ``vive ancora
        domani e viene pagato ugualmente''), ma lo studente, che perde
        l'occasione di mettere le mani su un argomento e diventare più abile;
  \item docenti e studenti devono trattarsi \textbf{da colleghi}, perché è
        molto probabile che un domani lo diventino davvero.
\end{itemize}

\subsection{Perché l'impegno conta più che mai}

Il lavoro del game designer è descritto come \emph{molto romantico, ma anche
molto duro}, e sempre più duro col passare degli anni. Durante il periodo del
Covid l'industria ha visto grandi investimenti; quei ``rubinetti'' sono stati
chiusi da tempo, il rapporto con i publisher è diventato più complesso e le
opportunità lavorative non sono più quelle di tre o quattro anni fa. Questo
significa:
\begin{itemize}
  \item più competizione per ogni posizione;
  \item aspettative più alte: ci si aspetta che i candidati siano ancora più
        preparati e sicuri delle proprie competenze;
  \item un unico modo per diventarlo: \textbf{lavorare}. Più tempo si dedica
        oggi ad apprendere il mestiere, più facilmente si sarà assumibili
        domani e più agevole sarà il lavoro di tutti i giorni.
\end{itemize}

\subsection{Sbagliare qui per non sbagliare fuori}

Il mantra del corso è: \textbf{quanto più provate qui, anche in modo goffo e
sbagliando, tanto meno sbaglierete là fuori}. In accademia, per quanto la
simulazione sia realistica, nessuno toglie lo stipendio e nessuno licenzia da
un momento all'altro. Nell'industria sì, e un errore può costare molto denaro
all'azienda, ancora di più se l'azienda è la propria.

\subsection{Il valore del fallimento}

Il fallimento, se affrontato con la mentalità corretta, è probabilmente
\emph{la cosa più vantaggiosa} che si possa fare durante la formazione:
\begin{itemize}
  \item un successo (``avete superato questa deadline'') conferma le proprie
        abilità, ma insegna relativamente poco;
  \item dopo $n$ fallimenti si impara moltissimo, e i successi arrivano;
  \item è fondamentale non demoralizzarsi: il feedback negativo non è mai un
        giudizio sulla persona, ma su ciò che è stato consegnato.
\end{itemize}

Il docente riconosce apertamente che dedicare tempo ed energia a un progetto e
vederlo rifiutato, o vederlo non uscire mai, è demotivante: è successo a lui e
lo ha visto accadere anche su progetti molto grandi, con team numerosi. Il
punto però non è ciò che si è ottenuto materialmente, ma \textbf{quanto si è
cresciuti}. Da ogni esperienza ci si portano a casa due cose:
\begin{enumerate}
  \item \textbf{l'esperienza}: quanto si è diventati bravi, veloci,
        praticamente abili nell'applicare un framework, quanto si è cresciuti
        come designer;
  \item \textbf{il rapporto umano}: nel lavoro si ``suda sangue'' fianco a
        fianco con altre persone, e questo crea legami molto forti che
        portano avanti l'industria nel senso più positivo del termine: nascono
        progetti privati, ci si porta da un'azienda all'altra, si formano team
        con basi solide.
\end{enumerate}

Va evitato anche l'eccesso opposto: la cultura del \emph{grinding} si è
appropriata di questi discorsi in modo distorto. Non bisogna diventare robot.

% ---------------------------------------------------------------------
\section{Organizzazione del corso e gerarchia}
\label{sec:gerarchia-corso}

\subsection{La struttura gerarchica}

La struttura dei referenti del corso ricalca volutamente una
\textbf{gerarchia tipica dell'industria}:
\begin{itemize}
  \item al vertice c'è il \textbf{direttore del corso} (Diego Ricciuti), che
        corrisponde al \emph{creative director}: supervisiona la visione del
        progetto (il concetto di ``visione'' verrà approfondito più avanti);
  \item sotto di lui ci sono i \textbf{referenti di dipartimento}, uno per
        disciplina, con un corrispettivo per ciascuna sede (Milano e Roma);
  \item in fondo ci sono gli studenti, che occupano il primo gradino di
        un'azienda: l'\emph{intern} (stagista) o il \emph{junior}, il passo
        che precede il lavoratore ``regolare''.
\end{itemize}

\begin{table}[htbp]
  \centering
  \begin{tabular}{@{}p{0.34\textwidth}p{0.29\textwidth}p{0.29\textwidth}@{}}
    \toprule
    \textbf{Disciplina} & \textbf{Milano} & \textbf{Roma} \\
    \midrule
    Direzione (\emph{creative director}) & \multicolumn{2}{p{0.58\textwidth}}{Diego Ricciuti} \\
    Game design      & Mattia Podini        & Giacomo ``Jack'' Danna \\
    Level design     & Stefano Zocchi       & Andrea Dispoli \\
    Narrative design & \multicolumn{2}{p{0.58\textwidth}}{Daniele Fusetto (lezioni condivise), Giacomo Masi} \\
    Programmazione (Unreal) & Alessandro Pregnolato & Matteo Velardi \\
    \bottomrule
  \end{tabular}
  \caption{Referenti del corso per disciplina e sede (nomi come riportati in lezione).}
  \label{tab:referenti}
\end{table}

Per il primo anno, il referente diretto degli studenti di game design è il
docente di game design: è a lui che ci si rivolge per dubbi e domande. Altre
figure utili sono il direttore didattico (severo, ``giustamente severo'', da
contattare solo per questioni particolarmente importanti), l'amministrazione
(contratti e documentazione) e la segreteria.

\subsection{I canali di comunicazione}

\begin{description}
  \item[Discord] è il canale per le comunicazioni ufficiali del corso e per
        le domande che non richiedono urgenza immediata; può avere qualche
        ritardo.
  \item[E-mail] è il canale per tutto ciò che è formale o ``ancora più
        ufficiale'' e per le comunicazioni rapide e immediate.
  \item[Telegram] esistono residui di vecchi canali precedenti a Discord, da
        evitare.
\end{description}

Le domande che valgono per tutti vanno fatte in classe (anche da un portavoce
a nome del gruppo); quelle singole e personali via Discord o, se formali, via
e-mail. Se il docente non risponde a una domanda c'è un motivo, ad esempio
ragioni legali, e lo dirà esplicitamente.

\subsection{I tre moduli e la settimana tipo}

Le tre discipline principali del primo anno (e non solo) sono:
\begin{itemize}
  \item \textbf{Game Design} (tendenzialmente il lunedì);
  \item \textbf{Level Design} (tendenzialmente il mercoledì);
  \item \textbf{Programmazione} (tendenzialmente il venerdì).
\end{itemize}
Nel primo mese le lezioni sono tutte di game design (lunedì, mercoledì e
venerdì). Sono previste anche lezioni di narrativa, ma il piano formativo si
aggiorna ogni anno: le ore di narrativa sono state ridotte a favore di quelle
di programmazione. Le poche lezioni di narrativa daranno quindi dei
\emph{cenni}, non la preparazione di un narrative designer. Ciò che serve
davvero è imparare \textbf{come comunicano gli altri dipartimenti} e come
lavorare con persone che fanno un lavoro diverso ma complementare al proprio.

\subsection{Le due anime delle lezioni di game design}

\begin{description}
  \item[Design methods (metodologie)] lezioni prevalentemente teoriche, ma con
        una buona componente pratica, in cui si studiano metodologie di design
        e i cosiddetti \emph{framework} dell'industria.
  \item[Documentazione e pitching] il cuore del lavoro quotidiano del
        designer.
\end{description}

\paragraph{La documentazione.}
Il docente si definisce ``il più grande nerd della documentazione'' che si
possa incontrare. È una parte dura, perché a pochi piace documentare, ma il
lavoro del designer \emph{è principalmente questo}: scrivere documenti e, da
quei documenti, passare alla pratica e all'implementazione. Chi non sa
documentare bene parte male. Quando non si è presenti, il documento è il
proprio \textbf{biglietto da visita}: qualcun altro lo leggerà, quindi va
scritto nel modo più chiaro e comprensibile possibile.

Un esempio concreto del perché la documentazione conta: CD Projekt Red ha
sviluppato un proprio engine proprietario, ma buona parte della sua
documentazione è stata ``buttata all'acqua''. Chi deve usarlo o lo sa già
usare, o ha competenze di un altro engine da trasferire, oppure è in
difficoltà. Si documenta \textbf{perché altri possano leggere}.

\paragraph{Il pitching.}
Il \emph{pitching} è l'esposizione, la presentazione orale. La parte creativa
del team presenta spesso:
\begin{itemize}
  \item le proprie idee ai colleghi;
  \item il gioco a convention ed eventi;
  \item il pitch del gioco agli investitori, per ottenere i fondi necessari a
        realizzarlo.
\end{itemize}
Saper parlare e comunicare chiaramente le idee è fondamentale. In sintesi:
\emph{il game designer parla}, e bisogna farsene una ragione.

\subsection{Altre componenti del corso}

\begin{description}
  \item[Gioco critico e playtest] la classe viene divisa in due gruppi: uno
        gioca, l'altro osserva. Chi osserva non applaude ``come le
        scimmiette'': annota, segna e analizza il \emph{playtest}, esattamente
        come si fa nell'industria. Il playtest fa parte delle mansioni del
        designer, ma esiste anche come specializzazione autonoma.
  \item[Pipeline] termine che indica tutto ciò che riguarda i processi
        formalizzati dell'industria; di solito si intende la \emph{pipeline di
        produzione}: come si divide il team, qual è la sua gerarchia, come il
        \emph{producer} gestisce le deadline. In breve: \emph{come} e
        \emph{quando} si fanno le cose. Al primo anno è meno rilevante (si
        dovrà ``nuotare'' un po' da soli); dal secondo anno si vedranno le
        metodologie effettivamente applicate ai progetti digitali. Si
        scoprirà quanto sia difficile lavorare a un progetto digitale, e
        ancor più fisico, senza un piano di sviluppo comprovato, come quelli
        usati per qualsiasi software in qualsiasi parte del mondo.
  \item[Engine] il motore di riferimento è \textbf{Unreal Engine}
        (sezione~\ref{sec:engine}).
  \item[Programmazione] principalmente tramite \emph{Blueprint} in Unreal,
        più eventuali sistemi di versionamento; il corso parte da zero.
  \item[Progetti pratici] probabilmente la cosa più importante. Al primo
        anno i progetti di gruppo restano circoscritti al corso di design,
        perché mancano ancora alcune nozioni (come comunicare, affrontare
        problemi, dare e ricevere feedback). In seguito si allargheranno a
        progetti digitali, molto più difficili di quelli fisici, con membri
        di altri team: un rapporto non unilaterale, ma ``covalente''.
\end{description}

% ---------------------------------------------------------------------
\section{Il metodo: l'inquiry-based learning}
\label{sec:inquiry}

\subsection{Che cos'è}

Il corso è costruito secondo l'\textbf{inquiry-based learning}
(apprendimento basato sull'indagine), un approccio pedagogico nato intorno
agli anni Sessanta, che si richiama a pensatori come Dewey e Piaget. In
generale, le discipline da cui il design attinge sono quasi sempre la
sociologia, la psicologia e, talvolta, l'economia.

Il funzionamento è un ciclo:
\begin{enumerate}
  \item il docente fornisce metodologie e framework dell'industria durante
        la lezione;
  \item lo studente, fuori dalle ore di lezione, lavora sui \emph{task}:
        mette a terra le nozioni, le maneggia, le assorbe meglio;
  \item da questo lavoro nascono nuovi dubbi;
  \item i dubbi tornano in aula: lo studente chiede, il docente risponde
        (ammesso che dare la risposta sia utile, e ammesso che una risposta
        esista).
\end{enumerate}

\subsection{Non sempre esiste una risposta giusta}

Non tutto ciò che si studia ha ``la risposta migliore possibile in assoluto
nel vuoto''. Esistono metodologie, limiti e framework che aiutano il lavoro,
ma non sono fissi: cambiano da azienda ad azienda, da lavoro a lavoro, da
progetto a progetto, e a volte vengono stravolti. Sta alla
\textbf{sensibilità del designer} capire quale strumento o framework usare in
un dato momento, e questa sensibilità si impara solo con l'esperienza. Nel
lavoro nessuno dirà ``usa questa cosa'' o ``è meglio quest'altra'': sarà il
designer a scegliere.

Anche l'esperienza del docente ha confini: sei-sette anni di lavoro, un modo
personale di approcciare il design e un certo insieme di progetti (ad esempio
nessun MMO). Oltre quei confini si cercano risposte nella ricerca, studiando
come aziende con esperienza specifica (ad esempio Riot) hanno affrontato
problemi simili e hanno condiviso le loro soluzioni, sempre valutando quali
metodi siano adatti alla dimensione del team, al tipo di progetto e al tempo
disponibile.

\subsection{Sensibilità, non memoria}

Non interessa imparare a memoria un framework o la teoria di un autore. Certe
terminologie vanno conosciute, ma per ragioni di comunicazione e chiarezza.
Ciò che conta è sviluppare la \textbf{capacità di farsi domande}, di farle
agli altri, di fare ricerca e di trovare una propria risposta. È esattamente
ciò che si fa nel lavoro: davanti a qualcosa che non si sa fare, l'abilità
del designer sta nel cercare come qualcun altro ha risolto un problema simile.

\subsection{Le lezioni come corsi di aggiornamento}

Le lezioni frontali simulano i \emph{corsi di aggiornamento} aziendali: le
sessioni che un lead di design tiene al suo team di design (o il lead di
level design al suo team, e così via). Alla fine dell'aggiornamento c'è quasi
sempre un lavoro da fare, un obiettivo da portare a terra:
\textbf{l'obiettivo della settimana}. Si lavora di settimana in settimana, e
il perché verrà chiarito in seguito.

Struttura tipica di una lezione:
\begin{enumerate}
  \item breve sessione di domande e risposte sulla lezione precedente;
  \item analisi di eventuali esercizi particolarmente critici;
  \item lezione vera e propria;
  \item discussione di gruppo, se c'è tempo;
  \item assegnazione del task della settimana.
\end{enumerate}
Il docente tende a ``parlare e straparlare'', ad aprire parentesi e a saltare
da un argomento all'altro: l'unico modo per fermarlo è che siano gli studenti
a interromperlo con domande.

\subsection{La gerarchia della competenza}

L'inquiry-based learning può essere rappresentato come una scala di
competenza:
\begin{enumerate}
  \item la sola lezione frontale porta a \textbf{conoscenze incomplete},
        perché manca la parte pratica e di ricerca;
  \item il lavoro a casa (task, ricerca) porta a una competenza \textbf{più
        completa}, non sbagliata in sé, ma ancora parziale e complementare ad
        altre;
  \item la \textbf{condivisione} dei task e dei feedback con gli altri porta
        idealmente alla competenza più completa possibile;
  \item il picco della competenza è la sua \textbf{applicazione nei
        progetti} (fisici o digitali, di gruppo): ``con le cose che avete
        imparato e ricercato, riuscite a fare un gioco?''. Quando la risposta
        diventa positiva, o comunque si inizia a navigare nei processi di
        design per realizzarne uno, quella è la vittoria.
\end{enumerate}
Non si parte subito creando giochi: si parte dall'analisi, dal design e dalle
regole generali di design, prima ``nel vuoto'', poi contestualizzando i
framework, infine unendo i puntini nei progetti.

Rimarranno comunque delle lacune, per tutti, docente compreso, perché il game
design \textbf{non è una scienza esatta}, anzi, a rigore non è affatto una
scienza. Prende in prestito metodologie dalle discipline scientifiche, ma
contiene una componente imprevedibile e incalcolabile: \textbf{la
produzione}. Si possono applicare perfettamente le regole di un framework, e
quel framework può comunque non essere adatto al gioco, al proprio modo di
sviluppare o al tempo di produzione disponibile (alcuni framework richiedono
una ricerca e un'attenzione che il tempo non concede).

C'è poi la componente creativa, la parte ``romantica'' del lavoro: esiste, ma
è \textbf{metà del lavoro}. L'altra metà è ricerca, metodo scientifico al
100\%. Dagli stessi dati $n$ designer diversi, con sensibilità diverse,
possono ottenere $n$ risposte diverse. Non tutte saranno valide: vanno
applicate, messe a terra, e questa scrematura fa parte del lavoro.

\subsection{I livelli dell'apprendimento}

Un secondo modello, a gradini, descrive i livelli di apprendimento (è la
cosiddetta tassonomia SOLO):
\begin{description}
  \item[Pre-strutturale] il docente dà il compito e lo studente lo risolve,
        possibilmente il prima possibile. È esattamente ciò che \emph{non} si
        vuole.
  \item[Uni-strutturale] si affronta una sola disciplina: è il primo passo
        necessario, ed è il motivo per cui nel primo mese ci sono solo lezioni
        di game design.
  \item[Multi-strutturale] entrano in gioco altri dipartimenti: level design,
        game design, programmazione e così via.
  \item[Relazionale] si capisce che alcune cose valgono per più
        dipartimenti: alcuni framework di narrativa derivano dal design o
        possono essere usati per fare design; le regole di documentazione si
        applicano anche ad altri dipartimenti; le regole di level design si
        mescolano con quelle di game design. Dopotutto, per fare un gioco
        servono molti dipartimenti.
  \item[Astratto esteso] il punto finale: muoversi autonomamente e capire
        cosa usare e quando, sulla base della propria esperienza. L'unico modo
        per arrivarci è fare pratica e fare progetti.
\end{description}

% ---------------------------------------------------------------------
\section{Task, feedback e regole del corso}
\label{sec:task-feedback}

\subsection{I task come opportunità}

I task non sono strettamente obbligatori (esiste un numero minimo di
consegne), ma sono \textbf{opportunità}. Quanto e come ci si lavora dipende
dallo studente. Un task fatto in cinque minuti, o in cui si è stati ``furbi'',
non verrà letto né riceverà feedback.

Domande da porsi davanti a ogni task:
\begin{enumerate}
  \item \textbf{Perché} mi è stato assegnato? Qual è l'obiettivo?
  \item Come ci lavoro per ottenere un risultato \textbf{che abbia senso}? Se
        la ricerca e i dati indicano una direzione, è inutile andare in
        un'altra.
\end{enumerate}

\subsection{Quanto lavorare: un calcolo indicativo}

Se il corso è una simulazione di lavoro, vale il carico di un lavoro normale:
\begin{itemize}
  \item un lavoro a tempo pieno comprende da 30 a 45 ore settimanali; si
        considera una media di \textbf{40 ore};
  \item si sottraggono circa 9 ore di lezione: restano $\approx 31$, diciamo
        \textbf{30 ore};
  \item con circa 3 task a settimana: $30 / 3 = $ \textbf{circa 10 ore per
        task}.
\end{itemize}
Il numero può variare, ma rende l'idea: dedicare un'ora a un task e poi fare
altro non fa guadagnare nulla a nessuno. Lo studente ``butta soldi'', il
docente risparmia tempo leggendo il task in un minuto (o non leggendolo).
Dedicando invece attenzione, ricerca, cura nella formattazione e
sperimentazione, e facendolo con cognizione di causa, si impara.

\subsection{Niente voti, solo feedback}

Ad eccezione del corso di programmazione, non ci sono voti (né ``ottimo'', né
10, né 100). Ci sono \textbf{feedback}:
\begin{itemize}
  \item in game design, generalmente sulla documentazione;
  \item in level design, sui progetti pratici.
\end{itemize}
Il docente indica cosa va bene e cosa va modificato, concentrandosi in
particolare su ciò che va cambiato. Il feedback \textbf{non è mai un giudizio
sulla persona}, ma su ciò che è stato consegnato nel concreto.

\subsection{La condivisione dei task}

Tutti i task consegnati vengono \textbf{condivisi} con la classe, e possono
essere letti solo dopo aver ricevuto i feedback del docente. Si cresce così
non solo con i propri feedback, ma anche con quelli dati agli altri. Regola
importante: ciò che il docente dice a uno studente \textbf{vale per tutti}. Se
a una persona viene detto ``questa cosa non va mai più fatta'', da quel
momento vale per l'intera classe.

\subsection{Non copiare}

Copiare i task degli altri è una delle cose peggiori che si possano fare, non
per una questione morale, ma perché annulla il senso del percorso. Si può
dedicare a ogni task il tempo e l'attenzione che si ritengono opportuni, ed è
normale che col tempo ci si specializzi (preferendo, ad esempio, il level
design o la programmazione tecnica) e si investa meno dove si sente meno
necessità. In questi casi l'alternativa corretta è:
\begin{itemize}
  \item non consegnare;
  \item studiare ciò che hanno fatto gli altri;
  \item eventualmente riconsegnare in seguito.
\end{itemize}
Ci saranno momenti in cui si verrà sommersi dalle cose da fare: bisogna
imparare a scegliere.

\subsection{Il ``non leggo''}

Il docente ha piena facoltà di non dare alcun feedback a un task (il
cosiddetto \textbf{``non leggo''}) se:
\begin{itemize}
  \item non si è lavorato con attenzione;
  \item non si rispettano le indicazioni date;
  \item è stata violata una regola estremamente importante.
\end{itemize}
Un task ``non letto'' equivale a un task non consegnato e non entra nel
conteggio. Attenzione: un task \emph{perfetto} e un task che \emph{fallisce}
nel comprendere o applicare un framework sono \textbf{ugualmente validi},
purché ci si sia impegnati e si sia dedicato tempo, pazienza e ``sudore''.

\paragraph{Riconsegna.} Un task non letto, in ritardo o che si vuole
migliorare può essere riconsegnato e verrà conteggiato, a patto che:
\begin{itemize}
  \item i problemi segnalati siano stati risolti;
  \item sia stato fatto ulteriore lavoro, integrando anche tutto ciò che nel
        frattempo è stato detto agli altri colleghi nei loro feedback.
\end{itemize}
Il task va sviluppato con le competenze che si hanno nel momento della
riconsegna, non con quelle del momento della prima consegna.

\subsection{Requisiti per superare l'anno}

\begin{itemize}
  \item consegnare almeno l'\textbf{80\% dei task} assegnati (soglia
        malleabile);
  \item partecipare almeno all'\textbf{80\% delle lezioni}, in presenza o
        da remoto, preferibilmente in presenza; i ritardi vengono valutati
        caso per caso;
  \item comportarsi da persone corrette e rispettose (è stata firmata una
        carta dello studente): non urlare, non alzare le mani, non
        vandalizzare, non compiere atti illegali.
\end{itemize}
Impegnarsi di più porta inoltre a essere selezionati per varie opportunità.

\subsection{Scadenze}

\begin{itemize}
  \item la scadenza di un task è \textbf{sempre prima dell'inizio della
        lezione} successiva indicata;
  \item se la classe, come team, mantiene comportamenti scorretti, il docente
        può \textbf{anticipare le deadline} per tutti.
\end{itemize}

\subsection{Accessibilità e DSA}

Gli studenti con disturbi specifici dell'apprendimento o neurodivergenze
(dislessia, ADHD, autismo, ecc.) sono invitati a segnalarlo. Quando si
affronteranno le regole di formattazione e documentazione (errori
grammaticali, mancanze, incoerenze) se ne terrà conto e si potranno trovare
dei compromessi, pur ricordando che nel mondo del lavoro nessuno farà
``sconti''. La documentazione può essere inviata all'amministrazione o alla
segreteria, che la comunicherà a tutti i docenti, ma il consiglio è parlarne
anche direttamente con ogni singolo docente.

% ---------------------------------------------------------------------
\section{Le competenze trasversali}
\label{sec:soft-skill}

\subsection{I giochi si fanno in gruppo}

Si possono fare giochi da soli, ed è un'impresa ammirevole; ma lavorando da
soli \textbf{tutto ciò che si guadagna in libertà si perde in produttività}.

\paragraph{Il caso Stardew Valley.} \emph{Stardew Valley} è stato realizzato
da una sola persona, Eric Barone, in anni di sviluppo e di sacrifici: faceva
un altro lavoro per mantenersi (in un cinema) e di notte programmava, faceva
l'arte e componeva la musica. Il risultato è stato straordinario e lo ha reso
ricco, ma si tratta di un caso estremamente raro. Quanti giochi simili, nati
nelle stesse condizioni, sono falliti? Molti di più.

\paragraph{Il rischio d'impresa.} Lavorare da soli non è sbagliato di per sé
e ha aspetti molto positivi: si ha il privilegio di creare esattamente ciò che
si vuole. Ma ci si fa carico anche di tutto il \textbf{rischio d'impresa}: se
si lavora a un gioco per cinque anni e poi il gioco va male, resta il
vantaggio di averlo pubblicato, ma il dopo è sperare di vivere di quella
release o di finanziare con essa il gioco successivo. Le grandi aziende hanno
gli stessi problemi, ma li distribuiscono su più persone e più responsabilità.

L'obiettivo di un professionista è \textbf{saper lavorare in diverse
condizioni}, e la più comune è quella del dipendente (o del freelance /
libero professionista con partita IVA): si viene assunti per le proprie
capacità da un'azienda, piccola o grande che sia, e si lavora con altre
persone.

\subsection{Essere persone con cui si vuole lavorare}

Se le persone non vogliono lavorare con voi, quello è \textbf{il primo motivo
di licenziamento}. E se non si viene licenziati, il clima nell'azienda
diventa terribile: la produttività crolla, la motivazione crolla. Le
\emph{soft skill}, cioè l'insieme di abilità che rendono una persona
piacevole con cui lavorare, vanno coltivate e diventano sempre più importanti
col passare degli anni (relativamente meno ora, molto di più quando i team
saranno più grandi). Anche chi è un ``lupo solitario'' deve imparare a:
\begin{itemize}
  \item lavorare con altre persone;
  \item comunicare con altre persone;
  \item trovare un punto in comune da cui partire per fare le cose.
\end{itemize}

\subsection{Indipendenza e proattività}

Sono le altre due parole chiave del corso:
\begin{description}
  \item[Indipendenza] spesso si verrà lasciati da soli, non per cattiveria ma
        per sfida: quando si cerca una risposta da soli, ciò che conta non è
        tanto la risposta quanto \textbf{il processo}. Una risposta data dal
        docente non vale niente.
  \item[Proattività] chiedere, domandare, agire da soli senza aspettare
        direttive, perché spesso quelle direttive non ci sono, a maggior
        ragione nei team piccoli.
\end{description}
Tutto questo si concretizza nel parlare, nel presentare e nel fare
\emph{post-mortem} (l'analisi a posteriori di un progetto concluso, che verrà
approfondita più avanti), e in generale nell'attenzione alla crescita
personale.

\subsection{Fare domande}

Il corso insiste sulle domande: alzare la mano, intervenire anche ``a
bomba'', fermare il docente ogni volta che qualcosa non è chiaro. Se una
persona non conosce qualcosa, molto probabilmente c'è almeno un'altra persona
che non la conosce. Il docente può dare per scontata la conoscenza di termini
che non lo sono, ed è del tutto normale, soprattutto all'inizio.

% ---------------------------------------------------------------------
\section{Strumenti: engine, Unreal e materiali}
\label{sec:engine}

\subsection{Unreal Engine e gli altri motori}

Il motore di riferimento del corso, usato in level design, è \textbf{Unreal
Engine}: ``un gran mattone'', ma che dà a un designer molte possibilità. Il
vantaggio principale è che permette a un designer non pratico di
programmazione di \textbf{prototipare tramite i Blueprint} (il sistema di
scripting visuale di Unreal). Questo non impedisce di passare a Unity o ad
altri engine, e viceversa: il docente stesso ha studiato e lavorato a lungo
con Unreal, ma la sua ultima esperienza professionale (come level designer) è
stata su Unity, un engine che non conosceva, a dimostrazione che il salto da
un engine all'altro è possibile.

Panorama dei motori citati:
\begin{itemize}
  \item \textbf{Unreal Engine} e \textbf{Unity}: i due più diffusi, con
        modelli di sorgente ``più o meno'' aperto;
  \item \textbf{Godot}: engine open source;
  \item \textbf{CryEngine}: proprietario di Crytek;
  \item \textbf{Snowdrop}: engine di Ubisoft;
  \item Rockstar usa un proprio engine proprietario;
  \item alcune aziende usano versioni \textbf{personalizzate} di engine
        esistenti: ad esempio Milestone, nei giochi di corse, usa un Unreal
        customizzato;
  \item CD Projekt Red ha un engine proprietario (con i problemi di
        documentazione già citati).
\end{itemize}
In sintesi: ogni azienda può avere un engine proprietario oppure una versione
modificata di un engine commerciale, ma si parte sempre da una base.

\subsection{Programmazione: cosa ci si aspetta da un designer}

Non è richiesto alcun bagaglio pregresso di programmazione: il corso parte da
zero. Gli studenti sono stati selezionati (il corso è a numero chiuso) per
competenze analitiche e critiche sopra la media. Anche nel mercato del lavoro,
per una posizione di game designer generalmente \textbf{non ci si aspetta
competenza di scripting} (C++, Java, C\#, ecc.), a meno di voler fare il
\emph{technical designer}, e anche lì varia da azienda ad azienda.

\subsection{Versione e preparazione}

Indicazioni pratiche emerse in aula:
\begin{itemize}
  \item scaricare Unreal Engine; la versione consigliata per il corso è
        stata indicata nella \textbf{5.6} (con conferma successiva). La
        stabilità dipende da cosa si deve fare, da quali tool e plugin si
        usano: i problemi nascono spesso dai plugin. Nel passaggio da Unreal 4
        a Unreal 5, ad esempio, il sistema \emph{Landscape} era inizialmente
        inutilizzabile, mentre in Unreal 4 funzionava molto bene;
  \item è consigliato installare tramite \textbf{Epic Games Launcher}, più
        comodo: permette anche di riscattare giochi gratuiti periodici e di
        accedere ai pacchetti di \textbf{Fab} (il marketplace di asset
        acquisito da Epic), che ogni mese offre alcuni asset gratuiti utili
        per esercitarsi nel level design (ad esempio per fare \emph{speed
        level design});
  \item molti pacchetti di UI o di sistemi di interazione sono utili per
        prototipare al volo, ma spesso inutilmente complicati;
  \item un plugin di \textbf{materiali di prototipazione a griglia} (adattivi,
        con metriche, allineabili alla griglia, con materiali già pronti e
        colori modificabili, dal costo di circa 5 euro) è indicato come
        ``salvavita'' per il level designer; se necessario, i pacchetti
        indispensabili vengono forniti;
  \item portare \textbf{sempre il portatile} alle lezioni di level design
        (da novembre, a partire dall'interfaccia dell'engine); alle lezioni
        di game design basta che lo porti almeno metà della classe, per i
        lavori a coppie. Serve comunque qualcosa per prendere appunti
        (``tavolette d'argilla e cannucce, se volete fare i babilonesi'',
        oppure carta e penna).
\end{itemize}
Chi ha un portatile poco potente dovrebbe comunque provare a installare
l'engine e aprire un template per vedere come gira: all'inizio non si faranno
cose pesanti. Nelle prime lezioni di level design è importante mettere subito
le mani sull'interfaccia: ascoltare dove si trova uno strumento (ad esempio lo
\emph{snapping}) senza provarlo porta a dimenticarlo.

% ---------------------------------------------------------------------
\section{Letture e risorse di riferimento}
\label{sec:risorse}

\subsection{I libri}

\begin{description}
  \item[\emph{The Art of Game Design: A Book of Lenses} (Jesse Schell)] il
        testo \textbf{davvero importante}. Alcune lezioni finali dell'anno si
        baseranno sulle sue terminologie. È utile perché apre la mente sulla
        possibilità di applicare metodologie diverse a giochi diversi, ma
        soprattutto per la \textbf{terminologia} comune. Esiste anche in
        audiolibro (circa 25 ore).
  \item[\emph{Challenges for Game Designers} (Brenda Romero e Ian Schreiber)]
        libro \emph{entry level} con moltissimi esercizi, molto caro al
        direttore del corso.
  \item[\emph{21st Century Game Design} (Chris Bateman)] utile come
        retrospettiva: è un libro visionario che ha previsto metodi e
        soprattutto \emph{archetipizzazioni} dei giocatori tuttora valide nel
        design e nel marketing. Servirà per una delle lezioni più importanti
        dell'anno. È utile leggerlo per vedere quanto aveva previsto, ma è
        utile ``fino a un certo punto''.
  \item[\emph{An Architectural Approach to Level Design} (Christopher
        Totten)] non strettamente necessario, ma di lettura semplice. Ha due
        anime: parte dall'architettura (babilonese, greca, ecc.) e la trasla
        su esempi di level design (come \emph{Doom} e \emph{Team Fortress}).
        Un tempo ne esisteva un riassunto su Gamasutra (oggi
        \emph{Game Developer}).
  \item[\emph{A Theory of Fun for Game Design} (Raph Koster)] utile fino a un
        certo punto: contiene diverse idee superate, che lo stesso Koster ha
        affrontato in vari talk reperibili nel GDC Vault. Verrà ripreso nelle
        prime lezioni, poi messo da parte.
\end{description}

\subsection{Canali e conferenze}

\begin{description}
  \item[GDC (Game Developers Conference)] il suo canale raccoglie talk e
        panel estremamente importanti: è ``uno strumento misterioso che ci
        servirà più tardi''. \textbf{Attenzione}: per molti talk mancano ancora
        le competenze per interpretarli, e alcuni potrebbero fuorviare o
        mettere in testa idee strane. La GDC si concentra spesso su
        produttività e mobile.
  \item[Il \emph{Dragon Speech} di Chris Crawford] il talk del fondatore della
        GDC: appartiene a un passato lontano, ma è molto motivante. Da vedere
        subito.
  \item[Game Maker's Toolkit] canale particolarmente utile per nozioni
        pratiche, con video brevi e alcune interviste.
  \item[Nordic Game] altra conferenza importante, con talk e panel di esperti
        del settore.
\end{description}

\subsection{Associazioni e premi: come leggerli}

Le associazioni di settore (ad esempio IIDEA in Italia, o la IGDA) non sono
il male, anzi: sono un canale importante per affacciarsi all'estero o per
accedere a fondi. Va però ricordato che \textbf{tutelano i propri associati}:
non esiste solo il loro palinsesto, e bisogna valutare quanto spazio occupano
nel panorama della visibilità.

Lo stesso vale per i \textbf{premi} (Game Awards, ``Best Mobile Game'',
``Games for Change'', ecc.): al designer \emph{non interessano}. Che un gioco
abbia vinto ``più di 90 premi'' non dice nulla di utile. Servono a dare
visibilità e lustro ai giochi ed è giusto che esistano, perché ampliano il
panorama e il numero di giocatori. Ma ciò che interessa al designer sono i
\textbf{dati}:
\begin{itemize}
  \item come sono state implementate certe cose;
  \item come è stato fatto il marketing e dove il gioco è stato presentato;
  \item non quante wishlist ha fatto, ma \emph{come} le ha fatte;
  \item in che periodo storico è uscito e quali erano i competitor in quel
        momento.
\end{itemize}

% ---------------------------------------------------------------------
\section{Che cosa fa un game designer}
\label{sec:cosa-fa}

Per arrivare a una prima definizione, la classe ha condotto un brainstorming:
ognuno doveva indicare una cosa che, secondo lui, fa il game designer (senza
ripetere quelle già dette). Ne è emersa la lista seguente, commentata punto
per punto.

\subsection{Progettare meccaniche e struttura del gioco}

Il ruolo più generico è quello del \textbf{core mechanic designer}, detto
anche \emph{gameplay designer} e, un tempo, \emph{black box designer}: si
occupa delle \textbf{meccaniche di gioco}. Si va dal piccolo al grande (in
funzione anche della dimensione dell'azienda):
\begin{itemize}
  \item nel \textbf{piccolo} ci sono le singole meccaniche;
  \item nel \textbf{grande} ci sono i \textbf{sistemi}. Nei giochi
        \emph{sistemici}, con molti sistemi interconnessi (ad esempio
        \emph{GTA}), esistono specializzazioni dedicate a bilanciare e
        supervisionare i sistemi: è il \emph{system design}.
\end{itemize}

\paragraph{Quando un gioco è ``sistemico''?}
Non dipende solo da come le meccaniche interagiscono con il giocatore (le
\emph{dinamiche}, secondo i framework che passano dalle meccaniche alle
dinamiche, che verranno studiati più avanti), ma anche da \emph{quante}
meccaniche ci sono e da quanto interagiscono tra loro. Non diremmo che
\emph{Super Mario} sia un gioco sistemico, eppure contiene dei sistemi: la
linea è sottile. Quando i sistemi aumentano molto, sono complessi e si
influenzano a vicenda, il gioco diventa sistemico. Esempio: in \emph{The
Legend of Zelda: Breath of the Wild} il vento modifica quanto il personaggio
sale o scende; il fuoco può incendiare l'erba; l'erba prende fuoco se colpita
da una freccia infuocata; e così via. Nel mercato moderno i giochi sistemici
sono spesso, ma non sempre, legati a mappe \emph{open world}.

\subsection{Le 3C: camera, character, control}

La base del gioco, da cui generalmente si inizia a lavorare (anche nei
prototipi), è la ``sacra trinità videoludica'' delle \textbf{3C}:
\begin{description}
  \item[Camera] che tipo di inquadratura ha il gioco: prima persona, terza
        persona, isometrica, ecc.
  \item[Character] quanto è grande il personaggio, che tipo di personaggio è
        e, di conseguenza, quali azioni può compiere.
  \item[Control] i controlli: che cosa il personaggio può fare, quali sono i
        \emph{verbi} di gioco.
\end{description}
Le 3C sono definite dal designer, ed esiste anche una specializzazione per
questo lavoro. Esistono varianti: quella del direttore del corso è
\textbf{4C + J}, cioè \emph{camera, character, control, conflict} (conflitto)
e \emph{juice} (il ``succo'', il feedback che rende piacevole l'interazione).
Verranno approfondite più avanti.

\subsection{Il personaggio}

Il designer definisce il personaggio dal punto di vista delle meccaniche
(dimensioni, azioni, capacità), all'interno delle 3C. L'aspetto estetico è
invece competenza di altri dipartimenti.

\subsection{Documentare le informazioni per gli altri dipartimenti}

Assolutamente vero: la documentazione è una delle prime attività del
designer. Più l'azienda è grande, più questo lavoro diventa prominente; più è
piccola, meno è rilevante. In \textbf{pre-produzione}, però, è sempre
fondamentale.

\subsection{Mediatore e leader ``al servizio''}

Il designer è un \textbf{mediatore} nei conflitti interni al team, una sorta
di collante che si relaziona con dipartimenti diversi (non sempre gli stessi,
a seconda della fase di produzione). A volte fa anche da leader, a seconda di
come sono gestite la pipeline e la gerarchia.

Questa leadership è difficile da definire perché \textbf{non è verticale}: è
una leadership \emph{al servizio} (\emph{servant leadership}). Il designer
``è asservito al programmatore'': il suo compito non è arrivare con un'idea e
scaricarla addosso agli altri, né portare la soluzione e lasciare il problema
a qualcun altro, ma \textbf{seguire il processo a basso livello}, sporcandosi
le mani.

\paragraph{I feature team.} Nei team medio-piccoli o medio-grandi possono
esistere sotto-team (\emph{feature team}) dedicati a un obiettivo specifico:
ad esempio un designer incaricato del character o della camera, affiancato da
uno o due programmatori e da un artista. In questi casi il designer
supervisiona il lavoro e gestisce la produzione del sotto-team.

\subsection{Gestire l'avanzamento del prodotto}

Esiste una specializzazione che si occupa solo di questo, il
\textbf{producer}, a un livello gerarchico completamente diverso (e che
potrebbe essere una strada percorribile al termine del percorso di studi).
All'interno di un feature team anche il designer può farlo, ma non è detto.
Può certamente \emph{dare la direzione}: nella sua documentazione può
indicare che per arrivare a uno stato giocabile del prototipo serve un certo
lavoro di arte e di programmazione, per lo stato successivo un altro, e così
via. In ambito produttivo esistono metodologie strutturate come
\textbf{Agile}, che verranno viste in seguito.

\subsection{Presentare le idee}

Il pitching (sezione~\ref{sec:gerarchia-corso}) è richiesto al designer sia
in pre-produzione sia in produzione, in misura che dipende dall'azienda. Ad
esempio Ubisoft lascia ai designer, e in realtà a qualsiasi membro, la
possibilità di \emph{pitchare} un'idea prima di una nuova produzione, perché
possa iniziare a prendere vita.

\subsection{Bozze e prototipi}

Le bozze vanno bene, ma devono essere \textbf{chiare}. C'è un'eterna lotta tra
documentazione e gioco giocabile: se si fa solo documentazione il gioco non
avanza; se si fa solo il gioco manca la documentazione e si perde la visione.
Sono due cose \textbf{complementari}.

Il modo per comunicare un'idea in forma grezza ma giocabile è il
\textbf{prototipo}. Un prototipo può essere qualsiasi cosa: di carta, in
Unreal con Blueprint ``brutti''. \emph{Non deve essere bello, deve essere
giocabile}.

\paragraph{Esempio: il prototipo di carta di Mario + Rabbids.} Durante lo
sviluppo di \emph{Mario + Rabbids}, i programmatori di Ubisoft erano impegnati
su altri progetti (come \emph{Ghost Recon}). I designer realizzarono allora
un \textbf{gioco da tavolo} che riproduceva le meccaniche del gioco, quelle
di uno strategico a turni sul modello di \emph{XCOM}. E funzionava.
L'obiettivo era \textbf{validare un'idea}, validare un gameplay.

Altri modi per rendere chiara un'idea sono i \emph{mock-up} nella
documentazione e le immagini modificate o i \emph{mash-up} di qualsiasi
natura. Lato documentazione l'importante è che sia chiaro; lato prototipo,
che sia giocabile.

\subsection{Fare ricerca per l'idea, e poi smontarla}

Il designer ha una componente creativa importante e fa ricerca per arrivare a
un'idea. Ma ciò che deve poi fare necessariamente è \textbf{smontarla}. Il
suo lavoro è scientifico, esattamente come il \textbf{metodo scientifico}, il
cui obiettivo è \emph{confutare un'ipotesi}: si arriva a un'ipotesi tramite
sperimentazione e poi si cerca $n$ volte di confutarla; se resiste, la si
assume come valida \emph{in un certo contesto} (le regole della fisica
classica, ad esempio, non valgono nella meccanica quantistica).

Allo stesso modo il designer arriva a un'idea (tramite ricerca o chiedendo
agli altri: le idee le hanno tutti, non sono un privilegio del designer),
cerca di smontarla e poi cerca di \textbf{metterla a terra}. Se l'idea non è
applicabile, non è implementabile, finché non si trova un modo per
realizzarla va scartata. Può essere l'idea più bella dell'universo, ma
\textbf{se non può essere messa a terra, non esiste}.

\paragraph{L'originalità.} Serve un'idea originale? Dipende dall'obiettivo e
dalla visione del progetto:
\begin{itemize}
  \item se l'obiettivo è vendere molto, un'idea originale non serve
        necessariamente;
  \item se si vuole creare un nuovo filone, serve un'idea, se non originale,
        almeno interessante: può essere un discorso di marketing o una
        meccanica innovativa, con il rischio di non vendere.
\end{itemize}
Il ruolo del \emph{director} e il concetto di \emph{visione} verranno
approfonditi in seguito.

\subsection{Analizzare i competitor}

Arrivati a un'idea (``vogliamo un certo tipo di combattimento''), si va a
vedere se qualcuno l'ha già fatta e \emph{come}. Le meccaniche non sono
privilegio di nessuno, salvo quando vengono brevettate, come il
\textbf{Nemesis System} di \emph{La Terra di Mezzo: L'Ombra di Mordor}.

\paragraph{Esempio: il combattimento da ridimensionare.} In un progetto
(presso Stormind Games) il team si trovò con un sistema di combattimento
troppo costoso: mancavano le risorse e gli animatori per realizzarlo. Poco
output lato gameplay significava non poter fare qualcosa di complesso e
profondo come in \emph{Elden Ring}; al contrario, il team aveva molti
artisti e poteva puntare su qualcosa di più artistico. Si iniziò così a
orbitare intorno al corpo a corpo di \emph{The Last of Us} (parte I e II).
Sul canale della GDC esistono diversi talk che raccontano cosa è stato fatto
e come, con successi, insuccessi e costi (ad esempio del \emph{motion
capture}): da quei talk si ricavarono moltissime informazioni utili.

I competitor possono essere:
\begin{itemize}
  \item \textbf{diretti}: un gioco molto simile al proprio;
  \item \textbf{di mercato}: un gioco che domina il mercato in cui ci si vuole
        affacciare;
  \item \textbf{indiretti}: un gioco con un sistema simile al proprio, ma
        tutto il resto diverso.
\end{itemize}
Regola d'oro: \textbf{là fuori, nascosto nell'ombra, c'è sicuramente un gioco
molto simile al vostro}. Va solo trovato. E se per qualche ragione non ha
venduto, bisogna capire perché.

\subsection{Lavorare in engine e fare tuning}

Il level designer lavora in engine, così come il technical designer; ma
qualsiasi game designer, a seconda di ciò che deve fare, può finire in
engine, soprattutto per il \textbf{bilanciamento} e il \textbf{tuning}.

Il \emph{tuning} è l'azione di bilanciare con precisione i dati. Esempio: in
\emph{Super Mario} si modifica quanto in alto salta il personaggio, con che
velocità, in quanto tempo raggiunge l'apice, e così via. Questo lavoro rende
il gioco:
\begin{itemize}
  \item ad alto livello, dotato della \textbf{progressione} giusta;
  \item a basso livello, \textbf{piacevole da giocare}.
\end{itemize}
Il \emph{fine tuning} è un lavoro di dettaglio, paziente e minuzioso.

\subsection{Bilanciare}

Il bilanciamento, emerso tardi nel brainstorming, è una delle attività
centrali del designer e ha anche una specializzazione dedicata
(l'\emph{economy designer}, sezione~\ref{sec:tipi-designer}).

\subsection{Analizzare il mercato}

In pre-produzione, di solito, sì: capire le tendenze, i generi che cambiano
nel tempo, la compatibilità del gioco con il momento storico, in vista ad
esempio di un possibile publisher. Va capito quanto, come e con quanto tempo
farlo, perché si è sempre limitati da molte cose.

\subsection{Sintesi}

\begin{quote}
Il designer parte dall'idea ma non si ferma lì: la confuta, lavora sui
dettagli, bilancia, fa \emph{fine tuning}; lavora con il team e fa da
\emph{manager} per diversi dipartimenti; e naviga costantemente tra i limiti.
\end{quote}

% ---------------------------------------------------------------------
\section{I limiti del design}
\label{sec:limiti}

I limiti ci sono sempre, e nel mondo del design sono spesso numerosi. I più
importanti entro cui navigare sono:
\begin{description}
  \item[Tempo] ad alto livello, l'idea deve essere realizzabile nel tempo di
        produzione disponibile; a basso livello, il proprio lavoro
        settimanale deve essere fattibile nel tempo che si ha. Non ci si può
        dividere tra mille feature team se realisticamente se ne può seguire
        uno solo.
  \item[Tecnologia] si può o non si può fare qualcosa? Ci sono i server per un
        multiplayer con un certo numero di giocatori? Ci si può permettere,
        ad esempio, una fauna che reagisce alla fisica come quella di
        \emph{GTA VI}, che richiede il meglio dei \emph{technical artist}?
  \item[Background e competenze] che cosa sa fare il team? Realisticamente,
        chi ha lavorato dieci anni su giochi horror viene assunto per un
        horror, perché ha le competenze per risolvere problemi già visti.
  \item[Persone e motivazione] il team ha la motivazione e le competenze per
        fare un certo tipo di gioco, e per lavorare insieme?
  \item[Denaro] tutto ruota intorno ai soldi. Che si lavori da soli o in mille,
        le persone vanno pagate e il gioco deve avere un ritorno economico.
        Bisogna decidere dove tracciare la linea: raggiungere il
        \textbf{break-even} (guadagnare esattamente quanto speso) o puntare a
        grandi guadagni. Con un \textbf{publisher}, inoltre, una parte
        consistente dei ricavi va a lui (ripartizioni dell'ordine di 80/20 o
        70/30, che dipendono dal contratto firmato).
\end{description}

\subsection{Non siete ``gli uomini delle idee''}

Per qualcuno può essere un problema scoprirlo, ma il designer non è
``l'uomo (o la donna) delle idee'': accanto alla parte creativa c'è una parte
tecnica, e si è fortemente limitati da molti fattori. Questi limiti, però,
sono un aiuto:
\begin{itemize}
  \item riducono il ventaglio di possibilità potenzialmente distruttive per il
        team (cose che non si riescono a fare, che non piacciono o per cui
        mancano le risorse);
  \item aiutano a mirare bene, trovando soluzioni creative in uno spazio più
        piccolo e più attuabile.
\end{itemize}

\begin{quote}
\centering
\textbf{Le idee sono nulla, l'implementazione è tutto.}
\end{quote}

% ---------------------------------------------------------------------
\section{Ruoli e specializzazioni}
\label{sec:tipi-designer}

\subsection{La scala gerarchica}

\begin{description}
  \item[Creative Director (CD)] mantiene la \textbf{visione} e dà l'idea, in
        modo che tutto segua una certa direzione. In teoria può diventarlo
        chiunque, ma con un'esperienza notevole, nell'ordine dei decenni: si
        assume che in quel tempo abbia visto come lavorano molti dipartimenti.
        Naturalmente un CD con un background artistico ha meno competenze di
        design, e viceversa.
  \item[Director di disciplina] uno per dipartimento: \emph{game director},
        \emph{level director}, \emph{art director}, ecc.
  \item[Lead] il capo del dipartimento, che gestisce i gruppi di persone. Chi
        lavora in un dipartimento risponde al proprio lead (ad esempio un game
        designer risponde al \emph{lead designer}, non al game director).
  \item[Senior / Principal] i profili più esperti (a seconda dell'azienda).
  \item[Junior] il profilo di inizio carriera.
  \item[Intern] lo stagista.
\end{description}
Lo schema si replica uno a uno per ogni dipartimento: arte, suono, animazione
(a volte arte e animazione insieme, o \emph{rigging} e animazione insieme),
programmazione, UX/UI e così via. A fare la differenza, a parità di ruolo, è
il \textbf{livello di seniority}.

\subsection{UX e UI}

UX e UI stanno a metà tra game design e arte: a seconda dell'azienda il ruolo
si chiama \emph{UX/UI Artist} o \emph{UX/UI Designer}.
\begin{description}
  \item[UX (User Experience)] come una persona interagisce con le cose. Esempio:
        su uno smartphone, un pulsante accende e spegne lo schermo; per
        digitare la password si scorre verso l'alto e compaiono i punti da
        toccare.
  \item[UI (User Interface)] come sono ``vestite'' queste cose: ad esempio
        l'aspetto della barra dei punti vita (HP), realizzata con asset 2D o
        3D.
\end{description}

\subsection{Le specializzazioni del game design}

Più l'azienda è grande, più la specializzazione è puntuale.
\begin{description}
  \item[Gameplay / Core mechanics designer] (sono la stessa cosa) si occupa
        delle meccaniche e delle 3C.
  \item[System designer] si occupa dei sistemi e delle loro interazioni.
  \item[Economy designer] si occupa del bilanciamento; lavora soprattutto con
        fogli di calcolo (Excel, Google Sheets).
  \item[Combat designer] si occupa del combattimento; negli annunci di lavoro
        indica quasi sempre un combattimento \emph{action}, con gameplay
        reattivo (non strategico), anche se non è una regola.
  \item[Multiplayer designer] si occupa delle interazioni multigiocatore.
  \item[Technical designer] fa da ponte tra game design e programmazione.
\end{description}

\subsection{Le specializzazioni del level design}

\begin{description}
  \item[Encounter designer] posiziona i nemici nella mappa in modo che gli
        scontri siano bilanciati e la progressione equilibrata come
        desiderato.
  \item[World / map designer] popola il mondo di gioco.
  \item[Quest designer] popola il mondo con le missioni e pensa alla
        progressione al loro interno, cioè a dove il giocatore deve andare (si
        vedano ad esempio gli annunci di lavoro di CD Projekt Red o di
        Larian).
  \item[Technical level designer] fa da ponte tra level design e
        programmazione.
\end{description}

\subsection{Discipline vicine: suono, UI, produzione}

Il monte ore del corso non permette di uscire come \emph{narrative designer}
(o \emph{sound designer}) completi. Ciò non toglie che si possano coltivare
queste discipline:
\begin{itemize}
  \item nei progetti di gruppo si può occuparsi anche del suono, della UI/UX,
        o provare ad applicare metodologie come Agile per fare
        \emph{producing} e gestire il team;
  \item il tempo dedicato a queste attività va però sottratto al game design,
        e dipende da quante persone ci sono nel team, da quanto lavoro c'è e
        da quanto è grande il progetto;
  \item avere sviluppato e rilasciato un progetto, anche piccolo (anche una
        \emph{game jam}), documentando i propri processi, è già un ottimo
        argomento da presentare in futuro;
  \item la via migliore resta lavorare nell'industria e, da lì, chiedere di
        imparare anche altre discipline.
\end{itemize}
Un gioco si può fare su qualsiasi cosa: sugli scarafaggi, sui piedini delle
sedie, sui pennarelli. Si parte avvantaggiati da ciò che è vicino alla propria
sensibilità: chi ama la musica può pensare a un gioco con meccaniche legate
alla musica, con un gameplay innovativo ma contenuto e realizzabile. Il design
è un cappello che comprende molte discipline.

% ---------------------------------------------------------------------
\section{Approfondimenti emersi in lezione}
\label{sec:approfondimenti}

\subsection{Roguelike e roguelite}

Un esempio di come il marketing conia generi: Supergiant Games, con
\emph{Hades}, rese celebre il termine \textbf{roguelite} al posto di
\emph{roguelike} (pur non essendo il primo roguelite in assoluto). La
differenza:
\begin{itemize}
  \item il \textbf{roguelike} si fonda su una serie di punti fissati in una
        sorta di \emph{manifesto} che definisce che cosa rende un gioco
        roguelike; uno su tutti: a ogni morte si ricomincia la \emph{run} da
        capo;
  \item il \textbf{roguelite} rompe alcuni di questi punti, in particolare
        introducendo una \textbf{progressione esterna alla run}: il giocatore
        può potenziarsi fuori dalla singola partita per vincere più
        facilmente in quelle successive. \emph{Hades} ne è il caso principe.
\end{itemize}
Analogamente, un mobile game sviluppato dal direttore del corso fu
promosso come ``il primo \emph{donation game} al mondo'': un gioco che
convertiva il tempo di gioco in donazioni reali a enti benefici per animali.
Anche qui un neologismo creato lato marketing per definire un genere.

\subsection{Applicazioni fuori dall'industria videoludica}

Le competenze sviluppate non servono solo ai videogiochi:
\begin{itemize}
  \item i \textbf{render architettonici} oggi possono essere realizzati con
        Unreal, e diversi professionisti formati nel game design lavorano in
        questo campo;
  \item la \textbf{gamification}, cioè l'applicazione di regole videoludiche
        a contesti non ludici (mostre interattive, progetti culturali,
        consulenze come quella per applicare questi concetti alle macchine da
        allenamento di un'azienda di fitness), è in Italia uno dei modi più
        redditizi di spendere queste competenze, spesso con finanziamenti
        pubblici.
\end{itemize}

% ---------------------------------------------------------------------
\section{Il primo task}
\label{sec:primo-task}

\begin{enumerate}
  \item Creare un indirizzo e-mail \textbf{professionale} (consigliato: un indirizzo Gmail nella forma
        \texttt{nome.cognome@gmail.com}; l'anno di nascita è accettabile).
        Va bene anche un indirizzo professionale già esistente usato per
        lavoro. Da evitare indirizzi poco seri (``non \texttt{blackfire67}'').
  \item Una persona della classe raccoglie tutti gli indirizzi in un unico
        documento e lo invia via e-mail al docente, che aggiungerà tutti al
        Drive condiviso del corso (dove verranno caricate lezioni, slide e
        materiali, per avere i documenti tutti nello stesso posto).
  \item Installare Discord e iscriversi al server del corso, dove passano le
        comunicazioni ufficiali. È consuetudine, e utile, creare anche un
        gruppo tra studenti.
\end{enumerate}
\textbf{Scadenza}: prima della lezione del venerdì.

% ---------------------------------------------------------------------
\section*{Punti chiave del capitolo}
\addcontentsline{toc}{section}{Punti chiave del capitolo}

\begin{itemize}
  \item Il corso è una \textbf{simulazione di lavoro}: ci si tratta da
        colleghi, con una gerarchia da studio di sviluppo.
  \item Si impara con l'\textbf{inquiry-based learning}: lezione, task,
        dubbi, domande, condivisione, progetti.
  \item Il fallimento affrontato con la mentalità giusta è la principale fonte
        di crescita; i feedback riguardano il lavoro, mai la persona.
  \item Il lavoro del designer è per metà creativo e per metà
        \textbf{ricerca e metodo scientifico}; il game design non è una
        scienza esatta.
  \item \textbf{Documentazione} e \textbf{pitching} sono il cuore del
        mestiere.
  \item Il designer progetta meccaniche e sistemi (partendo dalle
        \textbf{3C}), documenta, media, prototipa, confuta le idee, studia i
        competitor, bilancia e fa tuning.
  \item La leadership del designer è \textbf{al servizio} del team, non
        verticale.
  \item I limiti (tempo, tecnologia, competenze, persone, denaro) restringono
        e indirizzano la creatività.
  \item \textbf{Le idee sono nulla, l'implementazione è tutto.}
\end{itemize}

%    % ...
%
%  (Con hyperref aggiungere anche \renewcommand{\theHchapter}{z.\arabic{chapter}}
%  nella parte zero, per evitare ancore duplicate.)
%  Tutte le etichette di questo capitolo hanno il prefisso "z1:" per non
%  collidere con quelle dei capitoli del corso.
% =====================================================================

\chapter[La cultura del feedback]{Lavorare in squadra: la cultura del feedback}
\label{cap:z1-feedback}

Questo capitolo apre le \emph{lezioni zero}: incontri che precedono l'inizio
vero e proprio del corso di Game Design e che non riguardano una tecnica di
design in senso stretto, ma \textbf{la figura professionale} di chi lavora
nell'industria dei videogiochi. Prima di imparare a progettare un gioco
conviene infatti capire in che tipo di ambiente lo si progetterà: un ambiente
fatto di team, di gerarchie, di scadenze e, soprattutto, di persone che devono
riuscire a lavorare insieme.

La lezione è la prima giornata di un \textbf{workshop di due giorni} (circa
dieci ore in totale, pause pranzo comprese) tenuto da Stefano
Cozzi\footnote{Il cognome è riportato come compare nella trascrizione
automatica, che in più punti storpia i nomi propri; andrebbe verificato.
Nella stessa trascrizione lo studio per cui lavora appare come
``Stormine'', ``Stormline'' e ``Stormind'': si tratta di Stormind Games,
studio con sede a Catania, coerentemente con quanto detto in lezione.}, un
professionista con circa quindici anni di esperienza nel settore, tra
Blizzard Entertainment, Riot Games e, oggi, Stormind Games. Il workshop è
nato da un confronto con i coordinatori dell'accademia\footnote{Nella
trascrizione l'accademia è resa in modi diversi (``Hive'', ``a dive'',
``D.I.V.''); i coordinatori citati sono Enrico e Andrea, insieme al
direttore Diego Ricchiuti. Nel resto del capitolo si usa semplicemente
``l'accademia'', come negli altri capitoli.}: serviva un momento dedicato a
\emph{come ci si comporta in un ambiente professionale}, cosa che riguarda sia
il lavoro futuro sia i progetti di gruppo svolti durante il percorso di
studi, dove ci si deve già abituare a comportarsi ``come se si fosse in
un'azienda''.

La prima giornata ha un unico grande tema, sezionato in più parti:
\textbf{la cultura del feedback}. La seconda giornata (che sarà oggetto della
prossima lezione zero) sarà più pratica: \emph{time management},
\emph{brainstorming} e il ``guidare da dentro'' un team anche senza esserne
il leader (sezione~\ref{sec:z1-domani}).

Il percorso del capitolo segue quello della lezione:
\begin{enumerate}
  \item il formato del workshop e il percorso professionale del relatore
        (sezione~\ref{sec:z1-relatore});
  \item perché le aziende tengono tanto al lavoro di squadra e al
        \emph{fit} culturale (sezione~\ref{sec:z1-squadra});
  \item che cos'è il feedback e da che cosa dipende il modo di darlo
        (sezione~\ref{sec:z1-cose-feedback});
  \item le cinque fasi di vita di un team e il valore del
        \emph{post-mortem} (sezioni~\ref{sec:z1-fasi-team}
        e~\ref{sec:z1-qa-fasi});
  \item il confronto con i colleghi e alcune domande sul lavoro nei grandi
        studi (sezioni~\ref{sec:z1-rivalita} e~\ref{sec:z1-industria});
  \item il circolo virtuoso tra feedback, miglioramento e fiducia, e i suoi
        effetti sulla \emph{retention} e sulla cultura aziendale
        (sezioni~\ref{sec:z1-triangolo} e~\ref{sec:z1-cultura});
  \item il ``gioco del traditore'' e il principio del
        \textbf{default to trust} (sezioni~\ref{sec:z1-among-us}
        e~\ref{sec:z1-default-trust});
  \item le traiettorie di carriera e la valutazione delle performance
        (sezione~\ref{sec:z1-traiettorie});
  \item come si dà, come si riceve e quando si dà un feedback
        (sezioni~\ref{sec:z1-dare}, \ref{sec:z1-ricevere}
        e~\ref{sec:z1-quando});
  \item il feedback come regalo, le domande finali e le anticipazioni
        della seconda giornata (sezioni~\ref{sec:z1-regalo},
        \ref{sec:z1-qa-finale} e~\ref{sec:z1-domani}).
\end{enumerate}

Una premessa di metodo: il relatore ha alternato continuamente teoria,
domande alla classe e racconti della propria esperienza. I racconti non sono
digressioni: sono gli esempi concreti con cui i principi vengono resi
tangibili, e per questo sono riportati con cura.

% ---------------------------------------------------------------------
\section{Il workshop e il relatore}
\label{sec:z1-relatore}

\subsection{Formato e regole di partecipazione}

Il relatore ha cercato di ``riassumere quindici anni di esperienza in dieci
ore''. Ha materiale per circa il doppio delle ore disponibili: le cose
strettamente necessarie vengono trattate per prime, il resto è tenuto ``in
riserva'' come \emph{nice to have}. Le regole del workshop:
\begin{itemize}
  \item le domande si fanno \textbf{alla fine di ogni slide}, così da chiudere
        prima il discorso; non c'è un'unica sessione finale di domande e
        risposte (anche se ci sarà), perché il relatore preferisce fermarsi
        spesso e non fare ``uno spiegone di dieci ore'';
  \item le domande non devono per forza riguardare la slide appena vista: se
        una slide dà lo spunto per altro, se ne può parlare; al massimo sarà
        il relatore a rispondere in breve o a rimandare a una slide
        successiva;
  \item si può chiedere anche di esperienze passate, opinioni e sensazioni;
        se per NDA o per serietà professionale non può rispondere, lo dirà
        esplicitamente (``non lasciate che sia un \emph{blocker}'');
  \item su Zoom si usa l'alzata di mano o la chat, che il relatore tiene
        aperta sul secondo monitor;
  \item non ci sono pause ufficiali: la lezione è registrata e ognuno è
        libero di assentarsi quando serve, cercando però di non disturbare il
        flusso della lezione;
  \item le slide verranno fornite integralmente tramite la segreteria
        dell'accademia (Vanessa), insieme alle registrazioni: le slide senza
        il parlato servono poco, ma sono un buon promemoria.
\end{itemize}

\subsection{Il percorso professionale}

Il relatore si presenta (scherzosamente) come la ``guida escursionistica''
che accompagna la classe nella ``foresta incantata'' della cultura del
feedback, perché per molti è un territorio sconosciuto. Il suo percorso:
\begin{itemize}
  \item in Italia ha lavorato per aziende che fornivano servizi a terzi, tra
        cui il lavoro su \emph{Warhammer Online}, e per un canale televisivo
        sul digitale terrestre dedicato ai videogiochi\footnote{Nella
        trascrizione: ``Digital Game Network Channel''. Il nome esatto del
        canale andrebbe verificato.};
  \item a \textbf{Blizzard Entertainment} (circa sei anni, a Parigi) ha fatto
        il \emph{community manager} per l'Italia (\emph{StarCraft~II},
        \emph{World of Warcraft} da \emph{Mists of Pandaria} in poi,
        \emph{Hearthstone}, \emph{Diablo~III}, le espansioni di
        \emph{StarCraft}, il lancio di Battle.net), per poi diventare il
        manager dei community manager europei;
  \item a \textbf{Riot Games} (circa cinque anni, a Dublino) è stato
        brevemente community manager italiano e poi \emph{publishing
        manager} di un team di circa trenta persone per l'Europa. Il
        \emph{publishing} non è solo community: è l'unione di
        \textbf{community, marketing e PR};
  \item a \textbf{PG Esports} si è occupato di eventi dal vivo e online
        (ad esempio la Milano Games Week e tornei di \emph{Valorant} e
        \emph{League of Legends}), restando in contatto con Riot e Blizzard;
  \item da circa un anno è tornato nello sviluppo, a \textbf{Stormind
        Games}, dove oggi si occupa di risorse umane e organizzazione del
        lavoro; lo studio verrà presentato più ampiamente al \emph{Career
        Day} dell'accademia insieme alla sua HR manager;
  \item fa inoltre consulenze aziendali anche al di fuori dei videogiochi e
        insegna in master universitari (IULM di Milano).
\end{itemize}

Il suo ruolo ricorrente, sia a Blizzard sia a Riot, è stato quello di
\emph{Mr.~Wolf} (``per chi ha visto \emph{Pulp Fiction}''): quello che
risolve i problemi. Veniva inserito nei team più in difficoltà, ad esempio
dal punto di vista delle performance, per sistemarli e passare al successivo.
Era presente in alcuni dei momenti più critici, come il lancio di
\emph{Diablo~III} o il lancio di espansioni di \emph{World of
Warcraft}\footnote{La trascrizione riporta ``il lancio con i Gary's'' e
``World of Trenor'': si tratta verosimilmente di \emph{Mists of Pandaria} e
di \emph{Warlords of Draenor}.}, a fare \emph{damage control}.

La scoperta fondamentale di questi anni è che i problemi da risolvere
\textbf{quasi mai riguardavano il prodotto} (non era certo lui a sistemare i
server di \emph{Diablo~III}), ma i team che c'erano dietro: persone che
litigavano, non collaboravano, non riuscivano a stare nella stessa stanza
senza urlarsi contro. E nel 90\% dei casi tutto questo si riduce a
\textbf{problemi di comunicazione}: è raro che perfino le incompatibilità
caratteriali non siano risolvibili parlando.

\subsection{Struttura non vuol dire burocrazia}

A Stormind il relatore sta accompagnando un team passato da 30 a 95 persone
in meno di un anno, attraverso tutti i cambiamenti di cui parla il workshop.
Darsi una struttura viene spesso percepito male, come ``più processi, più
procedure'', ma bisogna distinguere:
\begin{itemize}
  \item un processo come ``timbrare il cartellino'' può essere davvero
        invasivo e noioso;
  \item definire \emph{come funziona un documento di game design, chi lo
        approva, chi lo visiona e chi no} è invece una struttura necessaria.
\end{itemize}
Uno studio di 30 persone nella stessa sede può farne a meno: quando c'è un
dubbio ci si alza e si va alla scrivania del collega. Uno studio di 95
persone sparse per l'Italia (e in parte per l'Europa) non può più farlo:
servono una struttura diversa, un modo di lavorare diverso e una
pianificazione del futuro diversa.

% ---------------------------------------------------------------------
\section{Perché il lavoro è sempre lavoro di squadra}
\label{sec:z1-squadra}

\subsection{Non si va in ufficio per fare amicizia}

Il primo messaggio è netto: \textbf{non si va in ufficio per fare amici} o
per trovare compagni di vita. Si va in ufficio perché c'è un lavoro da fare, e
lo scopo è riuscire a farlo \emph{anche con chi sta antipatico}. Con le
persone simpatiche è facile: diventeranno amici, e non sarà mai un problema.
È con chi sta antipatico che le cose si complicano, ed è lì che emergono i
\textbf{veri professionisti}: quelli che riescono a lavorare con chiunque
senza che l'antipatia pesi sul modo in cui lavorano e collaborano, mettendo il
prodotto e il team davanti al proprio ego e alle proprie preferenze personali.

La frase che fa da filo conduttore a tutto il workshop (e che il relatore
promette di riproporre ``fino alla nausea'') è:

\begin{quote}
  \emph{La capacità di un gruppo di persone di fare cose straordinarie
  dipende dal modo in cui queste persone operano insieme.}
\end{quote}

Le aziende raramente pensano che l'innesto di una singola persona possa
cambiare il destino di un progetto. Questo è tanto più vero quanto più grande
è il team, ma vale anche per i piccoli studi, anche per un gruppo di cinque
persone. Come negli sport di squadra: l'individuo conta fino a un certo punto,
e anche se si ha Dončić in squadra, se si gioca uno contro cinque si perde
comunque\footnote{Il riferimento è alle finali NBA in corso nei giorni del
workshop tra Dallas Mavericks e Boston Celtics (giugno 2024), il che
permette di datare approssimativamente la lezione.}.

\subsection{Lo \emph{useless genius} e il costo delle persone tossiche}

Molti dipartimenti HR usano la definizione di \textbf{useless
genius}\footnote{Il concetto è più noto in letteratura manageriale come
\emph{brilliant jerk}; in lezione è usata l'espressione \emph{useless
genius}.}: si può essere bravissimi, a volte i migliori nel proprio lavoro,
ma se non ci si sa comportare, relazionare e operare in un gruppo si è
fondamentalmente inutili ai fini del gruppo. Nel migliore dei casi inutili,
nel peggiore \emph{dannosi}.

In passato le aziende tendevano a ``chiudere un occhio'': la persona è
fondamentale, la teniamo. Oggi la domanda è diventata: \emph{quanto costa
all'azienda, oltre al suo stipendio, tenere una persona che crea un ambiente
tossico e peggiora chi le sta intorno?} Ci sono costi sommersi che le aziende
cominciano a vedere, e la tendenza è liberarsi di queste persone il prima
possibile. Il ragionamento è semplice:
\begin{itemize}
  \item le \textbf{competenze si possono trasmettere} o sostituire: si trova
        qualcuno che fa lo stesso lavoro magari nel doppio del tempo, oppure
        più persone che se lo dividono;
  \item il \textbf{carattere delle persone} invece cambia molto raramente.
\end{itemize}

\subsection{Il fit culturale}

Quando qualcuno \textbf{non è un \emph{fit} culturale} per un'azienda è
meglio separarsi presto. Ma non essere adatti a \emph{quell'}azienda non
significa non essere adatti a lavorare: è perfettamente legittimo che una
persona inadatta a un'azienda sia adattissima a un'altra.

Il relatore l'ha vissuto in prima persona. Dopo sei anni a Blizzard viene
contattato da Riot e affronta un processo di selezione lunghissimo (``quasi
undici colloqui''), al termine del quale viene scartato. Scrive allora una
mail che gli ha ``cambiato la traiettoria di carriera'': chiede
semplicemente perché non è stato ritenuto adatto. La risposta: in nessun
colloquio si era esposto abbastanza da far capire \emph{perché} volesse
lasciare Blizzard, e quindi pensavano che in realtà stesse bene dov'era. Il
motivo era che, per educazione professionale, non gli piace parlare male del
proprio posto di lavoro; portata all'eccesso, però, questa riservatezza era
diventata un ``è tutto bellissimo, ma vorrei provare Riot'', senza un perché.
Lo invitano comunque a un incontro dal vivo (``dal vivo è diverso''), e alla
fine viene assunto.

Una volta arrivato, i sei mesi di prova (\emph{probation}) sono serviti tutti
per smettere di ragionare e lavorare ``come a Blizzard''. La differenza
culturale era enorme, tanto che in Riot esisteva un gruppo di supporto
interno tra ex dipendenti Blizzard: un collega gliel'ha descritto come
``uscire da una relazione tossica e doversi riabituare a una relazione
normale''\footnote{Il relatore precisa che parla della Blizzard di allora,
che non esiste più: oggi è parte di Microsoft e nessuno dei suoi ex colleghi
vi lavora ancora, anche a causa dei licenziamenti di massa.}.
\begin{description}
  \item[Blizzard] era \emph{developer first}: il focus era chiaramente sugli
        sviluppatori, prima di tutti gli altri. Dopo sei anni c'erano ancora
        meeting a cui non poteva partecipare perché non abbastanza
        \emph{senior}, informazioni che non poteva sapere, piani che non
        poteva vedere.
  \item[Riot] era \emph{player first}: prima i giocatori, e i dipendenti come
        mezzo per trasferire il feedback dei giocatori dentro il gioco. Il
        secondo giorno lo portano in una stanza e gli mostrano il piano dei
        cinque anni successivi. Lui, sorpreso, pensa a un errore (``sono qui
        da oggi, non vorrei mettervi nei guai''); la risposta è: ``sei un
        dipendente Riot o no? Allora non ci sono livelli di accesso: devi
        sapere queste cose per fare meglio il tuo lavoro''.
\end{description}

Da questa esperienza ha imparato molto sulla \textbf{trasparenza}: condividere
le informazioni aiuta il lavoro. Non sempre, però: a volte non si condivide
qualcosa non per segretezza, ma perché non è ancora pronta per essere
presentata, esattamente come non si mostra un gioco quando è ancora solo
teoria (``te lo faccio vedere quando ho qualcosa di tangibile'').

La lezione per gli studenti è doppia: le aziende valuteranno il \emph{fit}
culturale sempre di più, ma capire se si è un \emph{fit} richiede anche
\textbf{lavoro da parte del candidato}. Bisogna avere la presenza di spirito
di dire: ``conosco i valori di questa azienda, non mando nemmeno il CV,
perché so che durerei un mese e mezzo''.

% ---------------------------------------------------------------------
\section{Che cos'è il feedback}
\label{sec:z1-cose-feedback}

\subsection{Le prime risposte della classe}

La prima domanda alla classe è: \emph{quando sentite la parola
``feedback'', a cosa pensate?} Le risposte raccolte, a voce e in chat:
\begin{itemize}
  \item un \textbf{riscontro} su ciò che si fa e su come lo si fa, sul lavoro
        o sul modo in cui ci si interfaccia con il gruppo; può essere
        lavorativo oppure comunicativo e comportamentale. Chi ha risposto lo
        vede come qualcosa che, se l'azienda non dà, va attivamente
        richiesto, perché è importante sapere come sta andando per
        migliorarsi;
  \item né positivo né negativo, ma \textbf{costruttivo}: in entrambi i casi
        serve a migliorarsi;
  \item uno \emph{scambio di opinioni}, un \emph{riscontro produttivo};
  \item un feedback ``di reazione'', non verbale: si fa una cosa e si
        osserva come cambia l'ambiente o il comportamento delle persone
        intorno (anche perché a voce la gente spesso si tiene le cose
        negative);
  \item per un designer, anche il \textbf{linguaggio del corpo}: osservando
        un giocatore giocare si può dedurre l'emozione che prova (frustrazione,
        incomprensione).
\end{itemize}

Sulla forma (scritto o orale?) le opinioni divergono: c'è chi lo preferisce
\emph{mille volte} scritto; chi pensa che dipenda dalle dimensioni
dell'azienda (orale nelle piccole, scritto in quelle con team di centinaia di
persone); chi pensa che dipenda dall'importanza del feedback (scritto se
tecnico o pieno di dati, da tenere sotto gli occhi; orale se è un consiglio
veloce).

Il relatore usa due esempi per mostrare che \emph{dipende}:
\begin{itemize}
  \item un feedback sulla programmazione ha senso darlo scritto, magari
        commentando direttamente il codice, più che ``andando a prendere un
        caffè'';
  \item ma se il feedback è ``ieri, durante il meeting, ho percepito male il
        fatto che mi interrompessi di continuo'', ce lo si aspetta
        \textbf{orale}.
\end{itemize}

\subsection{Da che cosa dipende il modo di dare feedback}

Lo \emph{spoiler} della giornata è che \textbf{non esiste una risposta univoca}
su come dare feedback. Il workshop non insegna ``il'' modo, ma fornisce
\emph{tutte le opzioni}, perché il modo giusto cambia in base a molti fattori:
\begin{itemize}
  \item a chi lo si dà: un collega, un superiore, il proprio capo, un proprio
        \emph{report} (un sottoposto);
  \item se è personale o professionale;
  \item se è la prima volta o la terza volta che si dà lo stesso feedback;
  \item se riguarda un progetto importante o secondario.
\end{itemize}

Manca però l'elemento fondamentale, che la classe individua subito:
\textbf{l'altro}. A parità di tutti gli altri fattori, il \emph{carattere
della persona che riceve il feedback} dovrebbe far cambiare il modo in cui lo
si dà. Per questo, prima della lezione, agli studenti era stato chiesto di
fare un test di personalità: non un test di psicologia avanzata, ma uno
strumento per rendersi conto di quanta varietà ci sia anche solo in una
classe di 53 persone (il che non vuol dire trovare 52 modi diversi di dare
feedback, ma lavorarci un po' sì).

Il relatore porta se stesso come esempio. Preferisce:
\begin{itemize}
  \item i feedback \textbf{positivi scritti}, perché è a disagio a ricevere
        complimenti in pubblico;
  \item i feedback sul proprio lavoro \textbf{diretti}: ``vai subito al
        punto e non girarmi intorno''.
\end{itemize}
È molto orgoglioso del suo lavoro e, quando qualcuno lo critica, la sua
primissima reazione istintiva è ``come ti permetti''. Negli anni ha imparato
a farla durare un secondo, ma chi gli dà feedback di persona dice di
vederla ancora negli occhi. Per mitigarla, dice esplicitamente a chi gli dà
feedback: \emph{``ignora il modo in cui reagisco nei primi cinque secondi,
vai avanti, poi mi passa''}. È considerato da tutti molto aperto al feedback
e non ha problemi a dire ``hai ragione, facciamo diversamente'': ma ci sono
reazioni istintive e caratteriali che è difficile controllare al cento per
cento.

\subsection{Con chi si lavora meglio}

Guardando le persone con cui si lavora abitualmente (anche nei progetti
dell'accademia), si noterà che si lavora meglio con due profili:
\begin{enumerate}
  \item chi è \textbf{esattamente come noi}, ma qui è un 50 e 50: alcuni
        caratteri vanno d'accordo con i propri simili, altri non vogliono
        nessuno con lo stesso carattere nel gruppo (e si potrebbe rientrare
        in questa categoria);
  \item chi è \textbf{affine}: ha qualcosa in comune con noi, ma non tutto,
        e proprio per questo può dare punti di vista differenti.
\end{enumerate}
Il lavoro della giornata è capire come valorizzare il secondo caso: ottenere
feedback da persone con caratteri, visioni e background diversi, per
completare una visione del proprio lavoro che è necessariamente limitata. È
come esplorare una mappa con la nebbia di guerra: un'altra persona, al di là
del muro, vede cose che noi non vediamo.

\subsection{Il feedback come somma di piccoli miglioramenti}

Per il relatore il feedback è questo: \textbf{applicare tanti piccoli
miglioramenti periodici}, che si sommano nel tempo. Migliorare dello 0,1\%
ogni giorno porta a fine anno a un miglioramento enorme, e questo vale per i
prodotti, per i giochi, per le persone e per i team.

\paragraph{L'esempio del ciclismo britannico.} La squadra di ciclismo
britannica, fino a una quindicina d'anni fa, non aveva vinto quasi nulla.
Cambiato il team manager, tutti si aspettavano grandi cambiamenti sportivi
(convocazioni, strategie di gara). Arrivarono invece \emph{minuscoli
cambiamenti non sportivi}: un altro chef in trasferta, materassi migliori nei
camper per riposare meglio, un sellino più comodo a fine tappa. Cambiamenti
marginali che da soli non avrebbero cambiato niente; eppure, dopo due o tre
anni di questo lavoro quotidiano, la squadra ha cominciato a vincere a
ripetizione, Olimpiadi comprese, cambiando il modo in cui il ciclismo è
visto in Gran Bretagna\footnote{È la nota strategia delle \emph{marginal
gains} associata a Dave Brailsford (British Cycling e Team Sky); il nome non
è citato in lezione.}. Certo, si è investito anche sulle persone giuste, ma
tutto è partito dai piccoli miglioramenti periodici.

Per gli studenti vale lo stesso: chiarire un conflitto con una persona con
cui si ha ancora poco a che fare, o una situazione che ha lasciato
``l'amaro in bocca'', sembra inutile oggi, ma darà i suoi frutti tra tre, sei,
nove mesi. Sono \emph{investimenti sul futuro}.

\paragraph{Peggiorare ogni giorno.} Esiste anche il percorso inverso:
lasciar correre, ``buttare sotto il tappeto'' le cose perché ``tanto è
inutile parlarne''. Queste cose tornano sempre fuori dopo otto mesi, come il
\emph{libro dei rancori} dei nani di Warhammer: ``ti ricordi quando, un anno
e mezzo fa, all'inizio del progetto, mi hai detto quella cosa? Mi sono
offeso''. A quel punto è quasi impossibile risolvere, perché nessuno ricorda
più con precisione cosa sia successo (``ma quale cena?'', ``forse era in
pizzeria'', ``no, eravamo da un'altra parte''). \textbf{Tenersi le cose nel
cassetto per tirarle fuori quando servono è sempre sconsigliato}: non servono
ed è impossibile usarle in modo costruttivo.

% ---------------------------------------------------------------------
\section{Le cinque fasi di vita di un team}
\label{sec:z1-fasi-team}

\subsection{Perché conoscerle}

Quasi tutti gli studenti lavorano già in gruppo (al primo anno su un gioco da
tavolo, negli anni successivi su un videogioco) e si saranno accorti che un
team attraversa delle fasi: il modo di lavorare insieme cambia nel tempo.
Spesso ci si fa prendere di sorpresa: ``eravamo così affiatati, cosa è
successo?''.

Un esempio da una consulenza fuori dai videogiochi: un'azienda farmaceutica
gli affida un team commerciale (rappresentanti che presentano medicinali a
ospedali, medici e privati). Il problema: ``prima andava tutto bene; ora che
siamo in diciotto dovremmo lavorare meno e meglio, e invece va peggio''. Va
peggio perché le persone nuove sono appena arrivate, mentre il management si
aspettava già di vedere i risultati dell'investimento.

Fino a pochi anni fa le fasi di un team venivano disegnate come una
\textbf{linea dritta verso l'alto}: le fasi erano le stesse, ma viste come un
miglioramento costante. Non si vedevano le parti \emph{fisiologiche} del
lavoro in team, cioè anche i peggioramenti. Accettarle è il primo passo per
accettare i fallimenti: nessun percorso personale o di team è lineare.

Il relatore ammette che il suo modo di lavorare è nato per il 70\% dal
proprio impegno a imparare, fare bene e rispettare le persone, e per il 30\%
dall'aver lavorato con ``incapaci colossali'' che gli hanno insegnato come
\emph{non} fare le cose. A Blizzard un manager gli fece capire che, pur non
sapendo ancora che tipo di \emph{people manager} sarebbe diventato, sapeva
benissimo quale tipo di manager \emph{non} voleva essere.

\subsection{Forming, storming, norming, performing, adjourning}

Le fasi sono cinque, un po' come le fasi del lutto. Molte aziende ne
considerano solo quattro (``e sono quelle che solitamente chiudono''); le
aziende al passo con i tempi ne considerano cinque\footnote{Il modello è
quello di Bruce Tuckman (1965), a cui la quinta fase, \emph{adjourning}, fu
aggiunta nel 1977. Il nome di Tuckman non è citato in lezione; la
trascrizione riporta la quinta fase come ``journing''.}.

\begin{description}
  \item[Forming] Si viene messi insieme e si comincia a lavorare. Non si sa
        nulla degli altri; spesso non si conoscono nemmeno i nomi, ed è
        normale.
  \item[Storming] Si cominciano a ``dare le testate come le capre in
        montagna per determinare il territorio'': questa cosa la faccio io,
        questa la vorrei fare così, chi decide? Si stabilisce chi fa cosa e
        come si divide il lavoro. Nel grafico è una \textbf{fase di
        peggioramento}: che il team litighi o discuta civilmente non cambia
        nulla, il team deve investire parte di tempo ed energie per risolvere
        queste questioni. È normale e \textbf{non è evitabile}.
  \item[Norming] Si comincia a capire come lavorano gli altri, cosa sanno
        fare e come far funzionare il gruppo; iniziano i miglioramenti. È una
        forma di \emph{accettazione} degli altri, che deriva dall'aver capito
        cosa devo fare io e cosa devono fare loro.
  \item[Performing] C'è uno sforzo \textbf{proattivo}: ``questa cosa
        toccherebbe a me, ma so che Alessio la fa molto meglio, quindi gli
        chiedo aiuto''. Si sfruttano i punti di forza degli altri, perché li
        si conosce e perché si sa che passare il testimone non avrà effetti
        negativi (il capo non dirà ``perché non l'hai fatta tu?'').
  \item[Adjourning] Il progetto finisce, si chiude o viene riorganizzato, il
        team viene sciolto. Quello che andrebbe fatto è analizzare
        \textbf{cosa ha funzionato, cosa no e cosa migliorare} per il
        progetto successivo.
\end{description}

\subsection{La fase che tutti saltano}

``Spoiler incredibile'': l'\emph{adjourning} è la fase che molte aziende
saltano, \textbf{perché c'è sempre il progetto dopo}.

A PG Esports, per la Milano Games Week (come per qualsiasi evento dal vivo a
cadenza annuale: Lucca Comics, Etna Comics, Comicon\dots), finita l'edizione
di un anno si cominciava dal giorno dopo a lavorare a quella successiva,
perché un anno sembra tanto ma non basta mai. Nei videogiochi è lo stesso:
finito un gioco, a parte le patch di mantenimento e del \emph{day one}, si
passa subito al successivo.

Il cambiamento introdotto dal relatore (in parte già praticato, ma da lui
strutturato) è stato: finita la Games Week, \textbf{tutti in ferie per una
settimana}; poi ci si ritrova per definire cosa ha funzionato e cosa no, con
un \textbf{post-mortem} (o \emph{debrief}); solo dopo si inizia a lavorare
all'edizione successiva. Senza questa fase manca la parte di
\emph{learning}, cioè di ciò che si è imparato dal progetto.

Il consiglio vale per gli studenti: quando si finisce il gioco da tavolo, e
poi il gioco digitale, bisogna trovare il tempo per il post-mortem,
altrimenti si butta via \textbf{il 90\% dell'utilità del progetto}. Una
struttura semplice per farlo è descritta nella sezione~\ref{sec:z1-qa-finale}.

% ---------------------------------------------------------------------
\section{Domande sulle fasi del team}
\label{sec:z1-qa-fasi}

\subsection{Lo storming che ritorna}

\paragraph{Capita che lo storming si ritagli spazio anche nelle altre fasi?}
Lo \emph{storming} è una fase precisa, che però molti manager, coordinatori e
team cercano di \textbf{saltare}: ``se io sono bravo come manager e loro sono
bravi come team, passeremo da forming a norming senza problemi''. È un errore,
perché senza una fase in cui si discute davvero non si sviluppa una
\emph{skill} fondamentale: \textbf{saper discutere}. Tutti conoscono persone
non abituate a farlo, per cui ogni minimo disaccordo diventa un affronto
personale (``avevamo deciso la pizza, perché andiamo al sushi? Perché mi fate
questo affronto?''): sono persone che non sono abituate ad avere conflitti
costruttivi.

Se si pensa di essere in \emph{norming} e continuano a emergere conflitti, il
punto è che \textbf{non si era davvero in norming}: lo storming non ``si
trascina'' nelle fasi successive, semplicemente non se n'è mai usciti.
Allo stesso modo, a chi chiede se ogni mini-fase di storming durante la vita
del gruppo sia frutto di uno storming fatto male o solo di carattere, la
risposta è già contenuta qui.

Un partecipante osserva che spesso, ancor prima di finire un progetto, si è
già messi su quello successivo: è il modo migliore per essere in costante
\emph{damage control}. Era l'esperienza del relatore a Blizzard: veniva
spostato fisicamente da un team all'altro, e finito il lancio di
\emph{Diablo~III} era già su \emph{World of Warcraft}, senza vivere il
``dopo''. Si è dovuto ritagliare da solo lo spazio per un debrief personale,
ed è lì che ha deciso che tipo di manager non voleva essere. Il consiglio:
\textbf{se non lo fa il team, fatelo voi}.

\subsection{Dove si passa la maggior parte del tempo}

\emph{In quale fase si passa più tempo nella propria carriera?} Gli studenti
rispondono \emph{storming} e \emph{norming}. Il relatore, in quindici anni,
calcola di aver passato in \emph{performing} circa \textbf{dieci mesi}. Non è
pessimismo: arrivare a performing deve essere sempre l'obiettivo, sia da
manager sia da membro del team; ma è difficile arrivarci, ed è normale che non
sempre ci si arrivi.

Lo storming può durare a lungo, ma con team molto esperti dura pochissimo,
perché c'è chi lo sa gestire. In realtà \textbf{la maggior parte del tempo si
passa in forming}, perché ogni equilibrio raggiunto vale per \emph{quelle}
persone, \emph{quel} progetto, \emph{quelle} tempistiche. Basta che cambi un
qualsiasi fattore (una deadline spostata, una persona nuova, qualcuno che si
dimette, un progetto cancellato) e si ricomincia, perché i bilanciamenti
precedenti non valgono più.

Esempio: un team di cinque persone affiatatissimo, in pieno performing. La
dirigenza, entusiasta, ``raddoppia le risorse'' e lo porta a dieci. O si
lascia tutto com'è (e i cinque nuovi non fanno niente), o bisogna ridiscutere
chi fa cosa e come i nuovi possono aiutare: che piaccia o no, si ricomincia
da forming.

Quindi la bravura di un manager, ma anche di un membro del team,
\textbf{non sta nel passare tanto tempo in performing, ma nell'accelerare il
passaggio da una fase all'altra}, e nel fare bene la fase di
\emph{adjourning}.

\subsection{Analizzare anche i successi}

C'è la tendenza a fare post-mortem solo dei progetti andati male: ``questa
cosa non ha funzionato, vediamo perché''. Va benissimo. Ma quando un progetto
va bene si dà per scontato che sia andato bene ``perché era buono''. Bisogna
invece chiedersi \emph{perché} è andato bene: era il prodotto? La
tempistica? Il modo in cui abbiamo fatto una cosa piuttosto che un'altra?
\textbf{Se non si capisce davvero il motivo del proprio successo, si
tenderanno a replicare le cose sbagliate.}

\paragraph{Blockbuster.} Quando arrivò la notizia che Netflix avrebbe
lanciato un servizio di streaming, Blockbuster (la catena di noleggio di
videocassette e DVD) aveva tutto il tempo per preparare un'offerta
concorrente. In una riunione del management raccontata in un libro (di cui il
relatore non ricorda il titolo), la direzione sembrava orientata a investire
nello streaming, finché un dirigente fece notare che il 12\% del fatturato
derivava dalle \emph{multe per riconsegna tardiva}, impossibili con lo
streaming, mentre secondo le stime lo streaming avrebbe prodotto il 10\% del
fatturato. ``Non possiamo rinunciare a questo 2\%'': niente streaming. Oggi
non c'è più un Blockbuster ``manco a cercarlo col lanternino''. Non avevano
capito qual era il vero motivo del loro successo, cioè il loro modello di
business\footnote{Le percentuali sono quelle riferite in lezione e non sono
state verificate.}.

\paragraph{The Founder.} Nel film sulla nascita di McDonald's c'è una scena in
cui il futuro socio del protagonista gli dice: ``tu pensi di essere nella
\emph{food industry}, e invece sei nella \emph{real estate industry}''. Non
dovrebbe investire sugli hamburger, ma sui terreni su cui costruisce i
ristoranti. Anche per un team vale lo stesso: bisogna capire che cosa
l'ha reso davvero di successo, oppure no.

\subsection{L'autoanalisi personale}

\paragraph{E quando si cambia azienda?} Uno studente chiede se il debrief sia
possibile anche quando si lascia un'azienda. Non si può tornare in ufficio a
discutere della propria esperienza, ma lo si può fare \textbf{per se stessi}.
Il relatore l'ha fatto, anche se all'epoca in modo poco consapevole, ed è
così che è passato dal \emph{community management} al \emph{people
management}. Analizzando cosa aveva fatto bene a Blizzard si era reso conto
che:
\begin{itemize}
  \item aveva costruito un ottimo rapporto di fiducia con i giocatori (con
        alcuni proprietari di fansite parla ancora oggi);
  \item pur senza la seniority, era diventato un punto di riferimento per il
        team, visto come imparziale ed equilibrato anche su temi che non lo
        riguardavano direttamente;
  \item sapeva \textbf{risolvere problemi}: durante il lancio di
        \emph{Diablo~III} era il più tranquillo di tutti, non perché
        sminuisse il problema, ma perché sapeva che si sarebbe risolto con i
        suoi tempi; il lavoro non era spegnere l'incendio ma ``far sfollare
        la gente'' e aspettare che tutto tornasse in ordine.
\end{itemize}
E aveva capito cosa sapeva fare meno bene, ad esempio \textbf{gestire il
proprio manager}. Avere un manager non significa che il manager gestisca te:
il manager gestisce il tuo lavoro, ma tu devi saper gestire i tuoi
\emph{stakeholder}. Spesso lo stakeholder è il proprio manager; da manager,
gli stakeholder saranno CEO, CFO e così via. Gestirli significa conoscerli,
sapere come portare loro le informazioni, come far presente che qualcosa non
va. Fa parte del lavoro, anche se a lui non piace, e ci ha lavorato.

Nel mese e mezzo di pausa prima di entrare in Riot si è concentrato su
questo: ``cosa porto come mio punto di forza?''. L'autoanalisi è culminata
nel chiedere espressamente alla persona che l'aveva assunto: ``prima mi avete
detto di no, ora mi avete preso: perché?''.

\paragraph{Come si capisce se un progetto è andato bene?} Dipende
dall'\textbf{obiettivo}. Se l'obiettivo era vendere milioni di copie e il
gioco le ha vendute, bisogna capire perché. Se, più ragionevolmente,
l'obiettivo era un gioco da tavolo che permettesse di mettere in pratica le
conoscenze del corso, va valutato per quello: ha permesso di farlo? Se sì,
perché? Se no, perché no? L'obiettivo va fissato \textbf{all'inizio del
progetto, non alla fine}, proprio per poterlo valutare dopo.

\subsection{Quando il team cambia di continuo}

\paragraph{Tutti questi cambiamenti non sono dannosi per il progetto?}
Assolutamente sì, ma quasi mai sono volontari. I progetti andati male sono
spesso progetti grandi.

Il relatore racconta di \emph{Cyberpunk 2077}. Conosceva persone in CD
Projekt Red (``il mondo è piccolissimo, le cose si sanno''). Alla Gamescom va
con Diego Ricchiuti, direttore dell'accademia, a vedere la presentazione del
primo gameplay, quella ``in cui regalavano il giubbottino da samurai''.
Finito il video, tra gli applausi, si guardano e capiscono entrambi che
\textbf{il gioco non è pronto}, a soli quattro mesi dall'uscita prevista:
Diego aveva notato il combattimento e la guida. Gli amici di CD Projekt Red
(uno aveva lavorato con lui a Blizzard per sei anni) gli confermano tra le
righe che ``ce la faremo, ma ci sono ancora cose da sistemare''. Il gioco fu
poi rimandato più volte\footnote{In lezione il relatore dice ``due, forse
tre volte''; l'uscita di \emph{Cyberpunk 2077} è stata effettivamente
rinviata tre volte prima del dicembre 2020.}.

Per deformazione professionale, chi ha sviluppato l'occhio critico non
riesce più a godersi un trailer: in un trailer di \emph{The Day Before} Diego
gli aveva fatto notare che l'effetto della fiammata di un fucile era stato
chiaramente aggiunto in postproduzione e non faceva parte del gameplay. ``È
il triste epilogo di questa storia'': non riescono più a godersi un trailer.

Il punto è che più il team è grande, più queste variabili vengono imposte, e
più è facile che il fallimento \textbf{non dipenda dal fatto che il prodotto
sia brutto} (\emph{Cyberpunk} è un bellissimo gioco sotto molti punti di
vista, e il relatore lo sta rigiocando ora che è stato corretto ed espanso).
Nei grandi gruppi qualcuno arriva, qualcuno si dimette, cambia il budget,
arriva il rinvio; nel caso di \emph{Anthem}, il team ha fatto così tante volte
forming--storming--forming--storming che ne è uscito ``un prodotto che
dovevamo espellere: l'abbiamo espulso, buttato fuori dal nostro organismo pur
di liberarcene''.

Tra i cambiamenti c'è anche la crescita verticale: se il capo viene promosso
perché lavora benissimo, bisogna scegliere un nuovo capo, ed è sempre un
caos (chi ci resta male, chi pensava di meritarlo, chi voleva il suo amico,
chi non sopporta il prescelto). Per questo il relatore è \textbf{contrario
alle promozioni durante i progetti}: si rischia di rovinare tutto il lavoro di
storming.

\subsection{Obiettivi chiari nei piccoli gruppi}

Nei gruppi piccoli, come quelli degli studenti, il problema deriva più
spesso dal \textbf{non avere chiaro che cosa si vuole ottenere}: un prodotto
puramente educativo e formativo, o un prodotto richiesto dall'accademia che,
se viene bene, si potrebbe anche vendere? Un gioco commerciale o un gioco
artistico? Va chiarito subito, perché ci saranno sempre momenti in cui
scegliere tra A e B, e a guidare la scelta sarà la domanda: \emph{quale delle
due opzioni ci porta più vicino all'obiettivo?} Se l'obiettivo non è chiaro,
non si sa cosa scegliere.

Questo porta molte aziende a fare \textbf{scelte di design che in realtà sono
scelte commerciali}. Le microtransazioni sono scelte commerciali: nessuno ha
mai pensato ``sai cosa piacerebbe tanto ai giocatori? Pagare le bustine''.
Prima le bustine si guadagnavano giocando; ora la tendenza è ``potresti
giocare, ma se paghi è meglio''. Si pensi alle polemiche su Ubisoft per
l'opzione di ``saltare il gioco'' pagando, o ai famosi \emph{boost} di livello
di \emph{World of Warcraft}: ``devo pagarti per saltare il gioco?''.

Una studentessa (Silvia) cita \emph{Final Fantasy XIV}, che a un certo punto fu
``resettato'' perché non funzionava: un esempio positivo. Il relatore aggiunge
quello negativo di \emph{Age of Conan}, che non è mai riuscito a fare quel
reset iniziale (``è una schifezza, ve lo ripresentiamo quando è pronto''). Ci
vuole consapevolezza, ma anche i mezzi: un conto è Square Enix, che può
permetterselo, un altro è un'azienda che non può chiudere tutto per un anno
continuando a pagare gli stipendi.

Riassumendo: nei piccoli gruppi i problemi sono più di
\textbf{comunicazione e chiarezza} (obiettivi non chiari o che si spostano);
nelle grandi aziende di più di \textbf{persone che vanno e vengono, obiettivi
e tempistiche che cambiano}. In accademia le scadenze sono fisse (il gioco da
tavolo si consegna a fine anno scolastico, a giugno, e giugno resta); in
azienda può succedere che una feature prevista per la \emph{vertical slice}
di luglio venga anticipata a maggio per portarla alla Gamescom, e un sistema
perfetto per consegnare a giugno può non esserlo più per consegnare a maggio.

\subsection{L'ambiente determina la qualità del prodotto}

\paragraph{È vero che chi passa per Electronic Arts finisce per fare brutti
giochi?} Uno studente cita BioWare, che faceva ottimi giochi finché era
autonoma. Il problema non è dei developer, che non sono ``improvvisamente
diventati incapaci'': è evidente che il problema è altrove, e che
\textbf{l'ambiente in cui si fa un gioco determina molto spesso la qualità del
prodotto}.

Un esempio locale positivo: Tommaso Bonanni, conosciuto dal relatore insieme a
Diego e già docente dell'accademia, ha fondato a Roma Caracal Games Studio,
applicando fin dall'inizio questi concetti di cultura aziendale. L'azienda è
rimasta piccola, ma ha molto chiaro cosa vuole fare e come. Quando si
predispone fin dall'inizio l'azienda e il team a lavorare in un certo modo è
fattibilissimo; è molto più difficile quando si è EA, si acquisisce
un'azienda e le si dice ``da oggi si fa così''.

Il relatore ha vissuto la transizione da Blizzard ad Activision Blizzard (uno
dei motivi per cui se n'è andato, volontariamente). Ha visto la cultura
interna cambiare per gradi:
\begin{enumerate}
  \item ``facciamo giochi fighissimi, perché i giochi devono essere
        fighissimi'';
  \item ``facciamo giochi fighi così vendono di più'';
  \item ``vendiamo di più facendo giochi fighi'' (sembra uguale, ma è molto
        diverso);
  \item ``facciamo più soldi''.
\end{enumerate}
Lo si capiva dalle scelte di design, osservando quale fosse il peso più
grande sulla bilancia. Lo stesso varrà per gli studenti, nei loro giochi e nei
colloqui di lavoro: capire subito che cosa vuole l'altra parte.

\paragraph{Leader e visione.} Nei gruppi dell'accademia, dal secondo anno, c'è
una persona che assume il ruolo di \emph{lead} o \emph{scrum master}. È
diverso dall'essere un \emph{game director} (``il gioco si fa come dico io,
aiutatemi a farlo con la mia visione''), come avviene nei grandi studi: da lì
nascono i nomi nei titoli, come \emph{Sid Meier's Civilization}, che indicano
la persona che ha fatto quel gioco, oppure un puro branding, come \emph{Tom
Clancy's Rainbow Six}.

\paragraph{Inclusion come jolly.} A margine, uno studente osserva che in molte
aziende ``la risposta giusta è sempre \emph{inclusion}'', usata come un
jolly. È comunissimo, conferma il relatore: \emph{diversity and inclusion}
viene spesso usata per dire ``così includiamo più clienti'', ma non è quello
il suo punto (``altro discorso'').

% ---------------------------------------------------------------------
\section{Confrontarsi con chi è più bravo}
\label{sec:z1-rivalita}

Arrivato a Riot dopo sei anni a Blizzard, il relatore pensava di avere un buon
bagaglio di esperienza. Si è trovato invece ``come un fante buttato in mezzo
ai Navy SEAL'': persone che avevano superato 11, 12, 13 colloqui, con un
livello molto alto.

Individua subito un collega francese che, nella sua testa, diventa ``la
versione migliore di me'': stesso ruolo (community manager), esperienza
simile (arrivava da Ubisoft, come quasi tutti i francesi di Riot), ma molto
più spigliato, sempre con la battuta pronta, ``l'anima della festa''. Il
relatore invece è introverso: un workshop come questo per lui è
\emph{emotivamente drenante}, non ama parlare in pubblico pur amando
condividere le proprie esperienze, e non riesce a sostenere conversazioni
come ``cosa hai fatto questo weekend?''. Lavorando raramente insieme, aveva
sviluppato una rivalità ``da allenatore di Pokémon'' (lui era Ash, il
relatore l'antagonista).

Dopo sei mesi, a un evento di \emph{team building} aziendale (a Tenerife), il
suo capo, molto attento a queste dinamiche, li mette apposta nello stesso
gruppo. A fine serata, complice l'open bar, il relatore gli confessa di averlo
mal sopportato dal primo giorno per invidia della sua estroversione. Il
collega ride e risponde che pochi giorni prima aveva detto al capo la stessa
cosa al contrario: invidiava a Stefano il suo essere sempre organizzato, sul
pezzo, presentarsi nella sala riunioni giusta (mentre lui finiva in quella
sbagliata), arrivare con gli appunti, con i dati a supporto e le analisi di
mercato, mentre lui spesso non aveva nemmeno letto la mail del meeting.
Un'invidia in due direzioni, che li ha portati a lavorare molto insieme e a
migliorarsi a vicenda.

Le lezioni:
\begin{itemize}
  \item non bisogna demoralizzarsi vedendo qualcuno che fa lo stesso lavoro
        ed è più avanti sotto alcuni aspetti: spesso c'è un bilanciamento. I
        ``punti caratteristica'' sono limitati: se li ha messi in carisma,
        non li avrà messi in saggezza;
  \item conviene \textbf{chiedere a queste persone come mai} sono così brave
        in qualcosa, come si fa con un progetto andato bene: magari non lo si
        farà mai bene quanto loro, ma si comincerà ad avere il
        \emph{mindset} giusto; e viceversa, loro saranno spinti a chiedere a
        noi;
  \item spesso \textbf{le proprie doti le scopriamo dagli altri}. Il relatore
        non si era mai accorto di essere ``quello che resta calmo durante le
        emergenze'' (il lancio di \emph{Diablo~III}, quello di
        \emph{Draenor}, i problemi dal vivo alla Games Week) finché non
        gliel'hanno fatto notare tutti: per lui era la normalità. A chi pensa
        ``non sono capace, forse non è il mio percorso'', può darsi che
        semplicemente nessuno gli abbia ancora detto ``questa cosa la sai
        fare davvero bene, concentrati su questo''.
\end{itemize}

% ---------------------------------------------------------------------
\section{Domande sul lavoro nell'industria}
\label{sec:z1-industria}

Diverse domande degli studenti hanno riguardato le condizioni di lavoro nei
grandi studi internazionali. Le risposte sono raccolte qui per argomento.

\subsection{Presenza, remoto e trasferimenti}

\paragraph{Grandi studi: quasi solo in presenza.} Sia a Blizzard sia a Riot si
lavorava \textbf{solo in presenza}, senza smart working: sei anni a Parigi e
cinque a Dublino. Si parlava solo inglese: anche nei meeting tra soli
italiani si tendeva a comunicare e a scrivere in inglese, per correttezza
verso chi si fosse aggiunto e per lasciare una traccia comprensibile a tutti.
Il francese serviva per fare la spesa e sopravvivere, non in ufficio.
L'ambiente era estremamente eterogeneo (francesi, polacchi, spagnoli e
molti altri): durante il colloquio dal vivo a Dublino, mentre si preoccupava
di essersi vestito in modo adeguato, vide passare dietro la parete a vetro un
collega svedese con i dreadlock, a piedi scalzi, con una pianta sotto un
braccio e un aquilone sotto l'altro, e capì di essere ``a posto così''.
Un'esposizione a molte culture e modi di pensare che il relatore considera
bellissima.

Anche oggi, tranne alcuni studi particolari, la maggior parte del lavoro si
svolge di persona: amici a Paradox lavorano a Stoccolma, i pochi rimasti a
Blizzard si sono dovuti trasferire a Irvine. Per posizioni \emph{junior} o
\emph{entry level} i grandi studi raramente concedono il lavoro da casa; per
profili senior può esserci qualche eccezione (ad esempio nove mesi da remoto
e tre in sede).

\paragraph{Relocation.} Spesso tocca spostarsi, ma il relatore non ha mai
pagato un trasloco di tasca propria: queste aziende offrono la
\emph{relocation} (aiuto a trovare casa, ad aprire il conto in banca e
supporto pratico), perché è un investimento sulla persona.

\paragraph{Colloqui mentre si lavora altrove.} Per i colloqui in presenza
prendeva ferie, senza dire dove andava. Il consiglio generale: per fare
colloqui \textbf{prendete ferie o permessi}, non datevi malati, ``che non è un
bel modo di comportarsi''.

\paragraph{Remoto ``doppio''.} Un collega lavora per l'azienda coreana che
gestisce \emph{PUBG}, con ufficio ad Amsterdam, ma lavora da Milano: un
``doppio remoto''. Si trova, ma in genere un ambiente anglofono con sedi in
Europa chiederà di trasferirsi in Europa.

\paragraph{Famiglia e trasferimenti.} Il relatore ha avuto la fortuna di
sposare una persona che fa lo stesso lavoro (conosciuta lavorando su
\emph{Warhammer Online}, poi con lui a Blizzard e a Riot, nella
localizzazione e nel doppiaggio, ad esempio di \emph{World of Warcraft} in
italiano): avendo competenze rare, spesso venivano assunti entrambi, al
massimo a pochi mesi di distanza. Per i developer la soluzione frequente è
che la famiglia segua la persona assunta, perché lo stipendio basta per
tutti. Ma non sempre c'è equilibrio, e a volte uno dei due deve sacrificarsi:
sua moglie ha lasciato Blizzard per seguirlo a Dublino (per poi essere
assunta da Riot); lui ha lasciato Riot per tornare in Italia quando lei ha
trovato lavoro qui. Ha dovuto rinunciare più volte a lavorare in Islanda per
CCP (\emph{EVE Online}), perché non era una scelta solo sua. Si fanno
compromessi, ma si riesce a far quadrare.

\paragraph{Ci si sposta per ogni progetto?} No: si sono spostati
esclusivamente per cambiare azienda. Cambiando progetto nella stessa azienda
è raro, quasi mai (``non l'ho mai visto succedere''), che venga chiesto di
trasferirsi. In Francia stavano benissimo (vivevano fuori Parigi; ``i
problemi con i francesi sono in realtà problemi con i parigini''), parlavano
francese ed erano vicini alla doppia cittadinanza: se ne sono andati perché
non vedevano un futuro a Blizzard, per il modo in cui l'azienda trattava i
colleghi.

\paragraph{Il mercato oggi.} Non esiste un \emph{silver bullet} per trovare
lavoro domani dove si vuole. Con tutti i \emph{layoff} recenti ci sono
migliaia di esperti del settore che cercano lavoro: serve pazienza, ma è
fattibile.

\paragraph{Studi indie e remoto.} Negli studi più piccoli il remoto è più
diffuso: a Stormind circa 60 persone su 95 lavorano da remoto, ma lo studio
si è strutturato fin dall'inizio per farlo. Serve consapevolezza. Bending
Spoons (che sviluppa app per mobile, non giochi, ed è l'azienda dell'app
Immuni), pur essendo grandissima e con sede a Milano, lavora da remoto e
lascia l'ufficio a chi vuole: lo smart working è posizionato quasi come un
\emph{benefit}, ma permette anche di avere un ufficio da 300 persone invece
che da 700. Le sedi fisiche costano decine di milioni all'anno (Blizzard ha
chiuso le sedi di Amsterdam e di Parigi). Tuttavia, per lavori molto
verticali e specifici gli studi preferiscono avere le persone dal vivo; i
ruoli più spesso remoti sono quelli che non fanno sviluppo (customer service,
PR\dots).

Anche gli artisti lavorano spesso in presenza: l'accademia può dare solo una
\emph{simulazione} di quanto stretto sarà il contatto tra programmazione,
design e art. ``Sarà ancora di più.''

\paragraph{La comunicazione da remoto.} È più difficile. Pur essendo
estremamente introverso, il relatore \textbf{obbliga tutti a usare la webcam}
nei meeting in cui conta: non in un aggiornamento con novanta persone
collegate, ma certamente in un colloquio a due sulla performance dell'ultimo
anno. La componente fisica è una parte enorme della
comunicazione\footnote{In lezione si parla di ``quasi l'ottanta per cento'';
il dato va preso come indicativo, poiché le percentuali di questo tipo
circolano in forme diverse e sono spesso citate fuori contesto.}: anche se si
è timidi, accendere la webcam è l'approssimazione più vicina all'essere dal
vivo. Negli ambienti dove si lavorava senza webcam la situazione era
``esponenzialmente peggiore''.

\subsection{Differenze culturali nel team}

\paragraph{Gestire sensibilità diverse.} A Riot i team erano divisi per
\emph{lingua}, non per nazione: il team di lingua spagnola includeva anche un
messicano, quello di lingua russa includeva russi e ucraini, proprio nel
periodo della guerra in Crimea. Tra i colleghi non c'era alcun problema (erano
amici), ma ci si è posti la domanda: \emph{come comunichiamo con una community
che sarà divisa su questo tema?} Nei giochi online la divisione era evidente:
in \emph{Rocket League} c'era chi, con un compagno di squadra con la bandiera
avversaria, si faceva otto autogol pur di farlo perdere. Allo stesso modo,
con molti americani mandati a Dublino a coprire posizioni scoperte, c'erano le
elezioni (Clinton--Trump, Biden--Trump) e varie guerre: serviva sempre una
certa sensibilità.

Imparare a navigare queste situazioni lo ha reso un professionista migliore:
gli ha insegnato a fermarsi e chiedersi ``io la percepisco così; gli altri
dieci, con background, storie e lingue diverse, come la penseranno?'', a
ragionare più \emph{long term} e più \emph{big picture}, con più empatia.

\paragraph{Stili di comunicazione.} Un partecipante che ha lavorato nel Regno
Unito chiede dello stile comunicativo italiano, a volte frainteso. Uno dei più
grandi complimenti ricevuti dal relatore è stato ``non sembri italiano''
(``gesticoli poco per essere italiano'', ``parli bene inglese per essere
italiano''). Ha imparato ad \textbf{adattare la comunicazione alla persona}:
con spagnoli e francesi la battuta sulla nazione funziona, perché c'è un
``odio--amore'' condiviso; con russi e giapponesi no. Doveva adattare anche
l'umorismo. Un collega metà russo e metà messicano diceva di aver preso ``il
lato peggiore di entrambe le nazionalità'', ed è stato un ottimo maestro per
imparare a leggere culture diverse.

Non bisogna generalizzare, ma gli stereotipi possono essere un punto di
partenza ragionevole, \emph{sapendo che troverete sempre le eccezioni}. E
vanno letti correttamente: gli americani tendono a essere più estroversi e ad
avere sempre qualcosa da dire non perché ``sono fatti così'', ma perché
vengono educati così (all'università tra i corsi c'è il \emph{debate}, saper
argomentare in pubblico); un collega giapponese può essere più introverso per
ragioni culturali. Nel dubbio si parte da lì, poi si conosce la persona.

\subsection{Negoziare, discriminazione e stipendi}

\paragraph{La famiglia come leva.} Avere un partner riduce le opzioni, ma le
aziende lo sanno: uno spostamento è per due. Si può usare come leva per
negoziare (un compenso diverso, un supporto per il trasloco): se un'azienda ha
deciso di volere quella persona, 2, 3, 4 o 10.000 euro in più all'anno, per
un'azienda che fattura centinaia di milioni o miliardi, non cambiano nulla.

\paragraph{Essere sposati è un motivo per non essere assunti?} No. In alcuni
paesi è proprio illegale discriminare su questo: in Francia era illegale
chiedere CV con data di nascita, foto, sesso o religione, e chiedere se si
vogliono figli. Purtroppo, dai racconti dei suoi studenti di master (in
settori come \emph{beauty and wellness} e \emph{food and wine}), in piccole
realtà italiane queste domande (``hai 25 anni, hai intenzione di
sposarti?'') vengono ancora fatte, e servono a filtrare i candidati. Nelle
grandi aziende ci sono sistemi per evitarlo; anzi, molte danno \emph{più}
benefit a chi ha famiglia (``porta pure la famiglia, è tutto pagato''),
perché le persone competenti hanno spesso una certa età e quindi,
ragionevolmente, una famiglia.

\paragraph{Gli stipendi.} Dipendono molto da dove si vive. Lo stipendio a Riot
non era paragonabile a quello di Blizzard perché in Francia era tassato al
52\% e in Irlanda al 17\%\footnote{Aliquote come riportate in lezione,
riferite alla situazione personale del relatore.}: nessuno stipendio
competitivo può compensare il 35\% di differenza di tasse. Negli Stati Uniti,
con pochissima previdenza pubblica, gli stipendi sono molto più alti perché
poi ci si paga tutto (``se hai un raffreddore paghi tu''). Bisogna valutare il
potere d'acquisto di quei soldi nel paese in cui si andrà a vivere.

\subsection{Freelance e \emph{work for hire}}

All'inizio della sessione pomeridiana uno studente che lavora nell'audio
chiede quali siano le condizioni per chi lavora come \emph{freelance}.
\begin{itemize}
  \item Dipende molto dal progetto: spesso il contratto è a progetto, per
        portare a termine una parte o un intero progetto.
  \item Non tutte le aziende usano freelance: Riot e Blizzard lo facevano
        pochissimo e preferivano assumere (magari trattando poi la persona
        ``come un freelance'', lasciandola a casa a fine progetto).
  \item Spesso i lavori offerti come freelance sono in realtà
        \textbf{sostitutivi di un full time} (``freelance, ma sei ore al
        giorno'').
  \item Si può fare, ma bisogna essere molto chiari nel definire l'impegno
        (quante ore, come vengono tracciate): non tutti specificano cosa
        significa davvero ``freelance''. Lavorando così per due o tre aziende
        si può arrivare a un tempo pieno.
  \item È più comune il \emph{work for hire}: aziende che forniscono personale
        ad altre aziende; si è assunti da X ma si lavora ``embeddati'' presso
        Y (una sorta di subappalto). Anche lo studio del relatore lo fa.
  \item Il puro freelance di quattro mesi esiste, ma di solito per quei
        quattro mesi diventa un lavoro a tempo pieno; è raro ``quattro mesi,
        due ore al giorno''. Nelle grandi aziende è improbabile lavorare nello
        stesso mese per Ubisoft, Riot e Blizzard tre ore ciascuna; è più
        realistico lavorare per più startup indie italiane, che non possono
        permettersi un'assunzione ma possono pagare due ore al giorno.
  \item Molto passa dai contatti: trovare le aziende giuste e rendere esplicite
        le condizioni.
\end{itemize}

% ---------------------------------------------------------------------
\section{Il circolo virtuoso: feedback, miglioramento, fiducia}
\label{sec:z1-triangolo}

\subsection{Il triangolo}

L'obiettivo a cui il workshop vuole arrivare è un \textbf{triangolo che si
autoalimenta}:
\begin{enumerate}
  \item si dà un \textbf{feedback} con l'obiettivo di migliorare le
        performance (inclusi elementi caratteriali);
  \item la persona \textbf{migliora} e vede di averne beneficiato;
  \item questo genera \textbf{fiducia}, che rende più facile dare e ricevere
        il feedback successivo.
\end{enumerate}
Togliendo uno qualsiasi dei tre elementi, il circolo non è altrettanto forte.

In qualsiasi fase si parta, il rispetto per i colleghi è inizialmente dettato
dall'\emph{educazione}, non dalla fiducia. La fiducia arriva quando i colleghi
cominciano a dire cose giuste e sensate e a \textbf{mantenere la parola}.
Dare e ricevere feedback è il modo in cui, in un certo senso,
``inganniamo il cervello'' per creare una \emph{case history} di fiducia: si
ha motivo di fidarsi di più, di ascoltare di più, di essere più ricettivi. E
si diventa sempre più bravi ad accettare il feedback e ad agire di
conseguenza.

\subsection{Accettare un feedback non vuol dire fare come dice l'altro}

Dieci anni fa il relatore non era ostile ai feedback, ma molto
\emph{difensivo}: ``sì, sì, va bene, ora ci guardo''; oppure, a chi gli diceva
``avresti dovuto farla così'', rispondeva spiegando perché l'aveva fatta in
quel modo. Oggi ha imparato a dire ``ok, grazie, ora ci ragiono'': non deve
difendere la propria posizione da chi viene solo a dire ``hai pensato a
questa cosa?''. Prima lo viveva come un attacco (``mi viene a dire che non so
fare il mio lavoro''), ora come un ``ottimo punto di vista, lo prendo in
considerazione''.

Con una postilla fondamentale: \textbf{accettare il feedback non significa
fare come dice l'altro}. ``Mi hai detto che preferisci questo mob verde. Ho
pensato anche a farlo verde? Sì. Penso comunque che giallo sia meglio? Sì, lo
faccio giallo.'' Il relatore accetta \emph{sempre} il feedback, valido o no:
garantisce di tenerlo a mente quando lavora. Che poi faccia come gli viene
detto è un altro discorso.

\subsection{Fiducia e scadenze}

La fiducia funziona allo stesso modo con gli impegni. Se si chiede a qualcuno
una cosa per venerdì e la consegna venerdì, la volta dopo gliela si chiederà
per venerdì. Se invece arriva il giovedì della settimana dopo, la volta dopo
gli si darà una \textbf{deadline anticipata}, come con chi arriva sempre in
ritardo agli appuntamenti (``se ti dico alle 7 arrivi alle 7 e mezza; la
prossima volta ti dico alle 6 e mezza'').

Al lavoro questo è dannosissimo, perché crea \textbf{deadline fittizie} che
non sono mai quelle vere, e nascono discussioni del tipo ``perché devo farlo
in due giorni se si fa in cinque?'' (``sì, se tu lo facessi in cinque''). La
conversazione giusta è un'altra: \emph{``che problemi hai avuto nel
consegnare venerdì?'' ``Questo, questo e questo.'' ``Cosa posso fare per
evitare che succeda?''}

Il relatore non vuole mai sentirsi dire ``non sono riuscito a farlo''.
Preferisce di gran lunga sentirsi dire, tre giorni prima, \textbf{``potrei non
riuscire a farlo, per questi motivi''}. Chi ha la visione d'insieme può dare
più budget, più persone, spostare la deadline, cambiare lo scopo del progetto:
può fare qualsiasi cosa per aiutare. Ma se lo si scopre il venerdì, cinque
minuti prima della deadline, non può fare più niente. Il messaggio è:
\textbf{``aiutami a poterti aiutare''}. Non penserà mai male di chi dice ``non
ce la faccio per questi motivi'', perché non sta dicendo ``non ho voglia'': sta
indicando i problemi, e lavorare su quei problemi è compito del manager. Non
può invece fidarsi di chi fino al giorno prima dice ``tranquillo, ce la
faccio'' e il giorno dopo ``non ce l'ho fatta''.

Questo è fondamentale per chi lavora in programmazione, design o art, perché
ci saranno sempre deadline: bisogna sempre \textbf{lavorare sulle aspettative},
dire al proprio manager (o al proprio report) cosa è lecito aspettarsi. Non
bisogna fare come i pinguini di \emph{Madagascar} (``sorridete e salutate, va
tutto bene''): se c'è un problema, va detto, altrimenti si peggiora solo la
propria situazione.

Da italiani si tende a non voler mai alzare la mano per chiedere aiuto
(``no, no, ce la facciamo, lo faccio io''). A volte non ce la si può fare, non
perché non si è capaci, ma perché vari fattori lo impediscono. Accettarlo è
importante: \textbf{è molto meglio avvisare e chiedere una mano che dire
``non ce l'ho fatta''}.

\subsection{Domande sul dare e ricevere feedback}

\paragraph{``Ne parliamo a fare, se poi fai come vuoi?''} Uno studente (Cristian)
riceve feedback che, dopo averli analizzati, non ritiene congrui; quando non
li segue gli viene detto ``ne parliamo a fare?''. C'è un problema di
\textbf{aspettative}: l'idea che accettare un feedback voglia dire fare come
dice l'altro. Accettarlo significa prendere in considerazione ciò che è stato
detto. Chi decide diversamente ha però un \textbf{onere}: argomentare perché
si fa comunque a modo proprio (``mi hai detto verde, la facciamo gialla
perché\dots''). Poi bisogna solo aspettare i risultati per vedere chi aveva
ragione.

La decisione spetta sempre \textbf{a chi ha la responsabilità di quella
decisione}. Se il game director deve scegliere tra A e B, anche se il
feedback dice B, la scelta è sua perché la responsabilità è sua. Non sarebbe
corretto essere ritenuti responsabili di una scelta delegata a chi dà
feedback. La risposta può essere semplicemente: ``l'ho tenuto in
considerazione, ma visto che ne sono responsabile io e le conseguenze sono
mie, ho scelto di fare X''.

\paragraph{E se l'altro fa ripicca?} Un partecipante chiede di chi, vedendo
il proprio feedback non accolto, si vendica ignorando a sua volta feedback
buoni. Per il relatore è un problema caratteriale: chi si comporta così ``è
già sulla strada per uscire''. Non c'è spazio per l'ego in team di questo
tipo. Il fenomeno nasce dall'idea che \emph{il feedback negativo faccia fare
brutta figura}, che il fatto che si potesse fare meglio tolga qualcosa alle
proprie capacità.

Non è così: \textbf{il feedback non è un giudizio}. Si può passare un anno a
Riot ricevendo solo feedback negativi e avere una valutazione annuale
stellare, perché ricevere molti feedback negativi significa essere una persona
approcciabile, che agisce sul feedback e invita gli altri a darlo. Il relatore
ha visto valutazioni molto peggiori per chi riceveva solo feedback positivi,
perché la gente aveva paura di dargliene di negativi: ma solo ``pacche sulle
spalle'' non fanno migliorare. Il problema va affrontato parlando con la
persona; a volte le conversazioni iniziano con \textbf{``ti do un feedback sul
feedback''}, cioè sul modo in cui l'altro ha dato il suo.

Esempi dalla sua esperienza:
\begin{itemize}
  \item lui chiede spesso di \textbf{non ricevere feedback senza esempi
        tangibili}, perché fa fatica ad astrarre: ``sei stato maleducato'' non
        significa nulla, ``ieri al meeting mi hai risposto male su questa
        cosa'' sì;
  \item ad altri serve l'opposto: ha ricevuto il feedback ``quando devi dirmi
        che ho sbagliato, puoi girarci un po' intorno e darmi del contesto?'',
        perché il suo stile diretto (``questa cosa non ha funzionato, dovremmo
        fare così'') feriva. Ha adattato il suo modo di dare feedback: con
        alcuni è diretto, con altri gira intorno. L'importante è saperlo.
\end{itemize}

\subsection{Non portare la soluzione}

Un consiglio centrale: \textbf{quando si dà un feedback, non si dà una
soluzione}. Non si dice ``io farei così'', ma ``secondo me questa cosa si
potrebbe fare diversamente, per questo e questo motivo: come potremmo farla
diversamente?''. È come chi smette di fumare perché vuole e chi smette perché
glielo dicono moglie e figli: se non si è convinti, spesso non si smette.

Quando il relatore diceva ``c'è questo problema, io farei così'', il meeting
successivo iniziava inevitabilmente con ``ci ho provato come hai detto tu, con
me non ha funzionato''. Sempre. Perché la soluzione era sua. La
responsabilità della soluzione deve essere invece \textbf{congiunta}: ogni
feedback deve creare un piccolo team che ha come obiettivo risolvere quel
problema. Questo toglie l'antagonismo (non ``ti dico che non sai fare il tuo
lavoro'', ma ``ti voglio aiutare a farlo meglio, lavoriamoci insieme'') e
sposta il focus: il problema non sei tu, è ciò che sta accadendo.

\paragraph{E se sono i piani alti a tenere all'oscuro dei problemi?} È una
questione di cultura aziendale. Il relatore fa sempre presente al proprio capo
il concetto di \textbf{service level agreement} (che verrà ripreso nella
seconda giornata): ``se mi dai questi input, ti garantisco questo output in
questi tempi''. Se il capo chiede delle cose senza dire budget, tempi e
perché, non può aspettarsi lo stesso risultato. È una comunicazione a due
vie: si dice al manager cosa serve per fare il lavoro, e lui deve fornirlo o
accettare che il risultato non sarà quello voluto.

\subsection{Amici, parenti e parole proibite}

\paragraph{``Non si lavora con amici e parenti''?} Per sua definizione
personale il relatore ha pochissimi amici (``stanno sulle dita di una mano'')
e una soglia molto alta per considerare amico un collega. Allo stesso tempo,
in quindici anni e con centinaia di colleghi, ha avuto solo \textbf{due}
esperienze in cui non è riuscito a lavorare con qualcuno: non serve essere
amici per far funzionare le cose. Con gli amici, anzi, gli viene più facile:
se dice a un amico che una cosa ``fa schifo'', l'amico non pensa ``come si
permette'', ma ``se perfino Stefano, che è mio amico, me lo dice, allora è
proprio così''. Diventa un rafforzativo, grazie alla fiducia costruita prima.
Il problema del lavorare con amici e parenti nasce dal \emph{non voler} dire
all'amico che sta sbagliando; se invece la confidenza permette di dirlo, si
può saltare la ``confezione'' del messaggio e passare direttamente al
contenuto.

\paragraph{Le frasi peggiori.} Alla domanda se esistano parole peggiori di
``ottimo lavoro'', il relatore indica la sua classifica: ``ottimo lavoro''
(detto in modo vuoto)\footnote{Il passaggio della trascrizione è
frammentario: si intende verosimilmente il complimento generico e di
circostanza, che non dice nulla di utile.}; poi la frase che lo ``fa uscire di
testa'', \textbf{``abbiamo sempre fatto così''} (``e infatti avete chiuso'',
come Blockbuster); infine \textbf{``siamo una famiglia''}, riferito al posto di
lavoro. Un'azienda non è mai una famiglia, perché non ci si può aspettare lo
stesso impegno da chi la possiede e da chi ci lavora: non c'è un
corrispettivo proporzionale tra impegno e ritorno. Per quanto inclusiva,
flessibile e razionale, non sarà mai una famiglia.

\subsection{Ego, muri e il peggior modo di scegliere un leader}

\paragraph{Chi mette davanti l'ego e alza muri.} Esistono, ma di solito durano
pochissimo in azienda. Non si tratta di avere un ego forte, ma di \textbf{non
saper ricevere feedback}. Spesso sono persone bravissime nel loro lavoro che
poi si chiedono perché tutti intorno diventino più bravi di loro: perché gli
altri accettano feedback e hanno più fonti da cui attingere, mentre loro si
sono rintanati nella loro torre.

Molti manager tendono ad \textbf{assumere persone un filo peggiori di loro},
per non fare brutta figura. Il relatore assume solo persone più brave di lui,
``perché non è una gara'': ha lavorato con molti community manager più bravi
di lui come community manager, ma lui era più bravo a fare il manager. Non
viene valutato su quanto sa fare il community manager, ma su quanto il suo
team sa farlo. Un partecipante chiede se la paura di avere qualcuno di
migliore esista anche in ambito accademico. Assolutamente sì: al relatore è
capitato di non essere preso come docente in un corso di formazione sui
videogiochi perché il responsabile aveva percepito che ne sapesse più di
lui\footnote{La trascrizione dice ``che non sapevo più di lui'', ma il senso
del passaggio richiede il contrario.}. ``Ma perché te ne importa? Tu sei quello che ha portato chi
sa fare il suo lavoro: ben venga.'' O si ha questa cultura, o si ha paura per
il proprio posto.

\paragraph{Promuovere il migliore.} Quando bisogna scegliere un leader per un
team (programmazione, design, art), molto spesso si sceglie \textbf{il migliore
tecnico}: il miglior programmatore, il miglior designer. È spesso la scelta
sbagliata: è come prendere un idraulico per fare l'elettricista. Fare il
\emph{people manager} non è come fare il designer o il programmatore; si può
essere un ottimo programmatore e un pessimo manager, e viceversa. Le aziende
l'hanno capito e lo fanno meno di prima. A Blizzard c'era un grave problema di
\emph{middle management}: molte persone in ruoli importanti erano state
promosse cercando le skill sbagliate, erano frustrate dal non saper fare il
proprio lavoro e schiacciavano chi stava sotto per paura di fare brutta
figura.

Un partecipante cita una teoria simile: una persona viene promossa finché non
raggiunge \emph{il proprio livello di incompetenza}, dove si blocca\footnote{È
il \emph{principio di Peter} (Laurence J.\ Peter, 1969), citato in lezione
come ``la teoria di Peter''.}. Un altro cita Michael Scott di \emph{The
Office}, su cui il relatore tornerà (sezione~\ref{sec:z1-quando}).

% ---------------------------------------------------------------------
\section{Perché la cultura conta: retention, costi e fiducia}
\label{sec:z1-cultura}

\subsection{L'ambiente di lavoro e la retention}

Perché tutto questo è così importante? Perché l'ambiente di lavoro conta
moltissimo per la \textbf{retention}, cioè per trattenere le persone e farle
lavorare in modo efficiente. Se contasse solo lo stipendio, molti non
lascerebbero mai Riot o Blizzard; e più si va avanti con l'età, più si impara
a dare valore anche al resto.

Inoltre, secondo gli studi citati nelle slide, chi vive il proprio lavoro in
modo negativo è \textbf{meno efficiente nel problem solving}: l'energia
mentale va verso ``come faccio a non farmi dare la colpa'' o ``come faccio a
fare bella figura'', non verso la soluzione del problema. In ambienti negativi,
dove ci si deve preoccupare della propria ``sopravvivenza'', si lavora
peggio.

\subsection{Il costo di chi se ne va}

Ogni volta che qualcuno si dimette o viene licenziato, il costo di trovare
una nuova persona, fare i colloqui, formarla e farle attraversare forming,
storming, norming\dots\ è circa \textbf{un terzo del suo stipendio}. Per questo
le aziende tendono a fare più \emph{retention} che \emph{acquisition}, almeno
con le figure considerate utili alla produzione.

Allo stesso tempo si vedono sempre più spesso \emph{layoff} e \emph{hiring
wave}: si licenziano 3.000 persone e sei mesi dopo se ne riassumono 2.500 o
3.500, ma tutte più junior. Il relatore ne ha scritto in un articolo su
Eurogamer.it, in cui spiega anche alcuni ``trucchi del mestiere'': perché i
licenziamenti avvengono sempre in certi periodi e sempre ``a gruppi'' (prima
Bethesda, poi Blizzard, poi un'altra\dots).

\paragraph{L'offerta di Riot.} A Riot esisteva una pratica (di cui il relatore
non ricorda il nome) per cui, firmato il contratto, si aveva \textbf{un anno
di tempo per dire ``non mi piace, non me la sento, me ne vado''} e ricevere
subito come buonuscita un anno di stipendio. Avevano stimato che fosse molto
meglio pagare subito un anno a qualcuno e lasciarlo andare che tenerlo due,
tre, quattro anni come \emph{underperformer}: la responsabilità della scelta
di assumere era dell'azienda\footnote{La pratica, nota come ``The Offer'', è
descritta in fonti pubbliche con condizioni diverse (un periodo più breve e
una somma pari a una frazione dello stipendio annuo). Si riportano le
condizioni come raccontate in lezione; andrebbero verificate.}. L'hanno usata
pochissimi, ma il solo fatto che esistesse faceva lavorare meglio anche lui da
manager, perché sapeva che chi non si trovava bene poteva andarsene con un
anno di tempo per trovare un altro lavoro.

\paragraph{Il ricircolo.} C'è un forte ricircolo nell'industria: molti ex
colleghi passano da Ubisoft a Paradox, da Paradox a Funcom, da Funcom a
Blizzard. Il relatore ha scelto di uscirne lavorando solo per aziende
\textbf{indipendenti}: Stormind non è un indie di cinque persone, ma è
indipendente, e non deve accontentare azionisti. Quando vede un'azienda con
utili miliardari licenziare migliaia di persone si chiede ``dov'è la linea''
(\emph{when enough is enough}): con azionisti che vogliono un ritorno, se
quello è l'obiettivo stanno facendo la cosa giusta, ma lui non condivide
l'obiettivo\footnote{In lezione si parla di ``6,5 miliardi di utile'' e di
``3.000 persone licenziate''; le cifre sono riportate come dette e non sono
state verificate.}.

\paragraph{L'open space vuoto.} A Riot, durante un'ondata di licenziamenti,
del suo team di trenta persone è rimasto solo lui: ``teniamo te come
coordinatore, poi ricostruirai il team''. Ha rifiutato di ricostruire un team
cancellato ``con un colpo di spugna'' e ha deciso che non avrebbe più lavorato
per una grande azienda. Nessuno aveva pensato di spostarlo, e ha lavorato per
tre mesi in un open space di cento metri quadri con la sua scrivania e altre
ventinove vuote. Su una lavagna aveva scritto i nomi dei suoi ragazzi e li
cancellava man mano che, anche grazie ai suoi \emph{referral}, trovavano
lavoro. Cancellato l'ultimo nome, si è dimesso: per lui il lavoro era finito.

\paragraph{Team che si dissolvono.} Uno studente cita un articolo su Larian
Studios, il cui management è contrario a ricostruire ogni volta il team a fine
progetto. È lo stesso problema: si pensa ai ``giochi Blizzard'' o ai ``giochi
BioWare'', ma le persone che li hanno fatti non ci sono più; è rimasto solo il
brand, e non ha senso avere le stesse aspettative da un gioco fatto da persone
diverse. E anche chi sopravvive ai layoff vive pensando che il prossimo
potrebbe essere lui: alcune persone, licenziate nell'ultima ondata a
Blizzard, gli hanno detto ``erano due anni che lo sapevo''. Vivevano a Cork,
``un paesotto irlandese dove o vai per lavorare a Blizzard o non ci vai'',
avevano comprato casa: non ci si può aspettare chissà che morale.

\subsection{La cultura è forte quando si lavora per gli altri}

\begin{quote}
  \emph{La cultura è forte quando le persone lavorano per gli altri e a favore
  degli altri, e non per se stesse.}
\end{quote}

Se si dà feedback per far migliorare l'altro e non per imporre il proprio
modo, se si aiuta per il bene del progetto e non per fare bella figura, la
cultura si rafforza. È dai \textbf{piccoli gesti} che emerge la cultura
aziendale: ad esempio, dal non dover chiedere ``perché mi stai aiutando?''.

\paragraph{L'incognita.} \textbf{L'incognita è spesso dannosa quanto le
cattive intenzioni.} In un esperimento (``mi pare a Stoccolma, cinque o sei
anni fa'') si distribuivano per strada banconote da 500 euro, e la maggior
parte delle persone le rifiutava chiedendosi ``dov'è la fregatura?''
\footnote{L'esperimento è riportato in modo vago in lezione e non è stato
possibile identificarlo con certezza.}. Lo stesso accadeva al relatore,
arrivato a PG e a Stormind da Riot: a ``vedo che sei in difficoltà, posso
aiutarti?'' la risposta era sempre ``perché?''. Alla fine rispondeva: ``perché
quando sarò in difficoltà io vorrei che tu mi aiutassi''.

Quando si comincia a lavorare per gli altri e si smette di pensare ``così avrò
la promozione, così avrò l'aumento'', promozioni, aumenti e rispetto arrivano
da soli. A Blizzard aveva fallito tre volte i colloqui per essere promosso a
manager (inteso come coordinatore del team); quando ha deciso di smettere di
pensarci e di fare solo il bene del proprio lavoro, la volta dopo è stato
promosso. Funziona così: chi aiuta i colleghi e fa andare bene il progetto
viene notato (``questa persona fa migliorare cinque persone, non solo se
stessa: vale di più'').

\subsection{Il lavoro è valutato come opera d'insieme}

Il relatore riconosce che buona parte del suo successo è venuto dal fatto che
il suo team lavorava bene, non dalla sua bravura personale. Anche i progetti
degli studenti, presentati a Diego, a Enrico, ad Andrea, verranno valutati
\textbf{come opera d'insieme}. Nessuno dice ``il \emph{combat system} di questo
gioco l'ha fatto X, che bravo'': si dice ``BioWare ha fatto un bel gioco'',
``Blizzard ha fatto un bel gioco''.

Lo ha replicato anche lato HR. In molte aziende ci sono \textbf{bonus
personali} legati a obiettivi (se il gioco vende X, se la feature è pronta in
tempo). Il relatore li ha dimezzati: metà restano personali, metà diventano
\textbf{bonus di team}, che tutti ricevono se il team raggiunge gli obiettivi.
Un incentivo concreto a collaborare: non solo si dice che il lavoro di gruppo è
importante, lo si dimostra dandogli più valore economico. Aveva proposto il
100\% di bonus di team, ma gli è stato concesso il 50 e 50; ha funzionato
comunque, anche se sarebbe stato ancora più potente. Ovviamente, alla
domanda ``perché mi aiuti?'' si può anche rispondere ``perché voglio il
bonus'': un aiuto dato ``per il motivo sbagliato'' era un compromesso che era
disposto ad accettare.

\subsection{Il coraggio di chiedere aiuto}

Un partecipante osserva che aiutarsi dovrebbe essere la norma. Il modello
diffuso, invece, è ``devi essere forte e indipendente, non chiedere mai una
mano'', e vale in tutte le direzioni (dal maschilismo alla ``donna forte e
indipendente che non ha bisogno di nessuno''): \textbf{non ammettere di avere
bisogno di aiuto è sempre uno dei motivi per cui le aziende vanno male}, che si
tratti di uomini, donne o persone non binarie. Il problema è non avere la
sicurezza di poter alzare la mano.

Il relatore ha cercato di togliere ogni conseguenza dall'aver alzato la mano:
se si \emph{depenalizza} qualcosa, le persone non sono più disincentivate a
farla. Ha detto espressamente che non avrebbe mai valutato nessuno per il
numero di volte in cui chiedeva aiuto o per il numero di feedback negativi
ricevuti, e questo ha funzionato.

\subsection{Perché si dà un feedback}

\emph{Perché avete dato il vostro ultimo feedback? Che cosa volevate ottenere?}
Le risposte della classe:
\begin{itemize}
  \item per \textbf{riconoscere il valore} di un collega, anche verso il
        manager, perché qualcuno che non aveva visibilità vedesse quanto
        quella persona fosse efficace;
  \item per \textbf{spingere una persona a superarsi};
  \item per aiutare un collega a rendere il prodotto all'altezza di ciò che
        voleva rappresentare;
  \item per \textbf{genuinità}: c'era qualcosa che poteva essere fatto meglio.
\end{itemize}
Il relatore aggiunge che la forma cambia col tempo: all'inizio era ``devo
darti un feedback, hai dieci minuti?''; ora è un messaggio informale su
Discord, perché ha costruito un rapporto di fiducia e sa che non verrà letto
come ``l'HR manager mi scrive, sono nei guai?'', ma come ``Stefano è così, mi
dà feedback, resterà tra noi e lo risolveremo insieme''.

Un altro motivo è \textbf{risolvere incompatibilità caratteriali}: le
conversazioni più difficili, del tipo ``io e te non stiamo funzionando, cosa
possiamo fare diversamente?''. Sono importanti perché, se due responsabili non
vanno d'accordo, il team se ne accorge, comincia a raccontarsi storie (``si
odiano, ci saranno problemi'') e, per difendere il proprio manager, comincia a
giudicare male il lavoro dell'altro gruppo. È la dinamica da fermare:
\textbf{valutare gli altri solo per difendere la propria fazione}.

Dare e ricevere feedback sono \textbf{due muscoli da allenare} (il relatore è
sempre stato più bravo a darlo che a riceverlo), e si impara col tempo. È come
fare colloqui: fare colloqui è una skill a parte, che non indica quanto si è
bravi nel proprio lavoro. Essere scartati non significa non saper fare quel
lavoro: spesso significa aver fatto male il colloquio o che è stato preferito
qualcun altro.

\paragraph{Feedback pubblici e privati.} Uno studente racconta di aver dato un
feedback negativo a un superiore, senza un motivo particolare, per dimostrare
a sé e ai colleghi che non si viene licenziati per questo. Il relatore coglie
l'occasione per precisare che \textbf{i feedback negativi si danno
individualmente}: lui dà moltissimi feedback di cui il resto dell'azienda non
sa nulla. Reagirebbe molto male se, in un meeting, qualcuno alzasse la mano per
dire ``comunque questa cosa non l'hai fatta bene'': quello non è feedback, è
discussione, accettabile solo se è un brainstorming e se l'opinione è stata
chiesta. Alcune aziende hanno piattaforme pubbliche di feedback, ma spesso
servono solo ad applaudire; lo stesso accade nelle aziende che
\emph{obbligano} a lasciare un certo numero di feedback a settimana, che
diventano ``lui è bravo, lei è competente, anche questa settimana l'abbiamo
fatto''. Bisognerebbe chiedere a chi le adotta cosa pensa di ottenere (forse
trasparenza), ma difficilmente una persona ``attaccata'' pubblicamente reagisce
bene.

% ---------------------------------------------------------------------
\section{Among Us: la caccia al traditore}
\label{sec:z1-among-us}

Prima della pausa pranzo il relatore mostra tre film di fantascienza di
qualche decennio fa e chiede che cosa li accomuni\footnote{I titoli non sono
leggibili nella trascrizione: viene citato \emph{Event Horizon} da uno
studente, e il film ``in mezzo'' è indicato come ``l'originale'',
verosimilmente \emph{La cosa} (\emph{The Thing}).}. In ognuno c'è un gruppo di
persone con un obiettivo e la presenza, reale o percepita, di qualcuno che
rema contro: \textbf{l'energia del gruppo viene spesa per trovare il
responsabile}.

\paragraph{L'esperimento degli alieni.} Nei master in presenza il relatore
propone un esercizio. Distribuisce segretamente cartellini plastificati con
scritto ``essere umano'' o ``alieno''; le carte alieno contengono istruzioni di
comportamento con cui gli alieni, che cercano di fingersi umani, si
tradirebbero (dare sempre ragione a chi ha parlato prima, finire ogni frase
con una domanda, parlare in un certo modo). Lo scopo: in dieci minuti formare
gruppi, e \textbf{vince il gruppo più numeroso senza alieni}; basta un alieno
per perdere. Si formano sempre gruppi piccoli (due, cinque, sei persone),
perché la tendenza è ``mi fido al massimo di voi due''.

Ma vincono i gruppi più grandi, per un motivo semplice: il relatore
\textbf{non distribuisce nessuna carta alieno}. Mostra a tutti entrambe le
carte, ma consegna solo carte ``umano''. Vuole simulare ciò che accade in
ufficio: la tendenza a percepire sempre qualcuno che rema contro (i manager,
l'altro team, quel collega), quando in realtà non c'è nessuno, ``è tutto nella
nostra testa''. Quando chiede le motivazioni, emergono sospetti
\emph{fabbricati}: ``faceva sempre domande'', ``parlava solo all'infinito''
(``no, parla così e basta''). Ci si fabbrica da soli i motivi per non fidarsi
degli altri, ed è comunissimo.

\paragraph{Among Us.} L'autosabotaggio è, per tornare ad \emph{Among Us},
\textbf{concentrarsi sul cercare cadaveri e traditori invece di fare le
task}: non si vince perché non si sono completate le task, mentre si cercava il
colpevole. Succede tantissimo negli uffici, ed è il motivo per cui l'80\% dei
problemi che il relatore risolveva erano di comunicazione. Entrando nei team
sentiva immancabilmente: ``facciamo il possibile, ma questa lotta continua con
l'altro gruppo è estenuante''. ``Perché c'è questo problema?'' ``Non lo so,
c'era già quando sono arrivato.'' Un paradosso: c'è un problema, nessuno sa
dire quale sia, ma c'è sempre stato e quindi lo si accetta.

\paragraph{L'esperimento delle scimmie.} È l'esempio, ``finto o reale'', delle
scimmie in gabbia\footnote{È il noto racconto delle ``cinque scimmie'', spesso
presentato come esperimento scientifico ma probabilmente apocrifo (lo stesso
relatore lo definisce ``finto barra reale''); nella trascrizione il numero
delle scimmie varia.}. In una gabbia con alcune scimmie pende una banana
elettrificata: chi la tocca fa prendere la scossa a tutte. Le scimmie imparano
e, quando una si avvicina alla banana, le altre la picchiano. Poi si
sostituisce una scimmia alla volta con una nuova: la nuova va verso la banana,
le altre la picchiano, e smette. Si continua finché nella gabbia ci sono solo
scimmie nuove, poi si toglie la corrente. Nessuna scimmia tocca la banana:
nessuna ha mai preso la scossa, ma tutte hanno imparato che chi tocca la
banana prende le botte.

Nei team è lo stesso: il motivo per cui si ``odia'' l'altro team si
trasforma, ma un vero motivo non c'è. Ne derivano la sfiducia, il non alzare
mai la mano per chiedere aiuto, il non dire mai ``bel lavoro'' agli altri
perché non sia mai che qualcuno ti veda fare un complimento ``a quelli''.
Risolvere queste dinamiche è lungo e tedioso, e sempre pieno di ``non mi
ricordo'', ``io non c'ero'' e ``abbiamo sempre fatto così'', la cosa più
frustrante per chi cerca di risolverle.

% ---------------------------------------------------------------------
\section{Il principio del default to trust}
\label{sec:z1-default-trust}

\subsection{La definizione}

Il relatore indica questa come la slide più importante della presentazione.
\textbf{Default to trust} è uno dei valori aziendali di Riot Games, che lui ha
tradotto, adattato e ``esportato'':

\begin{quote}
  \emph{Approccio ogni situazione di conflitto partendo dal presupposto che una
  persona abbia agito in un certo modo perché, con le informazioni a sua
  disposizione in quel momento, ha pensato che fosse la soluzione migliore, e
  non perché sia un'incapace o volesse danneggiare il progetto
  intenzionalmente.}
\end{quote}

In altre parole: \textbf{di default, mi fido}.

\subsection{L'esempio del ritardo al daily stand-up}

A Riot il relatore aveva un \emph{daily stand-up} dal vivo alle 15 con il suo
team di venti persone. Una ragazza arrivava regolarmente alle 15:15, per una,
due, tre settimane. Alla quarta ci ha parlato. \emph{Come iniziereste la
conversazione?} Le proposte della classe:
\begin{itemize}
  \item sottolineare l'importanza di esserci dall'inizio, perché potrebbe
        perdersi qualcosa che le interessa, e per rispetto dei colleghi
        presenti;
  \item ``devi arrivare in orario'', spiegando perché è importante;
  \item \textbf{chiedere perché}.
\end{itemize}

Il relatore osserva, anche dalle sue consulenze fuori dai videogiochi, che
\textbf{più si sale con la seniority}, più la risposta è ``devi smettere di
arrivare in ritardo'', spesso con un \emph{or else} finale (``altrimenti
succede quest'altra cosa''). Più si scende con l'età, e quindi con meno
abitudini consolidate, più la prima domanda diventa ``dimmi perché arrivi in
ritardo''. Già spiegare perché è importante è ``miglia avanti'' rispetto al
semplice ``smetti''.

Partire dal perché gli ha permesso di scoprire la causa: il meeting precedente
della ragazza, in un'altra zona del campus a un quarto d'ora di distanza,
finiva sempre tardi, e lei doveva attraversare il campus di corsa. Non aveva
la confidenza per dire all'altro manager, che insisteva perché si restasse
fino alla fine, che doveva andare; sacrificava quindi il meeting del relatore,
sapendo che era il più comprensivo dei due.

Il relatore si è offerto di \textbf{mediare}: ha detto all'altro manager che
o anticipava l'inizio del suo meeting, così da finire comunque in tempo,
oppure lasciava uscire la ragazza, perché quel meeting gli serviva. C'è anche
una questione di \textbf{economia di scala}, spesso sottovalutata: se venti
persone aspettano qualcuno per dieci minuti, si sono sprecati dieci minuti per
venti. Per lo stesso motivo il workshop è iniziato alle 10:05: iniziare alle
10:10 avrebbe significato far aspettare cinquanta persone cinque minuti in
più. A Riot c'era una scrivania vuota con un pupazzo gonfiabile, \textbf{Mr.\
Late}: sommando tutti i minuti di ritardo a meeting, call e conversazioni, in
un anno veniva fuori praticamente \emph{una persona a tempo pieno}. Serviva a
ricordare che anche cinque minuti contano, e che bisogna rispettare gli altri
arrivando in orario o almeno avvisando.

\paragraph{Gli effetti.} La conversazione era avvenuta nel primissimo periodo a
Riot, quando il team non aveva ancora motivo di fidarsi di lui. Risolvendo il
problema ha eliminato anche tutto lo strato di micro-dinamiche spiacevoli che
un manager impara a vedere (due colleghi che si guardano con la risatina
quando lei entra in ritardo). La ragazza ha cominciato a entrare tranquilla,
con una certa leggerezza: non le faceva piacere arrivare in ritardo (``ci sono
anche persone che lo fanno apposta, e quello è un problema da affrontare
diversamente''). Soprattutto, con i colleghi non parlava di lui come del capo
che minacciava, ma come di chi le aveva risolto un problema: era diventato
un \textbf{alleato}. La settimana dopo altre tre persone sono andate da lui
con un problema da risolvere. Ripetuto cinque, sei, sette, dodici volte, in un
paio di mesi poteva dire quasi qualsiasi cosa nei meeting sapendo che
richieste e critiche sarebbero state percepite come ``me lo chiede perché vuole
darmi una mano'', non per fare bella figura o per presentare i successi del
team ai suoi superiori.

\subsection{Condividere i successi, schermare i fallimenti}

Un altro elemento fondamentale (ripreso nella seconda giornata): il manager
\textbf{condivide i successi e protegge il team dai fallimenti}. ``Se va male
qualcosa è colpa mia; se va bene qualcosa, voglio che ne parliate voi.'' Il
lavoro del manager non dovrebbe essere valutato su ciò che presenta del lavoro
del team: tutti sanno comunque chi è il manager. Tanto vale lasciar presentare
il team, che potrà dire con orgoglio di aver presentato a Mike
Morhaime\footnote{Cofondatore ed ex presidente di Blizzard; nella
trascrizione ``Mike Morayme''.} la propria campagna europea per
\emph{StarCraft~II}. Ci sono manager che dicono ``come mi prendo le colpe, mi
prendo i meriti'', ma i meriti li hanno già: è inutile ostentarli. Con i
fallimenti invece il manager è la \textbf{rete di sicurezza} del team. Il
relatore ne ha fatto una questione di principio: \textbf{non scaricare mai la
colpa su qualcun altro}.

\paragraph{Colpevole o causa.} Un esempio che il relatore attribuisce a Pirate
Software, creatore di contenuti ed ex dipendente Blizzard\footnote{Il relatore
aggiunge che si tratterebbe del figlio dello sviluppatore a cui si sarebbe
ispirato il personaggio del giocatore ``senza vita'' nell'episodio di
\emph{South Park} dedicato a \emph{World of Warcraft}. L'informazione è
riportata come detta in lezione e non è verificata.}. Contattando il servizio
clienti Blizzard per motivi personali, aveva scoperto che i colleghi gli
avevano scritto nei ticket informazioni confidenziali che non avrebbero dovuto
condividere in quel modo, nemmeno con un collega. Non ha segnalato le persone
da punire: ha detto ``c'è un problema, evidentemente le policy non sono
chiare''. Quando gli hanno chiesto i nomi, ha rifiutato: ``se vi dico chi è,
punite quelli, e gli altri cinquecento continueranno a farlo; non avremo
cambiato la causa alla fonte''.

\begin{quote}
  \emph{È inutile trovare il colpevole. Voglio trovare la causa. Perché il
  colpevole cambia.}
\end{quote}

Il colpevole a volte è vittima della situazione: per questo la definizione
dice \emph{``con le informazioni a sua disposizione in quel momento''}. Se il
relatore avesse rimproverato la ragazza in ritardo, non solo non avrebbe
risolto il problema, ma avrebbe creato un'\textbf{ingiustizia}: lei, che
attraversava il campus di corsa per una colpa non sua e di cui non poteva
parlare, sarebbe stata anche rimproverata.

Nessuno fa un brutto lavoro apposta: nessuno è contento quando sbaglia o
perde. Perché al lavoro diamo per scontato che gli altri sbaglino
intenzionalmente? Se qualcuno ha fatto una cosa male, è perché in quel momento
credeva fosse la scelta migliore. La domanda giusta non è ``perché hai fatto
questa cosa che è stata un fallimento?'', ma:

\begin{quote}
  \emph{Che cosa ti ha fatto pensare che questa fosse la scelta giusta?}
\end{quote}

Magari erano limiti di budget, di tempo, di scopo: si agisce su quelli, non
sulla persona. Se si cambia la persona ma le si danno gli stessi fattori, è
molto probabile che arrivi alle stesse conclusioni, e ci si troverà a dover
``disciplinare'' anche lei.

\paragraph{La reputazione costruita sui rumor.} Durante il giro di
presentazione in un ufficio (non in nessuna delle aziende citate), al relatore
è stato detto: ``questo è il dipartimento di design; se ti serve qualcosa
chiedi a lei, non a lui, perché lui consegna sempre in ritardo''. Oltre a
trovarlo poco professionale, una settimana dopo ha chiesto: ``cosa ti ha fatto
questa persona?'' ``A me niente, però ho sentito che\dots''. Si torna alle
scimmie: magari quella persona non consegna in ritardo da mesi, ma ormai lo
dicono tutti. È un circolo vizioso da fermare subito.

\paragraph{Niente escalation inutili.} Lo stesso vale per il feedback: quando
qualcuno va dal relatore dicendo ``ho un problema con questa persona'', lui
chiede: \textbf{``perché ne stai parlando a me e non a lei?''}. ``Perché sei il
mio capo'' non basta. Passare da lui significa dare per scontato che l'altra
persona non abbia interesse a risolvere, mentre magari basta parlarle. È la
strada più facile (``lo dico a Stefano, ci pensa Stefano''), e proprio per
questo bisogna mettere dei paletti: il manager interviene quando serve, ma le
cose che si possono gestire da soli vanno gestite da soli. (Quando invece il
confronto diretto fallisce, l'escalation è legittima: si veda la
sezione~\ref{sec:z1-qa-finale}.)

% ---------------------------------------------------------------------
\section{Traiettorie di carriera e valutazione delle performance}
\label{sec:z1-traiettorie}

\subsection{Quale traiettoria vorreste?}

Il relatore mostra un grafico con quattro curve colorate, contrassegnate
anche da lettere (A, B, C, D) per eventuali persone daltoniche (``a me è
capitato che qualcuno scoprisse di esserlo durante una lezione''). Le curve
rappresentano la performance lavorativa nel tempo. Due domande: \emph{quale
vorreste essere? E quale pensate realisticamente di essere?} Un partecipante
aspira a B e pensa di essere C; quasi tutti scelgono B.

Perché non A, che sale più rapidamente? Un partecipante risponde che una
crescita così repentina può essere più dannosa che utile, come una ``rush'' di
studio molto faticosa che poi non si riesce a sostenere. Il relatore propone
una lettura suggerita da un amico psicologo, che spiega la forma del grafico:
di B, C e D \textbf{vediamo dove finiscono}; di A no. Potrebbe continuare a
salire oltre la scala del grafico, o crollare sotto D. Non saperlo ci spinge a
preferire B. (A chi obietta che anche le altre proseguono oltre il bordo, il
relatore precisa: il bordo destro è ``oggi'', e di A non sappiamo come sia
arrivata fin lì.)

\subsection{La carriera non è una linea retta}

Tutti immaginano la carriera come una linea: si parte da X e, attraverso
stage e posizioni, si arriva a Y. In realtà carriera e performance avranno un
andamento irregolare, con alti e bassi. Come non si accetta che un team possa
temporaneamente \emph{underperformare}, tanto meno lo si accetta da se stessi.
Bisogna accettare che ci saranno momenti bassi.

Per questo il relatore \textbf{detesta le valutazioni annuali delle
performance}: prendono un punto casuale della traiettoria, che se si è
fortunati è il picco più alto e se si è sfortunati il più basso, e lo usano
per valutare un intero anno. Insegna ai manager a fare più punti di controllo
e a guardare altre cose:
\begin{enumerate}
  \item \textbf{quanto tempo passa tra un picco e l'altro}, cioè quanto
        \emph{cooldown} serve a una persona. È il motivo per cui, dopo la
        Games Week, tutti staccavano una o due settimane. Il relatore stesso,
        che si impegna moltissimo e tende al burnout, non potrebbe rifare lo
        stesso workshop il lunedì e il martedì successivi con la stessa
        qualità;
  \item \textbf{la differenza tra la miglior performance e la peggiore}: che
        cosa ha portato a lavorare così bene, e così male? Con chi si
        lavorava, su cosa ci si concentrava, quanto tempo e budget c'erano?
        L'obiettivo è capire come \emph{replicare} le condizioni che fanno
        lavorare bene, invece del solito ``hai lavorato male, è colpa tua''.
        Il relatore ha avuto molti manager per cui quando lavorava male era
        colpa sua e quando lavorava bene ``che bravo'', ma sapeva che raramente
        era l'unico fattore della propria performance. È la fallacia del ``se
        sai fare 100, perché non fai sempre 100?'', come nella scena di
        \emph{Schindler's List} sull'operaio che produce una vite in pochi
        secondi: non sempre un ritmo è sostenibile;
  \item \textbf{il fallimento di oggi confrontato con quello di sei mesi fa}.
        Il relatore non è contento della presentazione di oggi (l'avrebbe
        migliorata per altri due anni), ma rispetto a quelle di cinque anni fa
        vede un grande miglioramento di chiarezza: il suo peggior fallimento di
        oggi equivale al suo primo successo di inizio carriera.
\end{enumerate}
Quello che conta è la \textbf{traiettoria} dall'inizio dell'anno a oggi, non
i singoli punti. E se si riesce a capire insieme quando e perché si ha
successo e quando e perché si fallisce, si può lavorare per replicare solo le
situazioni che fanno lavorare bene.

\paragraph{Lo stesso vale per i giochi.} È il motivo principale per cui i
giochi delle grandi case vanno bene o male. Quando si riescono a replicare le
condizioni che hanno prodotto \emph{Warcraft~II}, esce \emph{Warcraft~III};
quando non ci si riesce, esce ``il nuovo \emph{Dragon Age}'': la gente non è
arrabbiata perché il gioco non è bello, ma perché non è come gli altri
\emph{Dragon Age}; non ne replica il \emph{feeling}, anche se potrebbe
essere un gioco favoloso. L'eccezione positiva è il passaggio da \emph{Dune} a
\emph{Dune~II}: il primo era di fatto un'avventura grafica, il secondo ha
generato il genere degli strategici in tempo reale. Gli sviluppatori avevano
capito che il successo non stava nel tipo di gioco, ma nell'ambientazione che
attirava le persone: un bel gioco in quell'ambientazione sarebbe piaciuto
comunque. Il relatore comprò \emph{Dune~II} aspettandosi un seguito
dell'avventura e, pur sorpreso, lo finì.

\subsection{Lying, hiding, faking}

Quando si è valutati su un singolo punto nel tempo, e da quel punto dipendono
bonus, stipendio e il rischio di licenziamento, si comincia a fare
\textbf{lying, hiding e faking}:
\begin{description}
  \item[Lying] dire che è sempre andato tutto bene, che si finirà sempre in
        orario, che non c'è nessun problema;
  \item[Hiding] nascondere: ``non era un mio task'', ``non c'ero'', ``non ho
        visto niente'';
  \item[Faking] fingere, anche in cose apparentemente innocue come ``sì, sì,
        ho capito tutto'' quando si ha solo paura di dire ``me lo rispieghi?
        Non ho capito niente''.
\end{description}

\subsection{``Non ho capito'' e la conferma delle istruzioni}

Dire ``non ho capito, ripetimelo'' è ancora visto come una figura da sciocco.
Il relatore non si arrabbierà mai se qualcuno gli chiede di rispiegare
un'intera slide: le cose ovvie per lui possono non esserlo per altri. Più si
conosce qualcosa, \textbf{più è difficile spiegarla}, perché si danno per
scontate molte cose. Ne parlava proprio con Diego, che gli diceva che fare
l'insegnante diventa sempre più difficile man mano che si diventa esperti. È
come spiegare \emph{EVE Online}, uno dei giochi più complessi in circolazione:
chi lo conosce bene dà per scontato che ``ovviamente non puoi fare questa
cosa'', mentre per un nuovo giocatore non è affatto ovvio. A scuola chi dice
``non ho capito'' viene visto come lo ``scemo''; nessuno si chiede ``non è che
l'ho spiegato male?''.

Il racconto più ``traumatico'' dei ragazzi che fanno stage nelle grandi
aziende è sempre lo stesso: ``Cinzia, devi fare questa cosa.'' Cinzia la fa.
``Non te l'avevo chiesta così.'' Il problema è che nessuno ha
\textbf{confermato che Cinzia avesse capito}. Il consiglio, valido sempre:
\begin{itemize}
  \item quando si \emph{ricevono} istruzioni, ripeterle: ``ho capito bene?
        Vuoi che faccia questa cosa in questo modo?'';
  \item quando si \emph{danno} istruzioni, chiedere a chi le riceve di
        ripeterle, non perché lo si ritenga incapace, ma per capire che cosa
        ha capito dalla propria spiegazione (``quindi vuoi un quadrato?''
        ``No, ho detto un triangolo\dots'' ``Però mi hai detto una forma con
        quattro angoli.'' ``Hai ragione: voglio un quadrato'').
\end{itemize}
Frasi come ``non te lo ripeto due volte'' o ``stammi dietro'' portano solo a
fare le cose come si pensa che vadano fatte e a sentirsi dire a fine progetto
``non la volevo così'': che si tratti di una feature, di un design o di un
modello 3D. (Un partecipante osserva che lo stesso accade insegnando le lingue:
chi insegna dà per scontato, ad esempio, che si sappia dell'esistenza del
genere neutro in tedesco.)

\paragraph{Come convincere il team a chiedere?} Il relatore mette sempre come
premessa che \textbf{l'obiettivo non è trovare la persona più intelligente
nella stanza, ma fare un buon lavoro}. Che Gianni capisca alla prima e un altro
alla quinta non gli interessa, se alla fine entrambi lavorano bene: ripetere
cinque volte è comunque meno lavoro che far rifare le cose sette volte. È un
investimento a lungo termine, e non dissimula alcun fastidio: davvero non gli
dà fastidio. Ma ripetere le stesse parole non aiuta: ``non ho sentito'' si
risolve ripetendo; ``non ho capito'' si risolve cambiando spiegazione. Meglio
ancora: ``dimmi che cosa hai capito tu, e vediamo cosa non torna''. Se
intravedesse nel team uno sguardo del tipo ``è la quinta volta che lo
chiede'', interverrebbe: spesso chi giudica è proprio chi dice ``ho capito'' e
poi sbaglia. Preferisce di gran lunga chi chiede cinque volte.

\subsection{Il feedback più utile che abbiate ricevuto}

\emph{Qual è il feedback più utile che abbiate mai ricevuto?} Il relatore ne ha
almeno dieci che gli hanno cambiato il modo di lavorare. Le risposte:
\begin{itemize}
  \item un semplice \textbf{``grazie''} per un lavoro fatto, più utile di un
        ``questa cosa fa schifo'';
  \item un feedback sulla propria \textbf{impulsività} nel rispondere: fermarsi
        un attimo e capire di cosa si sta parlando prima di intervenire;
  \item \textbf{essere se stessi}, anche al lavoro, e non scusarsi
        continuamente per ogni cosa.
\end{itemize}
I feedback che toccano di più sono quelli \textbf{veri}. Il ``grazie''
significa ``ho visto quello che hai fatto e l'ho apprezzato''. Quelli negativi
che si ricordano sono quelli che hanno fatto realizzare ``questa cosa che
pensavo di saper fare, non la so fare'', oppure ``la so fare, ma questa volta
non mi è venuta'' (non è che se una cosa una volta non riesce non la si sa più
fare: può capitare di sbagliare una porta). Arrivano sempre da un
\textbf{punto cieco}: qualcosa che non si sapeva fosse percepito dagli altri.
Conta più ciò che le persone ti dicono di ciò che pensi che pensino di te; e
ciò che il relatore pensa che gli altri pensino di lui è sempre terribile,
mentre il loro giudizio si rivela molto più positivo del previsto.

\paragraph{Il denominatore comune siete voi.} Si potrebbe pensare che questi
concetti non siano applicabili da chi arriva in un'azienda senza poterne
cambiare la cultura. Ma c'è sempre un denominatore comune a cui applicarli:
\emph{se stessi}. La cultura aziendale può non spingere a dare feedback, può
vedere male chi chiede aiuto; ma il relatore ha visto che quando lui lo fa, le
persone intorno tendono a farlo, e un gruppo che lo fa lavora meglio. Perfino
i CEO più ostici a un certo punto gli chiedevano: ``perché il tuo team ora
performa meglio degli altri?''. La differenza è il default to trust: una
pratica che col tempo fa operare il team \emph{come un team} e non come un
insieme di individui. Un videogioco è un progetto corale: c'è chi fa giochi da
solo, ma non si fa da soli \emph{Red Dead Redemption~2}; e se il gruppo non è
affiatato e allineato non esce \emph{Red Dead Redemption~2}, esce
\emph{Anthem} (``oggi la mia fissazione è \emph{Anthem}, domani sarà
un'altra'').

% ---------------------------------------------------------------------
\section{Come dare un feedback}
\label{sec:z1-dare}

\subsection{Concentrarsi sull'azione e sui suoi effetti}

Il primo principio: \textbf{il feedback non riguarda mai la persona}. Non si
dice ``tu hai fatto questa cosa che non va bene'', ma ``questa cosa non va
bene: come facciamo a non farla più?''. Che sia un comportamento, un'azione,
una decisione o una strategia, ci si concentra su quello, e si
\textbf{esplicita l'effetto} che ha. Nell'esempio del ritardo l'effetto era
``quindici minuti di attesa moltiplicati per venti persone''. In un progetto
può essere una catena:
\begin{quote}
  se non rispetti le deadline, l'animazione che viene dopo resta ferma; se
  loro sono in ritardo, la \emph{vertical slice} è in ritardo; se la vertical
  slice non è pronta per la Gamescom, non andiamo alla Gamescom e non troviamo
  un publisher; se non troviamo un publisher, chiudiamo.
\end{quote}
Più la conseguenza è grave, più dovrebbe sembrare piccolo il gesto di dire
``datemi una mano'': è il \emph{male minore}. Poi, a cose fatte, ci si dà
feedback: cosa ha funzionato, cosa no, cosa fare di più, di meno, cosa smettere
di fare.

\subsection{Il pericolo degli ambienti solo positivi}

Gli ambienti dove ci si fa sempre grandi complimenti e nessuno parla dei
problemi creano difficoltà vere alle aziende. Il problema non è mai la
mancanza di complimenti, ma un \textbf{eccesso di ottimismo}: si comincia a
pensare che le cose andranno come si spera, perché tutti intorno danno solo
feedback positivi.

A Blizzard i feedback determinavano gli \textbf{aumenti annuali}: ognuno
doveva indicare cinque persone che avrebbero dato feedback su di lui, e
ovviamente tutti indicavano gli amici, o colleghi di progetti andati bene. Un
sistema nato con le migliori intenzioni diventava uno \emph{scambio di voti}
(``se mi metti follow ti metto follow''), senza alcun miglioramento.

Al contrario, \textbf{le conversazioni più brutte sono state le più utili}:
affrontare i licenziamenti a Riot, o la cancellazione di \emph{Project Titan}
a Blizzard, il progetto che doveva essere il nuovo MMO dopo \emph{World of
Warcraft} e da cui poi è nato \emph{Overwatch}\footnote{Nella trascrizione si
legge anche ``doveva essere Warcraft 2'', probabilmente un lapsus o un errore
di trascrizione.}. Andando a vedere perché lo si cancellava, c'erano solo
motivazioni razionali e pragmatiche, di business, chiare a tutti, non un ``Mike
Morhaime ci odia''. Sono anche le conversazioni in cui ha lavorato su se
stesso: è molto \textbf{avverso ai conflitti}, e questo a volte gli ha
procurato conseguenze molto più gravi di quanto sarebbe costato un piccolo
disagio. Dieci anni fa faceva fatica ad accettare questo difetto; ha dovuto
subire le conseguenze di conflitti evitati e poi cresciuti.

\subsection{Che cosa fare}

Le linee guida della slide, che il relatore legge rapidamente perché saranno
disponibili nel materiale:
\begin{description}
  \item[Attenzione al linguaggio del corpo] dare un feedback a braccia
        incrociate o guardando altrove è diverso dal guardare l'altro negli
        occhi.
  \item[Sii specifico] con esempi tangibili.
  \item[Esplicita gli effetti] come visto sopra.
  \item[Discuti il prima possibile] non aspettare un mese per parlare di
        qualcosa successo ieri, altrimenti nessuno ricorderà più bene cosa sia
        accaduto.
  \item[Concentrati sull'azione e sul problema, non sulla persona] il nemico è
        l'azione, non la persona che si ha davanti. Non bisogna far cambiare
        l'altro: bisogna evitare che il problema si ripresenti.
\end{description}

\paragraph{Gli insulti sono rivolti al ruolo.} Un esempio solo apparentemente
scollegato. Il relatore ha lavorato su \emph{League of Legends}, ``non
esattamente noto per avere una community da orsetti del cuore'', in varie
\emph{task force} sulla tossicità. A chi smetteva di giocare veniva inviato un
sondaggio sul perché; il 33\% indicava la tossicità della community. Dopo varie
campagne di sensibilizzazione e cambiamenti si è scesi al 30\%: sembra solo
un 3\% in meno, ma valeva circa \textbf{120 milioni di revenue in più
all'anno}, perché quei giocatori restavano esposti a skin e contenuti e
continuavano a comprare. Anche qui, piccoli guadagni che si sommano.

Lavorando con degli psicologi, aveva notato che gli insulti \textbf{non sono
quasi mai rivolti alla persona, ma al ruolo}: ``il mio jungler fa schifo'',
``il mio DPS non è capace'', ``il nostro portiere è una pippa''. Chiamare
qualcuno per ruolo lo \emph{depersonalizza} e rende più facile insultarlo. Lo
stesso nell'esempio del giro dell'ufficio: non ``non chiedere a Giulio'', ma
``non chiedere a quel grafico, consegna sempre in ritardo''. Concentrarsi
sull'azione, senza attaccare la persona, fa sì che il feedback non sia (e non
venga percepito come) un attacco personale.

\paragraph{Considera il carattere.} Qui tornano utili i test di personalità.
Servono più agli studenti che al relatore (che sa già che usciranno 53 profili
diversi, senza un pattern): bisogna conoscere il proprio profilo e
\textbf{condividerlo con chi lavora con noi}, chiedendo anche agli altri di
fare lo stesso. In ufficio il relatore ha usato una versione più avanzata, a
pagamento, con un report di circa trenta pagine: il \textbf{profilo
DISC}\footnote{Nella trascrizione ``disk''. Si tratta del modello DISC
(Dominance, Influence, Steadiness, Conscientiousness).}. Ha mappato tutta
l'azienda e ha fatto stampare il profilo di ciascuno sulla sua scrivania, così
che tutti sapessero con chi stavano parlando. Capire come ragionano gli altri
(``è introverso, non gliene parlo davanti a tutti'') è una delle cose che gli
sono tornate più utili: \emph{adattare ciò che si vuole dire a chi si ha
davanti}.

\paragraph{Assicurati di non disturbare.} Sotto il profilo di ognuno c'era un
piccolo \textbf{semaforo}: verde, ``parlami pure, interrompimi''; giallo,
``interrompimi solo se è urgente''; rosso, ``mi sto concentrando, non mi
parlare''.

\paragraph{Lascia finire il lavoro.} Durante gli eventi dal vivo il relatore si
era reso conto di quanto fosse grave dare feedback mentre qualcuno sta ancora
lavorando. Se una cosa sta andando male, dirlo in quel momento non aiuta: tra
la pressione di finire e quella del feedback, peggiorano entrambe le cose, e
anche il rapporto. Meglio lasciarla finire, anche male, e parlarne dopo.

\subsection{Che cosa non fare}

\begin{description}
  \item[Non basarti su supposizioni] ``mi hanno detto che ieri in un meeting
        hai fatto\dots'': o c'eri e l'hai visto, o non esiste.
  \item[Non giudicare la persona.]
  \item[Non fornire feedback vago.]
  \item[Non accettare e non fornire feedback di terze parti] Se qualcuno vi
        dicesse ``mi hanno detto che hai fatto un brutto lavoro'', il primo
        pensiero non sarebbe ``cosa ho fatto male?'', ma \emph{``chi te l'ha
        detto? Chi è la vipera che ho in seno? Chi è l'impostore?''}. Siamo
        fatti così. Ogni conversazione che inizia con ``mi hanno detto che'' è
        inutile: l'altro non ascolterà, penserà solo a chi dare la caccia, e
        spesso penserà alla persona sbagliata, creando antipatie tra persone
        che non c'entravano nulla. Quando al relatore dicono ``mi hanno detto
        che lui ha fatto\dots'', risponde ``vai a parlarne con lui''; se gli
        dicono ``mi hanno detto che tu hai fatto\dots'', risponde ``mandami la
        persona che te l'ha detto, altrimenti non ne parliamo''. Non per
        cattiveria: se chiedesse ``che cosa ho detto di preciso a Marco?'' e la
        risposta fosse ``non lo so, non c'ero'', di cosa si starebbe parlando?
        \textbf{Parlate con la fonte}, altrimenti fate come se quel feedback non
        esistesse.
  \item[Non dare consigli non richiesti] come visto, non portare la soluzione:
        la soluzione deve venire da chi riceve il feedback.
  \item[``Non mi piace'' non è un feedback] ``Penso che questo design non
        funzioni perché, visto che il nostro obiettivo è\dots'' ha argomenti;
        ``non mi piace il combat system'' è un'opinione. Anche ``io l'avrei
        fatto diversamente'' non è un feedback.
  \item[Accettare non vuol dire fare come proponi] deve essere chiaro a chi dà
        il feedback: l'aspettativa è essere ascoltati e presi in
        considerazione, ma se la responsabilità è dell'altro, deciderà
        l'altro.
  \item[Non offenderti] se il feedback non viene seguito. Restarci male
        significa pensare di saperne più dell'altro, cioè \emph{venire meno al
        default to trust}. ``Decidi tu'' significa ``mi fido di te, perché sei
        lo specialista, e mi fido che mi ascolterai e farai la cosa migliore
        per tutti'', che potrebbe anche non essere ciò che ho detto io. Se
        voglio che l'altro faccia sempre come dico io, mi sto fidando solo di
        me stesso.
\end{description}

\subsection{Il primo feedback e il debriefing}

La cosa più difficile è dare il \textbf{primo} feedback: è un muscolo da
allenare. Chi non l'ha mai fatto (soprattutto con questa struttura) dovrebbe
provarci, anche su minuzie (``abbiamo concordato di usare WhatsApp nel gruppo:
non potremmo usare Telegram?''), anche solo per vedere come va. Altre
indicazioni:
\begin{itemize}
  \item feedback \textbf{solo uno a uno}, mai davanti ad altri, soprattutto se
        negativo;
  \item \textbf{prepararsi} prima;
  \item fare il \textbf{debriefing del feedback}: dopo, fermarsi a pensare a
        com'è andato. Il relatore, quando media un feedback tra due persone,
        si sente spesso rispondere ``benissimo'' da una e ``malissimo''
        dall'altra: la percezione è molto diversa. Non si fa problemi a tornare
        dalla stessa persona: ``ieri ti ho dato un feedback: com'è andata? Sei
        soddisfatto? Come poteva andare meglio?''. Con le persone con cui si ha
        meno a che fare, questo aiuta molto la volta successiva.
\end{itemize}

% ---------------------------------------------------------------------
\section{Come ricevere un feedback}
\label{sec:z1-ricevere}

Per il relatore è la parte più difficile. Si riassume in quattro passi:
\begin{enumerate}
  \item \textbf{Ascolta}: si parte sempre ascoltando.
  \item \textbf{Chiedi dettagli}, se serve, ma per capire meglio, non per
        difenderti.
  \item \textbf{Accetta, a prescindere}: non si sta decidendo chi ha ragione;
        si sta ascoltando un'opinione. Non c'è nulla da cui difendersi.
  \item \textbf{Agisci}: fatti un piano. Può essere banale (``mi hanno detto
        che interrompo sempre: cerco di non interrompere''), oppure ``devo
        migliorare questa cosa: come?''. Pensa anche a \emph{chi} potrebbe
        aiutarti e a \emph{che supporto} ti serve: il piano non deve essere
        per forza ``devo risolvere tutto da solo''.
\end{enumerate}
Una domanda che molti sottovalutano: \emph{``sono junior in un'azienda: sarò
mica il primo ad avere questo problema?''}. Probabilmente qualcun altro l'ha
già avuto: basta chiedere ``tu come hai fatto?''.

\subsection{Il feedback positivo}

Tutto questo vale per i feedback costruttivi negativi, a cui si pensa di
default; ma esistono anche i \textbf{feedback positivi}, e vanno usati. Quando
si riceve un ``grazie'' o un ``bravo'', la domanda è: \emph{come faccio a
ricreare le stesse condizioni?}

Il relatore ha lavorato anche in aziende dove il feedback era solo negativo.
Il rischio? Si immagini l'impiegato ideale: vuole sempre migliorare, è sempre
pronto a dire ``magari sbaglio io''. Se nessuno gli dice cosa sta facendo
\emph{bene}, spinto dal desiderio di migliorare rischia di cambiare cose che
erano già buone, e di farle peggio. Il rinforzo positivo \textbf{non è solo una
gentilezza}: è un atto pratico, che evita che le persone peggiorino ciò che
sanno già fare bene.

\subsection{Il feedback accidentale}

Esiste infine il \textbf{feedback accidentale}, molto frequente negli uffici:
\begin{itemize}
  \item il ``chiedi a lei, non a quel grafico'' del giro dell'ufficio è un
        feedback negativo su una persona, dato accidentalmente, né richiesto
        né voluto;
  \item a Riot li chiamavano \textbf{drive-by}, come nelle scene dei film in
        cui passa un'auto, si abbassa il finestrino, si spara e si riparte. In
        molte aziende con una ``cultura del feedback'' in cui tutti devono
        dire la loro, si riceve una valanga di commenti del tipo ``comunque il
        sito a me non piace, ciao''. Non si sa perché, non si sa cosa fare di
        diverso, non si sa che rischi ci siano: \emph{a che cosa è servita
        questa conversazione?} A niente.
\end{itemize}

% ---------------------------------------------------------------------
\section{Quando, quanto e con chi}
\label{sec:z1-quando}

\subsection{Frequentemente e recentemente}

\emph{Quando avete ricevuto il vostro ultimo feedback?} Per gli studenti,
qualche giorno fa (nel corso) o qualche mese fa (al lavoro, dopo una
consegna). Come aveva osservato Silvia, il feedback andrebbe dato
spesso.

Il relatore usa l'esempio di una nave da crociera (``la Costa Concordia'') che
deve andare da A a B. La rotta ideale (azzurra) è quasi impossibile: nessun
progetto va esattamente come previsto. C'è la rotta ``panoramica'' (verde), che
gira intorno all'isola d'Elba; ci sono poi rotte (arancione e rosa) che a un
certo punto si allontanano molto e poi correggono bruscamente: qualcuno si è
accorto che si stava sbagliando, ``ma non è quello che avevo chiesto'',
``non è dove dovevamo andare''. Sono i famosi \textbf{sanity check} o
\emph{checkpoint}.

L'ideale è avere \textbf{tanti piccoli check per fare tante micro-correzioni}:
correggere la rotta di un paio di gradi è molto più rapido che farlo dopo
essersi allontanati. Bastano trenta secondi di ``ehi, vado bene?''. Spesso non
lo si fa per paura di sembrare insicuri o poco autonomi, ma per il relatore è
l'opposto: quando qualcuno gli chiede ``vado bene così?'' e lui risponde
``sì'', si alleggerisce anche del proprio carico mentale. Un capo che assegna
un compito e lo vede solo alla fine vive sempre col dubbio (``avranno capito?
Lo staranno facendo giusto? Arriverà in tempo?''); con un check di dieci
secondi al giorno (``tutto ok, Stefano'') è tranquillo, e la persona va avanti
sapendo che quanto fatto finora va bene. Non serve fare ogni giorno il giro del
dipartimento: anche ogni uno o due giorni basta a evitare la rotta verde.

Per il feedback è lo stesso: \textbf{va dato subito}. Dire oggi ``ieri hai
fatto questa cosa'' è molto diverso da dire ``negli ultimi sei mesi hai fatto
questa cosa''. Il relatore si è sentito dire più volte, alla valutazione di
dicembre, ``da maggio hai fatto questa cosa che non andava''. La sua risposta:
``e aspetti dicembre per dirmelo? Mi hai lasciato sbagliare per sei mesi?''. È
una mancanza di rispetto, e la conversazione non può più essere costruttiva.
Per questo non ama le valutazioni annuali: i manager crescono abituandosi a
prendere appunti delle cose negative per tirarle fuori alla valutazione
(``a marzo hai fatto questa cosa'' ``sì, però poi ho smesso'' ``sì, però a
marzo l'hai fatta'').

In sintesi:
\begin{description}
  \item[Frequentemente] riguarda l'intervallo tra un check e l'altro.
  \item[Recentemente] significa non lasciar passare un anno prima di dare il
        feedback, altrimenti non serve più a niente.
\end{description}

\subsection{Quanto dura un feedback}

Un feedback dura \textbf{da trenta secondi a quindici minuti}; raramente di
più. Chi non ne ha mai dato uno dovrebbe sforzarsi di darne almeno uno, magari
positivo per cominciare. Ogni feedback dato o ricevuto è un
\textbf{deposito nel proprio ``fondo fiduciario''} (\emph{trust fund}), che
costruisce la reputazione futura. Il relatore si è costruito la reputazione di
chi dà feedback il più possibile oggettivi e analitici: giusti o sbagliati che
siano, costringono a pensare. Non impone mai il suo punto di vista (``si fa
come dico io''), ma chiede ``tu come faresti? Come posso aiutarti?''. Il
feedback, per lui, serve a \emph{far crescere qualcun altro}.

\subsection{I manager peggiori: ``a che numero sto pensando?''}

I peggiori manager che il relatore abbia avuto erano quelli che volevano che si
facesse come volevano loro, anche quando chiedevano ``tu come la faresti?''.
In realtà la domanda era ``\emph{a che numero sto pensando?}'': se dicevi 5,
``no, dovremmo fare 9''; se dicevi 9, ``no, 5''. Lo stress costante di dover
indovinare cosa vuole il capo nasce dalla mancanza di fiducia (``se mi chiedi
come fare, fammelo fare'').

Lo stesso gli succede con i CEO che lo chiamano per ``sistemare la cultura
aziendale'' e poi dicono ``non sono sicuro, io farei diversamente''. La sua
risposta: ``lo so che faresti diversamente, per questo mi hai chiamato''. Se si
chiama un esperto, bisogna lasciargli fare il suo lavoro: non si va dal
meccanico a cambiare le gomme per poi dirgli ``non dovresti girare un po' di
più quel bullone?''. Chiedere un parere a un esperto significa dargli anche il
potere di scegliere. Il problema è il \emph{micromanagement}, o in generale il
non accettare punti di vista diversi dal proprio (magari perché si è il
fondatore e si pensa di essere l'unico a saper fare le cose).

\subsection{Dare feedback a un superiore egocentrico}

\paragraph{Come si dà un feedback negativo a un superiore con un grande ego?}
Il trucco è \textbf{non puntare sull'ego, ma sull'azione e sui suoi effetti}.
Non ``tu fai così e io mi sento così'': anzi, si alimenta il suo ego
rendendolo il \emph{main character} dell'avventura, l'eroe che risolve il
problema (``tu, grande manager, aiutaci a risolvere questo problema che
affligge noi che facciamo l'operatività''). Non è l'ideale (``io cambierei
manager''), ma pochissimi possono scegliersi i colleghi: il relatore stesso ha
quasi sempre ricevuto team preassemblati (``ecco le chiavi, questa è la tua
macchina''). Bisogna giocare con le carte che si hanno in mano.

\paragraph{Un esempio pratico.} Con un manager che fa micromanagement e mette
sempre ``il becco'' nel proprio lavoro, si può dire: ``se ti fidi di me per
questo lavoro, con le tue capacità potresti dedicarti a questa cosa più
strategica e di lungo termine: pensa quanto sarebbe più utile per l'azienda''.
In sostanza: ``guarda cosa potresti raggiungere se mi lasciassi fare il mio
lavoro''. Non come un ``tu fai il tuo e io faccio il mio'', ma come un
``ti sollevo da questa parte del carico, che so che ritieni importante: lascia
che ci pensi io, così puoi dedicarti ad altro'', che magari gli dà anche più
visibilità.

Si gioca su ciò che quella persona pensa di saper fare, e che probabilmente sa
anche fare. Anche Michael Scott di \emph{The Office}, per quanto preso in giro,
è raccontato come un venditore bravissimo, che ha chiuso i contratti più
importanti per l'azienda, e rispetto ai manager che gli succedono è forse
``uno dei meno peggiori''. Tutti, anche i manager peggiori, hanno punti di
forza: \textbf{giocate su quelli}.

% ---------------------------------------------------------------------
\section{Il feedback come regalo}
\label{sec:z1-regalo}

Per chiudere, il relatore risponde alla domanda d'apertura, che aveva lasciato
apposta in sospeso: che cos'è il feedback per lui?

\begin{quote}
  \emph{Il feedback è un regalo che mi viene fatto da qualcuno che so che vuole
  il mio interesse.}
\end{quote}

Ha costruito un team di persone che si fidano di lui e che non hanno motivo di
sabotarlo o di ferirlo: per questo interpreta anche le cose peggiori con il
default to trust. È un regalo, e cerca di prenderlo in considerazione il più
possibile.

Ci è arrivato partendo da se stesso. Oltre a essere avverso ai conflitti, si
definisce ``avverso alla comunicazione umana'': paradossalmente, una persona
introversa e timida finita a fare comunicazione, marketing e ora HR (``succede
spesso''). Prima di ordinare una pizza al telefono ha dieci minuti di
preparazione mentale, e se l'ordine delle domande cambia va in confusione. Per
lui un feedback di dieci minuti sono spesso \textbf{tre ore di preparazione}:
costruisce un \textbf{albero dei dialoghi}, come in un'avventura grafica
personale, con tutte le opzioni. Poi spesso finisce con ``hai fatto questa
cosa, andrebbe migliorata'' ``hai ragione, grazie'', e basta.

Così ha pensato: chissà quanti si fanno gli stessi problemi con lui, e passano
ore a prepararsi per dirgli ``questa cosa l'hai fatta male'', sapendo che
riceveranno comunque quell'occhiata fugace del ``come ti permetti''. Eppure
vengono a dirglielo. Non può far altro che rispondere: ``grazie, hai ragione,
cercherò di fare meglio''. È così che ha metabolizzato il ricevere feedback:
capendo che non c'è nessuna cattiva intenzione e che, se non proprio felice
(``ci ho messo una pietra sopra''), può almeno essere \emph{grato} che ci sia
ancora qualcuno che ha così a cuore il suo interesse da dargli feedback.

\paragraph{Il silenzio fa più paura.} Non bisogna spaventarsi dei feedback
negativi che emergono dai sondaggi aziendali: bisogna spaventarsi del
\textbf{silenzio}, quando nessuno dice più nulla. Significa che nessuno si
fida più che si sappiano risolvere i problemi.

\paragraph{La sindrome dell'impostore.} Il suo motivatore personale, all'inizio,
era soprattutto ``non voglio fallire'': per questo cercava di migliorare, e
migliorare è diventato accettare feedback, perché più feedback riceve, più
migliora. Il feedback è un \textbf{acceleratore}, e lui cerca di incentivare
gli altri a darglielo. Resta però sempre la vocina del ``e se si accorgono che
non sei capace?''. Con la \emph{impostor syndrome} convive da anni e
continuerà a conviverci; se ce l'avete, conviveteci anche voi. Avrete sempre
la sensazione che ``se ne accorgeranno'': forse sì, ma sapete fare anche
altre cose. Non concentratevi solo su ciò che non sapete fare. Gli errori
capitano: la perfezione non è di questo mondo.

\paragraph{Esercitarsi in accademia.} Per fare pratica con il feedback,
l'ambiente migliore è quello dei colleghi di corso: persone nella stessa
situazione, con esperienza simile e lo stesso obiettivo, in un contesto in cui
non si può essere licenziati e l'unico effetto a lungo termine è ``ho parlato
con questa persona, ora le cose vanno meglio''. E bisogna avere il coraggio di
parlare anche con i docenti (``dire a Diego: secondo me questa lezione potevi
farla in questo modo''; ``so che è una persona che lo accetta''), senza farsi
bloccare dalla gerarchia. Nell'azienda giusta (come quella in cui lavora ora) il
relatore va sempre dal CEO a dare feedback, sapendo che non verrà licenziato:
il CEO gli ha detto ``ti ho preso per darmi feedback''.

La frase di chiusura:

\begin{quote}
  \emph{Per superare le vostre difficoltà, tutto quello che vi viene chiesto è
  il coraggio di chiedere aiuto.}
\end{quote}

Chiedere aiuto non è una debolezza: è un segno di consapevolezza di sé
(\emph{self-awareness}) e di forza personale. Se tutti lo facessero, i
progetti andrebbero molto meglio; spesso invece ci si fa frenare dal ``non
voglio fare brutta figura'', ``non voglio passare per quello scemo'', e sono
proprio queste paure a rallentarci e a impedirci di raggiungere gli obiettivi.

% ---------------------------------------------------------------------
\section{Domande e risposte finali}
\label{sec:z1-qa-finale}

L'ultima mezz'ora è dedicata a domande libere.

\paragraph{E se una persona peggiora per uno o due anni di fila?} Se nessuno
gliel'ha mai detto, c'è un grosso problema, non suo ma dell'azienda. Se invece
c'è consapevolezza, il percorso è: ``cosa vuoi fare per risolvere?''; ``ho
provato questo, questo e questo''; ``visto che non ha funzionato, vuoi un mio
consiglio? Proviamo con quello?''; ``che supporto ti serve?''. Se, con tutto il
supporto, il tempo e le soluzioni concordate per mesi, il miglioramento non
arriva, per il relatore la soluzione è: \textbf{``dimmi un prezzo per andare
via''}. Preferisce pagare una buonuscita (\emph{severance package}) che
avviare la strada dei richiami formali sulla performance fino al
licenziamento per giusta causa. A un certo punto andare avanti non è un favore
per nessuno.

\paragraph{È vero che nei videogiochi si viene sfruttati? Meglio un piccolo
studio italiano?} Per il relatore si è quasi più a rischio a Blizzard che in un
piccolo studio di Roma o Torino. Nel piccolo studio si sa dove ci si va a
cacciare: non ci sarà il team building a Santo Domingo, ma ci si concentra su
un gioco alla volta e, se va bene, si fa il successivo. A Blizzard si può
lavorare su \emph{Overwatch}, che vende milioni, e perdere comunque il posto,
in modo totalmente scollegato dal proprio lavoro. Gli stipendi più alti
all'estero compensano anche questo rischio, che spesso non dipende da come va
il gioco (anzi, ``il gioco non ha fatto come volevamo'' è spesso una scusa per
licenziamenti che non avevano motivo).

\paragraph{Cambiare settore a trent'anni, da web developer full remote?} Lo
stipendio indicato dallo studente non è fuori dalla portata di vari studi
italiani, anzi è superabile. Non va esclusa la possibilità di lavorare da
remoto per aziende piccole ma estere. Alcuni amici del relatore, pagati con
stipendi italiani, si sono trasferiti a Bali o in paesi con un costo della vita
più basso; il collega di \emph{PUBG} percepisce uno stipendio olandese vivendo
a Milano, che per quanto cara non costa come i Paesi Bassi.

\paragraph{Come si diventa manager?} Dipende dall'azienda: alcune promuovono i
migliori (processo tendenzialmente sbagliato), altre propongono un percorso di
\emph{people management}. Diventare manager significa \textbf{abbandonare ciò
che si sa fare per imparare a fare altro}. Il relatore ha dovuto accettare che
non sarebbe più stato ``il miglior community manager italiano del mondo''
(``eravamo in pochi, quindi era facile''): dedicandosi alle persone non poteva
più fare il suo lavoro bene come prima. È stata una negoziazione con se
stesso: fare il manager significa non essere più l'esperto di qualcosa, ma
saper gestire e far crescere persone. Significa anche \textbf{cedere il
controllo della propria carriera}: se il team decidesse di non fare più nulla,
il manager ne pagherebbe le conseguenze. Il suo lavoro diventa scegliere le
persone giuste e farle lavorare bene. Lui lo esplicita sempre: sa di affidare
al team parte della propria carriera, e il team sa che lui combatterà per
loro, anche prendendosi le colpe.

\paragraph{L'inglese.} \textbf{Va assolutamente imparato}, anche lavorando in
Italia. Non saperlo preclude circa il 95\% del mercato (tutto quello fuori
dall'Italia) e soprattutto il \emph{networking} con persone che potrebbero dare
occasioni future: l'inglese è spesso l'unica lingua in comune con chi lavora
all'estero ed è oggi la lingua universale dei videogiochi.
\begin{itemize}
  \item \emph{Che livello serve?} Non importa B2 o C2: bisogna saper sostenere
        una conversazione, guardare almeno Netflix con sottotitoli in inglese,
        comunicare senza dodici pause e senza ``come si dice\dots''. Il
        certificato non verrà mai chiesto: verrà chiesto se si sa parlare
        inglese. Il relatore ha sentito persone con un C2 che non sapevano
        parlarlo.
  \item \emph{Come allenarlo senza occasioni di conversazione?} Netflix, Prime
        Video, Now TV, Udemy, Coursera, YouTube, Duolingo, siti di scambio
        linguistico con madrelingua, Discord: non sono mai esistiti così tanti
        mezzi gratuiti. Il relatore ha imparato l'inglese giocando (con una
        pronuncia sbagliata, perché leggeva senza sentire, su cui ha poi
        dovuto lavorare); guardare un film in inglese con i sottotitoli in
        italiano insegna anche la traduzione.
  \item Oltre all'ascolto passivo è fondamentale \textbf{parlare}: in una
        flotta di \emph{EVE Online} con 150 persone quasi tutte madrelingua
        inglese si è esposti a decine di accenti e modi di dire (mille modi di
        dire ``come stai'' oltre a \emph{how are you}); lavorando all'estero si
        impara a distinguere gli accenti (Manchester, Londra) e ad apprezzarli
        anche nei media.
\end{itemize}

\paragraph{Quanto deve sapere un junior?} Dipende da quanto è tecnico il
ruolo. Un game designer junior che non sa qualcosa può essere affiancato da un
senior; un programmatore che non sa programmare, in quel momento, non è utile
all'azienda.

\paragraph{Annunci ``junior con cinque anni di esperienza''.} ``Cercasi
ventenne con trent'anni di esperienza.'' Non è l'annuncio per un junior: è un
errore dell'azienda, che in realtà cerca un \emph{mid}. Riceverà profili
junior, si stupirà che non abbiano cinque anni di esperienza, rifarà l'annuncio
e avrà perso sei mesi. Il consiglio: tenere d'occhio quella pagina, perché
dopo qualche mese l'annuncio verrà corretto.

\paragraph{Una struttura per il post-mortem.} A chi chiede un formato di base,
il relatore anticipa quanto vedrà nella seconda giornata a proposito di
brainstorming: una lavagna \textbf{divisa in quattro aree}, su cui ognuno
attacca tutti i post-it che vuole:
\begin{center}
  \begin{tabular}{@{}p{0.44\textwidth}p{0.44\textwidth}@{}}
    \toprule
    \textbf{Smettere di fare} & \textbf{Iniziare a fare} \\
    ciò che non ha funzionato & ciò che non abbiamo fatto e andava fatto \\
    \midrule
    \textbf{Fare di meno} & \textbf{Fare di più} \\
    ciò che ha fatto perdere tempo rispetto all'impatto ottenuto &
    ciò che ha funzionato ma, fatto con troppa leggerezza, non ha dato i
    frutti sperati \\
    \bottomrule
  \end{tabular}
\end{center}
Con i post-it ognuno ha voce e tutto risulta catalogato; per il report basta
una foto della lavagna (o uno screenshot, se la lavagna è digitale, come Miro).

\paragraph{Come sbloccare la paralisi del feedback verso il management?} Uno
studente ha visto manager dire ``non abbiate paura, la mia porta è sempre
aperta'' senza che cambiasse nulla. C'è una dissonanza tra ciò che dicono e
ciò che fanno: è ``la mamma che dice vieni qua che non ti faccio niente''.
Qualcosa deve aver fatto pensare alle persone che non è sicuro. Prima va
sistemato quello; poi si può provare con i \textbf{feedback anonimi}.
L'obiettivo del relatore resta il feedback aperto e firmato, ma a Stormind ha
iniziato con feedback anonimi verso il management per verificare se
l'anonimato fosse un fattore. Non lo era (evidentemente non c'erano problemi a
parlare con il CEO), e la volta successiva intende abbandonarli. Ma è un buon
\emph{check}: se senza anonimato non arriva nulla e con l'anonimato arrivano
cento feedback, il problema è chiaro.

\paragraph{Quando è giusto rivolgersi a un superiore?} La strada giusta resta
parlare direttamente con l'interessato. Ma se si prova e non si trova una
soluzione, la cosa più ovvia è andare dal proprio manager insieme: ``ci
aiuti?''. Quando il relatore parlava di mediare i dialoghi, era proprio questo:
fare da moderatore, aiutare a sbloccare i punti fermi o, se non si riesce a
scegliere tra A e B, scegliere lui. L'escalation, quando necessaria, va fatta.

\paragraph{Il tempo libero veniva rispettato?} No, e se l'è dovuto imporre
lui. Il lavoro gli piaceva così tanto che continuare a casa non gli sembrava
lavorare. Gli si è ritorto contro con il lancio di \emph{Diablo~III}, fatto in
orario americano (le due o l'una di notte in Europa), con l'ufficio chiuso
perché a nessuno era stato chiesto di restare. Il lancio va malissimo, i
giocatori protestano, e lui, sveglio, risponde sui forum fino alle tre, quattro,
cinque del mattino: ``non posso risolvere il problema, ma resto qui ad
aggiornarvi con quello che riesco a sapere dal team di sviluppo''. Il giorno
dopo trova sulla scrivania un \textbf{richiamo formale dell'HR}: aver lavorato
fuori orario senza autorizzazione, condotta non professionale. Per lui è stato
un interruttore: ``ora me ne vado''. Due settimane dopo faceva colloqui con
Riot. Ancora oggi è contrario a quel richiamo, ma gli ha fatto capire come
ragiona un'azienda: ``hai un contratto 9--18, non mi interessa se l'hai fatto
con le migliori intenzioni''. Eppure per anni ha goduto di un'ottima
reputazione in quelle community come ``quello che è rimasto con noi fino alle
cinque del mattino'', e anche l'azienda ne ha beneficiato.

Ci sono poi aziende in cui, soprattutto per i developer, viene richiesto il
famigerato \textbf{crunch}: lavoro straordinario continuativo, senza pause.
Dipende molto dall'azienda; sta venendo (``grazie al cielo'') deprecato, perché
molte aziende capiscono che non è sano ed è visto male anche dal pubblico, ma
esiste ancora.

\paragraph{E se il tuo capo demotiva qualcuno del tuo team?} Qui ``diventa
personale'': il team del relatore si fida di lui perché sa che in questi casi
``diventa una bestia''. ``Te la puoi prendere con me quanto vuoi, ma il mio
team non lo tocchi.'' Lo gestisce direttamente con l'altro manager, per quanto
brutta possa diventare la conversazione. Lo sforzo più grande è \textbf{non
andarne poi a parlare con il team} (``ragazzi, ho fatto questa cosa''), perché
minerebbe la fiducia verso l'alto: non vuole un rapporto ottimo fino a lui e
pessimo oltre. L'obiettivo è che l'altro manager ``si faccia perdonare''
cambiando il proprio comportamento, e che il team percepisca da solo il
cambiamento.

\paragraph{Chiusura della giornata.} Il relatore chiede un feedback collettivo
in chat (``se domani ci dai \emph{more of the same} siamo contenti, oppure
ditemi cosa limare''), si scusa per i ripetuti crash della connessione e
ringrazia per la proattività nelle domande, che gli servono anche come
\emph{check}: se nessuno chiede nulla su un argomento, significa che può
andare avanti. Racconta di aver passato una settimana a decidere che cosa
mettere nelle slide e di averle poi preparate in poco tempo (rifacendone
alcune fino alle tre di notte). Ricorda infine di usare il test di
personalità per sé e per il proprio gruppo, facendolo fare a chi non l'ha
ancora fatto: tornerà molto utile per cooperare meglio.

% ---------------------------------------------------------------------
\section{Anticipazioni sulla seconda giornata}
\label{sec:z1-domani}

La seconda giornata del workshop sarà più pratica: ``oggi vi ho spiegato
perché, domani vi insegno come lo fanno queste aziende''. Gli argomenti
previsti:
\begin{description}
  \item[Time management] come organizzarsi una giornata professionale in
        azienda: come strutturare una \emph{task list}, cosa fare e cosa non
        fare, come e cosa prioritizzare.
  \item[Brainstorming] i gruppi dovranno fare moltissimi meeting (come facciamo
        il gioco, come lo chiamiamo, che scopo gli diamo). Se non si sanno fare
        bene, i brainstorming sono \textbf{dannosi ed eterni}: un brainstorming
        dovrebbe durare \textbf{al massimo due ore}. Verranno mostrati metodi
        che accorciano la durata, massimizzano l'effetto e, soprattutto,
        \textbf{fanno partecipare tutti}: non c'è niente di peggio di un
        brainstorming in cui parlano in due per il 90\% del tempo. Il classico
        ``ragazzi, non esistono idee sbagliate, sparate'' è il modo peggiore di
        farlo; non è colpa di nessuno, ma si può fare diversamente.
  \item[Guidare da dentro] come essere leader nel proprio team anche senza
        esserne il leader: sviluppare responsabilità personale e verso gli
        altri. Non essere il manager di qualcuno non significa non esserne
        responsabili: la persona accanto è responsabilità di tutti, perché si
        è un team e non un insieme di individui. ``Se qualcuno cade, tutti lo
        facciamo rialzare: non c'è l'addetto al rialzamento.'' È la versione
        \emph{advanced} della giornata: come comportarsi da manager senza
        averne il titolo.
\end{description}
Torneranno inoltre il \emph{service level agreement} e il tema della
condivisione dei successi e della protezione del team dai fallimenti.

% ---------------------------------------------------------------------
\section*{Punti chiave del capitolo}
\addcontentsline{toc}{section}{Punti chiave del capitolo}

\begin{itemize}
  \item ``La capacità di un gruppo di persone di fare cose straordinarie
        dipende dal modo in cui queste persone operano insieme.'' Il 90\% dei
        problemi dei team sono \textbf{problemi di comunicazione}.
  \item Le aziende si liberano sempre più presto dello \textbf{useless
        genius}: le competenze si trasmettono, il carattere no. Il
        \textbf{fit culturale} va valutato anche dal candidato.
  \item Non esiste un unico modo di dare feedback: dipende dal contesto e
        soprattutto da \textbf{chi lo riceve}. Il feedback è una somma di
        piccoli miglioramenti periodici; tenersi le cose nel cassetto non
        serve mai.
  \item I team attraversano \textbf{forming, storming, norming, performing e
        adjourning}. Lo storming non si salta; si passa la maggior parte del
        tempo in forming; la bravura sta nell'accelerare i passaggi e nel
        fare bene il \textbf{post-mortem}, anche dei successi.
  \item Gli obiettivi di un progetto vanno fissati \textbf{all'inizio}: solo
        così si può valutare se è andato bene e scegliere tra A e B.
  \item Il circolo virtuoso è \textbf{feedback $\rightarrow$ miglioramento
        $\rightarrow$ fiducia}. Accettare un feedback non significa fare come
        dice l'altro: decide chi ha la responsabilità, ma deve argomentare.
  \item ``Aiutami a poterti aiutare'': meglio dire in anticipo ``potrei non
        farcela'' che ``non ce l'ho fatta''. Chiedere aiuto non è una
        debolezza.
  \item Quando si dà feedback \textbf{non si porta la soluzione}: la si
        costruisce insieme a chi lo riceve.
  \item La cultura è forte quando si lavora \textbf{per gli altri}; il lavoro
        è valutato come opera d'insieme (anche con bonus di team).
  \item Come in \emph{Among Us}, i team sprecano energie a cercare il
        traditore invece di fare le task. \textbf{Default to trust}: chi ha
        sbagliato ha agito con le informazioni che aveva. Cercare
        \textbf{la causa, non il colpevole}: ``che cosa ti ha fatto pensare
        che fosse la scelta giusta?''.
  \item Carriere e performance non sono lineari: vanno valutate le
        \textbf{traiettorie}, non un punto dell'anno. Le valutazioni puntuali
        generano \emph{lying, hiding, faking}.
  \item Confermare sempre le istruzioni facendosele ripetere; ``non ho
        capito'' non è una colpa.
  \item Dare feedback: sull'\textbf{azione} e sui suoi \textbf{effetti}, specifico,
        presto, uno a uno, adattato al carattere, senza terze parti, senza
        ``non mi piace'', senza offendersi se non viene seguito; poi fare il
        debriefing.
  \item Ricevere feedback: \textbf{ascolta, chiedi dettagli, accetta, agisci}.
        Il feedback positivo serve a non peggiorare ciò che si fa bene; il
        feedback accidentale e il \emph{drive-by} non servono a nulla.
  \item Feedback \textbf{frequenti e recenti}: tanti piccoli \emph{sanity
        check} costano meno di una grande correzione di rotta.
  \item Il feedback è \textbf{un regalo} di chi ha a cuore il nostro
        interesse; fa paura il silenzio, non le critiche.
\end{itemize}

%    % ... altre lezioni zero ...
%    \part{Il corso di Game Design}
%    \renewcommand{\thechapter}{\arabic{chapter}}
%    \setcounter{chapter}{0}
%    % =====================================================================
%  Capitolo 1 — Introduzione al corso e al mestiere del game designer
%  Fonte: trascrizione della lezione GD1 (Game Design, primo anno)
%  Pacchetti richiesti nel preambolo del libro:
%    \usepackage[italian]{babel}, \usepackage[T1]{fontenc},
%    \usepackage[utf8]{inputenc} (solo pdfLaTeX), \usepackage{booktabs}
% =====================================================================

\chapter[Il mestiere del game designer]{Il mestiere del game designer: introduzione al corso}
\label{cap:introduzione-game-design}

Questo primo capitolo fa da soglia. Prima di affrontare metodi, framework e
documenti, è necessario capire \emph{in che modo} si impara il game design,
\emph{con chi} lo si fa, \emph{con quali strumenti} e, soprattutto,
\emph{che cosa fa concretamente} un game designer nel suo lavoro quotidiano.
Il capitolo segue l'impostazione della lezione inaugurale del corso: si parte
dal contesto formativo (il corso come simulazione di lavoro), si passa al
metodo di apprendimento (l'\emph{inquiry-based learning}), si elencano
letture e risorse di riferimento, e si arriva infine a una prima definizione
operativa del ruolo del designer, dei limiti entro cui lavora e delle sue
possibili specializzazioni.

Una premessa terminologica: il game design è una disciplina giovane e quasi
tutta la sua letteratura è scritta in inglese. Per questo motivo molti termini
tecnici verranno lasciati in lingua originale (\emph{pipeline},
\emph{pitching}, \emph{tuning}, \emph{feature team}, \dots): sono i termini che
si useranno nel lavoro, ed è importante imparare a riconoscerli e a usarli
correttamente fin da subito.

% ---------------------------------------------------------------------
\section{Il contesto: un corso come simulazione di lavoro}
\label{sec:simulazione-lavoro}

\subsection{Il principio fondamentale}

Il primo concetto da fissare è che il corso \textbf{si presenta come una
simulazione del lavoro}. Dal momento in cui si entra in aula si viene trattati
come \emph{potenziali professionisti dell'industria}: non soltanto studenti,
ma lavoratori effettivi di uno studio (l'accademia), con una gerarchia, dei
referenti, dei compiti e delle scadenze.

Ne derivano alcune conseguenze pratiche:
\begin{itemize}
  \item il docente non assegna un ``compitino'' da risolvere per accumulare
        punti o presenze: ogni \emph{task} è pensato per far dedicare tempo
        all'apprendimento di qualcosa di concreto;
  \item non consegnare un task non danneggia il docente (che ``vive ancora
        domani e viene pagato ugualmente''), ma lo studente, che perde
        l'occasione di mettere le mani su un argomento e diventare più abile;
  \item docenti e studenti devono trattarsi \textbf{da colleghi}, perché è
        molto probabile che un domani lo diventino davvero.
\end{itemize}

\subsection{Perché l'impegno conta più che mai}

Il lavoro del game designer è descritto come \emph{molto romantico, ma anche
molto duro}, e sempre più duro col passare degli anni. Durante il periodo del
Covid l'industria ha visto grandi investimenti; quei ``rubinetti'' sono stati
chiusi da tempo, il rapporto con i publisher è diventato più complesso e le
opportunità lavorative non sono più quelle di tre o quattro anni fa. Questo
significa:
\begin{itemize}
  \item più competizione per ogni posizione;
  \item aspettative più alte: ci si aspetta che i candidati siano ancora più
        preparati e sicuri delle proprie competenze;
  \item un unico modo per diventarlo: \textbf{lavorare}. Più tempo si dedica
        oggi ad apprendere il mestiere, più facilmente si sarà assumibili
        domani e più agevole sarà il lavoro di tutti i giorni.
\end{itemize}

\subsection{Sbagliare qui per non sbagliare fuori}

Il mantra del corso è: \textbf{quanto più provate qui, anche in modo goffo e
sbagliando, tanto meno sbaglierete là fuori}. In accademia, per quanto la
simulazione sia realistica, nessuno toglie lo stipendio e nessuno licenzia da
un momento all'altro. Nell'industria sì, e un errore può costare molto denaro
all'azienda, ancora di più se l'azienda è la propria.

\subsection{Il valore del fallimento}

Il fallimento, se affrontato con la mentalità corretta, è probabilmente
\emph{la cosa più vantaggiosa} che si possa fare durante la formazione:
\begin{itemize}
  \item un successo (``avete superato questa deadline'') conferma le proprie
        abilità, ma insegna relativamente poco;
  \item dopo $n$ fallimenti si impara moltissimo, e i successi arrivano;
  \item è fondamentale non demoralizzarsi: il feedback negativo non è mai un
        giudizio sulla persona, ma su ciò che è stato consegnato.
\end{itemize}

Il docente riconosce apertamente che dedicare tempo ed energia a un progetto e
vederlo rifiutato, o vederlo non uscire mai, è demotivante: è successo a lui e
lo ha visto accadere anche su progetti molto grandi, con team numerosi. Il
punto però non è ciò che si è ottenuto materialmente, ma \textbf{quanto si è
cresciuti}. Da ogni esperienza ci si portano a casa due cose:
\begin{enumerate}
  \item \textbf{l'esperienza}: quanto si è diventati bravi, veloci,
        praticamente abili nell'applicare un framework, quanto si è cresciuti
        come designer;
  \item \textbf{il rapporto umano}: nel lavoro si ``suda sangue'' fianco a
        fianco con altre persone, e questo crea legami molto forti che
        portano avanti l'industria nel senso più positivo del termine: nascono
        progetti privati, ci si porta da un'azienda all'altra, si formano team
        con basi solide.
\end{enumerate}

Va evitato anche l'eccesso opposto: la cultura del \emph{grinding} si è
appropriata di questi discorsi in modo distorto. Non bisogna diventare robot.

% ---------------------------------------------------------------------
\section{Organizzazione del corso e gerarchia}
\label{sec:gerarchia-corso}

\subsection{La struttura gerarchica}

La struttura dei referenti del corso ricalca volutamente una
\textbf{gerarchia tipica dell'industria}:
\begin{itemize}
  \item al vertice c'è il \textbf{direttore del corso} (Diego Ricciuti), che
        corrisponde al \emph{creative director}: supervisiona la visione del
        progetto (il concetto di ``visione'' verrà approfondito più avanti);
  \item sotto di lui ci sono i \textbf{referenti di dipartimento}, uno per
        disciplina, con un corrispettivo per ciascuna sede (Milano e Roma);
  \item in fondo ci sono gli studenti, che occupano il primo gradino di
        un'azienda: l'\emph{intern} (stagista) o il \emph{junior}, il passo
        che precede il lavoratore ``regolare''.
\end{itemize}

\begin{table}[htbp]
  \centering
  \begin{tabular}{@{}p{0.34\textwidth}p{0.29\textwidth}p{0.29\textwidth}@{}}
    \toprule
    \textbf{Disciplina} & \textbf{Milano} & \textbf{Roma} \\
    \midrule
    Direzione (\emph{creative director}) & \multicolumn{2}{p{0.58\textwidth}}{Diego Ricciuti} \\
    Game design      & Mattia Podini        & Giacomo ``Jack'' Danna \\
    Level design     & Stefano Zocchi       & Andrea Dispoli \\
    Narrative design & \multicolumn{2}{p{0.58\textwidth}}{Daniele Fusetto (lezioni condivise), Giacomo Masi} \\
    Programmazione (Unreal) & Alessandro Pregnolato & Matteo Velardi \\
    \bottomrule
  \end{tabular}
  \caption{Referenti del corso per disciplina e sede (nomi come riportati in lezione).}
  \label{tab:referenti}
\end{table}

Per il primo anno, il referente diretto degli studenti di game design è il
docente di game design: è a lui che ci si rivolge per dubbi e domande. Altre
figure utili sono il direttore didattico (severo, ``giustamente severo'', da
contattare solo per questioni particolarmente importanti), l'amministrazione
(contratti e documentazione) e la segreteria.

\subsection{I canali di comunicazione}

\begin{description}
  \item[Discord] è il canale per le comunicazioni ufficiali del corso e per
        le domande che non richiedono urgenza immediata; può avere qualche
        ritardo.
  \item[E-mail] è il canale per tutto ciò che è formale o ``ancora più
        ufficiale'' e per le comunicazioni rapide e immediate.
  \item[Telegram] esistono residui di vecchi canali precedenti a Discord, da
        evitare.
\end{description}

Le domande che valgono per tutti vanno fatte in classe (anche da un portavoce
a nome del gruppo); quelle singole e personali via Discord o, se formali, via
e-mail. Se il docente non risponde a una domanda c'è un motivo, ad esempio
ragioni legali, e lo dirà esplicitamente.

\subsection{I tre moduli e la settimana tipo}

Le tre discipline principali del primo anno (e non solo) sono:
\begin{itemize}
  \item \textbf{Game Design} (tendenzialmente il lunedì);
  \item \textbf{Level Design} (tendenzialmente il mercoledì);
  \item \textbf{Programmazione} (tendenzialmente il venerdì).
\end{itemize}
Nel primo mese le lezioni sono tutte di game design (lunedì, mercoledì e
venerdì). Sono previste anche lezioni di narrativa, ma il piano formativo si
aggiorna ogni anno: le ore di narrativa sono state ridotte a favore di quelle
di programmazione. Le poche lezioni di narrativa daranno quindi dei
\emph{cenni}, non la preparazione di un narrative designer. Ciò che serve
davvero è imparare \textbf{come comunicano gli altri dipartimenti} e come
lavorare con persone che fanno un lavoro diverso ma complementare al proprio.

\subsection{Le due anime delle lezioni di game design}

\begin{description}
  \item[Design methods (metodologie)] lezioni prevalentemente teoriche, ma con
        una buona componente pratica, in cui si studiano metodologie di design
        e i cosiddetti \emph{framework} dell'industria.
  \item[Documentazione e pitching] il cuore del lavoro quotidiano del
        designer.
\end{description}

\paragraph{La documentazione.}
Il docente si definisce ``il più grande nerd della documentazione'' che si
possa incontrare. È una parte dura, perché a pochi piace documentare, ma il
lavoro del designer \emph{è principalmente questo}: scrivere documenti e, da
quei documenti, passare alla pratica e all'implementazione. Chi non sa
documentare bene parte male. Quando non si è presenti, il documento è il
proprio \textbf{biglietto da visita}: qualcun altro lo leggerà, quindi va
scritto nel modo più chiaro e comprensibile possibile.

Un esempio concreto del perché la documentazione conta: CD Projekt Red ha
sviluppato un proprio engine proprietario, ma buona parte della sua
documentazione è stata ``buttata all'acqua''. Chi deve usarlo o lo sa già
usare, o ha competenze di un altro engine da trasferire, oppure è in
difficoltà. Si documenta \textbf{perché altri possano leggere}.

\paragraph{Il pitching.}
Il \emph{pitching} è l'esposizione, la presentazione orale. La parte creativa
del team presenta spesso:
\begin{itemize}
  \item le proprie idee ai colleghi;
  \item il gioco a convention ed eventi;
  \item il pitch del gioco agli investitori, per ottenere i fondi necessari a
        realizzarlo.
\end{itemize}
Saper parlare e comunicare chiaramente le idee è fondamentale. In sintesi:
\emph{il game designer parla}, e bisogna farsene una ragione.

\subsection{Altre componenti del corso}

\begin{description}
  \item[Gioco critico e playtest] la classe viene divisa in due gruppi: uno
        gioca, l'altro osserva. Chi osserva non applaude ``come le
        scimmiette'': annota, segna e analizza il \emph{playtest}, esattamente
        come si fa nell'industria. Il playtest fa parte delle mansioni del
        designer, ma esiste anche come specializzazione autonoma.
  \item[Pipeline] termine che indica tutto ciò che riguarda i processi
        formalizzati dell'industria; di solito si intende la \emph{pipeline di
        produzione}: come si divide il team, qual è la sua gerarchia, come il
        \emph{producer} gestisce le deadline. In breve: \emph{come} e
        \emph{quando} si fanno le cose. Al primo anno è meno rilevante (si
        dovrà ``nuotare'' un po' da soli); dal secondo anno si vedranno le
        metodologie effettivamente applicate ai progetti digitali. Si
        scoprirà quanto sia difficile lavorare a un progetto digitale, e
        ancor più fisico, senza un piano di sviluppo comprovato, come quelli
        usati per qualsiasi software in qualsiasi parte del mondo.
  \item[Engine] il motore di riferimento è \textbf{Unreal Engine}
        (sezione~\ref{sec:engine}).
  \item[Programmazione] principalmente tramite \emph{Blueprint} in Unreal,
        più eventuali sistemi di versionamento; il corso parte da zero.
  \item[Progetti pratici] probabilmente la cosa più importante. Al primo
        anno i progetti di gruppo restano circoscritti al corso di design,
        perché mancano ancora alcune nozioni (come comunicare, affrontare
        problemi, dare e ricevere feedback). In seguito si allargheranno a
        progetti digitali, molto più difficili di quelli fisici, con membri
        di altri team: un rapporto non unilaterale, ma ``covalente''.
\end{description}

% ---------------------------------------------------------------------
\section{Il metodo: l'inquiry-based learning}
\label{sec:inquiry}

\subsection{Che cos'è}

Il corso è costruito secondo l'\textbf{inquiry-based learning}
(apprendimento basato sull'indagine), un approccio pedagogico nato intorno
agli anni Sessanta, che si richiama a pensatori come Dewey e Piaget. In
generale, le discipline da cui il design attinge sono quasi sempre la
sociologia, la psicologia e, talvolta, l'economia.

Il funzionamento è un ciclo:
\begin{enumerate}
  \item il docente fornisce metodologie e framework dell'industria durante
        la lezione;
  \item lo studente, fuori dalle ore di lezione, lavora sui \emph{task}:
        mette a terra le nozioni, le maneggia, le assorbe meglio;
  \item da questo lavoro nascono nuovi dubbi;
  \item i dubbi tornano in aula: lo studente chiede, il docente risponde
        (ammesso che dare la risposta sia utile, e ammesso che una risposta
        esista).
\end{enumerate}

\subsection{Non sempre esiste una risposta giusta}

Non tutto ciò che si studia ha ``la risposta migliore possibile in assoluto
nel vuoto''. Esistono metodologie, limiti e framework che aiutano il lavoro,
ma non sono fissi: cambiano da azienda ad azienda, da lavoro a lavoro, da
progetto a progetto, e a volte vengono stravolti. Sta alla
\textbf{sensibilità del designer} capire quale strumento o framework usare in
un dato momento, e questa sensibilità si impara solo con l'esperienza. Nel
lavoro nessuno dirà ``usa questa cosa'' o ``è meglio quest'altra'': sarà il
designer a scegliere.

Anche l'esperienza del docente ha confini: sei-sette anni di lavoro, un modo
personale di approcciare il design e un certo insieme di progetti (ad esempio
nessun MMO). Oltre quei confini si cercano risposte nella ricerca, studiando
come aziende con esperienza specifica (ad esempio Riot) hanno affrontato
problemi simili e hanno condiviso le loro soluzioni, sempre valutando quali
metodi siano adatti alla dimensione del team, al tipo di progetto e al tempo
disponibile.

\subsection{Sensibilità, non memoria}

Non interessa imparare a memoria un framework o la teoria di un autore. Certe
terminologie vanno conosciute, ma per ragioni di comunicazione e chiarezza.
Ciò che conta è sviluppare la \textbf{capacità di farsi domande}, di farle
agli altri, di fare ricerca e di trovare una propria risposta. È esattamente
ciò che si fa nel lavoro: davanti a qualcosa che non si sa fare, l'abilità
del designer sta nel cercare come qualcun altro ha risolto un problema simile.

\subsection{Le lezioni come corsi di aggiornamento}

Le lezioni frontali simulano i \emph{corsi di aggiornamento} aziendali: le
sessioni che un lead di design tiene al suo team di design (o il lead di
level design al suo team, e così via). Alla fine dell'aggiornamento c'è quasi
sempre un lavoro da fare, un obiettivo da portare a terra:
\textbf{l'obiettivo della settimana}. Si lavora di settimana in settimana, e
il perché verrà chiarito in seguito.

Struttura tipica di una lezione:
\begin{enumerate}
  \item breve sessione di domande e risposte sulla lezione precedente;
  \item analisi di eventuali esercizi particolarmente critici;
  \item lezione vera e propria;
  \item discussione di gruppo, se c'è tempo;
  \item assegnazione del task della settimana.
\end{enumerate}
Il docente tende a ``parlare e straparlare'', ad aprire parentesi e a saltare
da un argomento all'altro: l'unico modo per fermarlo è che siano gli studenti
a interromperlo con domande.

\subsection{La gerarchia della competenza}

L'inquiry-based learning può essere rappresentato come una scala di
competenza:
\begin{enumerate}
  \item la sola lezione frontale porta a \textbf{conoscenze incomplete},
        perché manca la parte pratica e di ricerca;
  \item il lavoro a casa (task, ricerca) porta a una competenza \textbf{più
        completa}, non sbagliata in sé, ma ancora parziale e complementare ad
        altre;
  \item la \textbf{condivisione} dei task e dei feedback con gli altri porta
        idealmente alla competenza più completa possibile;
  \item il picco della competenza è la sua \textbf{applicazione nei
        progetti} (fisici o digitali, di gruppo): ``con le cose che avete
        imparato e ricercato, riuscite a fare un gioco?''. Quando la risposta
        diventa positiva, o comunque si inizia a navigare nei processi di
        design per realizzarne uno, quella è la vittoria.
\end{enumerate}
Non si parte subito creando giochi: si parte dall'analisi, dal design e dalle
regole generali di design, prima ``nel vuoto'', poi contestualizzando i
framework, infine unendo i puntini nei progetti.

Rimarranno comunque delle lacune, per tutti, docente compreso, perché il game
design \textbf{non è una scienza esatta}, anzi, a rigore non è affatto una
scienza. Prende in prestito metodologie dalle discipline scientifiche, ma
contiene una componente imprevedibile e incalcolabile: \textbf{la
produzione}. Si possono applicare perfettamente le regole di un framework, e
quel framework può comunque non essere adatto al gioco, al proprio modo di
sviluppare o al tempo di produzione disponibile (alcuni framework richiedono
una ricerca e un'attenzione che il tempo non concede).

C'è poi la componente creativa, la parte ``romantica'' del lavoro: esiste, ma
è \textbf{metà del lavoro}. L'altra metà è ricerca, metodo scientifico al
100\%. Dagli stessi dati $n$ designer diversi, con sensibilità diverse,
possono ottenere $n$ risposte diverse. Non tutte saranno valide: vanno
applicate, messe a terra, e questa scrematura fa parte del lavoro.

\subsection{I livelli dell'apprendimento}

Un secondo modello, a gradini, descrive i livelli di apprendimento (è la
cosiddetta tassonomia SOLO):
\begin{description}
  \item[Pre-strutturale] il docente dà il compito e lo studente lo risolve,
        possibilmente il prima possibile. È esattamente ciò che \emph{non} si
        vuole.
  \item[Uni-strutturale] si affronta una sola disciplina: è il primo passo
        necessario, ed è il motivo per cui nel primo mese ci sono solo lezioni
        di game design.
  \item[Multi-strutturale] entrano in gioco altri dipartimenti: level design,
        game design, programmazione e così via.
  \item[Relazionale] si capisce che alcune cose valgono per più
        dipartimenti: alcuni framework di narrativa derivano dal design o
        possono essere usati per fare design; le regole di documentazione si
        applicano anche ad altri dipartimenti; le regole di level design si
        mescolano con quelle di game design. Dopotutto, per fare un gioco
        servono molti dipartimenti.
  \item[Astratto esteso] il punto finale: muoversi autonomamente e capire
        cosa usare e quando, sulla base della propria esperienza. L'unico modo
        per arrivarci è fare pratica e fare progetti.
\end{description}

% ---------------------------------------------------------------------
\section{Task, feedback e regole del corso}
\label{sec:task-feedback}

\subsection{I task come opportunità}

I task non sono strettamente obbligatori (esiste un numero minimo di
consegne), ma sono \textbf{opportunità}. Quanto e come ci si lavora dipende
dallo studente. Un task fatto in cinque minuti, o in cui si è stati ``furbi'',
non verrà letto né riceverà feedback.

Domande da porsi davanti a ogni task:
\begin{enumerate}
  \item \textbf{Perché} mi è stato assegnato? Qual è l'obiettivo?
  \item Come ci lavoro per ottenere un risultato \textbf{che abbia senso}? Se
        la ricerca e i dati indicano una direzione, è inutile andare in
        un'altra.
\end{enumerate}

\subsection{Quanto lavorare: un calcolo indicativo}

Se il corso è una simulazione di lavoro, vale il carico di un lavoro normale:
\begin{itemize}
  \item un lavoro a tempo pieno comprende da 30 a 45 ore settimanali; si
        considera una media di \textbf{40 ore};
  \item si sottraggono circa 9 ore di lezione: restano $\approx 31$, diciamo
        \textbf{30 ore};
  \item con circa 3 task a settimana: $30 / 3 = $ \textbf{circa 10 ore per
        task}.
\end{itemize}
Il numero può variare, ma rende l'idea: dedicare un'ora a un task e poi fare
altro non fa guadagnare nulla a nessuno. Lo studente ``butta soldi'', il
docente risparmia tempo leggendo il task in un minuto (o non leggendolo).
Dedicando invece attenzione, ricerca, cura nella formattazione e
sperimentazione, e facendolo con cognizione di causa, si impara.

\subsection{Niente voti, solo feedback}

Ad eccezione del corso di programmazione, non ci sono voti (né ``ottimo'', né
10, né 100). Ci sono \textbf{feedback}:
\begin{itemize}
  \item in game design, generalmente sulla documentazione;
  \item in level design, sui progetti pratici.
\end{itemize}
Il docente indica cosa va bene e cosa va modificato, concentrandosi in
particolare su ciò che va cambiato. Il feedback \textbf{non è mai un giudizio
sulla persona}, ma su ciò che è stato consegnato nel concreto.

\subsection{La condivisione dei task}

Tutti i task consegnati vengono \textbf{condivisi} con la classe, e possono
essere letti solo dopo aver ricevuto i feedback del docente. Si cresce così
non solo con i propri feedback, ma anche con quelli dati agli altri. Regola
importante: ciò che il docente dice a uno studente \textbf{vale per tutti}. Se
a una persona viene detto ``questa cosa non va mai più fatta'', da quel
momento vale per l'intera classe.

\subsection{Non copiare}

Copiare i task degli altri è una delle cose peggiori che si possano fare, non
per una questione morale, ma perché annulla il senso del percorso. Si può
dedicare a ogni task il tempo e l'attenzione che si ritengono opportuni, ed è
normale che col tempo ci si specializzi (preferendo, ad esempio, il level
design o la programmazione tecnica) e si investa meno dove si sente meno
necessità. In questi casi l'alternativa corretta è:
\begin{itemize}
  \item non consegnare;
  \item studiare ciò che hanno fatto gli altri;
  \item eventualmente riconsegnare in seguito.
\end{itemize}
Ci saranno momenti in cui si verrà sommersi dalle cose da fare: bisogna
imparare a scegliere.

\subsection{Il ``non leggo''}

Il docente ha piena facoltà di non dare alcun feedback a un task (il
cosiddetto \textbf{``non leggo''}) se:
\begin{itemize}
  \item non si è lavorato con attenzione;
  \item non si rispettano le indicazioni date;
  \item è stata violata una regola estremamente importante.
\end{itemize}
Un task ``non letto'' equivale a un task non consegnato e non entra nel
conteggio. Attenzione: un task \emph{perfetto} e un task che \emph{fallisce}
nel comprendere o applicare un framework sono \textbf{ugualmente validi},
purché ci si sia impegnati e si sia dedicato tempo, pazienza e ``sudore''.

\paragraph{Riconsegna.} Un task non letto, in ritardo o che si vuole
migliorare può essere riconsegnato e verrà conteggiato, a patto che:
\begin{itemize}
  \item i problemi segnalati siano stati risolti;
  \item sia stato fatto ulteriore lavoro, integrando anche tutto ciò che nel
        frattempo è stato detto agli altri colleghi nei loro feedback.
\end{itemize}
Il task va sviluppato con le competenze che si hanno nel momento della
riconsegna, non con quelle del momento della prima consegna.

\subsection{Requisiti per superare l'anno}

\begin{itemize}
  \item consegnare almeno l'\textbf{80\% dei task} assegnati (soglia
        malleabile);
  \item partecipare almeno all'\textbf{80\% delle lezioni}, in presenza o
        da remoto, preferibilmente in presenza; i ritardi vengono valutati
        caso per caso;
  \item comportarsi da persone corrette e rispettose (è stata firmata una
        carta dello studente): non urlare, non alzare le mani, non
        vandalizzare, non compiere atti illegali.
\end{itemize}
Impegnarsi di più porta inoltre a essere selezionati per varie opportunità.

\subsection{Scadenze}

\begin{itemize}
  \item la scadenza di un task è \textbf{sempre prima dell'inizio della
        lezione} successiva indicata;
  \item se la classe, come team, mantiene comportamenti scorretti, il docente
        può \textbf{anticipare le deadline} per tutti.
\end{itemize}

\subsection{Accessibilità e DSA}

Gli studenti con disturbi specifici dell'apprendimento o neurodivergenze
(dislessia, ADHD, autismo, ecc.) sono invitati a segnalarlo. Quando si
affronteranno le regole di formattazione e documentazione (errori
grammaticali, mancanze, incoerenze) se ne terrà conto e si potranno trovare
dei compromessi, pur ricordando che nel mondo del lavoro nessuno farà
``sconti''. La documentazione può essere inviata all'amministrazione o alla
segreteria, che la comunicherà a tutti i docenti, ma il consiglio è parlarne
anche direttamente con ogni singolo docente.

% ---------------------------------------------------------------------
\section{Le competenze trasversali}
\label{sec:soft-skill}

\subsection{I giochi si fanno in gruppo}

Si possono fare giochi da soli, ed è un'impresa ammirevole; ma lavorando da
soli \textbf{tutto ciò che si guadagna in libertà si perde in produttività}.

\paragraph{Il caso Stardew Valley.} \emph{Stardew Valley} è stato realizzato
da una sola persona, Eric Barone, in anni di sviluppo e di sacrifici: faceva
un altro lavoro per mantenersi (in un cinema) e di notte programmava, faceva
l'arte e componeva la musica. Il risultato è stato straordinario e lo ha reso
ricco, ma si tratta di un caso estremamente raro. Quanti giochi simili, nati
nelle stesse condizioni, sono falliti? Molti di più.

\paragraph{Il rischio d'impresa.} Lavorare da soli non è sbagliato di per sé
e ha aspetti molto positivi: si ha il privilegio di creare esattamente ciò che
si vuole. Ma ci si fa carico anche di tutto il \textbf{rischio d'impresa}: se
si lavora a un gioco per cinque anni e poi il gioco va male, resta il
vantaggio di averlo pubblicato, ma il dopo è sperare di vivere di quella
release o di finanziare con essa il gioco successivo. Le grandi aziende hanno
gli stessi problemi, ma li distribuiscono su più persone e più responsabilità.

L'obiettivo di un professionista è \textbf{saper lavorare in diverse
condizioni}, e la più comune è quella del dipendente (o del freelance /
libero professionista con partita IVA): si viene assunti per le proprie
capacità da un'azienda, piccola o grande che sia, e si lavora con altre
persone.

\subsection{Essere persone con cui si vuole lavorare}

Se le persone non vogliono lavorare con voi, quello è \textbf{il primo motivo
di licenziamento}. E se non si viene licenziati, il clima nell'azienda
diventa terribile: la produttività crolla, la motivazione crolla. Le
\emph{soft skill}, cioè l'insieme di abilità che rendono una persona
piacevole con cui lavorare, vanno coltivate e diventano sempre più importanti
col passare degli anni (relativamente meno ora, molto di più quando i team
saranno più grandi). Anche chi è un ``lupo solitario'' deve imparare a:
\begin{itemize}
  \item lavorare con altre persone;
  \item comunicare con altre persone;
  \item trovare un punto in comune da cui partire per fare le cose.
\end{itemize}

\subsection{Indipendenza e proattività}

Sono le altre due parole chiave del corso:
\begin{description}
  \item[Indipendenza] spesso si verrà lasciati da soli, non per cattiveria ma
        per sfida: quando si cerca una risposta da soli, ciò che conta non è
        tanto la risposta quanto \textbf{il processo}. Una risposta data dal
        docente non vale niente.
  \item[Proattività] chiedere, domandare, agire da soli senza aspettare
        direttive, perché spesso quelle direttive non ci sono, a maggior
        ragione nei team piccoli.
\end{description}
Tutto questo si concretizza nel parlare, nel presentare e nel fare
\emph{post-mortem} (l'analisi a posteriori di un progetto concluso, che verrà
approfondita più avanti), e in generale nell'attenzione alla crescita
personale.

\subsection{Fare domande}

Il corso insiste sulle domande: alzare la mano, intervenire anche ``a
bomba'', fermare il docente ogni volta che qualcosa non è chiaro. Se una
persona non conosce qualcosa, molto probabilmente c'è almeno un'altra persona
che non la conosce. Il docente può dare per scontata la conoscenza di termini
che non lo sono, ed è del tutto normale, soprattutto all'inizio.

% ---------------------------------------------------------------------
\section{Strumenti: engine, Unreal e materiali}
\label{sec:engine}

\subsection{Unreal Engine e gli altri motori}

Il motore di riferimento del corso, usato in level design, è \textbf{Unreal
Engine}: ``un gran mattone'', ma che dà a un designer molte possibilità. Il
vantaggio principale è che permette a un designer non pratico di
programmazione di \textbf{prototipare tramite i Blueprint} (il sistema di
scripting visuale di Unreal). Questo non impedisce di passare a Unity o ad
altri engine, e viceversa: il docente stesso ha studiato e lavorato a lungo
con Unreal, ma la sua ultima esperienza professionale (come level designer) è
stata su Unity, un engine che non conosceva, a dimostrazione che il salto da
un engine all'altro è possibile.

Panorama dei motori citati:
\begin{itemize}
  \item \textbf{Unreal Engine} e \textbf{Unity}: i due più diffusi, con
        modelli di sorgente ``più o meno'' aperto;
  \item \textbf{Godot}: engine open source;
  \item \textbf{CryEngine}: proprietario di Crytek;
  \item \textbf{Snowdrop}: engine di Ubisoft;
  \item Rockstar usa un proprio engine proprietario;
  \item alcune aziende usano versioni \textbf{personalizzate} di engine
        esistenti: ad esempio Milestone, nei giochi di corse, usa un Unreal
        customizzato;
  \item CD Projekt Red ha un engine proprietario (con i problemi di
        documentazione già citati).
\end{itemize}
In sintesi: ogni azienda può avere un engine proprietario oppure una versione
modificata di un engine commerciale, ma si parte sempre da una base.

\subsection{Programmazione: cosa ci si aspetta da un designer}

Non è richiesto alcun bagaglio pregresso di programmazione: il corso parte da
zero. Gli studenti sono stati selezionati (il corso è a numero chiuso) per
competenze analitiche e critiche sopra la media. Anche nel mercato del lavoro,
per una posizione di game designer generalmente \textbf{non ci si aspetta
competenza di scripting} (C++, Java, C\#, ecc.), a meno di voler fare il
\emph{technical designer}, e anche lì varia da azienda ad azienda.

\subsection{Versione e preparazione}

Indicazioni pratiche emerse in aula:
\begin{itemize}
  \item scaricare Unreal Engine; la versione consigliata per il corso è
        stata indicata nella \textbf{5.6} (con conferma successiva). La
        stabilità dipende da cosa si deve fare, da quali tool e plugin si
        usano: i problemi nascono spesso dai plugin. Nel passaggio da Unreal 4
        a Unreal 5, ad esempio, il sistema \emph{Landscape} era inizialmente
        inutilizzabile, mentre in Unreal 4 funzionava molto bene;
  \item è consigliato installare tramite \textbf{Epic Games Launcher}, più
        comodo: permette anche di riscattare giochi gratuiti periodici e di
        accedere ai pacchetti di \textbf{Fab} (il marketplace di asset
        acquisito da Epic), che ogni mese offre alcuni asset gratuiti utili
        per esercitarsi nel level design (ad esempio per fare \emph{speed
        level design});
  \item molti pacchetti di UI o di sistemi di interazione sono utili per
        prototipare al volo, ma spesso inutilmente complicati;
  \item un plugin di \textbf{materiali di prototipazione a griglia} (adattivi,
        con metriche, allineabili alla griglia, con materiali già pronti e
        colori modificabili, dal costo di circa 5 euro) è indicato come
        ``salvavita'' per il level designer; se necessario, i pacchetti
        indispensabili vengono forniti;
  \item portare \textbf{sempre il portatile} alle lezioni di level design
        (da novembre, a partire dall'interfaccia dell'engine); alle lezioni
        di game design basta che lo porti almeno metà della classe, per i
        lavori a coppie. Serve comunque qualcosa per prendere appunti
        (``tavolette d'argilla e cannucce, se volete fare i babilonesi'',
        oppure carta e penna).
\end{itemize}
Chi ha un portatile poco potente dovrebbe comunque provare a installare
l'engine e aprire un template per vedere come gira: all'inizio non si faranno
cose pesanti. Nelle prime lezioni di level design è importante mettere subito
le mani sull'interfaccia: ascoltare dove si trova uno strumento (ad esempio lo
\emph{snapping}) senza provarlo porta a dimenticarlo.

% ---------------------------------------------------------------------
\section{Letture e risorse di riferimento}
\label{sec:risorse}

\subsection{I libri}

\begin{description}
  \item[\emph{The Art of Game Design: A Book of Lenses} (Jesse Schell)] il
        testo \textbf{davvero importante}. Alcune lezioni finali dell'anno si
        baseranno sulle sue terminologie. È utile perché apre la mente sulla
        possibilità di applicare metodologie diverse a giochi diversi, ma
        soprattutto per la \textbf{terminologia} comune. Esiste anche in
        audiolibro (circa 25 ore).
  \item[\emph{Challenges for Game Designers} (Brenda Romero e Ian Schreiber)]
        libro \emph{entry level} con moltissimi esercizi, molto caro al
        direttore del corso.
  \item[\emph{21st Century Game Design} (Chris Bateman)] utile come
        retrospettiva: è un libro visionario che ha previsto metodi e
        soprattutto \emph{archetipizzazioni} dei giocatori tuttora valide nel
        design e nel marketing. Servirà per una delle lezioni più importanti
        dell'anno. È utile leggerlo per vedere quanto aveva previsto, ma è
        utile ``fino a un certo punto''.
  \item[\emph{An Architectural Approach to Level Design} (Christopher
        Totten)] non strettamente necessario, ma di lettura semplice. Ha due
        anime: parte dall'architettura (babilonese, greca, ecc.) e la trasla
        su esempi di level design (come \emph{Doom} e \emph{Team Fortress}).
        Un tempo ne esisteva un riassunto su Gamasutra (oggi
        \emph{Game Developer}).
  \item[\emph{A Theory of Fun for Game Design} (Raph Koster)] utile fino a un
        certo punto: contiene diverse idee superate, che lo stesso Koster ha
        affrontato in vari talk reperibili nel GDC Vault. Verrà ripreso nelle
        prime lezioni, poi messo da parte.
\end{description}

\subsection{Canali e conferenze}

\begin{description}
  \item[GDC (Game Developers Conference)] il suo canale raccoglie talk e
        panel estremamente importanti: è ``uno strumento misterioso che ci
        servirà più tardi''. \textbf{Attenzione}: per molti talk mancano ancora
        le competenze per interpretarli, e alcuni potrebbero fuorviare o
        mettere in testa idee strane. La GDC si concentra spesso su
        produttività e mobile.
  \item[Il \emph{Dragon Speech} di Chris Crawford] il talk del fondatore della
        GDC: appartiene a un passato lontano, ma è molto motivante. Da vedere
        subito.
  \item[Game Maker's Toolkit] canale particolarmente utile per nozioni
        pratiche, con video brevi e alcune interviste.
  \item[Nordic Game] altra conferenza importante, con talk e panel di esperti
        del settore.
\end{description}

\subsection{Associazioni e premi: come leggerli}

Le associazioni di settore (ad esempio IIDEA in Italia, o la IGDA) non sono
il male, anzi: sono un canale importante per affacciarsi all'estero o per
accedere a fondi. Va però ricordato che \textbf{tutelano i propri associati}:
non esiste solo il loro palinsesto, e bisogna valutare quanto spazio occupano
nel panorama della visibilità.

Lo stesso vale per i \textbf{premi} (Game Awards, ``Best Mobile Game'',
``Games for Change'', ecc.): al designer \emph{non interessano}. Che un gioco
abbia vinto ``più di 90 premi'' non dice nulla di utile. Servono a dare
visibilità e lustro ai giochi ed è giusto che esistano, perché ampliano il
panorama e il numero di giocatori. Ma ciò che interessa al designer sono i
\textbf{dati}:
\begin{itemize}
  \item come sono state implementate certe cose;
  \item come è stato fatto il marketing e dove il gioco è stato presentato;
  \item non quante wishlist ha fatto, ma \emph{come} le ha fatte;
  \item in che periodo storico è uscito e quali erano i competitor in quel
        momento.
\end{itemize}

% ---------------------------------------------------------------------
\section{Che cosa fa un game designer}
\label{sec:cosa-fa}

Per arrivare a una prima definizione, la classe ha condotto un brainstorming:
ognuno doveva indicare una cosa che, secondo lui, fa il game designer (senza
ripetere quelle già dette). Ne è emersa la lista seguente, commentata punto
per punto.

\subsection{Progettare meccaniche e struttura del gioco}

Il ruolo più generico è quello del \textbf{core mechanic designer}, detto
anche \emph{gameplay designer} e, un tempo, \emph{black box designer}: si
occupa delle \textbf{meccaniche di gioco}. Si va dal piccolo al grande (in
funzione anche della dimensione dell'azienda):
\begin{itemize}
  \item nel \textbf{piccolo} ci sono le singole meccaniche;
  \item nel \textbf{grande} ci sono i \textbf{sistemi}. Nei giochi
        \emph{sistemici}, con molti sistemi interconnessi (ad esempio
        \emph{GTA}), esistono specializzazioni dedicate a bilanciare e
        supervisionare i sistemi: è il \emph{system design}.
\end{itemize}

\paragraph{Quando un gioco è ``sistemico''?}
Non dipende solo da come le meccaniche interagiscono con il giocatore (le
\emph{dinamiche}, secondo i framework che passano dalle meccaniche alle
dinamiche, che verranno studiati più avanti), ma anche da \emph{quante}
meccaniche ci sono e da quanto interagiscono tra loro. Non diremmo che
\emph{Super Mario} sia un gioco sistemico, eppure contiene dei sistemi: la
linea è sottile. Quando i sistemi aumentano molto, sono complessi e si
influenzano a vicenda, il gioco diventa sistemico. Esempio: in \emph{The
Legend of Zelda: Breath of the Wild} il vento modifica quanto il personaggio
sale o scende; il fuoco può incendiare l'erba; l'erba prende fuoco se colpita
da una freccia infuocata; e così via. Nel mercato moderno i giochi sistemici
sono spesso, ma non sempre, legati a mappe \emph{open world}.

\subsection{Le 3C: camera, character, control}

La base del gioco, da cui generalmente si inizia a lavorare (anche nei
prototipi), è la ``sacra trinità videoludica'' delle \textbf{3C}:
\begin{description}
  \item[Camera] che tipo di inquadratura ha il gioco: prima persona, terza
        persona, isometrica, ecc.
  \item[Character] quanto è grande il personaggio, che tipo di personaggio è
        e, di conseguenza, quali azioni può compiere.
  \item[Control] i controlli: che cosa il personaggio può fare, quali sono i
        \emph{verbi} di gioco.
\end{description}
Le 3C sono definite dal designer, ed esiste anche una specializzazione per
questo lavoro. Esistono varianti: quella del direttore del corso è
\textbf{4C + J}, cioè \emph{camera, character, control, conflict} (conflitto)
e \emph{juice} (il ``succo'', il feedback che rende piacevole l'interazione).
Verranno approfondite più avanti.

\subsection{Il personaggio}

Il designer definisce il personaggio dal punto di vista delle meccaniche
(dimensioni, azioni, capacità), all'interno delle 3C. L'aspetto estetico è
invece competenza di altri dipartimenti.

\subsection{Documentare le informazioni per gli altri dipartimenti}

Assolutamente vero: la documentazione è una delle prime attività del
designer. Più l'azienda è grande, più questo lavoro diventa prominente; più è
piccola, meno è rilevante. In \textbf{pre-produzione}, però, è sempre
fondamentale.

\subsection{Mediatore e leader ``al servizio''}

Il designer è un \textbf{mediatore} nei conflitti interni al team, una sorta
di collante che si relaziona con dipartimenti diversi (non sempre gli stessi,
a seconda della fase di produzione). A volte fa anche da leader, a seconda di
come sono gestite la pipeline e la gerarchia.

Questa leadership è difficile da definire perché \textbf{non è verticale}: è
una leadership \emph{al servizio} (\emph{servant leadership}). Il designer
``è asservito al programmatore'': il suo compito non è arrivare con un'idea e
scaricarla addosso agli altri, né portare la soluzione e lasciare il problema
a qualcun altro, ma \textbf{seguire il processo a basso livello}, sporcandosi
le mani.

\paragraph{I feature team.} Nei team medio-piccoli o medio-grandi possono
esistere sotto-team (\emph{feature team}) dedicati a un obiettivo specifico:
ad esempio un designer incaricato del character o della camera, affiancato da
uno o due programmatori e da un artista. In questi casi il designer
supervisiona il lavoro e gestisce la produzione del sotto-team.

\subsection{Gestire l'avanzamento del prodotto}

Esiste una specializzazione che si occupa solo di questo, il
\textbf{producer}, a un livello gerarchico completamente diverso (e che
potrebbe essere una strada percorribile al termine del percorso di studi).
All'interno di un feature team anche il designer può farlo, ma non è detto.
Può certamente \emph{dare la direzione}: nella sua documentazione può
indicare che per arrivare a uno stato giocabile del prototipo serve un certo
lavoro di arte e di programmazione, per lo stato successivo un altro, e così
via. In ambito produttivo esistono metodologie strutturate come
\textbf{Agile}, che verranno viste in seguito.

\subsection{Presentare le idee}

Il pitching (sezione~\ref{sec:gerarchia-corso}) è richiesto al designer sia
in pre-produzione sia in produzione, in misura che dipende dall'azienda. Ad
esempio Ubisoft lascia ai designer, e in realtà a qualsiasi membro, la
possibilità di \emph{pitchare} un'idea prima di una nuova produzione, perché
possa iniziare a prendere vita.

\subsection{Bozze e prototipi}

Le bozze vanno bene, ma devono essere \textbf{chiare}. C'è un'eterna lotta tra
documentazione e gioco giocabile: se si fa solo documentazione il gioco non
avanza; se si fa solo il gioco manca la documentazione e si perde la visione.
Sono due cose \textbf{complementari}.

Il modo per comunicare un'idea in forma grezza ma giocabile è il
\textbf{prototipo}. Un prototipo può essere qualsiasi cosa: di carta, in
Unreal con Blueprint ``brutti''. \emph{Non deve essere bello, deve essere
giocabile}.

\paragraph{Esempio: il prototipo di carta di Mario + Rabbids.} Durante lo
sviluppo di \emph{Mario + Rabbids}, i programmatori di Ubisoft erano impegnati
su altri progetti (come \emph{Ghost Recon}). I designer realizzarono allora
un \textbf{gioco da tavolo} che riproduceva le meccaniche del gioco, quelle
di uno strategico a turni sul modello di \emph{XCOM}. E funzionava.
L'obiettivo era \textbf{validare un'idea}, validare un gameplay.

Altri modi per rendere chiara un'idea sono i \emph{mock-up} nella
documentazione e le immagini modificate o i \emph{mash-up} di qualsiasi
natura. Lato documentazione l'importante è che sia chiaro; lato prototipo,
che sia giocabile.

\subsection{Fare ricerca per l'idea, e poi smontarla}

Il designer ha una componente creativa importante e fa ricerca per arrivare a
un'idea. Ma ciò che deve poi fare necessariamente è \textbf{smontarla}. Il
suo lavoro è scientifico, esattamente come il \textbf{metodo scientifico}, il
cui obiettivo è \emph{confutare un'ipotesi}: si arriva a un'ipotesi tramite
sperimentazione e poi si cerca $n$ volte di confutarla; se resiste, la si
assume come valida \emph{in un certo contesto} (le regole della fisica
classica, ad esempio, non valgono nella meccanica quantistica).

Allo stesso modo il designer arriva a un'idea (tramite ricerca o chiedendo
agli altri: le idee le hanno tutti, non sono un privilegio del designer),
cerca di smontarla e poi cerca di \textbf{metterla a terra}. Se l'idea non è
applicabile, non è implementabile, finché non si trova un modo per
realizzarla va scartata. Può essere l'idea più bella dell'universo, ma
\textbf{se non può essere messa a terra, non esiste}.

\paragraph{L'originalità.} Serve un'idea originale? Dipende dall'obiettivo e
dalla visione del progetto:
\begin{itemize}
  \item se l'obiettivo è vendere molto, un'idea originale non serve
        necessariamente;
  \item se si vuole creare un nuovo filone, serve un'idea, se non originale,
        almeno interessante: può essere un discorso di marketing o una
        meccanica innovativa, con il rischio di non vendere.
\end{itemize}
Il ruolo del \emph{director} e il concetto di \emph{visione} verranno
approfonditi in seguito.

\subsection{Analizzare i competitor}

Arrivati a un'idea (``vogliamo un certo tipo di combattimento''), si va a
vedere se qualcuno l'ha già fatta e \emph{come}. Le meccaniche non sono
privilegio di nessuno, salvo quando vengono brevettate, come il
\textbf{Nemesis System} di \emph{La Terra di Mezzo: L'Ombra di Mordor}.

\paragraph{Esempio: il combattimento da ridimensionare.} In un progetto
(presso Stormind Games) il team si trovò con un sistema di combattimento
troppo costoso: mancavano le risorse e gli animatori per realizzarlo. Poco
output lato gameplay significava non poter fare qualcosa di complesso e
profondo come in \emph{Elden Ring}; al contrario, il team aveva molti
artisti e poteva puntare su qualcosa di più artistico. Si iniziò così a
orbitare intorno al corpo a corpo di \emph{The Last of Us} (parte I e II).
Sul canale della GDC esistono diversi talk che raccontano cosa è stato fatto
e come, con successi, insuccessi e costi (ad esempio del \emph{motion
capture}): da quei talk si ricavarono moltissime informazioni utili.

I competitor possono essere:
\begin{itemize}
  \item \textbf{diretti}: un gioco molto simile al proprio;
  \item \textbf{di mercato}: un gioco che domina il mercato in cui ci si vuole
        affacciare;
  \item \textbf{indiretti}: un gioco con un sistema simile al proprio, ma
        tutto il resto diverso.
\end{itemize}
Regola d'oro: \textbf{là fuori, nascosto nell'ombra, c'è sicuramente un gioco
molto simile al vostro}. Va solo trovato. E se per qualche ragione non ha
venduto, bisogna capire perché.

\subsection{Lavorare in engine e fare tuning}

Il level designer lavora in engine, così come il technical designer; ma
qualsiasi game designer, a seconda di ciò che deve fare, può finire in
engine, soprattutto per il \textbf{bilanciamento} e il \textbf{tuning}.

Il \emph{tuning} è l'azione di bilanciare con precisione i dati. Esempio: in
\emph{Super Mario} si modifica quanto in alto salta il personaggio, con che
velocità, in quanto tempo raggiunge l'apice, e così via. Questo lavoro rende
il gioco:
\begin{itemize}
  \item ad alto livello, dotato della \textbf{progressione} giusta;
  \item a basso livello, \textbf{piacevole da giocare}.
\end{itemize}
Il \emph{fine tuning} è un lavoro di dettaglio, paziente e minuzioso.

\subsection{Bilanciare}

Il bilanciamento, emerso tardi nel brainstorming, è una delle attività
centrali del designer e ha anche una specializzazione dedicata
(l'\emph{economy designer}, sezione~\ref{sec:tipi-designer}).

\subsection{Analizzare il mercato}

In pre-produzione, di solito, sì: capire le tendenze, i generi che cambiano
nel tempo, la compatibilità del gioco con il momento storico, in vista ad
esempio di un possibile publisher. Va capito quanto, come e con quanto tempo
farlo, perché si è sempre limitati da molte cose.

\subsection{Sintesi}

\begin{quote}
Il designer parte dall'idea ma non si ferma lì: la confuta, lavora sui
dettagli, bilancia, fa \emph{fine tuning}; lavora con il team e fa da
\emph{manager} per diversi dipartimenti; e naviga costantemente tra i limiti.
\end{quote}

% ---------------------------------------------------------------------
\section{I limiti del design}
\label{sec:limiti}

I limiti ci sono sempre, e nel mondo del design sono spesso numerosi. I più
importanti entro cui navigare sono:
\begin{description}
  \item[Tempo] ad alto livello, l'idea deve essere realizzabile nel tempo di
        produzione disponibile; a basso livello, il proprio lavoro
        settimanale deve essere fattibile nel tempo che si ha. Non ci si può
        dividere tra mille feature team se realisticamente se ne può seguire
        uno solo.
  \item[Tecnologia] si può o non si può fare qualcosa? Ci sono i server per un
        multiplayer con un certo numero di giocatori? Ci si può permettere,
        ad esempio, una fauna che reagisce alla fisica come quella di
        \emph{GTA VI}, che richiede il meglio dei \emph{technical artist}?
  \item[Background e competenze] che cosa sa fare il team? Realisticamente,
        chi ha lavorato dieci anni su giochi horror viene assunto per un
        horror, perché ha le competenze per risolvere problemi già visti.
  \item[Persone e motivazione] il team ha la motivazione e le competenze per
        fare un certo tipo di gioco, e per lavorare insieme?
  \item[Denaro] tutto ruota intorno ai soldi. Che si lavori da soli o in mille,
        le persone vanno pagate e il gioco deve avere un ritorno economico.
        Bisogna decidere dove tracciare la linea: raggiungere il
        \textbf{break-even} (guadagnare esattamente quanto speso) o puntare a
        grandi guadagni. Con un \textbf{publisher}, inoltre, una parte
        consistente dei ricavi va a lui (ripartizioni dell'ordine di 80/20 o
        70/30, che dipendono dal contratto firmato).
\end{description}

\subsection{Non siete ``gli uomini delle idee''}

Per qualcuno può essere un problema scoprirlo, ma il designer non è
``l'uomo (o la donna) delle idee'': accanto alla parte creativa c'è una parte
tecnica, e si è fortemente limitati da molti fattori. Questi limiti, però,
sono un aiuto:
\begin{itemize}
  \item riducono il ventaglio di possibilità potenzialmente distruttive per il
        team (cose che non si riescono a fare, che non piacciono o per cui
        mancano le risorse);
  \item aiutano a mirare bene, trovando soluzioni creative in uno spazio più
        piccolo e più attuabile.
\end{itemize}

\begin{quote}
\centering
\textbf{Le idee sono nulla, l'implementazione è tutto.}
\end{quote}

% ---------------------------------------------------------------------
\section{Ruoli e specializzazioni}
\label{sec:tipi-designer}

\subsection{La scala gerarchica}

\begin{description}
  \item[Creative Director (CD)] mantiene la \textbf{visione} e dà l'idea, in
        modo che tutto segua una certa direzione. In teoria può diventarlo
        chiunque, ma con un'esperienza notevole, nell'ordine dei decenni: si
        assume che in quel tempo abbia visto come lavorano molti dipartimenti.
        Naturalmente un CD con un background artistico ha meno competenze di
        design, e viceversa.
  \item[Director di disciplina] uno per dipartimento: \emph{game director},
        \emph{level director}, \emph{art director}, ecc.
  \item[Lead] il capo del dipartimento, che gestisce i gruppi di persone. Chi
        lavora in un dipartimento risponde al proprio lead (ad esempio un game
        designer risponde al \emph{lead designer}, non al game director).
  \item[Senior / Principal] i profili più esperti (a seconda dell'azienda).
  \item[Junior] il profilo di inizio carriera.
  \item[Intern] lo stagista.
\end{description}
Lo schema si replica uno a uno per ogni dipartimento: arte, suono, animazione
(a volte arte e animazione insieme, o \emph{rigging} e animazione insieme),
programmazione, UX/UI e così via. A fare la differenza, a parità di ruolo, è
il \textbf{livello di seniority}.

\subsection{UX e UI}

UX e UI stanno a metà tra game design e arte: a seconda dell'azienda il ruolo
si chiama \emph{UX/UI Artist} o \emph{UX/UI Designer}.
\begin{description}
  \item[UX (User Experience)] come una persona interagisce con le cose. Esempio:
        su uno smartphone, un pulsante accende e spegne lo schermo; per
        digitare la password si scorre verso l'alto e compaiono i punti da
        toccare.
  \item[UI (User Interface)] come sono ``vestite'' queste cose: ad esempio
        l'aspetto della barra dei punti vita (HP), realizzata con asset 2D o
        3D.
\end{description}

\subsection{Le specializzazioni del game design}

Più l'azienda è grande, più la specializzazione è puntuale.
\begin{description}
  \item[Gameplay / Core mechanics designer] (sono la stessa cosa) si occupa
        delle meccaniche e delle 3C.
  \item[System designer] si occupa dei sistemi e delle loro interazioni.
  \item[Economy designer] si occupa del bilanciamento; lavora soprattutto con
        fogli di calcolo (Excel, Google Sheets).
  \item[Combat designer] si occupa del combattimento; negli annunci di lavoro
        indica quasi sempre un combattimento \emph{action}, con gameplay
        reattivo (non strategico), anche se non è una regola.
  \item[Multiplayer designer] si occupa delle interazioni multigiocatore.
  \item[Technical designer] fa da ponte tra game design e programmazione.
\end{description}

\subsection{Le specializzazioni del level design}

\begin{description}
  \item[Encounter designer] posiziona i nemici nella mappa in modo che gli
        scontri siano bilanciati e la progressione equilibrata come
        desiderato.
  \item[World / map designer] popola il mondo di gioco.
  \item[Quest designer] popola il mondo con le missioni e pensa alla
        progressione al loro interno, cioè a dove il giocatore deve andare (si
        vedano ad esempio gli annunci di lavoro di CD Projekt Red o di
        Larian).
  \item[Technical level designer] fa da ponte tra level design e
        programmazione.
\end{description}

\subsection{Discipline vicine: suono, UI, produzione}

Il monte ore del corso non permette di uscire come \emph{narrative designer}
(o \emph{sound designer}) completi. Ciò non toglie che si possano coltivare
queste discipline:
\begin{itemize}
  \item nei progetti di gruppo si può occuparsi anche del suono, della UI/UX,
        o provare ad applicare metodologie come Agile per fare
        \emph{producing} e gestire il team;
  \item il tempo dedicato a queste attività va però sottratto al game design,
        e dipende da quante persone ci sono nel team, da quanto lavoro c'è e
        da quanto è grande il progetto;
  \item avere sviluppato e rilasciato un progetto, anche piccolo (anche una
        \emph{game jam}), documentando i propri processi, è già un ottimo
        argomento da presentare in futuro;
  \item la via migliore resta lavorare nell'industria e, da lì, chiedere di
        imparare anche altre discipline.
\end{itemize}
Un gioco si può fare su qualsiasi cosa: sugli scarafaggi, sui piedini delle
sedie, sui pennarelli. Si parte avvantaggiati da ciò che è vicino alla propria
sensibilità: chi ama la musica può pensare a un gioco con meccaniche legate
alla musica, con un gameplay innovativo ma contenuto e realizzabile. Il design
è un cappello che comprende molte discipline.

% ---------------------------------------------------------------------
\section{Approfondimenti emersi in lezione}
\label{sec:approfondimenti}

\subsection{Roguelike e roguelite}

Un esempio di come il marketing conia generi: Supergiant Games, con
\emph{Hades}, rese celebre il termine \textbf{roguelite} al posto di
\emph{roguelike} (pur non essendo il primo roguelite in assoluto). La
differenza:
\begin{itemize}
  \item il \textbf{roguelike} si fonda su una serie di punti fissati in una
        sorta di \emph{manifesto} che definisce che cosa rende un gioco
        roguelike; uno su tutti: a ogni morte si ricomincia la \emph{run} da
        capo;
  \item il \textbf{roguelite} rompe alcuni di questi punti, in particolare
        introducendo una \textbf{progressione esterna alla run}: il giocatore
        può potenziarsi fuori dalla singola partita per vincere più
        facilmente in quelle successive. \emph{Hades} ne è il caso principe.
\end{itemize}
Analogamente, un mobile game sviluppato dal direttore del corso fu
promosso come ``il primo \emph{donation game} al mondo'': un gioco che
convertiva il tempo di gioco in donazioni reali a enti benefici per animali.
Anche qui un neologismo creato lato marketing per definire un genere.

\subsection{Applicazioni fuori dall'industria videoludica}

Le competenze sviluppate non servono solo ai videogiochi:
\begin{itemize}
  \item i \textbf{render architettonici} oggi possono essere realizzati con
        Unreal, e diversi professionisti formati nel game design lavorano in
        questo campo;
  \item la \textbf{gamification}, cioè l'applicazione di regole videoludiche
        a contesti non ludici (mostre interattive, progetti culturali,
        consulenze come quella per applicare questi concetti alle macchine da
        allenamento di un'azienda di fitness), è in Italia uno dei modi più
        redditizi di spendere queste competenze, spesso con finanziamenti
        pubblici.
\end{itemize}

% ---------------------------------------------------------------------
\section{Il primo task}
\label{sec:primo-task}

\begin{enumerate}
  \item Creare un indirizzo e-mail \textbf{professionale} (consigliato: un indirizzo Gmail nella forma
        \texttt{nome.cognome@gmail.com}; l'anno di nascita è accettabile).
        Va bene anche un indirizzo professionale già esistente usato per
        lavoro. Da evitare indirizzi poco seri (``non \texttt{blackfire67}'').
  \item Una persona della classe raccoglie tutti gli indirizzi in un unico
        documento e lo invia via e-mail al docente, che aggiungerà tutti al
        Drive condiviso del corso (dove verranno caricate lezioni, slide e
        materiali, per avere i documenti tutti nello stesso posto).
  \item Installare Discord e iscriversi al server del corso, dove passano le
        comunicazioni ufficiali. È consuetudine, e utile, creare anche un
        gruppo tra studenti.
\end{enumerate}
\textbf{Scadenza}: prima della lezione del venerdì.

% ---------------------------------------------------------------------
\section*{Punti chiave del capitolo}
\addcontentsline{toc}{section}{Punti chiave del capitolo}

\begin{itemize}
  \item Il corso è una \textbf{simulazione di lavoro}: ci si tratta da
        colleghi, con una gerarchia da studio di sviluppo.
  \item Si impara con l'\textbf{inquiry-based learning}: lezione, task,
        dubbi, domande, condivisione, progetti.
  \item Il fallimento affrontato con la mentalità giusta è la principale fonte
        di crescita; i feedback riguardano il lavoro, mai la persona.
  \item Il lavoro del designer è per metà creativo e per metà
        \textbf{ricerca e metodo scientifico}; il game design non è una
        scienza esatta.
  \item \textbf{Documentazione} e \textbf{pitching} sono il cuore del
        mestiere.
  \item Il designer progetta meccaniche e sistemi (partendo dalle
        \textbf{3C}), documenta, media, prototipa, confuta le idee, studia i
        competitor, bilancia e fa tuning.
  \item La leadership del designer è \textbf{al servizio} del team, non
        verticale.
  \item I limiti (tempo, tecnologia, competenze, persone, denaro) restringono
        e indirizzano la creatività.
  \item \textbf{Le idee sono nulla, l'implementazione è tutto.}
\end{itemize}

%    % ...
%
%  (Con hyperref aggiungere anche \renewcommand{\theHchapter}{z.\arabic{chapter}}
%  nella parte zero, per evitare ancore duplicate.)
%  Tutte le etichette di questo capitolo hanno il prefisso "z1:" per non
%  collidere con quelle dei capitoli del corso.
% =====================================================================

\chapter[La cultura del feedback]{Lavorare in squadra: la cultura del feedback}
\label{cap:z1-feedback}

Questo capitolo apre le \emph{lezioni zero}: incontri che precedono l'inizio
vero e proprio del corso di Game Design e che non riguardano una tecnica di
design in senso stretto, ma \textbf{la figura professionale} di chi lavora
nell'industria dei videogiochi. Prima di imparare a progettare un gioco
conviene infatti capire in che tipo di ambiente lo si progetterà: un ambiente
fatto di team, di gerarchie, di scadenze e, soprattutto, di persone che devono
riuscire a lavorare insieme.

La lezione è la prima giornata di un \textbf{workshop di due giorni} (circa
dieci ore in totale, pause pranzo comprese) tenuto da Stefano
Cozzi\footnote{Il cognome è riportato come compare nella trascrizione
automatica, che in più punti storpia i nomi propri; andrebbe verificato.
Nella stessa trascrizione lo studio per cui lavora appare come
``Stormine'', ``Stormline'' e ``Stormind'': si tratta di Stormind Games,
studio con sede a Catania, coerentemente con quanto detto in lezione.}, un
professionista con circa quindici anni di esperienza nel settore, tra
Blizzard Entertainment, Riot Games e, oggi, Stormind Games. Il workshop è
nato da un confronto con i coordinatori dell'accademia\footnote{Nella
trascrizione l'accademia è resa in modi diversi (``Hive'', ``a dive'',
``D.I.V.''); i coordinatori citati sono Enrico e Andrea, insieme al
direttore Diego Ricchiuti. Nel resto del capitolo si usa semplicemente
``l'accademia'', come negli altri capitoli.}: serviva un momento dedicato a
\emph{come ci si comporta in un ambiente professionale}, cosa che riguarda sia
il lavoro futuro sia i progetti di gruppo svolti durante il percorso di
studi, dove ci si deve già abituare a comportarsi ``come se si fosse in
un'azienda''.

La prima giornata ha un unico grande tema, sezionato in più parti:
\textbf{la cultura del feedback}. La seconda giornata (che sarà oggetto della
prossima lezione zero) sarà più pratica: \emph{time management},
\emph{brainstorming} e il ``guidare da dentro'' un team anche senza esserne
il leader (sezione~\ref{sec:z1-domani}).

Il percorso del capitolo segue quello della lezione:
\begin{enumerate}
  \item il formato del workshop e il percorso professionale del relatore
        (sezione~\ref{sec:z1-relatore});
  \item perché le aziende tengono tanto al lavoro di squadra e al
        \emph{fit} culturale (sezione~\ref{sec:z1-squadra});
  \item che cos'è il feedback e da che cosa dipende il modo di darlo
        (sezione~\ref{sec:z1-cose-feedback});
  \item le cinque fasi di vita di un team e il valore del
        \emph{post-mortem} (sezioni~\ref{sec:z1-fasi-team}
        e~\ref{sec:z1-qa-fasi});
  \item il confronto con i colleghi e alcune domande sul lavoro nei grandi
        studi (sezioni~\ref{sec:z1-rivalita} e~\ref{sec:z1-industria});
  \item il circolo virtuoso tra feedback, miglioramento e fiducia, e i suoi
        effetti sulla \emph{retention} e sulla cultura aziendale
        (sezioni~\ref{sec:z1-triangolo} e~\ref{sec:z1-cultura});
  \item il ``gioco del traditore'' e il principio del
        \textbf{default to trust} (sezioni~\ref{sec:z1-among-us}
        e~\ref{sec:z1-default-trust});
  \item le traiettorie di carriera e la valutazione delle performance
        (sezione~\ref{sec:z1-traiettorie});
  \item come si dà, come si riceve e quando si dà un feedback
        (sezioni~\ref{sec:z1-dare}, \ref{sec:z1-ricevere}
        e~\ref{sec:z1-quando});
  \item il feedback come regalo, le domande finali e le anticipazioni
        della seconda giornata (sezioni~\ref{sec:z1-regalo},
        \ref{sec:z1-qa-finale} e~\ref{sec:z1-domani}).
\end{enumerate}

Una premessa di metodo: il relatore ha alternato continuamente teoria,
domande alla classe e racconti della propria esperienza. I racconti non sono
digressioni: sono gli esempi concreti con cui i principi vengono resi
tangibili, e per questo sono riportati con cura.

% ---------------------------------------------------------------------
\section{Il workshop e il relatore}
\label{sec:z1-relatore}

\subsection{Formato e regole di partecipazione}

Il relatore ha cercato di ``riassumere quindici anni di esperienza in dieci
ore''. Ha materiale per circa il doppio delle ore disponibili: le cose
strettamente necessarie vengono trattate per prime, il resto è tenuto ``in
riserva'' come \emph{nice to have}. Le regole del workshop:
\begin{itemize}
  \item le domande si fanno \textbf{alla fine di ogni slide}, così da chiudere
        prima il discorso; non c'è un'unica sessione finale di domande e
        risposte (anche se ci sarà), perché il relatore preferisce fermarsi
        spesso e non fare ``uno spiegone di dieci ore'';
  \item le domande non devono per forza riguardare la slide appena vista: se
        una slide dà lo spunto per altro, se ne può parlare; al massimo sarà
        il relatore a rispondere in breve o a rimandare a una slide
        successiva;
  \item si può chiedere anche di esperienze passate, opinioni e sensazioni;
        se per NDA o per serietà professionale non può rispondere, lo dirà
        esplicitamente (``non lasciate che sia un \emph{blocker}'');
  \item su Zoom si usa l'alzata di mano o la chat, che il relatore tiene
        aperta sul secondo monitor;
  \item non ci sono pause ufficiali: la lezione è registrata e ognuno è
        libero di assentarsi quando serve, cercando però di non disturbare il
        flusso della lezione;
  \item le slide verranno fornite integralmente tramite la segreteria
        dell'accademia (Vanessa), insieme alle registrazioni: le slide senza
        il parlato servono poco, ma sono un buon promemoria.
\end{itemize}

\subsection{Il percorso professionale}

Il relatore si presenta (scherzosamente) come la ``guida escursionistica''
che accompagna la classe nella ``foresta incantata'' della cultura del
feedback, perché per molti è un territorio sconosciuto. Il suo percorso:
\begin{itemize}
  \item in Italia ha lavorato per aziende che fornivano servizi a terzi, tra
        cui il lavoro su \emph{Warhammer Online}, e per un canale televisivo
        sul digitale terrestre dedicato ai videogiochi\footnote{Nella
        trascrizione: ``Digital Game Network Channel''. Il nome esatto del
        canale andrebbe verificato.};
  \item a \textbf{Blizzard Entertainment} (circa sei anni, a Parigi) ha fatto
        il \emph{community manager} per l'Italia (\emph{StarCraft~II},
        \emph{World of Warcraft} da \emph{Mists of Pandaria} in poi,
        \emph{Hearthstone}, \emph{Diablo~III}, le espansioni di
        \emph{StarCraft}, il lancio di Battle.net), per poi diventare il
        manager dei community manager europei;
  \item a \textbf{Riot Games} (circa cinque anni, a Dublino) è stato
        brevemente community manager italiano e poi \emph{publishing
        manager} di un team di circa trenta persone per l'Europa. Il
        \emph{publishing} non è solo community: è l'unione di
        \textbf{community, marketing e PR};
  \item a \textbf{PG Esports} si è occupato di eventi dal vivo e online
        (ad esempio la Milano Games Week e tornei di \emph{Valorant} e
        \emph{League of Legends}), restando in contatto con Riot e Blizzard;
  \item da circa un anno è tornato nello sviluppo, a \textbf{Stormind
        Games}, dove oggi si occupa di risorse umane e organizzazione del
        lavoro; lo studio verrà presentato più ampiamente al \emph{Career
        Day} dell'accademia insieme alla sua HR manager;
  \item fa inoltre consulenze aziendali anche al di fuori dei videogiochi e
        insegna in master universitari (IULM di Milano).
\end{itemize}

Il suo ruolo ricorrente, sia a Blizzard sia a Riot, è stato quello di
\emph{Mr.~Wolf} (``per chi ha visto \emph{Pulp Fiction}''): quello che
risolve i problemi. Veniva inserito nei team più in difficoltà, ad esempio
dal punto di vista delle performance, per sistemarli e passare al successivo.
Era presente in alcuni dei momenti più critici, come il lancio di
\emph{Diablo~III} o il lancio di espansioni di \emph{World of
Warcraft}\footnote{La trascrizione riporta ``il lancio con i Gary's'' e
``World of Trenor'': si tratta verosimilmente di \emph{Mists of Pandaria} e
di \emph{Warlords of Draenor}.}, a fare \emph{damage control}.

La scoperta fondamentale di questi anni è che i problemi da risolvere
\textbf{quasi mai riguardavano il prodotto} (non era certo lui a sistemare i
server di \emph{Diablo~III}), ma i team che c'erano dietro: persone che
litigavano, non collaboravano, non riuscivano a stare nella stessa stanza
senza urlarsi contro. E nel 90\% dei casi tutto questo si riduce a
\textbf{problemi di comunicazione}: è raro che perfino le incompatibilità
caratteriali non siano risolvibili parlando.

\subsection{Struttura non vuol dire burocrazia}

A Stormind il relatore sta accompagnando un team passato da 30 a 95 persone
in meno di un anno, attraverso tutti i cambiamenti di cui parla il workshop.
Darsi una struttura viene spesso percepito male, come ``più processi, più
procedure'', ma bisogna distinguere:
\begin{itemize}
  \item un processo come ``timbrare il cartellino'' può essere davvero
        invasivo e noioso;
  \item definire \emph{come funziona un documento di game design, chi lo
        approva, chi lo visiona e chi no} è invece una struttura necessaria.
\end{itemize}
Uno studio di 30 persone nella stessa sede può farne a meno: quando c'è un
dubbio ci si alza e si va alla scrivania del collega. Uno studio di 95
persone sparse per l'Italia (e in parte per l'Europa) non può più farlo:
servono una struttura diversa, un modo di lavorare diverso e una
pianificazione del futuro diversa.

% ---------------------------------------------------------------------
\section{Perché il lavoro è sempre lavoro di squadra}
\label{sec:z1-squadra}

\subsection{Non si va in ufficio per fare amicizia}

Il primo messaggio è netto: \textbf{non si va in ufficio per fare amici} o
per trovare compagni di vita. Si va in ufficio perché c'è un lavoro da fare, e
lo scopo è riuscire a farlo \emph{anche con chi sta antipatico}. Con le
persone simpatiche è facile: diventeranno amici, e non sarà mai un problema.
È con chi sta antipatico che le cose si complicano, ed è lì che emergono i
\textbf{veri professionisti}: quelli che riescono a lavorare con chiunque
senza che l'antipatia pesi sul modo in cui lavorano e collaborano, mettendo il
prodotto e il team davanti al proprio ego e alle proprie preferenze personali.

La frase che fa da filo conduttore a tutto il workshop (e che il relatore
promette di riproporre ``fino alla nausea'') è:

\begin{quote}
  \emph{La capacità di un gruppo di persone di fare cose straordinarie
  dipende dal modo in cui queste persone operano insieme.}
\end{quote}

Le aziende raramente pensano che l'innesto di una singola persona possa
cambiare il destino di un progetto. Questo è tanto più vero quanto più grande
è il team, ma vale anche per i piccoli studi, anche per un gruppo di cinque
persone. Come negli sport di squadra: l'individuo conta fino a un certo punto,
e anche se si ha Dončić in squadra, se si gioca uno contro cinque si perde
comunque\footnote{Il riferimento è alle finali NBA in corso nei giorni del
workshop tra Dallas Mavericks e Boston Celtics (giugno 2024), il che
permette di datare approssimativamente la lezione.}.

\subsection{Lo \emph{useless genius} e il costo delle persone tossiche}

Molti dipartimenti HR usano la definizione di \textbf{useless
genius}\footnote{Il concetto è più noto in letteratura manageriale come
\emph{brilliant jerk}; in lezione è usata l'espressione \emph{useless
genius}.}: si può essere bravissimi, a volte i migliori nel proprio lavoro,
ma se non ci si sa comportare, relazionare e operare in un gruppo si è
fondamentalmente inutili ai fini del gruppo. Nel migliore dei casi inutili,
nel peggiore \emph{dannosi}.

In passato le aziende tendevano a ``chiudere un occhio'': la persona è
fondamentale, la teniamo. Oggi la domanda è diventata: \emph{quanto costa
all'azienda, oltre al suo stipendio, tenere una persona che crea un ambiente
tossico e peggiora chi le sta intorno?} Ci sono costi sommersi che le aziende
cominciano a vedere, e la tendenza è liberarsi di queste persone il prima
possibile. Il ragionamento è semplice:
\begin{itemize}
  \item le \textbf{competenze si possono trasmettere} o sostituire: si trova
        qualcuno che fa lo stesso lavoro magari nel doppio del tempo, oppure
        più persone che se lo dividono;
  \item il \textbf{carattere delle persone} invece cambia molto raramente.
\end{itemize}

\subsection{Il fit culturale}

Quando qualcuno \textbf{non è un \emph{fit} culturale} per un'azienda è
meglio separarsi presto. Ma non essere adatti a \emph{quell'}azienda non
significa non essere adatti a lavorare: è perfettamente legittimo che una
persona inadatta a un'azienda sia adattissima a un'altra.

Il relatore l'ha vissuto in prima persona. Dopo sei anni a Blizzard viene
contattato da Riot e affronta un processo di selezione lunghissimo (``quasi
undici colloqui''), al termine del quale viene scartato. Scrive allora una
mail che gli ha ``cambiato la traiettoria di carriera'': chiede
semplicemente perché non è stato ritenuto adatto. La risposta: in nessun
colloquio si era esposto abbastanza da far capire \emph{perché} volesse
lasciare Blizzard, e quindi pensavano che in realtà stesse bene dov'era. Il
motivo era che, per educazione professionale, non gli piace parlare male del
proprio posto di lavoro; portata all'eccesso, però, questa riservatezza era
diventata un ``è tutto bellissimo, ma vorrei provare Riot'', senza un perché.
Lo invitano comunque a un incontro dal vivo (``dal vivo è diverso''), e alla
fine viene assunto.

Una volta arrivato, i sei mesi di prova (\emph{probation}) sono serviti tutti
per smettere di ragionare e lavorare ``come a Blizzard''. La differenza
culturale era enorme, tanto che in Riot esisteva un gruppo di supporto
interno tra ex dipendenti Blizzard: un collega gliel'ha descritto come
``uscire da una relazione tossica e doversi riabituare a una relazione
normale''\footnote{Il relatore precisa che parla della Blizzard di allora,
che non esiste più: oggi è parte di Microsoft e nessuno dei suoi ex colleghi
vi lavora ancora, anche a causa dei licenziamenti di massa.}.
\begin{description}
  \item[Blizzard] era \emph{developer first}: il focus era chiaramente sugli
        sviluppatori, prima di tutti gli altri. Dopo sei anni c'erano ancora
        meeting a cui non poteva partecipare perché non abbastanza
        \emph{senior}, informazioni che non poteva sapere, piani che non
        poteva vedere.
  \item[Riot] era \emph{player first}: prima i giocatori, e i dipendenti come
        mezzo per trasferire il feedback dei giocatori dentro il gioco. Il
        secondo giorno lo portano in una stanza e gli mostrano il piano dei
        cinque anni successivi. Lui, sorpreso, pensa a un errore (``sono qui
        da oggi, non vorrei mettervi nei guai''); la risposta è: ``sei un
        dipendente Riot o no? Allora non ci sono livelli di accesso: devi
        sapere queste cose per fare meglio il tuo lavoro''.
\end{description}

Da questa esperienza ha imparato molto sulla \textbf{trasparenza}: condividere
le informazioni aiuta il lavoro. Non sempre, però: a volte non si condivide
qualcosa non per segretezza, ma perché non è ancora pronta per essere
presentata, esattamente come non si mostra un gioco quando è ancora solo
teoria (``te lo faccio vedere quando ho qualcosa di tangibile'').

La lezione per gli studenti è doppia: le aziende valuteranno il \emph{fit}
culturale sempre di più, ma capire se si è un \emph{fit} richiede anche
\textbf{lavoro da parte del candidato}. Bisogna avere la presenza di spirito
di dire: ``conosco i valori di questa azienda, non mando nemmeno il CV,
perché so che durerei un mese e mezzo''.

% ---------------------------------------------------------------------
\section{Che cos'è il feedback}
\label{sec:z1-cose-feedback}

\subsection{Le prime risposte della classe}

La prima domanda alla classe è: \emph{quando sentite la parola
``feedback'', a cosa pensate?} Le risposte raccolte, a voce e in chat:
\begin{itemize}
  \item un \textbf{riscontro} su ciò che si fa e su come lo si fa, sul lavoro
        o sul modo in cui ci si interfaccia con il gruppo; può essere
        lavorativo oppure comunicativo e comportamentale. Chi ha risposto lo
        vede come qualcosa che, se l'azienda non dà, va attivamente
        richiesto, perché è importante sapere come sta andando per
        migliorarsi;
  \item né positivo né negativo, ma \textbf{costruttivo}: in entrambi i casi
        serve a migliorarsi;
  \item uno \emph{scambio di opinioni}, un \emph{riscontro produttivo};
  \item un feedback ``di reazione'', non verbale: si fa una cosa e si
        osserva come cambia l'ambiente o il comportamento delle persone
        intorno (anche perché a voce la gente spesso si tiene le cose
        negative);
  \item per un designer, anche il \textbf{linguaggio del corpo}: osservando
        un giocatore giocare si può dedurre l'emozione che prova (frustrazione,
        incomprensione).
\end{itemize}

Sulla forma (scritto o orale?) le opinioni divergono: c'è chi lo preferisce
\emph{mille volte} scritto; chi pensa che dipenda dalle dimensioni
dell'azienda (orale nelle piccole, scritto in quelle con team di centinaia di
persone); chi pensa che dipenda dall'importanza del feedback (scritto se
tecnico o pieno di dati, da tenere sotto gli occhi; orale se è un consiglio
veloce).

Il relatore usa due esempi per mostrare che \emph{dipende}:
\begin{itemize}
  \item un feedback sulla programmazione ha senso darlo scritto, magari
        commentando direttamente il codice, più che ``andando a prendere un
        caffè'';
  \item ma se il feedback è ``ieri, durante il meeting, ho percepito male il
        fatto che mi interrompessi di continuo'', ce lo si aspetta
        \textbf{orale}.
\end{itemize}

\subsection{Da che cosa dipende il modo di dare feedback}

Lo \emph{spoiler} della giornata è che \textbf{non esiste una risposta univoca}
su come dare feedback. Il workshop non insegna ``il'' modo, ma fornisce
\emph{tutte le opzioni}, perché il modo giusto cambia in base a molti fattori:
\begin{itemize}
  \item a chi lo si dà: un collega, un superiore, il proprio capo, un proprio
        \emph{report} (un sottoposto);
  \item se è personale o professionale;
  \item se è la prima volta o la terza volta che si dà lo stesso feedback;
  \item se riguarda un progetto importante o secondario.
\end{itemize}

Manca però l'elemento fondamentale, che la classe individua subito:
\textbf{l'altro}. A parità di tutti gli altri fattori, il \emph{carattere
della persona che riceve il feedback} dovrebbe far cambiare il modo in cui lo
si dà. Per questo, prima della lezione, agli studenti era stato chiesto di
fare un test di personalità: non un test di psicologia avanzata, ma uno
strumento per rendersi conto di quanta varietà ci sia anche solo in una
classe di 53 persone (il che non vuol dire trovare 52 modi diversi di dare
feedback, ma lavorarci un po' sì).

Il relatore porta se stesso come esempio. Preferisce:
\begin{itemize}
  \item i feedback \textbf{positivi scritti}, perché è a disagio a ricevere
        complimenti in pubblico;
  \item i feedback sul proprio lavoro \textbf{diretti}: ``vai subito al
        punto e non girarmi intorno''.
\end{itemize}
È molto orgoglioso del suo lavoro e, quando qualcuno lo critica, la sua
primissima reazione istintiva è ``come ti permetti''. Negli anni ha imparato
a farla durare un secondo, ma chi gli dà feedback di persona dice di
vederla ancora negli occhi. Per mitigarla, dice esplicitamente a chi gli dà
feedback: \emph{``ignora il modo in cui reagisco nei primi cinque secondi,
vai avanti, poi mi passa''}. È considerato da tutti molto aperto al feedback
e non ha problemi a dire ``hai ragione, facciamo diversamente'': ma ci sono
reazioni istintive e caratteriali che è difficile controllare al cento per
cento.

\subsection{Con chi si lavora meglio}

Guardando le persone con cui si lavora abitualmente (anche nei progetti
dell'accademia), si noterà che si lavora meglio con due profili:
\begin{enumerate}
  \item chi è \textbf{esattamente come noi}, ma qui è un 50 e 50: alcuni
        caratteri vanno d'accordo con i propri simili, altri non vogliono
        nessuno con lo stesso carattere nel gruppo (e si potrebbe rientrare
        in questa categoria);
  \item chi è \textbf{affine}: ha qualcosa in comune con noi, ma non tutto,
        e proprio per questo può dare punti di vista differenti.
\end{enumerate}
Il lavoro della giornata è capire come valorizzare il secondo caso: ottenere
feedback da persone con caratteri, visioni e background diversi, per
completare una visione del proprio lavoro che è necessariamente limitata. È
come esplorare una mappa con la nebbia di guerra: un'altra persona, al di là
del muro, vede cose che noi non vediamo.

\subsection{Il feedback come somma di piccoli miglioramenti}

Per il relatore il feedback è questo: \textbf{applicare tanti piccoli
miglioramenti periodici}, che si sommano nel tempo. Migliorare dello 0,1\%
ogni giorno porta a fine anno a un miglioramento enorme, e questo vale per i
prodotti, per i giochi, per le persone e per i team.

\paragraph{L'esempio del ciclismo britannico.} La squadra di ciclismo
britannica, fino a una quindicina d'anni fa, non aveva vinto quasi nulla.
Cambiato il team manager, tutti si aspettavano grandi cambiamenti sportivi
(convocazioni, strategie di gara). Arrivarono invece \emph{minuscoli
cambiamenti non sportivi}: un altro chef in trasferta, materassi migliori nei
camper per riposare meglio, un sellino più comodo a fine tappa. Cambiamenti
marginali che da soli non avrebbero cambiato niente; eppure, dopo due o tre
anni di questo lavoro quotidiano, la squadra ha cominciato a vincere a
ripetizione, Olimpiadi comprese, cambiando il modo in cui il ciclismo è
visto in Gran Bretagna\footnote{È la nota strategia delle \emph{marginal
gains} associata a Dave Brailsford (British Cycling e Team Sky); il nome non
è citato in lezione.}. Certo, si è investito anche sulle persone giuste, ma
tutto è partito dai piccoli miglioramenti periodici.

Per gli studenti vale lo stesso: chiarire un conflitto con una persona con
cui si ha ancora poco a che fare, o una situazione che ha lasciato
``l'amaro in bocca'', sembra inutile oggi, ma darà i suoi frutti tra tre, sei,
nove mesi. Sono \emph{investimenti sul futuro}.

\paragraph{Peggiorare ogni giorno.} Esiste anche il percorso inverso:
lasciar correre, ``buttare sotto il tappeto'' le cose perché ``tanto è
inutile parlarne''. Queste cose tornano sempre fuori dopo otto mesi, come il
\emph{libro dei rancori} dei nani di Warhammer: ``ti ricordi quando, un anno
e mezzo fa, all'inizio del progetto, mi hai detto quella cosa? Mi sono
offeso''. A quel punto è quasi impossibile risolvere, perché nessuno ricorda
più con precisione cosa sia successo (``ma quale cena?'', ``forse era in
pizzeria'', ``no, eravamo da un'altra parte''). \textbf{Tenersi le cose nel
cassetto per tirarle fuori quando servono è sempre sconsigliato}: non servono
ed è impossibile usarle in modo costruttivo.

% ---------------------------------------------------------------------
\section{Le cinque fasi di vita di un team}
\label{sec:z1-fasi-team}

\subsection{Perché conoscerle}

Quasi tutti gli studenti lavorano già in gruppo (al primo anno su un gioco da
tavolo, negli anni successivi su un videogioco) e si saranno accorti che un
team attraversa delle fasi: il modo di lavorare insieme cambia nel tempo.
Spesso ci si fa prendere di sorpresa: ``eravamo così affiatati, cosa è
successo?''.

Un esempio da una consulenza fuori dai videogiochi: un'azienda farmaceutica
gli affida un team commerciale (rappresentanti che presentano medicinali a
ospedali, medici e privati). Il problema: ``prima andava tutto bene; ora che
siamo in diciotto dovremmo lavorare meno e meglio, e invece va peggio''. Va
peggio perché le persone nuove sono appena arrivate, mentre il management si
aspettava già di vedere i risultati dell'investimento.

Fino a pochi anni fa le fasi di un team venivano disegnate come una
\textbf{linea dritta verso l'alto}: le fasi erano le stesse, ma viste come un
miglioramento costante. Non si vedevano le parti \emph{fisiologiche} del
lavoro in team, cioè anche i peggioramenti. Accettarle è il primo passo per
accettare i fallimenti: nessun percorso personale o di team è lineare.

Il relatore ammette che il suo modo di lavorare è nato per il 70\% dal
proprio impegno a imparare, fare bene e rispettare le persone, e per il 30\%
dall'aver lavorato con ``incapaci colossali'' che gli hanno insegnato come
\emph{non} fare le cose. A Blizzard un manager gli fece capire che, pur non
sapendo ancora che tipo di \emph{people manager} sarebbe diventato, sapeva
benissimo quale tipo di manager \emph{non} voleva essere.

\subsection{Forming, storming, norming, performing, adjourning}

Le fasi sono cinque, un po' come le fasi del lutto. Molte aziende ne
considerano solo quattro (``e sono quelle che solitamente chiudono''); le
aziende al passo con i tempi ne considerano cinque\footnote{Il modello è
quello di Bruce Tuckman (1965), a cui la quinta fase, \emph{adjourning}, fu
aggiunta nel 1977. Il nome di Tuckman non è citato in lezione; la
trascrizione riporta la quinta fase come ``journing''.}.

\begin{description}
  \item[Forming] Si viene messi insieme e si comincia a lavorare. Non si sa
        nulla degli altri; spesso non si conoscono nemmeno i nomi, ed è
        normale.
  \item[Storming] Si cominciano a ``dare le testate come le capre in
        montagna per determinare il territorio'': questa cosa la faccio io,
        questa la vorrei fare così, chi decide? Si stabilisce chi fa cosa e
        come si divide il lavoro. Nel grafico è una \textbf{fase di
        peggioramento}: che il team litighi o discuta civilmente non cambia
        nulla, il team deve investire parte di tempo ed energie per risolvere
        queste questioni. È normale e \textbf{non è evitabile}.
  \item[Norming] Si comincia a capire come lavorano gli altri, cosa sanno
        fare e come far funzionare il gruppo; iniziano i miglioramenti. È una
        forma di \emph{accettazione} degli altri, che deriva dall'aver capito
        cosa devo fare io e cosa devono fare loro.
  \item[Performing] C'è uno sforzo \textbf{proattivo}: ``questa cosa
        toccherebbe a me, ma so che Alessio la fa molto meglio, quindi gli
        chiedo aiuto''. Si sfruttano i punti di forza degli altri, perché li
        si conosce e perché si sa che passare il testimone non avrà effetti
        negativi (il capo non dirà ``perché non l'hai fatta tu?'').
  \item[Adjourning] Il progetto finisce, si chiude o viene riorganizzato, il
        team viene sciolto. Quello che andrebbe fatto è analizzare
        \textbf{cosa ha funzionato, cosa no e cosa migliorare} per il
        progetto successivo.
\end{description}

\subsection{La fase che tutti saltano}

``Spoiler incredibile'': l'\emph{adjourning} è la fase che molte aziende
saltano, \textbf{perché c'è sempre il progetto dopo}.

A PG Esports, per la Milano Games Week (come per qualsiasi evento dal vivo a
cadenza annuale: Lucca Comics, Etna Comics, Comicon\dots), finita l'edizione
di un anno si cominciava dal giorno dopo a lavorare a quella successiva,
perché un anno sembra tanto ma non basta mai. Nei videogiochi è lo stesso:
finito un gioco, a parte le patch di mantenimento e del \emph{day one}, si
passa subito al successivo.

Il cambiamento introdotto dal relatore (in parte già praticato, ma da lui
strutturato) è stato: finita la Games Week, \textbf{tutti in ferie per una
settimana}; poi ci si ritrova per definire cosa ha funzionato e cosa no, con
un \textbf{post-mortem} (o \emph{debrief}); solo dopo si inizia a lavorare
all'edizione successiva. Senza questa fase manca la parte di
\emph{learning}, cioè di ciò che si è imparato dal progetto.

Il consiglio vale per gli studenti: quando si finisce il gioco da tavolo, e
poi il gioco digitale, bisogna trovare il tempo per il post-mortem,
altrimenti si butta via \textbf{il 90\% dell'utilità del progetto}. Una
struttura semplice per farlo è descritta nella sezione~\ref{sec:z1-qa-finale}.

% ---------------------------------------------------------------------
\section{Domande sulle fasi del team}
\label{sec:z1-qa-fasi}

\subsection{Lo storming che ritorna}

\paragraph{Capita che lo storming si ritagli spazio anche nelle altre fasi?}
Lo \emph{storming} è una fase precisa, che però molti manager, coordinatori e
team cercano di \textbf{saltare}: ``se io sono bravo come manager e loro sono
bravi come team, passeremo da forming a norming senza problemi''. È un errore,
perché senza una fase in cui si discute davvero non si sviluppa una
\emph{skill} fondamentale: \textbf{saper discutere}. Tutti conoscono persone
non abituate a farlo, per cui ogni minimo disaccordo diventa un affronto
personale (``avevamo deciso la pizza, perché andiamo al sushi? Perché mi fate
questo affronto?''): sono persone che non sono abituate ad avere conflitti
costruttivi.

Se si pensa di essere in \emph{norming} e continuano a emergere conflitti, il
punto è che \textbf{non si era davvero in norming}: lo storming non ``si
trascina'' nelle fasi successive, semplicemente non se n'è mai usciti.
Allo stesso modo, a chi chiede se ogni mini-fase di storming durante la vita
del gruppo sia frutto di uno storming fatto male o solo di carattere, la
risposta è già contenuta qui.

Un partecipante osserva che spesso, ancor prima di finire un progetto, si è
già messi su quello successivo: è il modo migliore per essere in costante
\emph{damage control}. Era l'esperienza del relatore a Blizzard: veniva
spostato fisicamente da un team all'altro, e finito il lancio di
\emph{Diablo~III} era già su \emph{World of Warcraft}, senza vivere il
``dopo''. Si è dovuto ritagliare da solo lo spazio per un debrief personale,
ed è lì che ha deciso che tipo di manager non voleva essere. Il consiglio:
\textbf{se non lo fa il team, fatelo voi}.

\subsection{Dove si passa la maggior parte del tempo}

\emph{In quale fase si passa più tempo nella propria carriera?} Gli studenti
rispondono \emph{storming} e \emph{norming}. Il relatore, in quindici anni,
calcola di aver passato in \emph{performing} circa \textbf{dieci mesi}. Non è
pessimismo: arrivare a performing deve essere sempre l'obiettivo, sia da
manager sia da membro del team; ma è difficile arrivarci, ed è normale che non
sempre ci si arrivi.

Lo storming può durare a lungo, ma con team molto esperti dura pochissimo,
perché c'è chi lo sa gestire. In realtà \textbf{la maggior parte del tempo si
passa in forming}, perché ogni equilibrio raggiunto vale per \emph{quelle}
persone, \emph{quel} progetto, \emph{quelle} tempistiche. Basta che cambi un
qualsiasi fattore (una deadline spostata, una persona nuova, qualcuno che si
dimette, un progetto cancellato) e si ricomincia, perché i bilanciamenti
precedenti non valgono più.

Esempio: un team di cinque persone affiatatissimo, in pieno performing. La
dirigenza, entusiasta, ``raddoppia le risorse'' e lo porta a dieci. O si
lascia tutto com'è (e i cinque nuovi non fanno niente), o bisogna ridiscutere
chi fa cosa e come i nuovi possono aiutare: che piaccia o no, si ricomincia
da forming.

Quindi la bravura di un manager, ma anche di un membro del team,
\textbf{non sta nel passare tanto tempo in performing, ma nell'accelerare il
passaggio da una fase all'altra}, e nel fare bene la fase di
\emph{adjourning}.

\subsection{Analizzare anche i successi}

C'è la tendenza a fare post-mortem solo dei progetti andati male: ``questa
cosa non ha funzionato, vediamo perché''. Va benissimo. Ma quando un progetto
va bene si dà per scontato che sia andato bene ``perché era buono''. Bisogna
invece chiedersi \emph{perché} è andato bene: era il prodotto? La
tempistica? Il modo in cui abbiamo fatto una cosa piuttosto che un'altra?
\textbf{Se non si capisce davvero il motivo del proprio successo, si
tenderanno a replicare le cose sbagliate.}

\paragraph{Blockbuster.} Quando arrivò la notizia che Netflix avrebbe
lanciato un servizio di streaming, Blockbuster (la catena di noleggio di
videocassette e DVD) aveva tutto il tempo per preparare un'offerta
concorrente. In una riunione del management raccontata in un libro (di cui il
relatore non ricorda il titolo), la direzione sembrava orientata a investire
nello streaming, finché un dirigente fece notare che il 12\% del fatturato
derivava dalle \emph{multe per riconsegna tardiva}, impossibili con lo
streaming, mentre secondo le stime lo streaming avrebbe prodotto il 10\% del
fatturato. ``Non possiamo rinunciare a questo 2\%'': niente streaming. Oggi
non c'è più un Blockbuster ``manco a cercarlo col lanternino''. Non avevano
capito qual era il vero motivo del loro successo, cioè il loro modello di
business\footnote{Le percentuali sono quelle riferite in lezione e non sono
state verificate.}.

\paragraph{The Founder.} Nel film sulla nascita di McDonald's c'è una scena in
cui il futuro socio del protagonista gli dice: ``tu pensi di essere nella
\emph{food industry}, e invece sei nella \emph{real estate industry}''. Non
dovrebbe investire sugli hamburger, ma sui terreni su cui costruisce i
ristoranti. Anche per un team vale lo stesso: bisogna capire che cosa
l'ha reso davvero di successo, oppure no.

\subsection{L'autoanalisi personale}

\paragraph{E quando si cambia azienda?} Uno studente chiede se il debrief sia
possibile anche quando si lascia un'azienda. Non si può tornare in ufficio a
discutere della propria esperienza, ma lo si può fare \textbf{per se stessi}.
Il relatore l'ha fatto, anche se all'epoca in modo poco consapevole, ed è
così che è passato dal \emph{community management} al \emph{people
management}. Analizzando cosa aveva fatto bene a Blizzard si era reso conto
che:
\begin{itemize}
  \item aveva costruito un ottimo rapporto di fiducia con i giocatori (con
        alcuni proprietari di fansite parla ancora oggi);
  \item pur senza la seniority, era diventato un punto di riferimento per il
        team, visto come imparziale ed equilibrato anche su temi che non lo
        riguardavano direttamente;
  \item sapeva \textbf{risolvere problemi}: durante il lancio di
        \emph{Diablo~III} era il più tranquillo di tutti, non perché
        sminuisse il problema, ma perché sapeva che si sarebbe risolto con i
        suoi tempi; il lavoro non era spegnere l'incendio ma ``far sfollare
        la gente'' e aspettare che tutto tornasse in ordine.
\end{itemize}
E aveva capito cosa sapeva fare meno bene, ad esempio \textbf{gestire il
proprio manager}. Avere un manager non significa che il manager gestisca te:
il manager gestisce il tuo lavoro, ma tu devi saper gestire i tuoi
\emph{stakeholder}. Spesso lo stakeholder è il proprio manager; da manager,
gli stakeholder saranno CEO, CFO e così via. Gestirli significa conoscerli,
sapere come portare loro le informazioni, come far presente che qualcosa non
va. Fa parte del lavoro, anche se a lui non piace, e ci ha lavorato.

Nel mese e mezzo di pausa prima di entrare in Riot si è concentrato su
questo: ``cosa porto come mio punto di forza?''. L'autoanalisi è culminata
nel chiedere espressamente alla persona che l'aveva assunto: ``prima mi avete
detto di no, ora mi avete preso: perché?''.

\paragraph{Come si capisce se un progetto è andato bene?} Dipende
dall'\textbf{obiettivo}. Se l'obiettivo era vendere milioni di copie e il
gioco le ha vendute, bisogna capire perché. Se, più ragionevolmente,
l'obiettivo era un gioco da tavolo che permettesse di mettere in pratica le
conoscenze del corso, va valutato per quello: ha permesso di farlo? Se sì,
perché? Se no, perché no? L'obiettivo va fissato \textbf{all'inizio del
progetto, non alla fine}, proprio per poterlo valutare dopo.

\subsection{Quando il team cambia di continuo}

\paragraph{Tutti questi cambiamenti non sono dannosi per il progetto?}
Assolutamente sì, ma quasi mai sono volontari. I progetti andati male sono
spesso progetti grandi.

Il relatore racconta di \emph{Cyberpunk 2077}. Conosceva persone in CD
Projekt Red (``il mondo è piccolissimo, le cose si sanno''). Alla Gamescom va
con Diego Ricchiuti, direttore dell'accademia, a vedere la presentazione del
primo gameplay, quella ``in cui regalavano il giubbottino da samurai''.
Finito il video, tra gli applausi, si guardano e capiscono entrambi che
\textbf{il gioco non è pronto}, a soli quattro mesi dall'uscita prevista:
Diego aveva notato il combattimento e la guida. Gli amici di CD Projekt Red
(uno aveva lavorato con lui a Blizzard per sei anni) gli confermano tra le
righe che ``ce la faremo, ma ci sono ancora cose da sistemare''. Il gioco fu
poi rimandato più volte\footnote{In lezione il relatore dice ``due, forse
tre volte''; l'uscita di \emph{Cyberpunk 2077} è stata effettivamente
rinviata tre volte prima del dicembre 2020.}.

Per deformazione professionale, chi ha sviluppato l'occhio critico non
riesce più a godersi un trailer: in un trailer di \emph{The Day Before} Diego
gli aveva fatto notare che l'effetto della fiammata di un fucile era stato
chiaramente aggiunto in postproduzione e non faceva parte del gameplay. ``È
il triste epilogo di questa storia'': non riescono più a godersi un trailer.

Il punto è che più il team è grande, più queste variabili vengono imposte, e
più è facile che il fallimento \textbf{non dipenda dal fatto che il prodotto
sia brutto} (\emph{Cyberpunk} è un bellissimo gioco sotto molti punti di
vista, e il relatore lo sta rigiocando ora che è stato corretto ed espanso).
Nei grandi gruppi qualcuno arriva, qualcuno si dimette, cambia il budget,
arriva il rinvio; nel caso di \emph{Anthem}, il team ha fatto così tante volte
forming--storming--forming--storming che ne è uscito ``un prodotto che
dovevamo espellere: l'abbiamo espulso, buttato fuori dal nostro organismo pur
di liberarcene''.

Tra i cambiamenti c'è anche la crescita verticale: se il capo viene promosso
perché lavora benissimo, bisogna scegliere un nuovo capo, ed è sempre un
caos (chi ci resta male, chi pensava di meritarlo, chi voleva il suo amico,
chi non sopporta il prescelto). Per questo il relatore è \textbf{contrario
alle promozioni durante i progetti}: si rischia di rovinare tutto il lavoro di
storming.

\subsection{Obiettivi chiari nei piccoli gruppi}

Nei gruppi piccoli, come quelli degli studenti, il problema deriva più
spesso dal \textbf{non avere chiaro che cosa si vuole ottenere}: un prodotto
puramente educativo e formativo, o un prodotto richiesto dall'accademia che,
se viene bene, si potrebbe anche vendere? Un gioco commerciale o un gioco
artistico? Va chiarito subito, perché ci saranno sempre momenti in cui
scegliere tra A e B, e a guidare la scelta sarà la domanda: \emph{quale delle
due opzioni ci porta più vicino all'obiettivo?} Se l'obiettivo non è chiaro,
non si sa cosa scegliere.

Questo porta molte aziende a fare \textbf{scelte di design che in realtà sono
scelte commerciali}. Le microtransazioni sono scelte commerciali: nessuno ha
mai pensato ``sai cosa piacerebbe tanto ai giocatori? Pagare le bustine''.
Prima le bustine si guadagnavano giocando; ora la tendenza è ``potresti
giocare, ma se paghi è meglio''. Si pensi alle polemiche su Ubisoft per
l'opzione di ``saltare il gioco'' pagando, o ai famosi \emph{boost} di livello
di \emph{World of Warcraft}: ``devo pagarti per saltare il gioco?''.

Una studentessa (Silvia) cita \emph{Final Fantasy XIV}, che a un certo punto fu
``resettato'' perché non funzionava: un esempio positivo. Il relatore aggiunge
quello negativo di \emph{Age of Conan}, che non è mai riuscito a fare quel
reset iniziale (``è una schifezza, ve lo ripresentiamo quando è pronto''). Ci
vuole consapevolezza, ma anche i mezzi: un conto è Square Enix, che può
permetterselo, un altro è un'azienda che non può chiudere tutto per un anno
continuando a pagare gli stipendi.

Riassumendo: nei piccoli gruppi i problemi sono più di
\textbf{comunicazione e chiarezza} (obiettivi non chiari o che si spostano);
nelle grandi aziende di più di \textbf{persone che vanno e vengono, obiettivi
e tempistiche che cambiano}. In accademia le scadenze sono fisse (il gioco da
tavolo si consegna a fine anno scolastico, a giugno, e giugno resta); in
azienda può succedere che una feature prevista per la \emph{vertical slice}
di luglio venga anticipata a maggio per portarla alla Gamescom, e un sistema
perfetto per consegnare a giugno può non esserlo più per consegnare a maggio.

\subsection{L'ambiente determina la qualità del prodotto}

\paragraph{È vero che chi passa per Electronic Arts finisce per fare brutti
giochi?} Uno studente cita BioWare, che faceva ottimi giochi finché era
autonoma. Il problema non è dei developer, che non sono ``improvvisamente
diventati incapaci'': è evidente che il problema è altrove, e che
\textbf{l'ambiente in cui si fa un gioco determina molto spesso la qualità del
prodotto}.

Un esempio locale positivo: Tommaso Bonanni, conosciuto dal relatore insieme a
Diego e già docente dell'accademia, ha fondato a Roma Caracal Games Studio,
applicando fin dall'inizio questi concetti di cultura aziendale. L'azienda è
rimasta piccola, ma ha molto chiaro cosa vuole fare e come. Quando si
predispone fin dall'inizio l'azienda e il team a lavorare in un certo modo è
fattibilissimo; è molto più difficile quando si è EA, si acquisisce
un'azienda e le si dice ``da oggi si fa così''.

Il relatore ha vissuto la transizione da Blizzard ad Activision Blizzard (uno
dei motivi per cui se n'è andato, volontariamente). Ha visto la cultura
interna cambiare per gradi:
\begin{enumerate}
  \item ``facciamo giochi fighissimi, perché i giochi devono essere
        fighissimi'';
  \item ``facciamo giochi fighi così vendono di più'';
  \item ``vendiamo di più facendo giochi fighi'' (sembra uguale, ma è molto
        diverso);
  \item ``facciamo più soldi''.
\end{enumerate}
Lo si capiva dalle scelte di design, osservando quale fosse il peso più
grande sulla bilancia. Lo stesso varrà per gli studenti, nei loro giochi e nei
colloqui di lavoro: capire subito che cosa vuole l'altra parte.

\paragraph{Leader e visione.} Nei gruppi dell'accademia, dal secondo anno, c'è
una persona che assume il ruolo di \emph{lead} o \emph{scrum master}. È
diverso dall'essere un \emph{game director} (``il gioco si fa come dico io,
aiutatemi a farlo con la mia visione''), come avviene nei grandi studi: da lì
nascono i nomi nei titoli, come \emph{Sid Meier's Civilization}, che indicano
la persona che ha fatto quel gioco, oppure un puro branding, come \emph{Tom
Clancy's Rainbow Six}.

\paragraph{Inclusion come jolly.} A margine, uno studente osserva che in molte
aziende ``la risposta giusta è sempre \emph{inclusion}'', usata come un
jolly. È comunissimo, conferma il relatore: \emph{diversity and inclusion}
viene spesso usata per dire ``così includiamo più clienti'', ma non è quello
il suo punto (``altro discorso'').

% ---------------------------------------------------------------------
\section{Confrontarsi con chi è più bravo}
\label{sec:z1-rivalita}

Arrivato a Riot dopo sei anni a Blizzard, il relatore pensava di avere un buon
bagaglio di esperienza. Si è trovato invece ``come un fante buttato in mezzo
ai Navy SEAL'': persone che avevano superato 11, 12, 13 colloqui, con un
livello molto alto.

Individua subito un collega francese che, nella sua testa, diventa ``la
versione migliore di me'': stesso ruolo (community manager), esperienza
simile (arrivava da Ubisoft, come quasi tutti i francesi di Riot), ma molto
più spigliato, sempre con la battuta pronta, ``l'anima della festa''. Il
relatore invece è introverso: un workshop come questo per lui è
\emph{emotivamente drenante}, non ama parlare in pubblico pur amando
condividere le proprie esperienze, e non riesce a sostenere conversazioni
come ``cosa hai fatto questo weekend?''. Lavorando raramente insieme, aveva
sviluppato una rivalità ``da allenatore di Pokémon'' (lui era Ash, il
relatore l'antagonista).

Dopo sei mesi, a un evento di \emph{team building} aziendale (a Tenerife), il
suo capo, molto attento a queste dinamiche, li mette apposta nello stesso
gruppo. A fine serata, complice l'open bar, il relatore gli confessa di averlo
mal sopportato dal primo giorno per invidia della sua estroversione. Il
collega ride e risponde che pochi giorni prima aveva detto al capo la stessa
cosa al contrario: invidiava a Stefano il suo essere sempre organizzato, sul
pezzo, presentarsi nella sala riunioni giusta (mentre lui finiva in quella
sbagliata), arrivare con gli appunti, con i dati a supporto e le analisi di
mercato, mentre lui spesso non aveva nemmeno letto la mail del meeting.
Un'invidia in due direzioni, che li ha portati a lavorare molto insieme e a
migliorarsi a vicenda.

Le lezioni:
\begin{itemize}
  \item non bisogna demoralizzarsi vedendo qualcuno che fa lo stesso lavoro
        ed è più avanti sotto alcuni aspetti: spesso c'è un bilanciamento. I
        ``punti caratteristica'' sono limitati: se li ha messi in carisma,
        non li avrà messi in saggezza;
  \item conviene \textbf{chiedere a queste persone come mai} sono così brave
        in qualcosa, come si fa con un progetto andato bene: magari non lo si
        farà mai bene quanto loro, ma si comincerà ad avere il
        \emph{mindset} giusto; e viceversa, loro saranno spinti a chiedere a
        noi;
  \item spesso \textbf{le proprie doti le scopriamo dagli altri}. Il relatore
        non si era mai accorto di essere ``quello che resta calmo durante le
        emergenze'' (il lancio di \emph{Diablo~III}, quello di
        \emph{Draenor}, i problemi dal vivo alla Games Week) finché non
        gliel'hanno fatto notare tutti: per lui era la normalità. A chi pensa
        ``non sono capace, forse non è il mio percorso'', può darsi che
        semplicemente nessuno gli abbia ancora detto ``questa cosa la sai
        fare davvero bene, concentrati su questo''.
\end{itemize}

% ---------------------------------------------------------------------
\section{Domande sul lavoro nell'industria}
\label{sec:z1-industria}

Diverse domande degli studenti hanno riguardato le condizioni di lavoro nei
grandi studi internazionali. Le risposte sono raccolte qui per argomento.

\subsection{Presenza, remoto e trasferimenti}

\paragraph{Grandi studi: quasi solo in presenza.} Sia a Blizzard sia a Riot si
lavorava \textbf{solo in presenza}, senza smart working: sei anni a Parigi e
cinque a Dublino. Si parlava solo inglese: anche nei meeting tra soli
italiani si tendeva a comunicare e a scrivere in inglese, per correttezza
verso chi si fosse aggiunto e per lasciare una traccia comprensibile a tutti.
Il francese serviva per fare la spesa e sopravvivere, non in ufficio.
L'ambiente era estremamente eterogeneo (francesi, polacchi, spagnoli e
molti altri): durante il colloquio dal vivo a Dublino, mentre si preoccupava
di essersi vestito in modo adeguato, vide passare dietro la parete a vetro un
collega svedese con i dreadlock, a piedi scalzi, con una pianta sotto un
braccio e un aquilone sotto l'altro, e capì di essere ``a posto così''.
Un'esposizione a molte culture e modi di pensare che il relatore considera
bellissima.

Anche oggi, tranne alcuni studi particolari, la maggior parte del lavoro si
svolge di persona: amici a Paradox lavorano a Stoccolma, i pochi rimasti a
Blizzard si sono dovuti trasferire a Irvine. Per posizioni \emph{junior} o
\emph{entry level} i grandi studi raramente concedono il lavoro da casa; per
profili senior può esserci qualche eccezione (ad esempio nove mesi da remoto
e tre in sede).

\paragraph{Relocation.} Spesso tocca spostarsi, ma il relatore non ha mai
pagato un trasloco di tasca propria: queste aziende offrono la
\emph{relocation} (aiuto a trovare casa, ad aprire il conto in banca e
supporto pratico), perché è un investimento sulla persona.

\paragraph{Colloqui mentre si lavora altrove.} Per i colloqui in presenza
prendeva ferie, senza dire dove andava. Il consiglio generale: per fare
colloqui \textbf{prendete ferie o permessi}, non datevi malati, ``che non è un
bel modo di comportarsi''.

\paragraph{Remoto ``doppio''.} Un collega lavora per l'azienda coreana che
gestisce \emph{PUBG}, con ufficio ad Amsterdam, ma lavora da Milano: un
``doppio remoto''. Si trova, ma in genere un ambiente anglofono con sedi in
Europa chiederà di trasferirsi in Europa.

\paragraph{Famiglia e trasferimenti.} Il relatore ha avuto la fortuna di
sposare una persona che fa lo stesso lavoro (conosciuta lavorando su
\emph{Warhammer Online}, poi con lui a Blizzard e a Riot, nella
localizzazione e nel doppiaggio, ad esempio di \emph{World of Warcraft} in
italiano): avendo competenze rare, spesso venivano assunti entrambi, al
massimo a pochi mesi di distanza. Per i developer la soluzione frequente è
che la famiglia segua la persona assunta, perché lo stipendio basta per
tutti. Ma non sempre c'è equilibrio, e a volte uno dei due deve sacrificarsi:
sua moglie ha lasciato Blizzard per seguirlo a Dublino (per poi essere
assunta da Riot); lui ha lasciato Riot per tornare in Italia quando lei ha
trovato lavoro qui. Ha dovuto rinunciare più volte a lavorare in Islanda per
CCP (\emph{EVE Online}), perché non era una scelta solo sua. Si fanno
compromessi, ma si riesce a far quadrare.

\paragraph{Ci si sposta per ogni progetto?} No: si sono spostati
esclusivamente per cambiare azienda. Cambiando progetto nella stessa azienda
è raro, quasi mai (``non l'ho mai visto succedere''), che venga chiesto di
trasferirsi. In Francia stavano benissimo (vivevano fuori Parigi; ``i
problemi con i francesi sono in realtà problemi con i parigini''), parlavano
francese ed erano vicini alla doppia cittadinanza: se ne sono andati perché
non vedevano un futuro a Blizzard, per il modo in cui l'azienda trattava i
colleghi.

\paragraph{Il mercato oggi.} Non esiste un \emph{silver bullet} per trovare
lavoro domani dove si vuole. Con tutti i \emph{layoff} recenti ci sono
migliaia di esperti del settore che cercano lavoro: serve pazienza, ma è
fattibile.

\paragraph{Studi indie e remoto.} Negli studi più piccoli il remoto è più
diffuso: a Stormind circa 60 persone su 95 lavorano da remoto, ma lo studio
si è strutturato fin dall'inizio per farlo. Serve consapevolezza. Bending
Spoons (che sviluppa app per mobile, non giochi, ed è l'azienda dell'app
Immuni), pur essendo grandissima e con sede a Milano, lavora da remoto e
lascia l'ufficio a chi vuole: lo smart working è posizionato quasi come un
\emph{benefit}, ma permette anche di avere un ufficio da 300 persone invece
che da 700. Le sedi fisiche costano decine di milioni all'anno (Blizzard ha
chiuso le sedi di Amsterdam e di Parigi). Tuttavia, per lavori molto
verticali e specifici gli studi preferiscono avere le persone dal vivo; i
ruoli più spesso remoti sono quelli che non fanno sviluppo (customer service,
PR\dots).

Anche gli artisti lavorano spesso in presenza: l'accademia può dare solo una
\emph{simulazione} di quanto stretto sarà il contatto tra programmazione,
design e art. ``Sarà ancora di più.''

\paragraph{La comunicazione da remoto.} È più difficile. Pur essendo
estremamente introverso, il relatore \textbf{obbliga tutti a usare la webcam}
nei meeting in cui conta: non in un aggiornamento con novanta persone
collegate, ma certamente in un colloquio a due sulla performance dell'ultimo
anno. La componente fisica è una parte enorme della
comunicazione\footnote{In lezione si parla di ``quasi l'ottanta per cento'';
il dato va preso come indicativo, poiché le percentuali di questo tipo
circolano in forme diverse e sono spesso citate fuori contesto.}: anche se si
è timidi, accendere la webcam è l'approssimazione più vicina all'essere dal
vivo. Negli ambienti dove si lavorava senza webcam la situazione era
``esponenzialmente peggiore''.

\subsection{Differenze culturali nel team}

\paragraph{Gestire sensibilità diverse.} A Riot i team erano divisi per
\emph{lingua}, non per nazione: il team di lingua spagnola includeva anche un
messicano, quello di lingua russa includeva russi e ucraini, proprio nel
periodo della guerra in Crimea. Tra i colleghi non c'era alcun problema (erano
amici), ma ci si è posti la domanda: \emph{come comunichiamo con una community
che sarà divisa su questo tema?} Nei giochi online la divisione era evidente:
in \emph{Rocket League} c'era chi, con un compagno di squadra con la bandiera
avversaria, si faceva otto autogol pur di farlo perdere. Allo stesso modo,
con molti americani mandati a Dublino a coprire posizioni scoperte, c'erano le
elezioni (Clinton--Trump, Biden--Trump) e varie guerre: serviva sempre una
certa sensibilità.

Imparare a navigare queste situazioni lo ha reso un professionista migliore:
gli ha insegnato a fermarsi e chiedersi ``io la percepisco così; gli altri
dieci, con background, storie e lingue diverse, come la penseranno?'', a
ragionare più \emph{long term} e più \emph{big picture}, con più empatia.

\paragraph{Stili di comunicazione.} Un partecipante che ha lavorato nel Regno
Unito chiede dello stile comunicativo italiano, a volte frainteso. Uno dei più
grandi complimenti ricevuti dal relatore è stato ``non sembri italiano''
(``gesticoli poco per essere italiano'', ``parli bene inglese per essere
italiano''). Ha imparato ad \textbf{adattare la comunicazione alla persona}:
con spagnoli e francesi la battuta sulla nazione funziona, perché c'è un
``odio--amore'' condiviso; con russi e giapponesi no. Doveva adattare anche
l'umorismo. Un collega metà russo e metà messicano diceva di aver preso ``il
lato peggiore di entrambe le nazionalità'', ed è stato un ottimo maestro per
imparare a leggere culture diverse.

Non bisogna generalizzare, ma gli stereotipi possono essere un punto di
partenza ragionevole, \emph{sapendo che troverete sempre le eccezioni}. E
vanno letti correttamente: gli americani tendono a essere più estroversi e ad
avere sempre qualcosa da dire non perché ``sono fatti così'', ma perché
vengono educati così (all'università tra i corsi c'è il \emph{debate}, saper
argomentare in pubblico); un collega giapponese può essere più introverso per
ragioni culturali. Nel dubbio si parte da lì, poi si conosce la persona.

\subsection{Negoziare, discriminazione e stipendi}

\paragraph{La famiglia come leva.} Avere un partner riduce le opzioni, ma le
aziende lo sanno: uno spostamento è per due. Si può usare come leva per
negoziare (un compenso diverso, un supporto per il trasloco): se un'azienda ha
deciso di volere quella persona, 2, 3, 4 o 10.000 euro in più all'anno, per
un'azienda che fattura centinaia di milioni o miliardi, non cambiano nulla.

\paragraph{Essere sposati è un motivo per non essere assunti?} No. In alcuni
paesi è proprio illegale discriminare su questo: in Francia era illegale
chiedere CV con data di nascita, foto, sesso o religione, e chiedere se si
vogliono figli. Purtroppo, dai racconti dei suoi studenti di master (in
settori come \emph{beauty and wellness} e \emph{food and wine}), in piccole
realtà italiane queste domande (``hai 25 anni, hai intenzione di
sposarti?'') vengono ancora fatte, e servono a filtrare i candidati. Nelle
grandi aziende ci sono sistemi per evitarlo; anzi, molte danno \emph{più}
benefit a chi ha famiglia (``porta pure la famiglia, è tutto pagato''),
perché le persone competenti hanno spesso una certa età e quindi,
ragionevolmente, una famiglia.

\paragraph{Gli stipendi.} Dipendono molto da dove si vive. Lo stipendio a Riot
non era paragonabile a quello di Blizzard perché in Francia era tassato al
52\% e in Irlanda al 17\%\footnote{Aliquote come riportate in lezione,
riferite alla situazione personale del relatore.}: nessuno stipendio
competitivo può compensare il 35\% di differenza di tasse. Negli Stati Uniti,
con pochissima previdenza pubblica, gli stipendi sono molto più alti perché
poi ci si paga tutto (``se hai un raffreddore paghi tu''). Bisogna valutare il
potere d'acquisto di quei soldi nel paese in cui si andrà a vivere.

\subsection{Freelance e \emph{work for hire}}

All'inizio della sessione pomeridiana uno studente che lavora nell'audio
chiede quali siano le condizioni per chi lavora come \emph{freelance}.
\begin{itemize}
  \item Dipende molto dal progetto: spesso il contratto è a progetto, per
        portare a termine una parte o un intero progetto.
  \item Non tutte le aziende usano freelance: Riot e Blizzard lo facevano
        pochissimo e preferivano assumere (magari trattando poi la persona
        ``come un freelance'', lasciandola a casa a fine progetto).
  \item Spesso i lavori offerti come freelance sono in realtà
        \textbf{sostitutivi di un full time} (``freelance, ma sei ore al
        giorno'').
  \item Si può fare, ma bisogna essere molto chiari nel definire l'impegno
        (quante ore, come vengono tracciate): non tutti specificano cosa
        significa davvero ``freelance''. Lavorando così per due o tre aziende
        si può arrivare a un tempo pieno.
  \item È più comune il \emph{work for hire}: aziende che forniscono personale
        ad altre aziende; si è assunti da X ma si lavora ``embeddati'' presso
        Y (una sorta di subappalto). Anche lo studio del relatore lo fa.
  \item Il puro freelance di quattro mesi esiste, ma di solito per quei
        quattro mesi diventa un lavoro a tempo pieno; è raro ``quattro mesi,
        due ore al giorno''. Nelle grandi aziende è improbabile lavorare nello
        stesso mese per Ubisoft, Riot e Blizzard tre ore ciascuna; è più
        realistico lavorare per più startup indie italiane, che non possono
        permettersi un'assunzione ma possono pagare due ore al giorno.
  \item Molto passa dai contatti: trovare le aziende giuste e rendere esplicite
        le condizioni.
\end{itemize}

% ---------------------------------------------------------------------
\section{Il circolo virtuoso: feedback, miglioramento, fiducia}
\label{sec:z1-triangolo}

\subsection{Il triangolo}

L'obiettivo a cui il workshop vuole arrivare è un \textbf{triangolo che si
autoalimenta}:
\begin{enumerate}
  \item si dà un \textbf{feedback} con l'obiettivo di migliorare le
        performance (inclusi elementi caratteriali);
  \item la persona \textbf{migliora} e vede di averne beneficiato;
  \item questo genera \textbf{fiducia}, che rende più facile dare e ricevere
        il feedback successivo.
\end{enumerate}
Togliendo uno qualsiasi dei tre elementi, il circolo non è altrettanto forte.

In qualsiasi fase si parta, il rispetto per i colleghi è inizialmente dettato
dall'\emph{educazione}, non dalla fiducia. La fiducia arriva quando i colleghi
cominciano a dire cose giuste e sensate e a \textbf{mantenere la parola}.
Dare e ricevere feedback è il modo in cui, in un certo senso,
``inganniamo il cervello'' per creare una \emph{case history} di fiducia: si
ha motivo di fidarsi di più, di ascoltare di più, di essere più ricettivi. E
si diventa sempre più bravi ad accettare il feedback e ad agire di
conseguenza.

\subsection{Accettare un feedback non vuol dire fare come dice l'altro}

Dieci anni fa il relatore non era ostile ai feedback, ma molto
\emph{difensivo}: ``sì, sì, va bene, ora ci guardo''; oppure, a chi gli diceva
``avresti dovuto farla così'', rispondeva spiegando perché l'aveva fatta in
quel modo. Oggi ha imparato a dire ``ok, grazie, ora ci ragiono'': non deve
difendere la propria posizione da chi viene solo a dire ``hai pensato a
questa cosa?''. Prima lo viveva come un attacco (``mi viene a dire che non so
fare il mio lavoro''), ora come un ``ottimo punto di vista, lo prendo in
considerazione''.

Con una postilla fondamentale: \textbf{accettare il feedback non significa
fare come dice l'altro}. ``Mi hai detto che preferisci questo mob verde. Ho
pensato anche a farlo verde? Sì. Penso comunque che giallo sia meglio? Sì, lo
faccio giallo.'' Il relatore accetta \emph{sempre} il feedback, valido o no:
garantisce di tenerlo a mente quando lavora. Che poi faccia come gli viene
detto è un altro discorso.

\subsection{Fiducia e scadenze}

La fiducia funziona allo stesso modo con gli impegni. Se si chiede a qualcuno
una cosa per venerdì e la consegna venerdì, la volta dopo gliela si chiederà
per venerdì. Se invece arriva il giovedì della settimana dopo, la volta dopo
gli si darà una \textbf{deadline anticipata}, come con chi arriva sempre in
ritardo agli appuntamenti (``se ti dico alle 7 arrivi alle 7 e mezza; la
prossima volta ti dico alle 6 e mezza'').

Al lavoro questo è dannosissimo, perché crea \textbf{deadline fittizie} che
non sono mai quelle vere, e nascono discussioni del tipo ``perché devo farlo
in due giorni se si fa in cinque?'' (``sì, se tu lo facessi in cinque''). La
conversazione giusta è un'altra: \emph{``che problemi hai avuto nel
consegnare venerdì?'' ``Questo, questo e questo.'' ``Cosa posso fare per
evitare che succeda?''}

Il relatore non vuole mai sentirsi dire ``non sono riuscito a farlo''.
Preferisce di gran lunga sentirsi dire, tre giorni prima, \textbf{``potrei non
riuscire a farlo, per questi motivi''}. Chi ha la visione d'insieme può dare
più budget, più persone, spostare la deadline, cambiare lo scopo del progetto:
può fare qualsiasi cosa per aiutare. Ma se lo si scopre il venerdì, cinque
minuti prima della deadline, non può fare più niente. Il messaggio è:
\textbf{``aiutami a poterti aiutare''}. Non penserà mai male di chi dice ``non
ce la faccio per questi motivi'', perché non sta dicendo ``non ho voglia'': sta
indicando i problemi, e lavorare su quei problemi è compito del manager. Non
può invece fidarsi di chi fino al giorno prima dice ``tranquillo, ce la
faccio'' e il giorno dopo ``non ce l'ho fatta''.

Questo è fondamentale per chi lavora in programmazione, design o art, perché
ci saranno sempre deadline: bisogna sempre \textbf{lavorare sulle aspettative},
dire al proprio manager (o al proprio report) cosa è lecito aspettarsi. Non
bisogna fare come i pinguini di \emph{Madagascar} (``sorridete e salutate, va
tutto bene''): se c'è un problema, va detto, altrimenti si peggiora solo la
propria situazione.

Da italiani si tende a non voler mai alzare la mano per chiedere aiuto
(``no, no, ce la facciamo, lo faccio io''). A volte non ce la si può fare, non
perché non si è capaci, ma perché vari fattori lo impediscono. Accettarlo è
importante: \textbf{è molto meglio avvisare e chiedere una mano che dire
``non ce l'ho fatta''}.

\subsection{Domande sul dare e ricevere feedback}

\paragraph{``Ne parliamo a fare, se poi fai come vuoi?''} Uno studente (Cristian)
riceve feedback che, dopo averli analizzati, non ritiene congrui; quando non
li segue gli viene detto ``ne parliamo a fare?''. C'è un problema di
\textbf{aspettative}: l'idea che accettare un feedback voglia dire fare come
dice l'altro. Accettarlo significa prendere in considerazione ciò che è stato
detto. Chi decide diversamente ha però un \textbf{onere}: argomentare perché
si fa comunque a modo proprio (``mi hai detto verde, la facciamo gialla
perché\dots''). Poi bisogna solo aspettare i risultati per vedere chi aveva
ragione.

La decisione spetta sempre \textbf{a chi ha la responsabilità di quella
decisione}. Se il game director deve scegliere tra A e B, anche se il
feedback dice B, la scelta è sua perché la responsabilità è sua. Non sarebbe
corretto essere ritenuti responsabili di una scelta delegata a chi dà
feedback. La risposta può essere semplicemente: ``l'ho tenuto in
considerazione, ma visto che ne sono responsabile io e le conseguenze sono
mie, ho scelto di fare X''.

\paragraph{E se l'altro fa ripicca?} Un partecipante chiede di chi, vedendo
il proprio feedback non accolto, si vendica ignorando a sua volta feedback
buoni. Per il relatore è un problema caratteriale: chi si comporta così ``è
già sulla strada per uscire''. Non c'è spazio per l'ego in team di questo
tipo. Il fenomeno nasce dall'idea che \emph{il feedback negativo faccia fare
brutta figura}, che il fatto che si potesse fare meglio tolga qualcosa alle
proprie capacità.

Non è così: \textbf{il feedback non è un giudizio}. Si può passare un anno a
Riot ricevendo solo feedback negativi e avere una valutazione annuale
stellare, perché ricevere molti feedback negativi significa essere una persona
approcciabile, che agisce sul feedback e invita gli altri a darlo. Il relatore
ha visto valutazioni molto peggiori per chi riceveva solo feedback positivi,
perché la gente aveva paura di dargliene di negativi: ma solo ``pacche sulle
spalle'' non fanno migliorare. Il problema va affrontato parlando con la
persona; a volte le conversazioni iniziano con \textbf{``ti do un feedback sul
feedback''}, cioè sul modo in cui l'altro ha dato il suo.

Esempi dalla sua esperienza:
\begin{itemize}
  \item lui chiede spesso di \textbf{non ricevere feedback senza esempi
        tangibili}, perché fa fatica ad astrarre: ``sei stato maleducato'' non
        significa nulla, ``ieri al meeting mi hai risposto male su questa
        cosa'' sì;
  \item ad altri serve l'opposto: ha ricevuto il feedback ``quando devi dirmi
        che ho sbagliato, puoi girarci un po' intorno e darmi del contesto?'',
        perché il suo stile diretto (``questa cosa non ha funzionato, dovremmo
        fare così'') feriva. Ha adattato il suo modo di dare feedback: con
        alcuni è diretto, con altri gira intorno. L'importante è saperlo.
\end{itemize}

\subsection{Non portare la soluzione}

Un consiglio centrale: \textbf{quando si dà un feedback, non si dà una
soluzione}. Non si dice ``io farei così'', ma ``secondo me questa cosa si
potrebbe fare diversamente, per questo e questo motivo: come potremmo farla
diversamente?''. È come chi smette di fumare perché vuole e chi smette perché
glielo dicono moglie e figli: se non si è convinti, spesso non si smette.

Quando il relatore diceva ``c'è questo problema, io farei così'', il meeting
successivo iniziava inevitabilmente con ``ci ho provato come hai detto tu, con
me non ha funzionato''. Sempre. Perché la soluzione era sua. La
responsabilità della soluzione deve essere invece \textbf{congiunta}: ogni
feedback deve creare un piccolo team che ha come obiettivo risolvere quel
problema. Questo toglie l'antagonismo (non ``ti dico che non sai fare il tuo
lavoro'', ma ``ti voglio aiutare a farlo meglio, lavoriamoci insieme'') e
sposta il focus: il problema non sei tu, è ciò che sta accadendo.

\paragraph{E se sono i piani alti a tenere all'oscuro dei problemi?} È una
questione di cultura aziendale. Il relatore fa sempre presente al proprio capo
il concetto di \textbf{service level agreement} (che verrà ripreso nella
seconda giornata): ``se mi dai questi input, ti garantisco questo output in
questi tempi''. Se il capo chiede delle cose senza dire budget, tempi e
perché, non può aspettarsi lo stesso risultato. È una comunicazione a due
vie: si dice al manager cosa serve per fare il lavoro, e lui deve fornirlo o
accettare che il risultato non sarà quello voluto.

\subsection{Amici, parenti e parole proibite}

\paragraph{``Non si lavora con amici e parenti''?} Per sua definizione
personale il relatore ha pochissimi amici (``stanno sulle dita di una mano'')
e una soglia molto alta per considerare amico un collega. Allo stesso tempo,
in quindici anni e con centinaia di colleghi, ha avuto solo \textbf{due}
esperienze in cui non è riuscito a lavorare con qualcuno: non serve essere
amici per far funzionare le cose. Con gli amici, anzi, gli viene più facile:
se dice a un amico che una cosa ``fa schifo'', l'amico non pensa ``come si
permette'', ma ``se perfino Stefano, che è mio amico, me lo dice, allora è
proprio così''. Diventa un rafforzativo, grazie alla fiducia costruita prima.
Il problema del lavorare con amici e parenti nasce dal \emph{non voler} dire
all'amico che sta sbagliando; se invece la confidenza permette di dirlo, si
può saltare la ``confezione'' del messaggio e passare direttamente al
contenuto.

\paragraph{Le frasi peggiori.} Alla domanda se esistano parole peggiori di
``ottimo lavoro'', il relatore indica la sua classifica: ``ottimo lavoro''
(detto in modo vuoto)\footnote{Il passaggio della trascrizione è
frammentario: si intende verosimilmente il complimento generico e di
circostanza, che non dice nulla di utile.}; poi la frase che lo ``fa uscire di
testa'', \textbf{``abbiamo sempre fatto così''} (``e infatti avete chiuso'',
come Blockbuster); infine \textbf{``siamo una famiglia''}, riferito al posto di
lavoro. Un'azienda non è mai una famiglia, perché non ci si può aspettare lo
stesso impegno da chi la possiede e da chi ci lavora: non c'è un
corrispettivo proporzionale tra impegno e ritorno. Per quanto inclusiva,
flessibile e razionale, non sarà mai una famiglia.

\subsection{Ego, muri e il peggior modo di scegliere un leader}

\paragraph{Chi mette davanti l'ego e alza muri.} Esistono, ma di solito durano
pochissimo in azienda. Non si tratta di avere un ego forte, ma di \textbf{non
saper ricevere feedback}. Spesso sono persone bravissime nel loro lavoro che
poi si chiedono perché tutti intorno diventino più bravi di loro: perché gli
altri accettano feedback e hanno più fonti da cui attingere, mentre loro si
sono rintanati nella loro torre.

Molti manager tendono ad \textbf{assumere persone un filo peggiori di loro},
per non fare brutta figura. Il relatore assume solo persone più brave di lui,
``perché non è una gara'': ha lavorato con molti community manager più bravi
di lui come community manager, ma lui era più bravo a fare il manager. Non
viene valutato su quanto sa fare il community manager, ma su quanto il suo
team sa farlo. Un partecipante chiede se la paura di avere qualcuno di
migliore esista anche in ambito accademico. Assolutamente sì: al relatore è
capitato di non essere preso come docente in un corso di formazione sui
videogiochi perché il responsabile aveva percepito che ne sapesse più di
lui\footnote{La trascrizione dice ``che non sapevo più di lui'', ma il senso
del passaggio richiede il contrario.}. ``Ma perché te ne importa? Tu sei quello che ha portato chi
sa fare il suo lavoro: ben venga.'' O si ha questa cultura, o si ha paura per
il proprio posto.

\paragraph{Promuovere il migliore.} Quando bisogna scegliere un leader per un
team (programmazione, design, art), molto spesso si sceglie \textbf{il migliore
tecnico}: il miglior programmatore, il miglior designer. È spesso la scelta
sbagliata: è come prendere un idraulico per fare l'elettricista. Fare il
\emph{people manager} non è come fare il designer o il programmatore; si può
essere un ottimo programmatore e un pessimo manager, e viceversa. Le aziende
l'hanno capito e lo fanno meno di prima. A Blizzard c'era un grave problema di
\emph{middle management}: molte persone in ruoli importanti erano state
promosse cercando le skill sbagliate, erano frustrate dal non saper fare il
proprio lavoro e schiacciavano chi stava sotto per paura di fare brutta
figura.

Un partecipante cita una teoria simile: una persona viene promossa finché non
raggiunge \emph{il proprio livello di incompetenza}, dove si blocca\footnote{È
il \emph{principio di Peter} (Laurence J.\ Peter, 1969), citato in lezione
come ``la teoria di Peter''.}. Un altro cita Michael Scott di \emph{The
Office}, su cui il relatore tornerà (sezione~\ref{sec:z1-quando}).

% ---------------------------------------------------------------------
\section{Perché la cultura conta: retention, costi e fiducia}
\label{sec:z1-cultura}

\subsection{L'ambiente di lavoro e la retention}

Perché tutto questo è così importante? Perché l'ambiente di lavoro conta
moltissimo per la \textbf{retention}, cioè per trattenere le persone e farle
lavorare in modo efficiente. Se contasse solo lo stipendio, molti non
lascerebbero mai Riot o Blizzard; e più si va avanti con l'età, più si impara
a dare valore anche al resto.

Inoltre, secondo gli studi citati nelle slide, chi vive il proprio lavoro in
modo negativo è \textbf{meno efficiente nel problem solving}: l'energia
mentale va verso ``come faccio a non farmi dare la colpa'' o ``come faccio a
fare bella figura'', non verso la soluzione del problema. In ambienti negativi,
dove ci si deve preoccupare della propria ``sopravvivenza'', si lavora
peggio.

\subsection{Il costo di chi se ne va}

Ogni volta che qualcuno si dimette o viene licenziato, il costo di trovare
una nuova persona, fare i colloqui, formarla e farle attraversare forming,
storming, norming\dots\ è circa \textbf{un terzo del suo stipendio}. Per questo
le aziende tendono a fare più \emph{retention} che \emph{acquisition}, almeno
con le figure considerate utili alla produzione.

Allo stesso tempo si vedono sempre più spesso \emph{layoff} e \emph{hiring
wave}: si licenziano 3.000 persone e sei mesi dopo se ne riassumono 2.500 o
3.500, ma tutte più junior. Il relatore ne ha scritto in un articolo su
Eurogamer.it, in cui spiega anche alcuni ``trucchi del mestiere'': perché i
licenziamenti avvengono sempre in certi periodi e sempre ``a gruppi'' (prima
Bethesda, poi Blizzard, poi un'altra\dots).

\paragraph{L'offerta di Riot.} A Riot esisteva una pratica (di cui il relatore
non ricorda il nome) per cui, firmato il contratto, si aveva \textbf{un anno
di tempo per dire ``non mi piace, non me la sento, me ne vado''} e ricevere
subito come buonuscita un anno di stipendio. Avevano stimato che fosse molto
meglio pagare subito un anno a qualcuno e lasciarlo andare che tenerlo due,
tre, quattro anni come \emph{underperformer}: la responsabilità della scelta
di assumere era dell'azienda\footnote{La pratica, nota come ``The Offer'', è
descritta in fonti pubbliche con condizioni diverse (un periodo più breve e
una somma pari a una frazione dello stipendio annuo). Si riportano le
condizioni come raccontate in lezione; andrebbero verificate.}. L'hanno usata
pochissimi, ma il solo fatto che esistesse faceva lavorare meglio anche lui da
manager, perché sapeva che chi non si trovava bene poteva andarsene con un
anno di tempo per trovare un altro lavoro.

\paragraph{Il ricircolo.} C'è un forte ricircolo nell'industria: molti ex
colleghi passano da Ubisoft a Paradox, da Paradox a Funcom, da Funcom a
Blizzard. Il relatore ha scelto di uscirne lavorando solo per aziende
\textbf{indipendenti}: Stormind non è un indie di cinque persone, ma è
indipendente, e non deve accontentare azionisti. Quando vede un'azienda con
utili miliardari licenziare migliaia di persone si chiede ``dov'è la linea''
(\emph{when enough is enough}): con azionisti che vogliono un ritorno, se
quello è l'obiettivo stanno facendo la cosa giusta, ma lui non condivide
l'obiettivo\footnote{In lezione si parla di ``6,5 miliardi di utile'' e di
``3.000 persone licenziate''; le cifre sono riportate come dette e non sono
state verificate.}.

\paragraph{L'open space vuoto.} A Riot, durante un'ondata di licenziamenti,
del suo team di trenta persone è rimasto solo lui: ``teniamo te come
coordinatore, poi ricostruirai il team''. Ha rifiutato di ricostruire un team
cancellato ``con un colpo di spugna'' e ha deciso che non avrebbe più lavorato
per una grande azienda. Nessuno aveva pensato di spostarlo, e ha lavorato per
tre mesi in un open space di cento metri quadri con la sua scrivania e altre
ventinove vuote. Su una lavagna aveva scritto i nomi dei suoi ragazzi e li
cancellava man mano che, anche grazie ai suoi \emph{referral}, trovavano
lavoro. Cancellato l'ultimo nome, si è dimesso: per lui il lavoro era finito.

\paragraph{Team che si dissolvono.} Uno studente cita un articolo su Larian
Studios, il cui management è contrario a ricostruire ogni volta il team a fine
progetto. È lo stesso problema: si pensa ai ``giochi Blizzard'' o ai ``giochi
BioWare'', ma le persone che li hanno fatti non ci sono più; è rimasto solo il
brand, e non ha senso avere le stesse aspettative da un gioco fatto da persone
diverse. E anche chi sopravvive ai layoff vive pensando che il prossimo
potrebbe essere lui: alcune persone, licenziate nell'ultima ondata a
Blizzard, gli hanno detto ``erano due anni che lo sapevo''. Vivevano a Cork,
``un paesotto irlandese dove o vai per lavorare a Blizzard o non ci vai'',
avevano comprato casa: non ci si può aspettare chissà che morale.

\subsection{La cultura è forte quando si lavora per gli altri}

\begin{quote}
  \emph{La cultura è forte quando le persone lavorano per gli altri e a favore
  degli altri, e non per se stesse.}
\end{quote}

Se si dà feedback per far migliorare l'altro e non per imporre il proprio
modo, se si aiuta per il bene del progetto e non per fare bella figura, la
cultura si rafforza. È dai \textbf{piccoli gesti} che emerge la cultura
aziendale: ad esempio, dal non dover chiedere ``perché mi stai aiutando?''.

\paragraph{L'incognita.} \textbf{L'incognita è spesso dannosa quanto le
cattive intenzioni.} In un esperimento (``mi pare a Stoccolma, cinque o sei
anni fa'') si distribuivano per strada banconote da 500 euro, e la maggior
parte delle persone le rifiutava chiedendosi ``dov'è la fregatura?''
\footnote{L'esperimento è riportato in modo vago in lezione e non è stato
possibile identificarlo con certezza.}. Lo stesso accadeva al relatore,
arrivato a PG e a Stormind da Riot: a ``vedo che sei in difficoltà, posso
aiutarti?'' la risposta era sempre ``perché?''. Alla fine rispondeva: ``perché
quando sarò in difficoltà io vorrei che tu mi aiutassi''.

Quando si comincia a lavorare per gli altri e si smette di pensare ``così avrò
la promozione, così avrò l'aumento'', promozioni, aumenti e rispetto arrivano
da soli. A Blizzard aveva fallito tre volte i colloqui per essere promosso a
manager (inteso come coordinatore del team); quando ha deciso di smettere di
pensarci e di fare solo il bene del proprio lavoro, la volta dopo è stato
promosso. Funziona così: chi aiuta i colleghi e fa andare bene il progetto
viene notato (``questa persona fa migliorare cinque persone, non solo se
stessa: vale di più'').

\subsection{Il lavoro è valutato come opera d'insieme}

Il relatore riconosce che buona parte del suo successo è venuto dal fatto che
il suo team lavorava bene, non dalla sua bravura personale. Anche i progetti
degli studenti, presentati a Diego, a Enrico, ad Andrea, verranno valutati
\textbf{come opera d'insieme}. Nessuno dice ``il \emph{combat system} di questo
gioco l'ha fatto X, che bravo'': si dice ``BioWare ha fatto un bel gioco'',
``Blizzard ha fatto un bel gioco''.

Lo ha replicato anche lato HR. In molte aziende ci sono \textbf{bonus
personali} legati a obiettivi (se il gioco vende X, se la feature è pronta in
tempo). Il relatore li ha dimezzati: metà restano personali, metà diventano
\textbf{bonus di team}, che tutti ricevono se il team raggiunge gli obiettivi.
Un incentivo concreto a collaborare: non solo si dice che il lavoro di gruppo è
importante, lo si dimostra dandogli più valore economico. Aveva proposto il
100\% di bonus di team, ma gli è stato concesso il 50 e 50; ha funzionato
comunque, anche se sarebbe stato ancora più potente. Ovviamente, alla
domanda ``perché mi aiuti?'' si può anche rispondere ``perché voglio il
bonus'': un aiuto dato ``per il motivo sbagliato'' era un compromesso che era
disposto ad accettare.

\subsection{Il coraggio di chiedere aiuto}

Un partecipante osserva che aiutarsi dovrebbe essere la norma. Il modello
diffuso, invece, è ``devi essere forte e indipendente, non chiedere mai una
mano'', e vale in tutte le direzioni (dal maschilismo alla ``donna forte e
indipendente che non ha bisogno di nessuno''): \textbf{non ammettere di avere
bisogno di aiuto è sempre uno dei motivi per cui le aziende vanno male}, che si
tratti di uomini, donne o persone non binarie. Il problema è non avere la
sicurezza di poter alzare la mano.

Il relatore ha cercato di togliere ogni conseguenza dall'aver alzato la mano:
se si \emph{depenalizza} qualcosa, le persone non sono più disincentivate a
farla. Ha detto espressamente che non avrebbe mai valutato nessuno per il
numero di volte in cui chiedeva aiuto o per il numero di feedback negativi
ricevuti, e questo ha funzionato.

\subsection{Perché si dà un feedback}

\emph{Perché avete dato il vostro ultimo feedback? Che cosa volevate ottenere?}
Le risposte della classe:
\begin{itemize}
  \item per \textbf{riconoscere il valore} di un collega, anche verso il
        manager, perché qualcuno che non aveva visibilità vedesse quanto
        quella persona fosse efficace;
  \item per \textbf{spingere una persona a superarsi};
  \item per aiutare un collega a rendere il prodotto all'altezza di ciò che
        voleva rappresentare;
  \item per \textbf{genuinità}: c'era qualcosa che poteva essere fatto meglio.
\end{itemize}
Il relatore aggiunge che la forma cambia col tempo: all'inizio era ``devo
darti un feedback, hai dieci minuti?''; ora è un messaggio informale su
Discord, perché ha costruito un rapporto di fiducia e sa che non verrà letto
come ``l'HR manager mi scrive, sono nei guai?'', ma come ``Stefano è così, mi
dà feedback, resterà tra noi e lo risolveremo insieme''.

Un altro motivo è \textbf{risolvere incompatibilità caratteriali}: le
conversazioni più difficili, del tipo ``io e te non stiamo funzionando, cosa
possiamo fare diversamente?''. Sono importanti perché, se due responsabili non
vanno d'accordo, il team se ne accorge, comincia a raccontarsi storie (``si
odiano, ci saranno problemi'') e, per difendere il proprio manager, comincia a
giudicare male il lavoro dell'altro gruppo. È la dinamica da fermare:
\textbf{valutare gli altri solo per difendere la propria fazione}.

Dare e ricevere feedback sono \textbf{due muscoli da allenare} (il relatore è
sempre stato più bravo a darlo che a riceverlo), e si impara col tempo. È come
fare colloqui: fare colloqui è una skill a parte, che non indica quanto si è
bravi nel proprio lavoro. Essere scartati non significa non saper fare quel
lavoro: spesso significa aver fatto male il colloquio o che è stato preferito
qualcun altro.

\paragraph{Feedback pubblici e privati.} Uno studente racconta di aver dato un
feedback negativo a un superiore, senza un motivo particolare, per dimostrare
a sé e ai colleghi che non si viene licenziati per questo. Il relatore coglie
l'occasione per precisare che \textbf{i feedback negativi si danno
individualmente}: lui dà moltissimi feedback di cui il resto dell'azienda non
sa nulla. Reagirebbe molto male se, in un meeting, qualcuno alzasse la mano per
dire ``comunque questa cosa non l'hai fatta bene'': quello non è feedback, è
discussione, accettabile solo se è un brainstorming e se l'opinione è stata
chiesta. Alcune aziende hanno piattaforme pubbliche di feedback, ma spesso
servono solo ad applaudire; lo stesso accade nelle aziende che
\emph{obbligano} a lasciare un certo numero di feedback a settimana, che
diventano ``lui è bravo, lei è competente, anche questa settimana l'abbiamo
fatto''. Bisognerebbe chiedere a chi le adotta cosa pensa di ottenere (forse
trasparenza), ma difficilmente una persona ``attaccata'' pubblicamente reagisce
bene.

% ---------------------------------------------------------------------
\section{Among Us: la caccia al traditore}
\label{sec:z1-among-us}

Prima della pausa pranzo il relatore mostra tre film di fantascienza di
qualche decennio fa e chiede che cosa li accomuni\footnote{I titoli non sono
leggibili nella trascrizione: viene citato \emph{Event Horizon} da uno
studente, e il film ``in mezzo'' è indicato come ``l'originale'',
verosimilmente \emph{La cosa} (\emph{The Thing}).}. In ognuno c'è un gruppo di
persone con un obiettivo e la presenza, reale o percepita, di qualcuno che
rema contro: \textbf{l'energia del gruppo viene spesa per trovare il
responsabile}.

\paragraph{L'esperimento degli alieni.} Nei master in presenza il relatore
propone un esercizio. Distribuisce segretamente cartellini plastificati con
scritto ``essere umano'' o ``alieno''; le carte alieno contengono istruzioni di
comportamento con cui gli alieni, che cercano di fingersi umani, si
tradirebbero (dare sempre ragione a chi ha parlato prima, finire ogni frase
con una domanda, parlare in un certo modo). Lo scopo: in dieci minuti formare
gruppi, e \textbf{vince il gruppo più numeroso senza alieni}; basta un alieno
per perdere. Si formano sempre gruppi piccoli (due, cinque, sei persone),
perché la tendenza è ``mi fido al massimo di voi due''.

Ma vincono i gruppi più grandi, per un motivo semplice: il relatore
\textbf{non distribuisce nessuna carta alieno}. Mostra a tutti entrambe le
carte, ma consegna solo carte ``umano''. Vuole simulare ciò che accade in
ufficio: la tendenza a percepire sempre qualcuno che rema contro (i manager,
l'altro team, quel collega), quando in realtà non c'è nessuno, ``è tutto nella
nostra testa''. Quando chiede le motivazioni, emergono sospetti
\emph{fabbricati}: ``faceva sempre domande'', ``parlava solo all'infinito''
(``no, parla così e basta''). Ci si fabbrica da soli i motivi per non fidarsi
degli altri, ed è comunissimo.

\paragraph{Among Us.} L'autosabotaggio è, per tornare ad \emph{Among Us},
\textbf{concentrarsi sul cercare cadaveri e traditori invece di fare le
task}: non si vince perché non si sono completate le task, mentre si cercava il
colpevole. Succede tantissimo negli uffici, ed è il motivo per cui l'80\% dei
problemi che il relatore risolveva erano di comunicazione. Entrando nei team
sentiva immancabilmente: ``facciamo il possibile, ma questa lotta continua con
l'altro gruppo è estenuante''. ``Perché c'è questo problema?'' ``Non lo so,
c'era già quando sono arrivato.'' Un paradosso: c'è un problema, nessuno sa
dire quale sia, ma c'è sempre stato e quindi lo si accetta.

\paragraph{L'esperimento delle scimmie.} È l'esempio, ``finto o reale'', delle
scimmie in gabbia\footnote{È il noto racconto delle ``cinque scimmie'', spesso
presentato come esperimento scientifico ma probabilmente apocrifo (lo stesso
relatore lo definisce ``finto barra reale''); nella trascrizione il numero
delle scimmie varia.}. In una gabbia con alcune scimmie pende una banana
elettrificata: chi la tocca fa prendere la scossa a tutte. Le scimmie imparano
e, quando una si avvicina alla banana, le altre la picchiano. Poi si
sostituisce una scimmia alla volta con una nuova: la nuova va verso la banana,
le altre la picchiano, e smette. Si continua finché nella gabbia ci sono solo
scimmie nuove, poi si toglie la corrente. Nessuna scimmia tocca la banana:
nessuna ha mai preso la scossa, ma tutte hanno imparato che chi tocca la
banana prende le botte.

Nei team è lo stesso: il motivo per cui si ``odia'' l'altro team si
trasforma, ma un vero motivo non c'è. Ne derivano la sfiducia, il non alzare
mai la mano per chiedere aiuto, il non dire mai ``bel lavoro'' agli altri
perché non sia mai che qualcuno ti veda fare un complimento ``a quelli''.
Risolvere queste dinamiche è lungo e tedioso, e sempre pieno di ``non mi
ricordo'', ``io non c'ero'' e ``abbiamo sempre fatto così'', la cosa più
frustrante per chi cerca di risolverle.

% ---------------------------------------------------------------------
\section{Il principio del default to trust}
\label{sec:z1-default-trust}

\subsection{La definizione}

Il relatore indica questa come la slide più importante della presentazione.
\textbf{Default to trust} è uno dei valori aziendali di Riot Games, che lui ha
tradotto, adattato e ``esportato'':

\begin{quote}
  \emph{Approccio ogni situazione di conflitto partendo dal presupposto che una
  persona abbia agito in un certo modo perché, con le informazioni a sua
  disposizione in quel momento, ha pensato che fosse la soluzione migliore, e
  non perché sia un'incapace o volesse danneggiare il progetto
  intenzionalmente.}
\end{quote}

In altre parole: \textbf{di default, mi fido}.

\subsection{L'esempio del ritardo al daily stand-up}

A Riot il relatore aveva un \emph{daily stand-up} dal vivo alle 15 con il suo
team di venti persone. Una ragazza arrivava regolarmente alle 15:15, per una,
due, tre settimane. Alla quarta ci ha parlato. \emph{Come iniziereste la
conversazione?} Le proposte della classe:
\begin{itemize}
  \item sottolineare l'importanza di esserci dall'inizio, perché potrebbe
        perdersi qualcosa che le interessa, e per rispetto dei colleghi
        presenti;
  \item ``devi arrivare in orario'', spiegando perché è importante;
  \item \textbf{chiedere perché}.
\end{itemize}

Il relatore osserva, anche dalle sue consulenze fuori dai videogiochi, che
\textbf{più si sale con la seniority}, più la risposta è ``devi smettere di
arrivare in ritardo'', spesso con un \emph{or else} finale (``altrimenti
succede quest'altra cosa''). Più si scende con l'età, e quindi con meno
abitudini consolidate, più la prima domanda diventa ``dimmi perché arrivi in
ritardo''. Già spiegare perché è importante è ``miglia avanti'' rispetto al
semplice ``smetti''.

Partire dal perché gli ha permesso di scoprire la causa: il meeting precedente
della ragazza, in un'altra zona del campus a un quarto d'ora di distanza,
finiva sempre tardi, e lei doveva attraversare il campus di corsa. Non aveva
la confidenza per dire all'altro manager, che insisteva perché si restasse
fino alla fine, che doveva andare; sacrificava quindi il meeting del relatore,
sapendo che era il più comprensivo dei due.

Il relatore si è offerto di \textbf{mediare}: ha detto all'altro manager che
o anticipava l'inizio del suo meeting, così da finire comunque in tempo,
oppure lasciava uscire la ragazza, perché quel meeting gli serviva. C'è anche
una questione di \textbf{economia di scala}, spesso sottovalutata: se venti
persone aspettano qualcuno per dieci minuti, si sono sprecati dieci minuti per
venti. Per lo stesso motivo il workshop è iniziato alle 10:05: iniziare alle
10:10 avrebbe significato far aspettare cinquanta persone cinque minuti in
più. A Riot c'era una scrivania vuota con un pupazzo gonfiabile, \textbf{Mr.\
Late}: sommando tutti i minuti di ritardo a meeting, call e conversazioni, in
un anno veniva fuori praticamente \emph{una persona a tempo pieno}. Serviva a
ricordare che anche cinque minuti contano, e che bisogna rispettare gli altri
arrivando in orario o almeno avvisando.

\paragraph{Gli effetti.} La conversazione era avvenuta nel primissimo periodo a
Riot, quando il team non aveva ancora motivo di fidarsi di lui. Risolvendo il
problema ha eliminato anche tutto lo strato di micro-dinamiche spiacevoli che
un manager impara a vedere (due colleghi che si guardano con la risatina
quando lei entra in ritardo). La ragazza ha cominciato a entrare tranquilla,
con una certa leggerezza: non le faceva piacere arrivare in ritardo (``ci sono
anche persone che lo fanno apposta, e quello è un problema da affrontare
diversamente''). Soprattutto, con i colleghi non parlava di lui come del capo
che minacciava, ma come di chi le aveva risolto un problema: era diventato
un \textbf{alleato}. La settimana dopo altre tre persone sono andate da lui
con un problema da risolvere. Ripetuto cinque, sei, sette, dodici volte, in un
paio di mesi poteva dire quasi qualsiasi cosa nei meeting sapendo che
richieste e critiche sarebbero state percepite come ``me lo chiede perché vuole
darmi una mano'', non per fare bella figura o per presentare i successi del
team ai suoi superiori.

\subsection{Condividere i successi, schermare i fallimenti}

Un altro elemento fondamentale (ripreso nella seconda giornata): il manager
\textbf{condivide i successi e protegge il team dai fallimenti}. ``Se va male
qualcosa è colpa mia; se va bene qualcosa, voglio che ne parliate voi.'' Il
lavoro del manager non dovrebbe essere valutato su ciò che presenta del lavoro
del team: tutti sanno comunque chi è il manager. Tanto vale lasciar presentare
il team, che potrà dire con orgoglio di aver presentato a Mike
Morhaime\footnote{Cofondatore ed ex presidente di Blizzard; nella
trascrizione ``Mike Morayme''.} la propria campagna europea per
\emph{StarCraft~II}. Ci sono manager che dicono ``come mi prendo le colpe, mi
prendo i meriti'', ma i meriti li hanno già: è inutile ostentarli. Con i
fallimenti invece il manager è la \textbf{rete di sicurezza} del team. Il
relatore ne ha fatto una questione di principio: \textbf{non scaricare mai la
colpa su qualcun altro}.

\paragraph{Colpevole o causa.} Un esempio che il relatore attribuisce a Pirate
Software, creatore di contenuti ed ex dipendente Blizzard\footnote{Il relatore
aggiunge che si tratterebbe del figlio dello sviluppatore a cui si sarebbe
ispirato il personaggio del giocatore ``senza vita'' nell'episodio di
\emph{South Park} dedicato a \emph{World of Warcraft}. L'informazione è
riportata come detta in lezione e non è verificata.}. Contattando il servizio
clienti Blizzard per motivi personali, aveva scoperto che i colleghi gli
avevano scritto nei ticket informazioni confidenziali che non avrebbero dovuto
condividere in quel modo, nemmeno con un collega. Non ha segnalato le persone
da punire: ha detto ``c'è un problema, evidentemente le policy non sono
chiare''. Quando gli hanno chiesto i nomi, ha rifiutato: ``se vi dico chi è,
punite quelli, e gli altri cinquecento continueranno a farlo; non avremo
cambiato la causa alla fonte''.

\begin{quote}
  \emph{È inutile trovare il colpevole. Voglio trovare la causa. Perché il
  colpevole cambia.}
\end{quote}

Il colpevole a volte è vittima della situazione: per questo la definizione
dice \emph{``con le informazioni a sua disposizione in quel momento''}. Se il
relatore avesse rimproverato la ragazza in ritardo, non solo non avrebbe
risolto il problema, ma avrebbe creato un'\textbf{ingiustizia}: lei, che
attraversava il campus di corsa per una colpa non sua e di cui non poteva
parlare, sarebbe stata anche rimproverata.

Nessuno fa un brutto lavoro apposta: nessuno è contento quando sbaglia o
perde. Perché al lavoro diamo per scontato che gli altri sbaglino
intenzionalmente? Se qualcuno ha fatto una cosa male, è perché in quel momento
credeva fosse la scelta migliore. La domanda giusta non è ``perché hai fatto
questa cosa che è stata un fallimento?'', ma:

\begin{quote}
  \emph{Che cosa ti ha fatto pensare che questa fosse la scelta giusta?}
\end{quote}

Magari erano limiti di budget, di tempo, di scopo: si agisce su quelli, non
sulla persona. Se si cambia la persona ma le si danno gli stessi fattori, è
molto probabile che arrivi alle stesse conclusioni, e ci si troverà a dover
``disciplinare'' anche lei.

\paragraph{La reputazione costruita sui rumor.} Durante il giro di
presentazione in un ufficio (non in nessuna delle aziende citate), al relatore
è stato detto: ``questo è il dipartimento di design; se ti serve qualcosa
chiedi a lei, non a lui, perché lui consegna sempre in ritardo''. Oltre a
trovarlo poco professionale, una settimana dopo ha chiesto: ``cosa ti ha fatto
questa persona?'' ``A me niente, però ho sentito che\dots''. Si torna alle
scimmie: magari quella persona non consegna in ritardo da mesi, ma ormai lo
dicono tutti. È un circolo vizioso da fermare subito.

\paragraph{Niente escalation inutili.} Lo stesso vale per il feedback: quando
qualcuno va dal relatore dicendo ``ho un problema con questa persona'', lui
chiede: \textbf{``perché ne stai parlando a me e non a lei?''}. ``Perché sei il
mio capo'' non basta. Passare da lui significa dare per scontato che l'altra
persona non abbia interesse a risolvere, mentre magari basta parlarle. È la
strada più facile (``lo dico a Stefano, ci pensa Stefano''), e proprio per
questo bisogna mettere dei paletti: il manager interviene quando serve, ma le
cose che si possono gestire da soli vanno gestite da soli. (Quando invece il
confronto diretto fallisce, l'escalation è legittima: si veda la
sezione~\ref{sec:z1-qa-finale}.)

% ---------------------------------------------------------------------
\section{Traiettorie di carriera e valutazione delle performance}
\label{sec:z1-traiettorie}

\subsection{Quale traiettoria vorreste?}

Il relatore mostra un grafico con quattro curve colorate, contrassegnate
anche da lettere (A, B, C, D) per eventuali persone daltoniche (``a me è
capitato che qualcuno scoprisse di esserlo durante una lezione''). Le curve
rappresentano la performance lavorativa nel tempo. Due domande: \emph{quale
vorreste essere? E quale pensate realisticamente di essere?} Un partecipante
aspira a B e pensa di essere C; quasi tutti scelgono B.

Perché non A, che sale più rapidamente? Un partecipante risponde che una
crescita così repentina può essere più dannosa che utile, come una ``rush'' di
studio molto faticosa che poi non si riesce a sostenere. Il relatore propone
una lettura suggerita da un amico psicologo, che spiega la forma del grafico:
di B, C e D \textbf{vediamo dove finiscono}; di A no. Potrebbe continuare a
salire oltre la scala del grafico, o crollare sotto D. Non saperlo ci spinge a
preferire B. (A chi obietta che anche le altre proseguono oltre il bordo, il
relatore precisa: il bordo destro è ``oggi'', e di A non sappiamo come sia
arrivata fin lì.)

\subsection{La carriera non è una linea retta}

Tutti immaginano la carriera come una linea: si parte da X e, attraverso
stage e posizioni, si arriva a Y. In realtà carriera e performance avranno un
andamento irregolare, con alti e bassi. Come non si accetta che un team possa
temporaneamente \emph{underperformare}, tanto meno lo si accetta da se stessi.
Bisogna accettare che ci saranno momenti bassi.

Per questo il relatore \textbf{detesta le valutazioni annuali delle
performance}: prendono un punto casuale della traiettoria, che se si è
fortunati è il picco più alto e se si è sfortunati il più basso, e lo usano
per valutare un intero anno. Insegna ai manager a fare più punti di controllo
e a guardare altre cose:
\begin{enumerate}
  \item \textbf{quanto tempo passa tra un picco e l'altro}, cioè quanto
        \emph{cooldown} serve a una persona. È il motivo per cui, dopo la
        Games Week, tutti staccavano una o due settimane. Il relatore stesso,
        che si impegna moltissimo e tende al burnout, non potrebbe rifare lo
        stesso workshop il lunedì e il martedì successivi con la stessa
        qualità;
  \item \textbf{la differenza tra la miglior performance e la peggiore}: che
        cosa ha portato a lavorare così bene, e così male? Con chi si
        lavorava, su cosa ci si concentrava, quanto tempo e budget c'erano?
        L'obiettivo è capire come \emph{replicare} le condizioni che fanno
        lavorare bene, invece del solito ``hai lavorato male, è colpa tua''.
        Il relatore ha avuto molti manager per cui quando lavorava male era
        colpa sua e quando lavorava bene ``che bravo'', ma sapeva che raramente
        era l'unico fattore della propria performance. È la fallacia del ``se
        sai fare 100, perché non fai sempre 100?'', come nella scena di
        \emph{Schindler's List} sull'operaio che produce una vite in pochi
        secondi: non sempre un ritmo è sostenibile;
  \item \textbf{il fallimento di oggi confrontato con quello di sei mesi fa}.
        Il relatore non è contento della presentazione di oggi (l'avrebbe
        migliorata per altri due anni), ma rispetto a quelle di cinque anni fa
        vede un grande miglioramento di chiarezza: il suo peggior fallimento di
        oggi equivale al suo primo successo di inizio carriera.
\end{enumerate}
Quello che conta è la \textbf{traiettoria} dall'inizio dell'anno a oggi, non
i singoli punti. E se si riesce a capire insieme quando e perché si ha
successo e quando e perché si fallisce, si può lavorare per replicare solo le
situazioni che fanno lavorare bene.

\paragraph{Lo stesso vale per i giochi.} È il motivo principale per cui i
giochi delle grandi case vanno bene o male. Quando si riescono a replicare le
condizioni che hanno prodotto \emph{Warcraft~II}, esce \emph{Warcraft~III};
quando non ci si riesce, esce ``il nuovo \emph{Dragon Age}'': la gente non è
arrabbiata perché il gioco non è bello, ma perché non è come gli altri
\emph{Dragon Age}; non ne replica il \emph{feeling}, anche se potrebbe
essere un gioco favoloso. L'eccezione positiva è il passaggio da \emph{Dune} a
\emph{Dune~II}: il primo era di fatto un'avventura grafica, il secondo ha
generato il genere degli strategici in tempo reale. Gli sviluppatori avevano
capito che il successo non stava nel tipo di gioco, ma nell'ambientazione che
attirava le persone: un bel gioco in quell'ambientazione sarebbe piaciuto
comunque. Il relatore comprò \emph{Dune~II} aspettandosi un seguito
dell'avventura e, pur sorpreso, lo finì.

\subsection{Lying, hiding, faking}

Quando si è valutati su un singolo punto nel tempo, e da quel punto dipendono
bonus, stipendio e il rischio di licenziamento, si comincia a fare
\textbf{lying, hiding e faking}:
\begin{description}
  \item[Lying] dire che è sempre andato tutto bene, che si finirà sempre in
        orario, che non c'è nessun problema;
  \item[Hiding] nascondere: ``non era un mio task'', ``non c'ero'', ``non ho
        visto niente'';
  \item[Faking] fingere, anche in cose apparentemente innocue come ``sì, sì,
        ho capito tutto'' quando si ha solo paura di dire ``me lo rispieghi?
        Non ho capito niente''.
\end{description}

\subsection{``Non ho capito'' e la conferma delle istruzioni}

Dire ``non ho capito, ripetimelo'' è ancora visto come una figura da sciocco.
Il relatore non si arrabbierà mai se qualcuno gli chiede di rispiegare
un'intera slide: le cose ovvie per lui possono non esserlo per altri. Più si
conosce qualcosa, \textbf{più è difficile spiegarla}, perché si danno per
scontate molte cose. Ne parlava proprio con Diego, che gli diceva che fare
l'insegnante diventa sempre più difficile man mano che si diventa esperti. È
come spiegare \emph{EVE Online}, uno dei giochi più complessi in circolazione:
chi lo conosce bene dà per scontato che ``ovviamente non puoi fare questa
cosa'', mentre per un nuovo giocatore non è affatto ovvio. A scuola chi dice
``non ho capito'' viene visto come lo ``scemo''; nessuno si chiede ``non è che
l'ho spiegato male?''.

Il racconto più ``traumatico'' dei ragazzi che fanno stage nelle grandi
aziende è sempre lo stesso: ``Cinzia, devi fare questa cosa.'' Cinzia la fa.
``Non te l'avevo chiesta così.'' Il problema è che nessuno ha
\textbf{confermato che Cinzia avesse capito}. Il consiglio, valido sempre:
\begin{itemize}
  \item quando si \emph{ricevono} istruzioni, ripeterle: ``ho capito bene?
        Vuoi che faccia questa cosa in questo modo?'';
  \item quando si \emph{danno} istruzioni, chiedere a chi le riceve di
        ripeterle, non perché lo si ritenga incapace, ma per capire che cosa
        ha capito dalla propria spiegazione (``quindi vuoi un quadrato?''
        ``No, ho detto un triangolo\dots'' ``Però mi hai detto una forma con
        quattro angoli.'' ``Hai ragione: voglio un quadrato'').
\end{itemize}
Frasi come ``non te lo ripeto due volte'' o ``stammi dietro'' portano solo a
fare le cose come si pensa che vadano fatte e a sentirsi dire a fine progetto
``non la volevo così'': che si tratti di una feature, di un design o di un
modello 3D. (Un partecipante osserva che lo stesso accade insegnando le lingue:
chi insegna dà per scontato, ad esempio, che si sappia dell'esistenza del
genere neutro in tedesco.)

\paragraph{Come convincere il team a chiedere?} Il relatore mette sempre come
premessa che \textbf{l'obiettivo non è trovare la persona più intelligente
nella stanza, ma fare un buon lavoro}. Che Gianni capisca alla prima e un altro
alla quinta non gli interessa, se alla fine entrambi lavorano bene: ripetere
cinque volte è comunque meno lavoro che far rifare le cose sette volte. È un
investimento a lungo termine, e non dissimula alcun fastidio: davvero non gli
dà fastidio. Ma ripetere le stesse parole non aiuta: ``non ho sentito'' si
risolve ripetendo; ``non ho capito'' si risolve cambiando spiegazione. Meglio
ancora: ``dimmi che cosa hai capito tu, e vediamo cosa non torna''. Se
intravedesse nel team uno sguardo del tipo ``è la quinta volta che lo
chiede'', interverrebbe: spesso chi giudica è proprio chi dice ``ho capito'' e
poi sbaglia. Preferisce di gran lunga chi chiede cinque volte.

\subsection{Il feedback più utile che abbiate ricevuto}

\emph{Qual è il feedback più utile che abbiate mai ricevuto?} Il relatore ne ha
almeno dieci che gli hanno cambiato il modo di lavorare. Le risposte:
\begin{itemize}
  \item un semplice \textbf{``grazie''} per un lavoro fatto, più utile di un
        ``questa cosa fa schifo'';
  \item un feedback sulla propria \textbf{impulsività} nel rispondere: fermarsi
        un attimo e capire di cosa si sta parlando prima di intervenire;
  \item \textbf{essere se stessi}, anche al lavoro, e non scusarsi
        continuamente per ogni cosa.
\end{itemize}
I feedback che toccano di più sono quelli \textbf{veri}. Il ``grazie''
significa ``ho visto quello che hai fatto e l'ho apprezzato''. Quelli negativi
che si ricordano sono quelli che hanno fatto realizzare ``questa cosa che
pensavo di saper fare, non la so fare'', oppure ``la so fare, ma questa volta
non mi è venuta'' (non è che se una cosa una volta non riesce non la si sa più
fare: può capitare di sbagliare una porta). Arrivano sempre da un
\textbf{punto cieco}: qualcosa che non si sapeva fosse percepito dagli altri.
Conta più ciò che le persone ti dicono di ciò che pensi che pensino di te; e
ciò che il relatore pensa che gli altri pensino di lui è sempre terribile,
mentre il loro giudizio si rivela molto più positivo del previsto.

\paragraph{Il denominatore comune siete voi.} Si potrebbe pensare che questi
concetti non siano applicabili da chi arriva in un'azienda senza poterne
cambiare la cultura. Ma c'è sempre un denominatore comune a cui applicarli:
\emph{se stessi}. La cultura aziendale può non spingere a dare feedback, può
vedere male chi chiede aiuto; ma il relatore ha visto che quando lui lo fa, le
persone intorno tendono a farlo, e un gruppo che lo fa lavora meglio. Perfino
i CEO più ostici a un certo punto gli chiedevano: ``perché il tuo team ora
performa meglio degli altri?''. La differenza è il default to trust: una
pratica che col tempo fa operare il team \emph{come un team} e non come un
insieme di individui. Un videogioco è un progetto corale: c'è chi fa giochi da
solo, ma non si fa da soli \emph{Red Dead Redemption~2}; e se il gruppo non è
affiatato e allineato non esce \emph{Red Dead Redemption~2}, esce
\emph{Anthem} (``oggi la mia fissazione è \emph{Anthem}, domani sarà
un'altra'').

% ---------------------------------------------------------------------
\section{Come dare un feedback}
\label{sec:z1-dare}

\subsection{Concentrarsi sull'azione e sui suoi effetti}

Il primo principio: \textbf{il feedback non riguarda mai la persona}. Non si
dice ``tu hai fatto questa cosa che non va bene'', ma ``questa cosa non va
bene: come facciamo a non farla più?''. Che sia un comportamento, un'azione,
una decisione o una strategia, ci si concentra su quello, e si
\textbf{esplicita l'effetto} che ha. Nell'esempio del ritardo l'effetto era
``quindici minuti di attesa moltiplicati per venti persone''. In un progetto
può essere una catena:
\begin{quote}
  se non rispetti le deadline, l'animazione che viene dopo resta ferma; se
  loro sono in ritardo, la \emph{vertical slice} è in ritardo; se la vertical
  slice non è pronta per la Gamescom, non andiamo alla Gamescom e non troviamo
  un publisher; se non troviamo un publisher, chiudiamo.
\end{quote}
Più la conseguenza è grave, più dovrebbe sembrare piccolo il gesto di dire
``datemi una mano'': è il \emph{male minore}. Poi, a cose fatte, ci si dà
feedback: cosa ha funzionato, cosa no, cosa fare di più, di meno, cosa smettere
di fare.

\subsection{Il pericolo degli ambienti solo positivi}

Gli ambienti dove ci si fa sempre grandi complimenti e nessuno parla dei
problemi creano difficoltà vere alle aziende. Il problema non è mai la
mancanza di complimenti, ma un \textbf{eccesso di ottimismo}: si comincia a
pensare che le cose andranno come si spera, perché tutti intorno danno solo
feedback positivi.

A Blizzard i feedback determinavano gli \textbf{aumenti annuali}: ognuno
doveva indicare cinque persone che avrebbero dato feedback su di lui, e
ovviamente tutti indicavano gli amici, o colleghi di progetti andati bene. Un
sistema nato con le migliori intenzioni diventava uno \emph{scambio di voti}
(``se mi metti follow ti metto follow''), senza alcun miglioramento.

Al contrario, \textbf{le conversazioni più brutte sono state le più utili}:
affrontare i licenziamenti a Riot, o la cancellazione di \emph{Project Titan}
a Blizzard, il progetto che doveva essere il nuovo MMO dopo \emph{World of
Warcraft} e da cui poi è nato \emph{Overwatch}\footnote{Nella trascrizione si
legge anche ``doveva essere Warcraft 2'', probabilmente un lapsus o un errore
di trascrizione.}. Andando a vedere perché lo si cancellava, c'erano solo
motivazioni razionali e pragmatiche, di business, chiare a tutti, non un ``Mike
Morhaime ci odia''. Sono anche le conversazioni in cui ha lavorato su se
stesso: è molto \textbf{avverso ai conflitti}, e questo a volte gli ha
procurato conseguenze molto più gravi di quanto sarebbe costato un piccolo
disagio. Dieci anni fa faceva fatica ad accettare questo difetto; ha dovuto
subire le conseguenze di conflitti evitati e poi cresciuti.

\subsection{Che cosa fare}

Le linee guida della slide, che il relatore legge rapidamente perché saranno
disponibili nel materiale:
\begin{description}
  \item[Attenzione al linguaggio del corpo] dare un feedback a braccia
        incrociate o guardando altrove è diverso dal guardare l'altro negli
        occhi.
  \item[Sii specifico] con esempi tangibili.
  \item[Esplicita gli effetti] come visto sopra.
  \item[Discuti il prima possibile] non aspettare un mese per parlare di
        qualcosa successo ieri, altrimenti nessuno ricorderà più bene cosa sia
        accaduto.
  \item[Concentrati sull'azione e sul problema, non sulla persona] il nemico è
        l'azione, non la persona che si ha davanti. Non bisogna far cambiare
        l'altro: bisogna evitare che il problema si ripresenti.
\end{description}

\paragraph{Gli insulti sono rivolti al ruolo.} Un esempio solo apparentemente
scollegato. Il relatore ha lavorato su \emph{League of Legends}, ``non
esattamente noto per avere una community da orsetti del cuore'', in varie
\emph{task force} sulla tossicità. A chi smetteva di giocare veniva inviato un
sondaggio sul perché; il 33\% indicava la tossicità della community. Dopo varie
campagne di sensibilizzazione e cambiamenti si è scesi al 30\%: sembra solo
un 3\% in meno, ma valeva circa \textbf{120 milioni di revenue in più
all'anno}, perché quei giocatori restavano esposti a skin e contenuti e
continuavano a comprare. Anche qui, piccoli guadagni che si sommano.

Lavorando con degli psicologi, aveva notato che gli insulti \textbf{non sono
quasi mai rivolti alla persona, ma al ruolo}: ``il mio jungler fa schifo'',
``il mio DPS non è capace'', ``il nostro portiere è una pippa''. Chiamare
qualcuno per ruolo lo \emph{depersonalizza} e rende più facile insultarlo. Lo
stesso nell'esempio del giro dell'ufficio: non ``non chiedere a Giulio'', ma
``non chiedere a quel grafico, consegna sempre in ritardo''. Concentrarsi
sull'azione, senza attaccare la persona, fa sì che il feedback non sia (e non
venga percepito come) un attacco personale.

\paragraph{Considera il carattere.} Qui tornano utili i test di personalità.
Servono più agli studenti che al relatore (che sa già che usciranno 53 profili
diversi, senza un pattern): bisogna conoscere il proprio profilo e
\textbf{condividerlo con chi lavora con noi}, chiedendo anche agli altri di
fare lo stesso. In ufficio il relatore ha usato una versione più avanzata, a
pagamento, con un report di circa trenta pagine: il \textbf{profilo
DISC}\footnote{Nella trascrizione ``disk''. Si tratta del modello DISC
(Dominance, Influence, Steadiness, Conscientiousness).}. Ha mappato tutta
l'azienda e ha fatto stampare il profilo di ciascuno sulla sua scrivania, così
che tutti sapessero con chi stavano parlando. Capire come ragionano gli altri
(``è introverso, non gliene parlo davanti a tutti'') è una delle cose che gli
sono tornate più utili: \emph{adattare ciò che si vuole dire a chi si ha
davanti}.

\paragraph{Assicurati di non disturbare.} Sotto il profilo di ognuno c'era un
piccolo \textbf{semaforo}: verde, ``parlami pure, interrompimi''; giallo,
``interrompimi solo se è urgente''; rosso, ``mi sto concentrando, non mi
parlare''.

\paragraph{Lascia finire il lavoro.} Durante gli eventi dal vivo il relatore si
era reso conto di quanto fosse grave dare feedback mentre qualcuno sta ancora
lavorando. Se una cosa sta andando male, dirlo in quel momento non aiuta: tra
la pressione di finire e quella del feedback, peggiorano entrambe le cose, e
anche il rapporto. Meglio lasciarla finire, anche male, e parlarne dopo.

\subsection{Che cosa non fare}

\begin{description}
  \item[Non basarti su supposizioni] ``mi hanno detto che ieri in un meeting
        hai fatto\dots'': o c'eri e l'hai visto, o non esiste.
  \item[Non giudicare la persona.]
  \item[Non fornire feedback vago.]
  \item[Non accettare e non fornire feedback di terze parti] Se qualcuno vi
        dicesse ``mi hanno detto che hai fatto un brutto lavoro'', il primo
        pensiero non sarebbe ``cosa ho fatto male?'', ma \emph{``chi te l'ha
        detto? Chi è la vipera che ho in seno? Chi è l'impostore?''}. Siamo
        fatti così. Ogni conversazione che inizia con ``mi hanno detto che'' è
        inutile: l'altro non ascolterà, penserà solo a chi dare la caccia, e
        spesso penserà alla persona sbagliata, creando antipatie tra persone
        che non c'entravano nulla. Quando al relatore dicono ``mi hanno detto
        che lui ha fatto\dots'', risponde ``vai a parlarne con lui''; se gli
        dicono ``mi hanno detto che tu hai fatto\dots'', risponde ``mandami la
        persona che te l'ha detto, altrimenti non ne parliamo''. Non per
        cattiveria: se chiedesse ``che cosa ho detto di preciso a Marco?'' e la
        risposta fosse ``non lo so, non c'ero'', di cosa si starebbe parlando?
        \textbf{Parlate con la fonte}, altrimenti fate come se quel feedback non
        esistesse.
  \item[Non dare consigli non richiesti] come visto, non portare la soluzione:
        la soluzione deve venire da chi riceve il feedback.
  \item[``Non mi piace'' non è un feedback] ``Penso che questo design non
        funzioni perché, visto che il nostro obiettivo è\dots'' ha argomenti;
        ``non mi piace il combat system'' è un'opinione. Anche ``io l'avrei
        fatto diversamente'' non è un feedback.
  \item[Accettare non vuol dire fare come proponi] deve essere chiaro a chi dà
        il feedback: l'aspettativa è essere ascoltati e presi in
        considerazione, ma se la responsabilità è dell'altro, deciderà
        l'altro.
  \item[Non offenderti] se il feedback non viene seguito. Restarci male
        significa pensare di saperne più dell'altro, cioè \emph{venire meno al
        default to trust}. ``Decidi tu'' significa ``mi fido di te, perché sei
        lo specialista, e mi fido che mi ascolterai e farai la cosa migliore
        per tutti'', che potrebbe anche non essere ciò che ho detto io. Se
        voglio che l'altro faccia sempre come dico io, mi sto fidando solo di
        me stesso.
\end{description}

\subsection{Il primo feedback e il debriefing}

La cosa più difficile è dare il \textbf{primo} feedback: è un muscolo da
allenare. Chi non l'ha mai fatto (soprattutto con questa struttura) dovrebbe
provarci, anche su minuzie (``abbiamo concordato di usare WhatsApp nel gruppo:
non potremmo usare Telegram?''), anche solo per vedere come va. Altre
indicazioni:
\begin{itemize}
  \item feedback \textbf{solo uno a uno}, mai davanti ad altri, soprattutto se
        negativo;
  \item \textbf{prepararsi} prima;
  \item fare il \textbf{debriefing del feedback}: dopo, fermarsi a pensare a
        com'è andato. Il relatore, quando media un feedback tra due persone,
        si sente spesso rispondere ``benissimo'' da una e ``malissimo''
        dall'altra: la percezione è molto diversa. Non si fa problemi a tornare
        dalla stessa persona: ``ieri ti ho dato un feedback: com'è andata? Sei
        soddisfatto? Come poteva andare meglio?''. Con le persone con cui si ha
        meno a che fare, questo aiuta molto la volta successiva.
\end{itemize}

% ---------------------------------------------------------------------
\section{Come ricevere un feedback}
\label{sec:z1-ricevere}

Per il relatore è la parte più difficile. Si riassume in quattro passi:
\begin{enumerate}
  \item \textbf{Ascolta}: si parte sempre ascoltando.
  \item \textbf{Chiedi dettagli}, se serve, ma per capire meglio, non per
        difenderti.
  \item \textbf{Accetta, a prescindere}: non si sta decidendo chi ha ragione;
        si sta ascoltando un'opinione. Non c'è nulla da cui difendersi.
  \item \textbf{Agisci}: fatti un piano. Può essere banale (``mi hanno detto
        che interrompo sempre: cerco di non interrompere''), oppure ``devo
        migliorare questa cosa: come?''. Pensa anche a \emph{chi} potrebbe
        aiutarti e a \emph{che supporto} ti serve: il piano non deve essere
        per forza ``devo risolvere tutto da solo''.
\end{enumerate}
Una domanda che molti sottovalutano: \emph{``sono junior in un'azienda: sarò
mica il primo ad avere questo problema?''}. Probabilmente qualcun altro l'ha
già avuto: basta chiedere ``tu come hai fatto?''.

\subsection{Il feedback positivo}

Tutto questo vale per i feedback costruttivi negativi, a cui si pensa di
default; ma esistono anche i \textbf{feedback positivi}, e vanno usati. Quando
si riceve un ``grazie'' o un ``bravo'', la domanda è: \emph{come faccio a
ricreare le stesse condizioni?}

Il relatore ha lavorato anche in aziende dove il feedback era solo negativo.
Il rischio? Si immagini l'impiegato ideale: vuole sempre migliorare, è sempre
pronto a dire ``magari sbaglio io''. Se nessuno gli dice cosa sta facendo
\emph{bene}, spinto dal desiderio di migliorare rischia di cambiare cose che
erano già buone, e di farle peggio. Il rinforzo positivo \textbf{non è solo una
gentilezza}: è un atto pratico, che evita che le persone peggiorino ciò che
sanno già fare bene.

\subsection{Il feedback accidentale}

Esiste infine il \textbf{feedback accidentale}, molto frequente negli uffici:
\begin{itemize}
  \item il ``chiedi a lei, non a quel grafico'' del giro dell'ufficio è un
        feedback negativo su una persona, dato accidentalmente, né richiesto
        né voluto;
  \item a Riot li chiamavano \textbf{drive-by}, come nelle scene dei film in
        cui passa un'auto, si abbassa il finestrino, si spara e si riparte. In
        molte aziende con una ``cultura del feedback'' in cui tutti devono
        dire la loro, si riceve una valanga di commenti del tipo ``comunque il
        sito a me non piace, ciao''. Non si sa perché, non si sa cosa fare di
        diverso, non si sa che rischi ci siano: \emph{a che cosa è servita
        questa conversazione?} A niente.
\end{itemize}

% ---------------------------------------------------------------------
\section{Quando, quanto e con chi}
\label{sec:z1-quando}

\subsection{Frequentemente e recentemente}

\emph{Quando avete ricevuto il vostro ultimo feedback?} Per gli studenti,
qualche giorno fa (nel corso) o qualche mese fa (al lavoro, dopo una
consegna). Come aveva osservato Silvia, il feedback andrebbe dato
spesso.

Il relatore usa l'esempio di una nave da crociera (``la Costa Concordia'') che
deve andare da A a B. La rotta ideale (azzurra) è quasi impossibile: nessun
progetto va esattamente come previsto. C'è la rotta ``panoramica'' (verde), che
gira intorno all'isola d'Elba; ci sono poi rotte (arancione e rosa) che a un
certo punto si allontanano molto e poi correggono bruscamente: qualcuno si è
accorto che si stava sbagliando, ``ma non è quello che avevo chiesto'',
``non è dove dovevamo andare''. Sono i famosi \textbf{sanity check} o
\emph{checkpoint}.

L'ideale è avere \textbf{tanti piccoli check per fare tante micro-correzioni}:
correggere la rotta di un paio di gradi è molto più rapido che farlo dopo
essersi allontanati. Bastano trenta secondi di ``ehi, vado bene?''. Spesso non
lo si fa per paura di sembrare insicuri o poco autonomi, ma per il relatore è
l'opposto: quando qualcuno gli chiede ``vado bene così?'' e lui risponde
``sì'', si alleggerisce anche del proprio carico mentale. Un capo che assegna
un compito e lo vede solo alla fine vive sempre col dubbio (``avranno capito?
Lo staranno facendo giusto? Arriverà in tempo?''); con un check di dieci
secondi al giorno (``tutto ok, Stefano'') è tranquillo, e la persona va avanti
sapendo che quanto fatto finora va bene. Non serve fare ogni giorno il giro del
dipartimento: anche ogni uno o due giorni basta a evitare la rotta verde.

Per il feedback è lo stesso: \textbf{va dato subito}. Dire oggi ``ieri hai
fatto questa cosa'' è molto diverso da dire ``negli ultimi sei mesi hai fatto
questa cosa''. Il relatore si è sentito dire più volte, alla valutazione di
dicembre, ``da maggio hai fatto questa cosa che non andava''. La sua risposta:
``e aspetti dicembre per dirmelo? Mi hai lasciato sbagliare per sei mesi?''. È
una mancanza di rispetto, e la conversazione non può più essere costruttiva.
Per questo non ama le valutazioni annuali: i manager crescono abituandosi a
prendere appunti delle cose negative per tirarle fuori alla valutazione
(``a marzo hai fatto questa cosa'' ``sì, però poi ho smesso'' ``sì, però a
marzo l'hai fatta'').

In sintesi:
\begin{description}
  \item[Frequentemente] riguarda l'intervallo tra un check e l'altro.
  \item[Recentemente] significa non lasciar passare un anno prima di dare il
        feedback, altrimenti non serve più a niente.
\end{description}

\subsection{Quanto dura un feedback}

Un feedback dura \textbf{da trenta secondi a quindici minuti}; raramente di
più. Chi non ne ha mai dato uno dovrebbe sforzarsi di darne almeno uno, magari
positivo per cominciare. Ogni feedback dato o ricevuto è un
\textbf{deposito nel proprio ``fondo fiduciario''} (\emph{trust fund}), che
costruisce la reputazione futura. Il relatore si è costruito la reputazione di
chi dà feedback il più possibile oggettivi e analitici: giusti o sbagliati che
siano, costringono a pensare. Non impone mai il suo punto di vista (``si fa
come dico io''), ma chiede ``tu come faresti? Come posso aiutarti?''. Il
feedback, per lui, serve a \emph{far crescere qualcun altro}.

\subsection{I manager peggiori: ``a che numero sto pensando?''}

I peggiori manager che il relatore abbia avuto erano quelli che volevano che si
facesse come volevano loro, anche quando chiedevano ``tu come la faresti?''.
In realtà la domanda era ``\emph{a che numero sto pensando?}'': se dicevi 5,
``no, dovremmo fare 9''; se dicevi 9, ``no, 5''. Lo stress costante di dover
indovinare cosa vuole il capo nasce dalla mancanza di fiducia (``se mi chiedi
come fare, fammelo fare'').

Lo stesso gli succede con i CEO che lo chiamano per ``sistemare la cultura
aziendale'' e poi dicono ``non sono sicuro, io farei diversamente''. La sua
risposta: ``lo so che faresti diversamente, per questo mi hai chiamato''. Se si
chiama un esperto, bisogna lasciargli fare il suo lavoro: non si va dal
meccanico a cambiare le gomme per poi dirgli ``non dovresti girare un po' di
più quel bullone?''. Chiedere un parere a un esperto significa dargli anche il
potere di scegliere. Il problema è il \emph{micromanagement}, o in generale il
non accettare punti di vista diversi dal proprio (magari perché si è il
fondatore e si pensa di essere l'unico a saper fare le cose).

\subsection{Dare feedback a un superiore egocentrico}

\paragraph{Come si dà un feedback negativo a un superiore con un grande ego?}
Il trucco è \textbf{non puntare sull'ego, ma sull'azione e sui suoi effetti}.
Non ``tu fai così e io mi sento così'': anzi, si alimenta il suo ego
rendendolo il \emph{main character} dell'avventura, l'eroe che risolve il
problema (``tu, grande manager, aiutaci a risolvere questo problema che
affligge noi che facciamo l'operatività''). Non è l'ideale (``io cambierei
manager''), ma pochissimi possono scegliersi i colleghi: il relatore stesso ha
quasi sempre ricevuto team preassemblati (``ecco le chiavi, questa è la tua
macchina''). Bisogna giocare con le carte che si hanno in mano.

\paragraph{Un esempio pratico.} Con un manager che fa micromanagement e mette
sempre ``il becco'' nel proprio lavoro, si può dire: ``se ti fidi di me per
questo lavoro, con le tue capacità potresti dedicarti a questa cosa più
strategica e di lungo termine: pensa quanto sarebbe più utile per l'azienda''.
In sostanza: ``guarda cosa potresti raggiungere se mi lasciassi fare il mio
lavoro''. Non come un ``tu fai il tuo e io faccio il mio'', ma come un
``ti sollevo da questa parte del carico, che so che ritieni importante: lascia
che ci pensi io, così puoi dedicarti ad altro'', che magari gli dà anche più
visibilità.

Si gioca su ciò che quella persona pensa di saper fare, e che probabilmente sa
anche fare. Anche Michael Scott di \emph{The Office}, per quanto preso in giro,
è raccontato come un venditore bravissimo, che ha chiuso i contratti più
importanti per l'azienda, e rispetto ai manager che gli succedono è forse
``uno dei meno peggiori''. Tutti, anche i manager peggiori, hanno punti di
forza: \textbf{giocate su quelli}.

% ---------------------------------------------------------------------
\section{Il feedback come regalo}
\label{sec:z1-regalo}

Per chiudere, il relatore risponde alla domanda d'apertura, che aveva lasciato
apposta in sospeso: che cos'è il feedback per lui?

\begin{quote}
  \emph{Il feedback è un regalo che mi viene fatto da qualcuno che so che vuole
  il mio interesse.}
\end{quote}

Ha costruito un team di persone che si fidano di lui e che non hanno motivo di
sabotarlo o di ferirlo: per questo interpreta anche le cose peggiori con il
default to trust. È un regalo, e cerca di prenderlo in considerazione il più
possibile.

Ci è arrivato partendo da se stesso. Oltre a essere avverso ai conflitti, si
definisce ``avverso alla comunicazione umana'': paradossalmente, una persona
introversa e timida finita a fare comunicazione, marketing e ora HR (``succede
spesso''). Prima di ordinare una pizza al telefono ha dieci minuti di
preparazione mentale, e se l'ordine delle domande cambia va in confusione. Per
lui un feedback di dieci minuti sono spesso \textbf{tre ore di preparazione}:
costruisce un \textbf{albero dei dialoghi}, come in un'avventura grafica
personale, con tutte le opzioni. Poi spesso finisce con ``hai fatto questa
cosa, andrebbe migliorata'' ``hai ragione, grazie'', e basta.

Così ha pensato: chissà quanti si fanno gli stessi problemi con lui, e passano
ore a prepararsi per dirgli ``questa cosa l'hai fatta male'', sapendo che
riceveranno comunque quell'occhiata fugace del ``come ti permetti''. Eppure
vengono a dirglielo. Non può far altro che rispondere: ``grazie, hai ragione,
cercherò di fare meglio''. È così che ha metabolizzato il ricevere feedback:
capendo che non c'è nessuna cattiva intenzione e che, se non proprio felice
(``ci ho messo una pietra sopra''), può almeno essere \emph{grato} che ci sia
ancora qualcuno che ha così a cuore il suo interesse da dargli feedback.

\paragraph{Il silenzio fa più paura.} Non bisogna spaventarsi dei feedback
negativi che emergono dai sondaggi aziendali: bisogna spaventarsi del
\textbf{silenzio}, quando nessuno dice più nulla. Significa che nessuno si
fida più che si sappiano risolvere i problemi.

\paragraph{La sindrome dell'impostore.} Il suo motivatore personale, all'inizio,
era soprattutto ``non voglio fallire'': per questo cercava di migliorare, e
migliorare è diventato accettare feedback, perché più feedback riceve, più
migliora. Il feedback è un \textbf{acceleratore}, e lui cerca di incentivare
gli altri a darglielo. Resta però sempre la vocina del ``e se si accorgono che
non sei capace?''. Con la \emph{impostor syndrome} convive da anni e
continuerà a conviverci; se ce l'avete, conviveteci anche voi. Avrete sempre
la sensazione che ``se ne accorgeranno'': forse sì, ma sapete fare anche
altre cose. Non concentratevi solo su ciò che non sapete fare. Gli errori
capitano: la perfezione non è di questo mondo.

\paragraph{Esercitarsi in accademia.} Per fare pratica con il feedback,
l'ambiente migliore è quello dei colleghi di corso: persone nella stessa
situazione, con esperienza simile e lo stesso obiettivo, in un contesto in cui
non si può essere licenziati e l'unico effetto a lungo termine è ``ho parlato
con questa persona, ora le cose vanno meglio''. E bisogna avere il coraggio di
parlare anche con i docenti (``dire a Diego: secondo me questa lezione potevi
farla in questo modo''; ``so che è una persona che lo accetta''), senza farsi
bloccare dalla gerarchia. Nell'azienda giusta (come quella in cui lavora ora) il
relatore va sempre dal CEO a dare feedback, sapendo che non verrà licenziato:
il CEO gli ha detto ``ti ho preso per darmi feedback''.

La frase di chiusura:

\begin{quote}
  \emph{Per superare le vostre difficoltà, tutto quello che vi viene chiesto è
  il coraggio di chiedere aiuto.}
\end{quote}

Chiedere aiuto non è una debolezza: è un segno di consapevolezza di sé
(\emph{self-awareness}) e di forza personale. Se tutti lo facessero, i
progetti andrebbero molto meglio; spesso invece ci si fa frenare dal ``non
voglio fare brutta figura'', ``non voglio passare per quello scemo'', e sono
proprio queste paure a rallentarci e a impedirci di raggiungere gli obiettivi.

% ---------------------------------------------------------------------
\section{Domande e risposte finali}
\label{sec:z1-qa-finale}

L'ultima mezz'ora è dedicata a domande libere.

\paragraph{E se una persona peggiora per uno o due anni di fila?} Se nessuno
gliel'ha mai detto, c'è un grosso problema, non suo ma dell'azienda. Se invece
c'è consapevolezza, il percorso è: ``cosa vuoi fare per risolvere?''; ``ho
provato questo, questo e questo''; ``visto che non ha funzionato, vuoi un mio
consiglio? Proviamo con quello?''; ``che supporto ti serve?''. Se, con tutto il
supporto, il tempo e le soluzioni concordate per mesi, il miglioramento non
arriva, per il relatore la soluzione è: \textbf{``dimmi un prezzo per andare
via''}. Preferisce pagare una buonuscita (\emph{severance package}) che
avviare la strada dei richiami formali sulla performance fino al
licenziamento per giusta causa. A un certo punto andare avanti non è un favore
per nessuno.

\paragraph{È vero che nei videogiochi si viene sfruttati? Meglio un piccolo
studio italiano?} Per il relatore si è quasi più a rischio a Blizzard che in un
piccolo studio di Roma o Torino. Nel piccolo studio si sa dove ci si va a
cacciare: non ci sarà il team building a Santo Domingo, ma ci si concentra su
un gioco alla volta e, se va bene, si fa il successivo. A Blizzard si può
lavorare su \emph{Overwatch}, che vende milioni, e perdere comunque il posto,
in modo totalmente scollegato dal proprio lavoro. Gli stipendi più alti
all'estero compensano anche questo rischio, che spesso non dipende da come va
il gioco (anzi, ``il gioco non ha fatto come volevamo'' è spesso una scusa per
licenziamenti che non avevano motivo).

\paragraph{Cambiare settore a trent'anni, da web developer full remote?} Lo
stipendio indicato dallo studente non è fuori dalla portata di vari studi
italiani, anzi è superabile. Non va esclusa la possibilità di lavorare da
remoto per aziende piccole ma estere. Alcuni amici del relatore, pagati con
stipendi italiani, si sono trasferiti a Bali o in paesi con un costo della vita
più basso; il collega di \emph{PUBG} percepisce uno stipendio olandese vivendo
a Milano, che per quanto cara non costa come i Paesi Bassi.

\paragraph{Come si diventa manager?} Dipende dall'azienda: alcune promuovono i
migliori (processo tendenzialmente sbagliato), altre propongono un percorso di
\emph{people management}. Diventare manager significa \textbf{abbandonare ciò
che si sa fare per imparare a fare altro}. Il relatore ha dovuto accettare che
non sarebbe più stato ``il miglior community manager italiano del mondo''
(``eravamo in pochi, quindi era facile''): dedicandosi alle persone non poteva
più fare il suo lavoro bene come prima. È stata una negoziazione con se
stesso: fare il manager significa non essere più l'esperto di qualcosa, ma
saper gestire e far crescere persone. Significa anche \textbf{cedere il
controllo della propria carriera}: se il team decidesse di non fare più nulla,
il manager ne pagherebbe le conseguenze. Il suo lavoro diventa scegliere le
persone giuste e farle lavorare bene. Lui lo esplicita sempre: sa di affidare
al team parte della propria carriera, e il team sa che lui combatterà per
loro, anche prendendosi le colpe.

\paragraph{L'inglese.} \textbf{Va assolutamente imparato}, anche lavorando in
Italia. Non saperlo preclude circa il 95\% del mercato (tutto quello fuori
dall'Italia) e soprattutto il \emph{networking} con persone che potrebbero dare
occasioni future: l'inglese è spesso l'unica lingua in comune con chi lavora
all'estero ed è oggi la lingua universale dei videogiochi.
\begin{itemize}
  \item \emph{Che livello serve?} Non importa B2 o C2: bisogna saper sostenere
        una conversazione, guardare almeno Netflix con sottotitoli in inglese,
        comunicare senza dodici pause e senza ``come si dice\dots''. Il
        certificato non verrà mai chiesto: verrà chiesto se si sa parlare
        inglese. Il relatore ha sentito persone con un C2 che non sapevano
        parlarlo.
  \item \emph{Come allenarlo senza occasioni di conversazione?} Netflix, Prime
        Video, Now TV, Udemy, Coursera, YouTube, Duolingo, siti di scambio
        linguistico con madrelingua, Discord: non sono mai esistiti così tanti
        mezzi gratuiti. Il relatore ha imparato l'inglese giocando (con una
        pronuncia sbagliata, perché leggeva senza sentire, su cui ha poi
        dovuto lavorare); guardare un film in inglese con i sottotitoli in
        italiano insegna anche la traduzione.
  \item Oltre all'ascolto passivo è fondamentale \textbf{parlare}: in una
        flotta di \emph{EVE Online} con 150 persone quasi tutte madrelingua
        inglese si è esposti a decine di accenti e modi di dire (mille modi di
        dire ``come stai'' oltre a \emph{how are you}); lavorando all'estero si
        impara a distinguere gli accenti (Manchester, Londra) e ad apprezzarli
        anche nei media.
\end{itemize}

\paragraph{Quanto deve sapere un junior?} Dipende da quanto è tecnico il
ruolo. Un game designer junior che non sa qualcosa può essere affiancato da un
senior; un programmatore che non sa programmare, in quel momento, non è utile
all'azienda.

\paragraph{Annunci ``junior con cinque anni di esperienza''.} ``Cercasi
ventenne con trent'anni di esperienza.'' Non è l'annuncio per un junior: è un
errore dell'azienda, che in realtà cerca un \emph{mid}. Riceverà profili
junior, si stupirà che non abbiano cinque anni di esperienza, rifarà l'annuncio
e avrà perso sei mesi. Il consiglio: tenere d'occhio quella pagina, perché
dopo qualche mese l'annuncio verrà corretto.

\paragraph{Una struttura per il post-mortem.} A chi chiede un formato di base,
il relatore anticipa quanto vedrà nella seconda giornata a proposito di
brainstorming: una lavagna \textbf{divisa in quattro aree}, su cui ognuno
attacca tutti i post-it che vuole:
\begin{center}
  \begin{tabular}{@{}p{0.44\textwidth}p{0.44\textwidth}@{}}
    \toprule
    \textbf{Smettere di fare} & \textbf{Iniziare a fare} \\
    ciò che non ha funzionato & ciò che non abbiamo fatto e andava fatto \\
    \midrule
    \textbf{Fare di meno} & \textbf{Fare di più} \\
    ciò che ha fatto perdere tempo rispetto all'impatto ottenuto &
    ciò che ha funzionato ma, fatto con troppa leggerezza, non ha dato i
    frutti sperati \\
    \bottomrule
  \end{tabular}
\end{center}
Con i post-it ognuno ha voce e tutto risulta catalogato; per il report basta
una foto della lavagna (o uno screenshot, se la lavagna è digitale, come Miro).

\paragraph{Come sbloccare la paralisi del feedback verso il management?} Uno
studente ha visto manager dire ``non abbiate paura, la mia porta è sempre
aperta'' senza che cambiasse nulla. C'è una dissonanza tra ciò che dicono e
ciò che fanno: è ``la mamma che dice vieni qua che non ti faccio niente''.
Qualcosa deve aver fatto pensare alle persone che non è sicuro. Prima va
sistemato quello; poi si può provare con i \textbf{feedback anonimi}.
L'obiettivo del relatore resta il feedback aperto e firmato, ma a Stormind ha
iniziato con feedback anonimi verso il management per verificare se
l'anonimato fosse un fattore. Non lo era (evidentemente non c'erano problemi a
parlare con il CEO), e la volta successiva intende abbandonarli. Ma è un buon
\emph{check}: se senza anonimato non arriva nulla e con l'anonimato arrivano
cento feedback, il problema è chiaro.

\paragraph{Quando è giusto rivolgersi a un superiore?} La strada giusta resta
parlare direttamente con l'interessato. Ma se si prova e non si trova una
soluzione, la cosa più ovvia è andare dal proprio manager insieme: ``ci
aiuti?''. Quando il relatore parlava di mediare i dialoghi, era proprio questo:
fare da moderatore, aiutare a sbloccare i punti fermi o, se non si riesce a
scegliere tra A e B, scegliere lui. L'escalation, quando necessaria, va fatta.

\paragraph{Il tempo libero veniva rispettato?} No, e se l'è dovuto imporre
lui. Il lavoro gli piaceva così tanto che continuare a casa non gli sembrava
lavorare. Gli si è ritorto contro con il lancio di \emph{Diablo~III}, fatto in
orario americano (le due o l'una di notte in Europa), con l'ufficio chiuso
perché a nessuno era stato chiesto di restare. Il lancio va malissimo, i
giocatori protestano, e lui, sveglio, risponde sui forum fino alle tre, quattro,
cinque del mattino: ``non posso risolvere il problema, ma resto qui ad
aggiornarvi con quello che riesco a sapere dal team di sviluppo''. Il giorno
dopo trova sulla scrivania un \textbf{richiamo formale dell'HR}: aver lavorato
fuori orario senza autorizzazione, condotta non professionale. Per lui è stato
un interruttore: ``ora me ne vado''. Due settimane dopo faceva colloqui con
Riot. Ancora oggi è contrario a quel richiamo, ma gli ha fatto capire come
ragiona un'azienda: ``hai un contratto 9--18, non mi interessa se l'hai fatto
con le migliori intenzioni''. Eppure per anni ha goduto di un'ottima
reputazione in quelle community come ``quello che è rimasto con noi fino alle
cinque del mattino'', e anche l'azienda ne ha beneficiato.

Ci sono poi aziende in cui, soprattutto per i developer, viene richiesto il
famigerato \textbf{crunch}: lavoro straordinario continuativo, senza pause.
Dipende molto dall'azienda; sta venendo (``grazie al cielo'') deprecato, perché
molte aziende capiscono che non è sano ed è visto male anche dal pubblico, ma
esiste ancora.

\paragraph{E se il tuo capo demotiva qualcuno del tuo team?} Qui ``diventa
personale'': il team del relatore si fida di lui perché sa che in questi casi
``diventa una bestia''. ``Te la puoi prendere con me quanto vuoi, ma il mio
team non lo tocchi.'' Lo gestisce direttamente con l'altro manager, per quanto
brutta possa diventare la conversazione. Lo sforzo più grande è \textbf{non
andarne poi a parlare con il team} (``ragazzi, ho fatto questa cosa''), perché
minerebbe la fiducia verso l'alto: non vuole un rapporto ottimo fino a lui e
pessimo oltre. L'obiettivo è che l'altro manager ``si faccia perdonare''
cambiando il proprio comportamento, e che il team percepisca da solo il
cambiamento.

\paragraph{Chiusura della giornata.} Il relatore chiede un feedback collettivo
in chat (``se domani ci dai \emph{more of the same} siamo contenti, oppure
ditemi cosa limare''), si scusa per i ripetuti crash della connessione e
ringrazia per la proattività nelle domande, che gli servono anche come
\emph{check}: se nessuno chiede nulla su un argomento, significa che può
andare avanti. Racconta di aver passato una settimana a decidere che cosa
mettere nelle slide e di averle poi preparate in poco tempo (rifacendone
alcune fino alle tre di notte). Ricorda infine di usare il test di
personalità per sé e per il proprio gruppo, facendolo fare a chi non l'ha
ancora fatto: tornerà molto utile per cooperare meglio.

% ---------------------------------------------------------------------
\section{Anticipazioni sulla seconda giornata}
\label{sec:z1-domani}

La seconda giornata del workshop sarà più pratica: ``oggi vi ho spiegato
perché, domani vi insegno come lo fanno queste aziende''. Gli argomenti
previsti:
\begin{description}
  \item[Time management] come organizzarsi una giornata professionale in
        azienda: come strutturare una \emph{task list}, cosa fare e cosa non
        fare, come e cosa prioritizzare.
  \item[Brainstorming] i gruppi dovranno fare moltissimi meeting (come facciamo
        il gioco, come lo chiamiamo, che scopo gli diamo). Se non si sanno fare
        bene, i brainstorming sono \textbf{dannosi ed eterni}: un brainstorming
        dovrebbe durare \textbf{al massimo due ore}. Verranno mostrati metodi
        che accorciano la durata, massimizzano l'effetto e, soprattutto,
        \textbf{fanno partecipare tutti}: non c'è niente di peggio di un
        brainstorming in cui parlano in due per il 90\% del tempo. Il classico
        ``ragazzi, non esistono idee sbagliate, sparate'' è il modo peggiore di
        farlo; non è colpa di nessuno, ma si può fare diversamente.
  \item[Guidare da dentro] come essere leader nel proprio team anche senza
        esserne il leader: sviluppare responsabilità personale e verso gli
        altri. Non essere il manager di qualcuno non significa non esserne
        responsabili: la persona accanto è responsabilità di tutti, perché si
        è un team e non un insieme di individui. ``Se qualcuno cade, tutti lo
        facciamo rialzare: non c'è l'addetto al rialzamento.'' È la versione
        \emph{advanced} della giornata: come comportarsi da manager senza
        averne il titolo.
\end{description}
Torneranno inoltre il \emph{service level agreement} e il tema della
condivisione dei successi e della protezione del team dai fallimenti.

% ---------------------------------------------------------------------
\section*{Punti chiave del capitolo}
\addcontentsline{toc}{section}{Punti chiave del capitolo}

\begin{itemize}
  \item ``La capacità di un gruppo di persone di fare cose straordinarie
        dipende dal modo in cui queste persone operano insieme.'' Il 90\% dei
        problemi dei team sono \textbf{problemi di comunicazione}.
  \item Le aziende si liberano sempre più presto dello \textbf{useless
        genius}: le competenze si trasmettono, il carattere no. Il
        \textbf{fit culturale} va valutato anche dal candidato.
  \item Non esiste un unico modo di dare feedback: dipende dal contesto e
        soprattutto da \textbf{chi lo riceve}. Il feedback è una somma di
        piccoli miglioramenti periodici; tenersi le cose nel cassetto non
        serve mai.
  \item I team attraversano \textbf{forming, storming, norming, performing e
        adjourning}. Lo storming non si salta; si passa la maggior parte del
        tempo in forming; la bravura sta nell'accelerare i passaggi e nel
        fare bene il \textbf{post-mortem}, anche dei successi.
  \item Gli obiettivi di un progetto vanno fissati \textbf{all'inizio}: solo
        così si può valutare se è andato bene e scegliere tra A e B.
  \item Il circolo virtuoso è \textbf{feedback $\rightarrow$ miglioramento
        $\rightarrow$ fiducia}. Accettare un feedback non significa fare come
        dice l'altro: decide chi ha la responsabilità, ma deve argomentare.
  \item ``Aiutami a poterti aiutare'': meglio dire in anticipo ``potrei non
        farcela'' che ``non ce l'ho fatta''. Chiedere aiuto non è una
        debolezza.
  \item Quando si dà feedback \textbf{non si porta la soluzione}: la si
        costruisce insieme a chi lo riceve.
  \item La cultura è forte quando si lavora \textbf{per gli altri}; il lavoro
        è valutato come opera d'insieme (anche con bonus di team).
  \item Come in \emph{Among Us}, i team sprecano energie a cercare il
        traditore invece di fare le task. \textbf{Default to trust}: chi ha
        sbagliato ha agito con le informazioni che aveva. Cercare
        \textbf{la causa, non il colpevole}: ``che cosa ti ha fatto pensare
        che fosse la scelta giusta?''.
  \item Carriere e performance non sono lineari: vanno valutate le
        \textbf{traiettorie}, non un punto dell'anno. Le valutazioni puntuali
        generano \emph{lying, hiding, faking}.
  \item Confermare sempre le istruzioni facendosele ripetere; ``non ho
        capito'' non è una colpa.
  \item Dare feedback: sull'\textbf{azione} e sui suoi \textbf{effetti}, specifico,
        presto, uno a uno, adattato al carattere, senza terze parti, senza
        ``non mi piace'', senza offendersi se non viene seguito; poi fare il
        debriefing.
  \item Ricevere feedback: \textbf{ascolta, chiedi dettagli, accetta, agisci}.
        Il feedback positivo serve a non peggiorare ciò che si fa bene; il
        feedback accidentale e il \emph{drive-by} non servono a nulla.
  \item Feedback \textbf{frequenti e recenti}: tanti piccoli \emph{sanity
        check} costano meno di una grande correzione di rotta.
  \item Il feedback è \textbf{un regalo} di chi ha a cuore il nostro
        interesse; fa paura il silenzio, non le critiche.
\end{itemize}

%    % ... altre lezioni zero ...
%    \part{Il corso di Game Design}
%    \renewcommand{\thechapter}{\arabic{chapter}}
%    \setcounter{chapter}{0}
%    % =====================================================================
%  Capitolo 1 — Introduzione al corso e al mestiere del game designer
%  Fonte: trascrizione della lezione GD1 (Game Design, primo anno)
%  Pacchetti richiesti nel preambolo del libro:
%    \usepackage[italian]{babel}, \usepackage[T1]{fontenc},
%    \usepackage[utf8]{inputenc} (solo pdfLaTeX), \usepackage{booktabs}
% =====================================================================

\chapter[Il mestiere del game designer]{Il mestiere del game designer: introduzione al corso}
\label{cap:introduzione-game-design}

Questo primo capitolo fa da soglia. Prima di affrontare metodi, framework e
documenti, è necessario capire \emph{in che modo} si impara il game design,
\emph{con chi} lo si fa, \emph{con quali strumenti} e, soprattutto,
\emph{che cosa fa concretamente} un game designer nel suo lavoro quotidiano.
Il capitolo segue l'impostazione della lezione inaugurale del corso: si parte
dal contesto formativo (il corso come simulazione di lavoro), si passa al
metodo di apprendimento (l'\emph{inquiry-based learning}), si elencano
letture e risorse di riferimento, e si arriva infine a una prima definizione
operativa del ruolo del designer, dei limiti entro cui lavora e delle sue
possibili specializzazioni.

Una premessa terminologica: il game design è una disciplina giovane e quasi
tutta la sua letteratura è scritta in inglese. Per questo motivo molti termini
tecnici verranno lasciati in lingua originale (\emph{pipeline},
\emph{pitching}, \emph{tuning}, \emph{feature team}, \dots): sono i termini che
si useranno nel lavoro, ed è importante imparare a riconoscerli e a usarli
correttamente fin da subito.

% ---------------------------------------------------------------------
\section{Il contesto: un corso come simulazione di lavoro}
\label{sec:simulazione-lavoro}

\subsection{Il principio fondamentale}

Il primo concetto da fissare è che il corso \textbf{si presenta come una
simulazione del lavoro}. Dal momento in cui si entra in aula si viene trattati
come \emph{potenziali professionisti dell'industria}: non soltanto studenti,
ma lavoratori effettivi di uno studio (l'accademia), con una gerarchia, dei
referenti, dei compiti e delle scadenze.

Ne derivano alcune conseguenze pratiche:
\begin{itemize}
  \item il docente non assegna un ``compitino'' da risolvere per accumulare
        punti o presenze: ogni \emph{task} è pensato per far dedicare tempo
        all'apprendimento di qualcosa di concreto;
  \item non consegnare un task non danneggia il docente (che ``vive ancora
        domani e viene pagato ugualmente''), ma lo studente, che perde
        l'occasione di mettere le mani su un argomento e diventare più abile;
  \item docenti e studenti devono trattarsi \textbf{da colleghi}, perché è
        molto probabile che un domani lo diventino davvero.
\end{itemize}

\subsection{Perché l'impegno conta più che mai}

Il lavoro del game designer è descritto come \emph{molto romantico, ma anche
molto duro}, e sempre più duro col passare degli anni. Durante il periodo del
Covid l'industria ha visto grandi investimenti; quei ``rubinetti'' sono stati
chiusi da tempo, il rapporto con i publisher è diventato più complesso e le
opportunità lavorative non sono più quelle di tre o quattro anni fa. Questo
significa:
\begin{itemize}
  \item più competizione per ogni posizione;
  \item aspettative più alte: ci si aspetta che i candidati siano ancora più
        preparati e sicuri delle proprie competenze;
  \item un unico modo per diventarlo: \textbf{lavorare}. Più tempo si dedica
        oggi ad apprendere il mestiere, più facilmente si sarà assumibili
        domani e più agevole sarà il lavoro di tutti i giorni.
\end{itemize}

\subsection{Sbagliare qui per non sbagliare fuori}

Il mantra del corso è: \textbf{quanto più provate qui, anche in modo goffo e
sbagliando, tanto meno sbaglierete là fuori}. In accademia, per quanto la
simulazione sia realistica, nessuno toglie lo stipendio e nessuno licenzia da
un momento all'altro. Nell'industria sì, e un errore può costare molto denaro
all'azienda, ancora di più se l'azienda è la propria.

\subsection{Il valore del fallimento}

Il fallimento, se affrontato con la mentalità corretta, è probabilmente
\emph{la cosa più vantaggiosa} che si possa fare durante la formazione:
\begin{itemize}
  \item un successo (``avete superato questa deadline'') conferma le proprie
        abilità, ma insegna relativamente poco;
  \item dopo $n$ fallimenti si impara moltissimo, e i successi arrivano;
  \item è fondamentale non demoralizzarsi: il feedback negativo non è mai un
        giudizio sulla persona, ma su ciò che è stato consegnato.
\end{itemize}

Il docente riconosce apertamente che dedicare tempo ed energia a un progetto e
vederlo rifiutato, o vederlo non uscire mai, è demotivante: è successo a lui e
lo ha visto accadere anche su progetti molto grandi, con team numerosi. Il
punto però non è ciò che si è ottenuto materialmente, ma \textbf{quanto si è
cresciuti}. Da ogni esperienza ci si portano a casa due cose:
\begin{enumerate}
  \item \textbf{l'esperienza}: quanto si è diventati bravi, veloci,
        praticamente abili nell'applicare un framework, quanto si è cresciuti
        come designer;
  \item \textbf{il rapporto umano}: nel lavoro si ``suda sangue'' fianco a
        fianco con altre persone, e questo crea legami molto forti che
        portano avanti l'industria nel senso più positivo del termine: nascono
        progetti privati, ci si porta da un'azienda all'altra, si formano team
        con basi solide.
\end{enumerate}

Va evitato anche l'eccesso opposto: la cultura del \emph{grinding} si è
appropriata di questi discorsi in modo distorto. Non bisogna diventare robot.

% ---------------------------------------------------------------------
\section{Organizzazione del corso e gerarchia}
\label{sec:gerarchia-corso}

\subsection{La struttura gerarchica}

La struttura dei referenti del corso ricalca volutamente una
\textbf{gerarchia tipica dell'industria}:
\begin{itemize}
  \item al vertice c'è il \textbf{direttore del corso} (Diego Ricciuti), che
        corrisponde al \emph{creative director}: supervisiona la visione del
        progetto (il concetto di ``visione'' verrà approfondito più avanti);
  \item sotto di lui ci sono i \textbf{referenti di dipartimento}, uno per
        disciplina, con un corrispettivo per ciascuna sede (Milano e Roma);
  \item in fondo ci sono gli studenti, che occupano il primo gradino di
        un'azienda: l'\emph{intern} (stagista) o il \emph{junior}, il passo
        che precede il lavoratore ``regolare''.
\end{itemize}

\begin{table}[htbp]
  \centering
  \begin{tabular}{@{}p{0.34\textwidth}p{0.29\textwidth}p{0.29\textwidth}@{}}
    \toprule
    \textbf{Disciplina} & \textbf{Milano} & \textbf{Roma} \\
    \midrule
    Direzione (\emph{creative director}) & \multicolumn{2}{p{0.58\textwidth}}{Diego Ricciuti} \\
    Game design      & Mattia Podini        & Giacomo ``Jack'' Danna \\
    Level design     & Stefano Zocchi       & Andrea Dispoli \\
    Narrative design & \multicolumn{2}{p{0.58\textwidth}}{Daniele Fusetto (lezioni condivise), Giacomo Masi} \\
    Programmazione (Unreal) & Alessandro Pregnolato & Matteo Velardi \\
    \bottomrule
  \end{tabular}
  \caption{Referenti del corso per disciplina e sede (nomi come riportati in lezione).}
  \label{tab:referenti}
\end{table}

Per il primo anno, il referente diretto degli studenti di game design è il
docente di game design: è a lui che ci si rivolge per dubbi e domande. Altre
figure utili sono il direttore didattico (severo, ``giustamente severo'', da
contattare solo per questioni particolarmente importanti), l'amministrazione
(contratti e documentazione) e la segreteria.

\subsection{I canali di comunicazione}

\begin{description}
  \item[Discord] è il canale per le comunicazioni ufficiali del corso e per
        le domande che non richiedono urgenza immediata; può avere qualche
        ritardo.
  \item[E-mail] è il canale per tutto ciò che è formale o ``ancora più
        ufficiale'' e per le comunicazioni rapide e immediate.
  \item[Telegram] esistono residui di vecchi canali precedenti a Discord, da
        evitare.
\end{description}

Le domande che valgono per tutti vanno fatte in classe (anche da un portavoce
a nome del gruppo); quelle singole e personali via Discord o, se formali, via
e-mail. Se il docente non risponde a una domanda c'è un motivo, ad esempio
ragioni legali, e lo dirà esplicitamente.

\subsection{I tre moduli e la settimana tipo}

Le tre discipline principali del primo anno (e non solo) sono:
\begin{itemize}
  \item \textbf{Game Design} (tendenzialmente il lunedì);
  \item \textbf{Level Design} (tendenzialmente il mercoledì);
  \item \textbf{Programmazione} (tendenzialmente il venerdì).
\end{itemize}
Nel primo mese le lezioni sono tutte di game design (lunedì, mercoledì e
venerdì). Sono previste anche lezioni di narrativa, ma il piano formativo si
aggiorna ogni anno: le ore di narrativa sono state ridotte a favore di quelle
di programmazione. Le poche lezioni di narrativa daranno quindi dei
\emph{cenni}, non la preparazione di un narrative designer. Ciò che serve
davvero è imparare \textbf{come comunicano gli altri dipartimenti} e come
lavorare con persone che fanno un lavoro diverso ma complementare al proprio.

\subsection{Le due anime delle lezioni di game design}

\begin{description}
  \item[Design methods (metodologie)] lezioni prevalentemente teoriche, ma con
        una buona componente pratica, in cui si studiano metodologie di design
        e i cosiddetti \emph{framework} dell'industria.
  \item[Documentazione e pitching] il cuore del lavoro quotidiano del
        designer.
\end{description}

\paragraph{La documentazione.}
Il docente si definisce ``il più grande nerd della documentazione'' che si
possa incontrare. È una parte dura, perché a pochi piace documentare, ma il
lavoro del designer \emph{è principalmente questo}: scrivere documenti e, da
quei documenti, passare alla pratica e all'implementazione. Chi non sa
documentare bene parte male. Quando non si è presenti, il documento è il
proprio \textbf{biglietto da visita}: qualcun altro lo leggerà, quindi va
scritto nel modo più chiaro e comprensibile possibile.

Un esempio concreto del perché la documentazione conta: CD Projekt Red ha
sviluppato un proprio engine proprietario, ma buona parte della sua
documentazione è stata ``buttata all'acqua''. Chi deve usarlo o lo sa già
usare, o ha competenze di un altro engine da trasferire, oppure è in
difficoltà. Si documenta \textbf{perché altri possano leggere}.

\paragraph{Il pitching.}
Il \emph{pitching} è l'esposizione, la presentazione orale. La parte creativa
del team presenta spesso:
\begin{itemize}
  \item le proprie idee ai colleghi;
  \item il gioco a convention ed eventi;
  \item il pitch del gioco agli investitori, per ottenere i fondi necessari a
        realizzarlo.
\end{itemize}
Saper parlare e comunicare chiaramente le idee è fondamentale. In sintesi:
\emph{il game designer parla}, e bisogna farsene una ragione.

\subsection{Altre componenti del corso}

\begin{description}
  \item[Gioco critico e playtest] la classe viene divisa in due gruppi: uno
        gioca, l'altro osserva. Chi osserva non applaude ``come le
        scimmiette'': annota, segna e analizza il \emph{playtest}, esattamente
        come si fa nell'industria. Il playtest fa parte delle mansioni del
        designer, ma esiste anche come specializzazione autonoma.
  \item[Pipeline] termine che indica tutto ciò che riguarda i processi
        formalizzati dell'industria; di solito si intende la \emph{pipeline di
        produzione}: come si divide il team, qual è la sua gerarchia, come il
        \emph{producer} gestisce le deadline. In breve: \emph{come} e
        \emph{quando} si fanno le cose. Al primo anno è meno rilevante (si
        dovrà ``nuotare'' un po' da soli); dal secondo anno si vedranno le
        metodologie effettivamente applicate ai progetti digitali. Si
        scoprirà quanto sia difficile lavorare a un progetto digitale, e
        ancor più fisico, senza un piano di sviluppo comprovato, come quelli
        usati per qualsiasi software in qualsiasi parte del mondo.
  \item[Engine] il motore di riferimento è \textbf{Unreal Engine}
        (sezione~\ref{sec:engine}).
  \item[Programmazione] principalmente tramite \emph{Blueprint} in Unreal,
        più eventuali sistemi di versionamento; il corso parte da zero.
  \item[Progetti pratici] probabilmente la cosa più importante. Al primo
        anno i progetti di gruppo restano circoscritti al corso di design,
        perché mancano ancora alcune nozioni (come comunicare, affrontare
        problemi, dare e ricevere feedback). In seguito si allargheranno a
        progetti digitali, molto più difficili di quelli fisici, con membri
        di altri team: un rapporto non unilaterale, ma ``covalente''.
\end{description}

% ---------------------------------------------------------------------
\section{Il metodo: l'inquiry-based learning}
\label{sec:inquiry}

\subsection{Che cos'è}

Il corso è costruito secondo l'\textbf{inquiry-based learning}
(apprendimento basato sull'indagine), un approccio pedagogico nato intorno
agli anni Sessanta, che si richiama a pensatori come Dewey e Piaget. In
generale, le discipline da cui il design attinge sono quasi sempre la
sociologia, la psicologia e, talvolta, l'economia.

Il funzionamento è un ciclo:
\begin{enumerate}
  \item il docente fornisce metodologie e framework dell'industria durante
        la lezione;
  \item lo studente, fuori dalle ore di lezione, lavora sui \emph{task}:
        mette a terra le nozioni, le maneggia, le assorbe meglio;
  \item da questo lavoro nascono nuovi dubbi;
  \item i dubbi tornano in aula: lo studente chiede, il docente risponde
        (ammesso che dare la risposta sia utile, e ammesso che una risposta
        esista).
\end{enumerate}

\subsection{Non sempre esiste una risposta giusta}

Non tutto ciò che si studia ha ``la risposta migliore possibile in assoluto
nel vuoto''. Esistono metodologie, limiti e framework che aiutano il lavoro,
ma non sono fissi: cambiano da azienda ad azienda, da lavoro a lavoro, da
progetto a progetto, e a volte vengono stravolti. Sta alla
\textbf{sensibilità del designer} capire quale strumento o framework usare in
un dato momento, e questa sensibilità si impara solo con l'esperienza. Nel
lavoro nessuno dirà ``usa questa cosa'' o ``è meglio quest'altra'': sarà il
designer a scegliere.

Anche l'esperienza del docente ha confini: sei-sette anni di lavoro, un modo
personale di approcciare il design e un certo insieme di progetti (ad esempio
nessun MMO). Oltre quei confini si cercano risposte nella ricerca, studiando
come aziende con esperienza specifica (ad esempio Riot) hanno affrontato
problemi simili e hanno condiviso le loro soluzioni, sempre valutando quali
metodi siano adatti alla dimensione del team, al tipo di progetto e al tempo
disponibile.

\subsection{Sensibilità, non memoria}

Non interessa imparare a memoria un framework o la teoria di un autore. Certe
terminologie vanno conosciute, ma per ragioni di comunicazione e chiarezza.
Ciò che conta è sviluppare la \textbf{capacità di farsi domande}, di farle
agli altri, di fare ricerca e di trovare una propria risposta. È esattamente
ciò che si fa nel lavoro: davanti a qualcosa che non si sa fare, l'abilità
del designer sta nel cercare come qualcun altro ha risolto un problema simile.

\subsection{Le lezioni come corsi di aggiornamento}

Le lezioni frontali simulano i \emph{corsi di aggiornamento} aziendali: le
sessioni che un lead di design tiene al suo team di design (o il lead di
level design al suo team, e così via). Alla fine dell'aggiornamento c'è quasi
sempre un lavoro da fare, un obiettivo da portare a terra:
\textbf{l'obiettivo della settimana}. Si lavora di settimana in settimana, e
il perché verrà chiarito in seguito.

Struttura tipica di una lezione:
\begin{enumerate}
  \item breve sessione di domande e risposte sulla lezione precedente;
  \item analisi di eventuali esercizi particolarmente critici;
  \item lezione vera e propria;
  \item discussione di gruppo, se c'è tempo;
  \item assegnazione del task della settimana.
\end{enumerate}
Il docente tende a ``parlare e straparlare'', ad aprire parentesi e a saltare
da un argomento all'altro: l'unico modo per fermarlo è che siano gli studenti
a interromperlo con domande.

\subsection{La gerarchia della competenza}

L'inquiry-based learning può essere rappresentato come una scala di
competenza:
\begin{enumerate}
  \item la sola lezione frontale porta a \textbf{conoscenze incomplete},
        perché manca la parte pratica e di ricerca;
  \item il lavoro a casa (task, ricerca) porta a una competenza \textbf{più
        completa}, non sbagliata in sé, ma ancora parziale e complementare ad
        altre;
  \item la \textbf{condivisione} dei task e dei feedback con gli altri porta
        idealmente alla competenza più completa possibile;
  \item il picco della competenza è la sua \textbf{applicazione nei
        progetti} (fisici o digitali, di gruppo): ``con le cose che avete
        imparato e ricercato, riuscite a fare un gioco?''. Quando la risposta
        diventa positiva, o comunque si inizia a navigare nei processi di
        design per realizzarne uno, quella è la vittoria.
\end{enumerate}
Non si parte subito creando giochi: si parte dall'analisi, dal design e dalle
regole generali di design, prima ``nel vuoto'', poi contestualizzando i
framework, infine unendo i puntini nei progetti.

Rimarranno comunque delle lacune, per tutti, docente compreso, perché il game
design \textbf{non è una scienza esatta}, anzi, a rigore non è affatto una
scienza. Prende in prestito metodologie dalle discipline scientifiche, ma
contiene una componente imprevedibile e incalcolabile: \textbf{la
produzione}. Si possono applicare perfettamente le regole di un framework, e
quel framework può comunque non essere adatto al gioco, al proprio modo di
sviluppare o al tempo di produzione disponibile (alcuni framework richiedono
una ricerca e un'attenzione che il tempo non concede).

C'è poi la componente creativa, la parte ``romantica'' del lavoro: esiste, ma
è \textbf{metà del lavoro}. L'altra metà è ricerca, metodo scientifico al
100\%. Dagli stessi dati $n$ designer diversi, con sensibilità diverse,
possono ottenere $n$ risposte diverse. Non tutte saranno valide: vanno
applicate, messe a terra, e questa scrematura fa parte del lavoro.

\subsection{I livelli dell'apprendimento}

Un secondo modello, a gradini, descrive i livelli di apprendimento (è la
cosiddetta tassonomia SOLO):
\begin{description}
  \item[Pre-strutturale] il docente dà il compito e lo studente lo risolve,
        possibilmente il prima possibile. È esattamente ciò che \emph{non} si
        vuole.
  \item[Uni-strutturale] si affronta una sola disciplina: è il primo passo
        necessario, ed è il motivo per cui nel primo mese ci sono solo lezioni
        di game design.
  \item[Multi-strutturale] entrano in gioco altri dipartimenti: level design,
        game design, programmazione e così via.
  \item[Relazionale] si capisce che alcune cose valgono per più
        dipartimenti: alcuni framework di narrativa derivano dal design o
        possono essere usati per fare design; le regole di documentazione si
        applicano anche ad altri dipartimenti; le regole di level design si
        mescolano con quelle di game design. Dopotutto, per fare un gioco
        servono molti dipartimenti.
  \item[Astratto esteso] il punto finale: muoversi autonomamente e capire
        cosa usare e quando, sulla base della propria esperienza. L'unico modo
        per arrivarci è fare pratica e fare progetti.
\end{description}

% ---------------------------------------------------------------------
\section{Task, feedback e regole del corso}
\label{sec:task-feedback}

\subsection{I task come opportunità}

I task non sono strettamente obbligatori (esiste un numero minimo di
consegne), ma sono \textbf{opportunità}. Quanto e come ci si lavora dipende
dallo studente. Un task fatto in cinque minuti, o in cui si è stati ``furbi'',
non verrà letto né riceverà feedback.

Domande da porsi davanti a ogni task:
\begin{enumerate}
  \item \textbf{Perché} mi è stato assegnato? Qual è l'obiettivo?
  \item Come ci lavoro per ottenere un risultato \textbf{che abbia senso}? Se
        la ricerca e i dati indicano una direzione, è inutile andare in
        un'altra.
\end{enumerate}

\subsection{Quanto lavorare: un calcolo indicativo}

Se il corso è una simulazione di lavoro, vale il carico di un lavoro normale:
\begin{itemize}
  \item un lavoro a tempo pieno comprende da 30 a 45 ore settimanali; si
        considera una media di \textbf{40 ore};
  \item si sottraggono circa 9 ore di lezione: restano $\approx 31$, diciamo
        \textbf{30 ore};
  \item con circa 3 task a settimana: $30 / 3 = $ \textbf{circa 10 ore per
        task}.
\end{itemize}
Il numero può variare, ma rende l'idea: dedicare un'ora a un task e poi fare
altro non fa guadagnare nulla a nessuno. Lo studente ``butta soldi'', il
docente risparmia tempo leggendo il task in un minuto (o non leggendolo).
Dedicando invece attenzione, ricerca, cura nella formattazione e
sperimentazione, e facendolo con cognizione di causa, si impara.

\subsection{Niente voti, solo feedback}

Ad eccezione del corso di programmazione, non ci sono voti (né ``ottimo'', né
10, né 100). Ci sono \textbf{feedback}:
\begin{itemize}
  \item in game design, generalmente sulla documentazione;
  \item in level design, sui progetti pratici.
\end{itemize}
Il docente indica cosa va bene e cosa va modificato, concentrandosi in
particolare su ciò che va cambiato. Il feedback \textbf{non è mai un giudizio
sulla persona}, ma su ciò che è stato consegnato nel concreto.

\subsection{La condivisione dei task}

Tutti i task consegnati vengono \textbf{condivisi} con la classe, e possono
essere letti solo dopo aver ricevuto i feedback del docente. Si cresce così
non solo con i propri feedback, ma anche con quelli dati agli altri. Regola
importante: ciò che il docente dice a uno studente \textbf{vale per tutti}. Se
a una persona viene detto ``questa cosa non va mai più fatta'', da quel
momento vale per l'intera classe.

\subsection{Non copiare}

Copiare i task degli altri è una delle cose peggiori che si possano fare, non
per una questione morale, ma perché annulla il senso del percorso. Si può
dedicare a ogni task il tempo e l'attenzione che si ritengono opportuni, ed è
normale che col tempo ci si specializzi (preferendo, ad esempio, il level
design o la programmazione tecnica) e si investa meno dove si sente meno
necessità. In questi casi l'alternativa corretta è:
\begin{itemize}
  \item non consegnare;
  \item studiare ciò che hanno fatto gli altri;
  \item eventualmente riconsegnare in seguito.
\end{itemize}
Ci saranno momenti in cui si verrà sommersi dalle cose da fare: bisogna
imparare a scegliere.

\subsection{Il ``non leggo''}

Il docente ha piena facoltà di non dare alcun feedback a un task (il
cosiddetto \textbf{``non leggo''}) se:
\begin{itemize}
  \item non si è lavorato con attenzione;
  \item non si rispettano le indicazioni date;
  \item è stata violata una regola estremamente importante.
\end{itemize}
Un task ``non letto'' equivale a un task non consegnato e non entra nel
conteggio. Attenzione: un task \emph{perfetto} e un task che \emph{fallisce}
nel comprendere o applicare un framework sono \textbf{ugualmente validi},
purché ci si sia impegnati e si sia dedicato tempo, pazienza e ``sudore''.

\paragraph{Riconsegna.} Un task non letto, in ritardo o che si vuole
migliorare può essere riconsegnato e verrà conteggiato, a patto che:
\begin{itemize}
  \item i problemi segnalati siano stati risolti;
  \item sia stato fatto ulteriore lavoro, integrando anche tutto ciò che nel
        frattempo è stato detto agli altri colleghi nei loro feedback.
\end{itemize}
Il task va sviluppato con le competenze che si hanno nel momento della
riconsegna, non con quelle del momento della prima consegna.

\subsection{Requisiti per superare l'anno}

\begin{itemize}
  \item consegnare almeno l'\textbf{80\% dei task} assegnati (soglia
        malleabile);
  \item partecipare almeno all'\textbf{80\% delle lezioni}, in presenza o
        da remoto, preferibilmente in presenza; i ritardi vengono valutati
        caso per caso;
  \item comportarsi da persone corrette e rispettose (è stata firmata una
        carta dello studente): non urlare, non alzare le mani, non
        vandalizzare, non compiere atti illegali.
\end{itemize}
Impegnarsi di più porta inoltre a essere selezionati per varie opportunità.

\subsection{Scadenze}

\begin{itemize}
  \item la scadenza di un task è \textbf{sempre prima dell'inizio della
        lezione} successiva indicata;
  \item se la classe, come team, mantiene comportamenti scorretti, il docente
        può \textbf{anticipare le deadline} per tutti.
\end{itemize}

\subsection{Accessibilità e DSA}

Gli studenti con disturbi specifici dell'apprendimento o neurodivergenze
(dislessia, ADHD, autismo, ecc.) sono invitati a segnalarlo. Quando si
affronteranno le regole di formattazione e documentazione (errori
grammaticali, mancanze, incoerenze) se ne terrà conto e si potranno trovare
dei compromessi, pur ricordando che nel mondo del lavoro nessuno farà
``sconti''. La documentazione può essere inviata all'amministrazione o alla
segreteria, che la comunicherà a tutti i docenti, ma il consiglio è parlarne
anche direttamente con ogni singolo docente.

% ---------------------------------------------------------------------
\section{Le competenze trasversali}
\label{sec:soft-skill}

\subsection{I giochi si fanno in gruppo}

Si possono fare giochi da soli, ed è un'impresa ammirevole; ma lavorando da
soli \textbf{tutto ciò che si guadagna in libertà si perde in produttività}.

\paragraph{Il caso Stardew Valley.} \emph{Stardew Valley} è stato realizzato
da una sola persona, Eric Barone, in anni di sviluppo e di sacrifici: faceva
un altro lavoro per mantenersi (in un cinema) e di notte programmava, faceva
l'arte e componeva la musica. Il risultato è stato straordinario e lo ha reso
ricco, ma si tratta di un caso estremamente raro. Quanti giochi simili, nati
nelle stesse condizioni, sono falliti? Molti di più.

\paragraph{Il rischio d'impresa.} Lavorare da soli non è sbagliato di per sé
e ha aspetti molto positivi: si ha il privilegio di creare esattamente ciò che
si vuole. Ma ci si fa carico anche di tutto il \textbf{rischio d'impresa}: se
si lavora a un gioco per cinque anni e poi il gioco va male, resta il
vantaggio di averlo pubblicato, ma il dopo è sperare di vivere di quella
release o di finanziare con essa il gioco successivo. Le grandi aziende hanno
gli stessi problemi, ma li distribuiscono su più persone e più responsabilità.

L'obiettivo di un professionista è \textbf{saper lavorare in diverse
condizioni}, e la più comune è quella del dipendente (o del freelance /
libero professionista con partita IVA): si viene assunti per le proprie
capacità da un'azienda, piccola o grande che sia, e si lavora con altre
persone.

\subsection{Essere persone con cui si vuole lavorare}

Se le persone non vogliono lavorare con voi, quello è \textbf{il primo motivo
di licenziamento}. E se non si viene licenziati, il clima nell'azienda
diventa terribile: la produttività crolla, la motivazione crolla. Le
\emph{soft skill}, cioè l'insieme di abilità che rendono una persona
piacevole con cui lavorare, vanno coltivate e diventano sempre più importanti
col passare degli anni (relativamente meno ora, molto di più quando i team
saranno più grandi). Anche chi è un ``lupo solitario'' deve imparare a:
\begin{itemize}
  \item lavorare con altre persone;
  \item comunicare con altre persone;
  \item trovare un punto in comune da cui partire per fare le cose.
\end{itemize}

\subsection{Indipendenza e proattività}

Sono le altre due parole chiave del corso:
\begin{description}
  \item[Indipendenza] spesso si verrà lasciati da soli, non per cattiveria ma
        per sfida: quando si cerca una risposta da soli, ciò che conta non è
        tanto la risposta quanto \textbf{il processo}. Una risposta data dal
        docente non vale niente.
  \item[Proattività] chiedere, domandare, agire da soli senza aspettare
        direttive, perché spesso quelle direttive non ci sono, a maggior
        ragione nei team piccoli.
\end{description}
Tutto questo si concretizza nel parlare, nel presentare e nel fare
\emph{post-mortem} (l'analisi a posteriori di un progetto concluso, che verrà
approfondita più avanti), e in generale nell'attenzione alla crescita
personale.

\subsection{Fare domande}

Il corso insiste sulle domande: alzare la mano, intervenire anche ``a
bomba'', fermare il docente ogni volta che qualcosa non è chiaro. Se una
persona non conosce qualcosa, molto probabilmente c'è almeno un'altra persona
che non la conosce. Il docente può dare per scontata la conoscenza di termini
che non lo sono, ed è del tutto normale, soprattutto all'inizio.

% ---------------------------------------------------------------------
\section{Strumenti: engine, Unreal e materiali}
\label{sec:engine}

\subsection{Unreal Engine e gli altri motori}

Il motore di riferimento del corso, usato in level design, è \textbf{Unreal
Engine}: ``un gran mattone'', ma che dà a un designer molte possibilità. Il
vantaggio principale è che permette a un designer non pratico di
programmazione di \textbf{prototipare tramite i Blueprint} (il sistema di
scripting visuale di Unreal). Questo non impedisce di passare a Unity o ad
altri engine, e viceversa: il docente stesso ha studiato e lavorato a lungo
con Unreal, ma la sua ultima esperienza professionale (come level designer) è
stata su Unity, un engine che non conosceva, a dimostrazione che il salto da
un engine all'altro è possibile.

Panorama dei motori citati:
\begin{itemize}
  \item \textbf{Unreal Engine} e \textbf{Unity}: i due più diffusi, con
        modelli di sorgente ``più o meno'' aperto;
  \item \textbf{Godot}: engine open source;
  \item \textbf{CryEngine}: proprietario di Crytek;
  \item \textbf{Snowdrop}: engine di Ubisoft;
  \item Rockstar usa un proprio engine proprietario;
  \item alcune aziende usano versioni \textbf{personalizzate} di engine
        esistenti: ad esempio Milestone, nei giochi di corse, usa un Unreal
        customizzato;
  \item CD Projekt Red ha un engine proprietario (con i problemi di
        documentazione già citati).
\end{itemize}
In sintesi: ogni azienda può avere un engine proprietario oppure una versione
modificata di un engine commerciale, ma si parte sempre da una base.

\subsection{Programmazione: cosa ci si aspetta da un designer}

Non è richiesto alcun bagaglio pregresso di programmazione: il corso parte da
zero. Gli studenti sono stati selezionati (il corso è a numero chiuso) per
competenze analitiche e critiche sopra la media. Anche nel mercato del lavoro,
per una posizione di game designer generalmente \textbf{non ci si aspetta
competenza di scripting} (C++, Java, C\#, ecc.), a meno di voler fare il
\emph{technical designer}, e anche lì varia da azienda ad azienda.

\subsection{Versione e preparazione}

Indicazioni pratiche emerse in aula:
\begin{itemize}
  \item scaricare Unreal Engine; la versione consigliata per il corso è
        stata indicata nella \textbf{5.6} (con conferma successiva). La
        stabilità dipende da cosa si deve fare, da quali tool e plugin si
        usano: i problemi nascono spesso dai plugin. Nel passaggio da Unreal 4
        a Unreal 5, ad esempio, il sistema \emph{Landscape} era inizialmente
        inutilizzabile, mentre in Unreal 4 funzionava molto bene;
  \item è consigliato installare tramite \textbf{Epic Games Launcher}, più
        comodo: permette anche di riscattare giochi gratuiti periodici e di
        accedere ai pacchetti di \textbf{Fab} (il marketplace di asset
        acquisito da Epic), che ogni mese offre alcuni asset gratuiti utili
        per esercitarsi nel level design (ad esempio per fare \emph{speed
        level design});
  \item molti pacchetti di UI o di sistemi di interazione sono utili per
        prototipare al volo, ma spesso inutilmente complicati;
  \item un plugin di \textbf{materiali di prototipazione a griglia} (adattivi,
        con metriche, allineabili alla griglia, con materiali già pronti e
        colori modificabili, dal costo di circa 5 euro) è indicato come
        ``salvavita'' per il level designer; se necessario, i pacchetti
        indispensabili vengono forniti;
  \item portare \textbf{sempre il portatile} alle lezioni di level design
        (da novembre, a partire dall'interfaccia dell'engine); alle lezioni
        di game design basta che lo porti almeno metà della classe, per i
        lavori a coppie. Serve comunque qualcosa per prendere appunti
        (``tavolette d'argilla e cannucce, se volete fare i babilonesi'',
        oppure carta e penna).
\end{itemize}
Chi ha un portatile poco potente dovrebbe comunque provare a installare
l'engine e aprire un template per vedere come gira: all'inizio non si faranno
cose pesanti. Nelle prime lezioni di level design è importante mettere subito
le mani sull'interfaccia: ascoltare dove si trova uno strumento (ad esempio lo
\emph{snapping}) senza provarlo porta a dimenticarlo.

% ---------------------------------------------------------------------
\section{Letture e risorse di riferimento}
\label{sec:risorse}

\subsection{I libri}

\begin{description}
  \item[\emph{The Art of Game Design: A Book of Lenses} (Jesse Schell)] il
        testo \textbf{davvero importante}. Alcune lezioni finali dell'anno si
        baseranno sulle sue terminologie. È utile perché apre la mente sulla
        possibilità di applicare metodologie diverse a giochi diversi, ma
        soprattutto per la \textbf{terminologia} comune. Esiste anche in
        audiolibro (circa 25 ore).
  \item[\emph{Challenges for Game Designers} (Brenda Romero e Ian Schreiber)]
        libro \emph{entry level} con moltissimi esercizi, molto caro al
        direttore del corso.
  \item[\emph{21st Century Game Design} (Chris Bateman)] utile come
        retrospettiva: è un libro visionario che ha previsto metodi e
        soprattutto \emph{archetipizzazioni} dei giocatori tuttora valide nel
        design e nel marketing. Servirà per una delle lezioni più importanti
        dell'anno. È utile leggerlo per vedere quanto aveva previsto, ma è
        utile ``fino a un certo punto''.
  \item[\emph{An Architectural Approach to Level Design} (Christopher
        Totten)] non strettamente necessario, ma di lettura semplice. Ha due
        anime: parte dall'architettura (babilonese, greca, ecc.) e la trasla
        su esempi di level design (come \emph{Doom} e \emph{Team Fortress}).
        Un tempo ne esisteva un riassunto su Gamasutra (oggi
        \emph{Game Developer}).
  \item[\emph{A Theory of Fun for Game Design} (Raph Koster)] utile fino a un
        certo punto: contiene diverse idee superate, che lo stesso Koster ha
        affrontato in vari talk reperibili nel GDC Vault. Verrà ripreso nelle
        prime lezioni, poi messo da parte.
\end{description}

\subsection{Canali e conferenze}

\begin{description}
  \item[GDC (Game Developers Conference)] il suo canale raccoglie talk e
        panel estremamente importanti: è ``uno strumento misterioso che ci
        servirà più tardi''. \textbf{Attenzione}: per molti talk mancano ancora
        le competenze per interpretarli, e alcuni potrebbero fuorviare o
        mettere in testa idee strane. La GDC si concentra spesso su
        produttività e mobile.
  \item[Il \emph{Dragon Speech} di Chris Crawford] il talk del fondatore della
        GDC: appartiene a un passato lontano, ma è molto motivante. Da vedere
        subito.
  \item[Game Maker's Toolkit] canale particolarmente utile per nozioni
        pratiche, con video brevi e alcune interviste.
  \item[Nordic Game] altra conferenza importante, con talk e panel di esperti
        del settore.
\end{description}

\subsection{Associazioni e premi: come leggerli}

Le associazioni di settore (ad esempio IIDEA in Italia, o la IGDA) non sono
il male, anzi: sono un canale importante per affacciarsi all'estero o per
accedere a fondi. Va però ricordato che \textbf{tutelano i propri associati}:
non esiste solo il loro palinsesto, e bisogna valutare quanto spazio occupano
nel panorama della visibilità.

Lo stesso vale per i \textbf{premi} (Game Awards, ``Best Mobile Game'',
``Games for Change'', ecc.): al designer \emph{non interessano}. Che un gioco
abbia vinto ``più di 90 premi'' non dice nulla di utile. Servono a dare
visibilità e lustro ai giochi ed è giusto che esistano, perché ampliano il
panorama e il numero di giocatori. Ma ciò che interessa al designer sono i
\textbf{dati}:
\begin{itemize}
  \item come sono state implementate certe cose;
  \item come è stato fatto il marketing e dove il gioco è stato presentato;
  \item non quante wishlist ha fatto, ma \emph{come} le ha fatte;
  \item in che periodo storico è uscito e quali erano i competitor in quel
        momento.
\end{itemize}

% ---------------------------------------------------------------------
\section{Che cosa fa un game designer}
\label{sec:cosa-fa}

Per arrivare a una prima definizione, la classe ha condotto un brainstorming:
ognuno doveva indicare una cosa che, secondo lui, fa il game designer (senza
ripetere quelle già dette). Ne è emersa la lista seguente, commentata punto
per punto.

\subsection{Progettare meccaniche e struttura del gioco}

Il ruolo più generico è quello del \textbf{core mechanic designer}, detto
anche \emph{gameplay designer} e, un tempo, \emph{black box designer}: si
occupa delle \textbf{meccaniche di gioco}. Si va dal piccolo al grande (in
funzione anche della dimensione dell'azienda):
\begin{itemize}
  \item nel \textbf{piccolo} ci sono le singole meccaniche;
  \item nel \textbf{grande} ci sono i \textbf{sistemi}. Nei giochi
        \emph{sistemici}, con molti sistemi interconnessi (ad esempio
        \emph{GTA}), esistono specializzazioni dedicate a bilanciare e
        supervisionare i sistemi: è il \emph{system design}.
\end{itemize}

\paragraph{Quando un gioco è ``sistemico''?}
Non dipende solo da come le meccaniche interagiscono con il giocatore (le
\emph{dinamiche}, secondo i framework che passano dalle meccaniche alle
dinamiche, che verranno studiati più avanti), ma anche da \emph{quante}
meccaniche ci sono e da quanto interagiscono tra loro. Non diremmo che
\emph{Super Mario} sia un gioco sistemico, eppure contiene dei sistemi: la
linea è sottile. Quando i sistemi aumentano molto, sono complessi e si
influenzano a vicenda, il gioco diventa sistemico. Esempio: in \emph{The
Legend of Zelda: Breath of the Wild} il vento modifica quanto il personaggio
sale o scende; il fuoco può incendiare l'erba; l'erba prende fuoco se colpita
da una freccia infuocata; e così via. Nel mercato moderno i giochi sistemici
sono spesso, ma non sempre, legati a mappe \emph{open world}.

\subsection{Le 3C: camera, character, control}

La base del gioco, da cui generalmente si inizia a lavorare (anche nei
prototipi), è la ``sacra trinità videoludica'' delle \textbf{3C}:
\begin{description}
  \item[Camera] che tipo di inquadratura ha il gioco: prima persona, terza
        persona, isometrica, ecc.
  \item[Character] quanto è grande il personaggio, che tipo di personaggio è
        e, di conseguenza, quali azioni può compiere.
  \item[Control] i controlli: che cosa il personaggio può fare, quali sono i
        \emph{verbi} di gioco.
\end{description}
Le 3C sono definite dal designer, ed esiste anche una specializzazione per
questo lavoro. Esistono varianti: quella del direttore del corso è
\textbf{4C + J}, cioè \emph{camera, character, control, conflict} (conflitto)
e \emph{juice} (il ``succo'', il feedback che rende piacevole l'interazione).
Verranno approfondite più avanti.

\subsection{Il personaggio}

Il designer definisce il personaggio dal punto di vista delle meccaniche
(dimensioni, azioni, capacità), all'interno delle 3C. L'aspetto estetico è
invece competenza di altri dipartimenti.

\subsection{Documentare le informazioni per gli altri dipartimenti}

Assolutamente vero: la documentazione è una delle prime attività del
designer. Più l'azienda è grande, più questo lavoro diventa prominente; più è
piccola, meno è rilevante. In \textbf{pre-produzione}, però, è sempre
fondamentale.

\subsection{Mediatore e leader ``al servizio''}

Il designer è un \textbf{mediatore} nei conflitti interni al team, una sorta
di collante che si relaziona con dipartimenti diversi (non sempre gli stessi,
a seconda della fase di produzione). A volte fa anche da leader, a seconda di
come sono gestite la pipeline e la gerarchia.

Questa leadership è difficile da definire perché \textbf{non è verticale}: è
una leadership \emph{al servizio} (\emph{servant leadership}). Il designer
``è asservito al programmatore'': il suo compito non è arrivare con un'idea e
scaricarla addosso agli altri, né portare la soluzione e lasciare il problema
a qualcun altro, ma \textbf{seguire il processo a basso livello}, sporcandosi
le mani.

\paragraph{I feature team.} Nei team medio-piccoli o medio-grandi possono
esistere sotto-team (\emph{feature team}) dedicati a un obiettivo specifico:
ad esempio un designer incaricato del character o della camera, affiancato da
uno o due programmatori e da un artista. In questi casi il designer
supervisiona il lavoro e gestisce la produzione del sotto-team.

\subsection{Gestire l'avanzamento del prodotto}

Esiste una specializzazione che si occupa solo di questo, il
\textbf{producer}, a un livello gerarchico completamente diverso (e che
potrebbe essere una strada percorribile al termine del percorso di studi).
All'interno di un feature team anche il designer può farlo, ma non è detto.
Può certamente \emph{dare la direzione}: nella sua documentazione può
indicare che per arrivare a uno stato giocabile del prototipo serve un certo
lavoro di arte e di programmazione, per lo stato successivo un altro, e così
via. In ambito produttivo esistono metodologie strutturate come
\textbf{Agile}, che verranno viste in seguito.

\subsection{Presentare le idee}

Il pitching (sezione~\ref{sec:gerarchia-corso}) è richiesto al designer sia
in pre-produzione sia in produzione, in misura che dipende dall'azienda. Ad
esempio Ubisoft lascia ai designer, e in realtà a qualsiasi membro, la
possibilità di \emph{pitchare} un'idea prima di una nuova produzione, perché
possa iniziare a prendere vita.

\subsection{Bozze e prototipi}

Le bozze vanno bene, ma devono essere \textbf{chiare}. C'è un'eterna lotta tra
documentazione e gioco giocabile: se si fa solo documentazione il gioco non
avanza; se si fa solo il gioco manca la documentazione e si perde la visione.
Sono due cose \textbf{complementari}.

Il modo per comunicare un'idea in forma grezza ma giocabile è il
\textbf{prototipo}. Un prototipo può essere qualsiasi cosa: di carta, in
Unreal con Blueprint ``brutti''. \emph{Non deve essere bello, deve essere
giocabile}.

\paragraph{Esempio: il prototipo di carta di Mario + Rabbids.} Durante lo
sviluppo di \emph{Mario + Rabbids}, i programmatori di Ubisoft erano impegnati
su altri progetti (come \emph{Ghost Recon}). I designer realizzarono allora
un \textbf{gioco da tavolo} che riproduceva le meccaniche del gioco, quelle
di uno strategico a turni sul modello di \emph{XCOM}. E funzionava.
L'obiettivo era \textbf{validare un'idea}, validare un gameplay.

Altri modi per rendere chiara un'idea sono i \emph{mock-up} nella
documentazione e le immagini modificate o i \emph{mash-up} di qualsiasi
natura. Lato documentazione l'importante è che sia chiaro; lato prototipo,
che sia giocabile.

\subsection{Fare ricerca per l'idea, e poi smontarla}

Il designer ha una componente creativa importante e fa ricerca per arrivare a
un'idea. Ma ciò che deve poi fare necessariamente è \textbf{smontarla}. Il
suo lavoro è scientifico, esattamente come il \textbf{metodo scientifico}, il
cui obiettivo è \emph{confutare un'ipotesi}: si arriva a un'ipotesi tramite
sperimentazione e poi si cerca $n$ volte di confutarla; se resiste, la si
assume come valida \emph{in un certo contesto} (le regole della fisica
classica, ad esempio, non valgono nella meccanica quantistica).

Allo stesso modo il designer arriva a un'idea (tramite ricerca o chiedendo
agli altri: le idee le hanno tutti, non sono un privilegio del designer),
cerca di smontarla e poi cerca di \textbf{metterla a terra}. Se l'idea non è
applicabile, non è implementabile, finché non si trova un modo per
realizzarla va scartata. Può essere l'idea più bella dell'universo, ma
\textbf{se non può essere messa a terra, non esiste}.

\paragraph{L'originalità.} Serve un'idea originale? Dipende dall'obiettivo e
dalla visione del progetto:
\begin{itemize}
  \item se l'obiettivo è vendere molto, un'idea originale non serve
        necessariamente;
  \item se si vuole creare un nuovo filone, serve un'idea, se non originale,
        almeno interessante: può essere un discorso di marketing o una
        meccanica innovativa, con il rischio di non vendere.
\end{itemize}
Il ruolo del \emph{director} e il concetto di \emph{visione} verranno
approfonditi in seguito.

\subsection{Analizzare i competitor}

Arrivati a un'idea (``vogliamo un certo tipo di combattimento''), si va a
vedere se qualcuno l'ha già fatta e \emph{come}. Le meccaniche non sono
privilegio di nessuno, salvo quando vengono brevettate, come il
\textbf{Nemesis System} di \emph{La Terra di Mezzo: L'Ombra di Mordor}.

\paragraph{Esempio: il combattimento da ridimensionare.} In un progetto
(presso Stormind Games) il team si trovò con un sistema di combattimento
troppo costoso: mancavano le risorse e gli animatori per realizzarlo. Poco
output lato gameplay significava non poter fare qualcosa di complesso e
profondo come in \emph{Elden Ring}; al contrario, il team aveva molti
artisti e poteva puntare su qualcosa di più artistico. Si iniziò così a
orbitare intorno al corpo a corpo di \emph{The Last of Us} (parte I e II).
Sul canale della GDC esistono diversi talk che raccontano cosa è stato fatto
e come, con successi, insuccessi e costi (ad esempio del \emph{motion
capture}): da quei talk si ricavarono moltissime informazioni utili.

I competitor possono essere:
\begin{itemize}
  \item \textbf{diretti}: un gioco molto simile al proprio;
  \item \textbf{di mercato}: un gioco che domina il mercato in cui ci si vuole
        affacciare;
  \item \textbf{indiretti}: un gioco con un sistema simile al proprio, ma
        tutto il resto diverso.
\end{itemize}
Regola d'oro: \textbf{là fuori, nascosto nell'ombra, c'è sicuramente un gioco
molto simile al vostro}. Va solo trovato. E se per qualche ragione non ha
venduto, bisogna capire perché.

\subsection{Lavorare in engine e fare tuning}

Il level designer lavora in engine, così come il technical designer; ma
qualsiasi game designer, a seconda di ciò che deve fare, può finire in
engine, soprattutto per il \textbf{bilanciamento} e il \textbf{tuning}.

Il \emph{tuning} è l'azione di bilanciare con precisione i dati. Esempio: in
\emph{Super Mario} si modifica quanto in alto salta il personaggio, con che
velocità, in quanto tempo raggiunge l'apice, e così via. Questo lavoro rende
il gioco:
\begin{itemize}
  \item ad alto livello, dotato della \textbf{progressione} giusta;
  \item a basso livello, \textbf{piacevole da giocare}.
\end{itemize}
Il \emph{fine tuning} è un lavoro di dettaglio, paziente e minuzioso.

\subsection{Bilanciare}

Il bilanciamento, emerso tardi nel brainstorming, è una delle attività
centrali del designer e ha anche una specializzazione dedicata
(l'\emph{economy designer}, sezione~\ref{sec:tipi-designer}).

\subsection{Analizzare il mercato}

In pre-produzione, di solito, sì: capire le tendenze, i generi che cambiano
nel tempo, la compatibilità del gioco con il momento storico, in vista ad
esempio di un possibile publisher. Va capito quanto, come e con quanto tempo
farlo, perché si è sempre limitati da molte cose.

\subsection{Sintesi}

\begin{quote}
Il designer parte dall'idea ma non si ferma lì: la confuta, lavora sui
dettagli, bilancia, fa \emph{fine tuning}; lavora con il team e fa da
\emph{manager} per diversi dipartimenti; e naviga costantemente tra i limiti.
\end{quote}

% ---------------------------------------------------------------------
\section{I limiti del design}
\label{sec:limiti}

I limiti ci sono sempre, e nel mondo del design sono spesso numerosi. I più
importanti entro cui navigare sono:
\begin{description}
  \item[Tempo] ad alto livello, l'idea deve essere realizzabile nel tempo di
        produzione disponibile; a basso livello, il proprio lavoro
        settimanale deve essere fattibile nel tempo che si ha. Non ci si può
        dividere tra mille feature team se realisticamente se ne può seguire
        uno solo.
  \item[Tecnologia] si può o non si può fare qualcosa? Ci sono i server per un
        multiplayer con un certo numero di giocatori? Ci si può permettere,
        ad esempio, una fauna che reagisce alla fisica come quella di
        \emph{GTA VI}, che richiede il meglio dei \emph{technical artist}?
  \item[Background e competenze] che cosa sa fare il team? Realisticamente,
        chi ha lavorato dieci anni su giochi horror viene assunto per un
        horror, perché ha le competenze per risolvere problemi già visti.
  \item[Persone e motivazione] il team ha la motivazione e le competenze per
        fare un certo tipo di gioco, e per lavorare insieme?
  \item[Denaro] tutto ruota intorno ai soldi. Che si lavori da soli o in mille,
        le persone vanno pagate e il gioco deve avere un ritorno economico.
        Bisogna decidere dove tracciare la linea: raggiungere il
        \textbf{break-even} (guadagnare esattamente quanto speso) o puntare a
        grandi guadagni. Con un \textbf{publisher}, inoltre, una parte
        consistente dei ricavi va a lui (ripartizioni dell'ordine di 80/20 o
        70/30, che dipendono dal contratto firmato).
\end{description}

\subsection{Non siete ``gli uomini delle idee''}

Per qualcuno può essere un problema scoprirlo, ma il designer non è
``l'uomo (o la donna) delle idee'': accanto alla parte creativa c'è una parte
tecnica, e si è fortemente limitati da molti fattori. Questi limiti, però,
sono un aiuto:
\begin{itemize}
  \item riducono il ventaglio di possibilità potenzialmente distruttive per il
        team (cose che non si riescono a fare, che non piacciono o per cui
        mancano le risorse);
  \item aiutano a mirare bene, trovando soluzioni creative in uno spazio più
        piccolo e più attuabile.
\end{itemize}

\begin{quote}
\centering
\textbf{Le idee sono nulla, l'implementazione è tutto.}
\end{quote}

% ---------------------------------------------------------------------
\section{Ruoli e specializzazioni}
\label{sec:tipi-designer}

\subsection{La scala gerarchica}

\begin{description}
  \item[Creative Director (CD)] mantiene la \textbf{visione} e dà l'idea, in
        modo che tutto segua una certa direzione. In teoria può diventarlo
        chiunque, ma con un'esperienza notevole, nell'ordine dei decenni: si
        assume che in quel tempo abbia visto come lavorano molti dipartimenti.
        Naturalmente un CD con un background artistico ha meno competenze di
        design, e viceversa.
  \item[Director di disciplina] uno per dipartimento: \emph{game director},
        \emph{level director}, \emph{art director}, ecc.
  \item[Lead] il capo del dipartimento, che gestisce i gruppi di persone. Chi
        lavora in un dipartimento risponde al proprio lead (ad esempio un game
        designer risponde al \emph{lead designer}, non al game director).
  \item[Senior / Principal] i profili più esperti (a seconda dell'azienda).
  \item[Junior] il profilo di inizio carriera.
  \item[Intern] lo stagista.
\end{description}
Lo schema si replica uno a uno per ogni dipartimento: arte, suono, animazione
(a volte arte e animazione insieme, o \emph{rigging} e animazione insieme),
programmazione, UX/UI e così via. A fare la differenza, a parità di ruolo, è
il \textbf{livello di seniority}.

\subsection{UX e UI}

UX e UI stanno a metà tra game design e arte: a seconda dell'azienda il ruolo
si chiama \emph{UX/UI Artist} o \emph{UX/UI Designer}.
\begin{description}
  \item[UX (User Experience)] come una persona interagisce con le cose. Esempio:
        su uno smartphone, un pulsante accende e spegne lo schermo; per
        digitare la password si scorre verso l'alto e compaiono i punti da
        toccare.
  \item[UI (User Interface)] come sono ``vestite'' queste cose: ad esempio
        l'aspetto della barra dei punti vita (HP), realizzata con asset 2D o
        3D.
\end{description}

\subsection{Le specializzazioni del game design}

Più l'azienda è grande, più la specializzazione è puntuale.
\begin{description}
  \item[Gameplay / Core mechanics designer] (sono la stessa cosa) si occupa
        delle meccaniche e delle 3C.
  \item[System designer] si occupa dei sistemi e delle loro interazioni.
  \item[Economy designer] si occupa del bilanciamento; lavora soprattutto con
        fogli di calcolo (Excel, Google Sheets).
  \item[Combat designer] si occupa del combattimento; negli annunci di lavoro
        indica quasi sempre un combattimento \emph{action}, con gameplay
        reattivo (non strategico), anche se non è una regola.
  \item[Multiplayer designer] si occupa delle interazioni multigiocatore.
  \item[Technical designer] fa da ponte tra game design e programmazione.
\end{description}

\subsection{Le specializzazioni del level design}

\begin{description}
  \item[Encounter designer] posiziona i nemici nella mappa in modo che gli
        scontri siano bilanciati e la progressione equilibrata come
        desiderato.
  \item[World / map designer] popola il mondo di gioco.
  \item[Quest designer] popola il mondo con le missioni e pensa alla
        progressione al loro interno, cioè a dove il giocatore deve andare (si
        vedano ad esempio gli annunci di lavoro di CD Projekt Red o di
        Larian).
  \item[Technical level designer] fa da ponte tra level design e
        programmazione.
\end{description}

\subsection{Discipline vicine: suono, UI, produzione}

Il monte ore del corso non permette di uscire come \emph{narrative designer}
(o \emph{sound designer}) completi. Ciò non toglie che si possano coltivare
queste discipline:
\begin{itemize}
  \item nei progetti di gruppo si può occuparsi anche del suono, della UI/UX,
        o provare ad applicare metodologie come Agile per fare
        \emph{producing} e gestire il team;
  \item il tempo dedicato a queste attività va però sottratto al game design,
        e dipende da quante persone ci sono nel team, da quanto lavoro c'è e
        da quanto è grande il progetto;
  \item avere sviluppato e rilasciato un progetto, anche piccolo (anche una
        \emph{game jam}), documentando i propri processi, è già un ottimo
        argomento da presentare in futuro;
  \item la via migliore resta lavorare nell'industria e, da lì, chiedere di
        imparare anche altre discipline.
\end{itemize}
Un gioco si può fare su qualsiasi cosa: sugli scarafaggi, sui piedini delle
sedie, sui pennarelli. Si parte avvantaggiati da ciò che è vicino alla propria
sensibilità: chi ama la musica può pensare a un gioco con meccaniche legate
alla musica, con un gameplay innovativo ma contenuto e realizzabile. Il design
è un cappello che comprende molte discipline.

% ---------------------------------------------------------------------
\section{Approfondimenti emersi in lezione}
\label{sec:approfondimenti}

\subsection{Roguelike e roguelite}

Un esempio di come il marketing conia generi: Supergiant Games, con
\emph{Hades}, rese celebre il termine \textbf{roguelite} al posto di
\emph{roguelike} (pur non essendo il primo roguelite in assoluto). La
differenza:
\begin{itemize}
  \item il \textbf{roguelike} si fonda su una serie di punti fissati in una
        sorta di \emph{manifesto} che definisce che cosa rende un gioco
        roguelike; uno su tutti: a ogni morte si ricomincia la \emph{run} da
        capo;
  \item il \textbf{roguelite} rompe alcuni di questi punti, in particolare
        introducendo una \textbf{progressione esterna alla run}: il giocatore
        può potenziarsi fuori dalla singola partita per vincere più
        facilmente in quelle successive. \emph{Hades} ne è il caso principe.
\end{itemize}
Analogamente, un mobile game sviluppato dal direttore del corso fu
promosso come ``il primo \emph{donation game} al mondo'': un gioco che
convertiva il tempo di gioco in donazioni reali a enti benefici per animali.
Anche qui un neologismo creato lato marketing per definire un genere.

\subsection{Applicazioni fuori dall'industria videoludica}

Le competenze sviluppate non servono solo ai videogiochi:
\begin{itemize}
  \item i \textbf{render architettonici} oggi possono essere realizzati con
        Unreal, e diversi professionisti formati nel game design lavorano in
        questo campo;
  \item la \textbf{gamification}, cioè l'applicazione di regole videoludiche
        a contesti non ludici (mostre interattive, progetti culturali,
        consulenze come quella per applicare questi concetti alle macchine da
        allenamento di un'azienda di fitness), è in Italia uno dei modi più
        redditizi di spendere queste competenze, spesso con finanziamenti
        pubblici.
\end{itemize}

% ---------------------------------------------------------------------
\section{Il primo task}
\label{sec:primo-task}

\begin{enumerate}
  \item Creare un indirizzo e-mail \textbf{professionale} (consigliato: un indirizzo Gmail nella forma
        \texttt{nome.cognome@gmail.com}; l'anno di nascita è accettabile).
        Va bene anche un indirizzo professionale già esistente usato per
        lavoro. Da evitare indirizzi poco seri (``non \texttt{blackfire67}'').
  \item Una persona della classe raccoglie tutti gli indirizzi in un unico
        documento e lo invia via e-mail al docente, che aggiungerà tutti al
        Drive condiviso del corso (dove verranno caricate lezioni, slide e
        materiali, per avere i documenti tutti nello stesso posto).
  \item Installare Discord e iscriversi al server del corso, dove passano le
        comunicazioni ufficiali. È consuetudine, e utile, creare anche un
        gruppo tra studenti.
\end{enumerate}
\textbf{Scadenza}: prima della lezione del venerdì.

% ---------------------------------------------------------------------
\section*{Punti chiave del capitolo}
\addcontentsline{toc}{section}{Punti chiave del capitolo}

\begin{itemize}
  \item Il corso è una \textbf{simulazione di lavoro}: ci si tratta da
        colleghi, con una gerarchia da studio di sviluppo.
  \item Si impara con l'\textbf{inquiry-based learning}: lezione, task,
        dubbi, domande, condivisione, progetti.
  \item Il fallimento affrontato con la mentalità giusta è la principale fonte
        di crescita; i feedback riguardano il lavoro, mai la persona.
  \item Il lavoro del designer è per metà creativo e per metà
        \textbf{ricerca e metodo scientifico}; il game design non è una
        scienza esatta.
  \item \textbf{Documentazione} e \textbf{pitching} sono il cuore del
        mestiere.
  \item Il designer progetta meccaniche e sistemi (partendo dalle
        \textbf{3C}), documenta, media, prototipa, confuta le idee, studia i
        competitor, bilancia e fa tuning.
  \item La leadership del designer è \textbf{al servizio} del team, non
        verticale.
  \item I limiti (tempo, tecnologia, competenze, persone, denaro) restringono
        e indirizzano la creatività.
  \item \textbf{Le idee sono nulla, l'implementazione è tutto.}
\end{itemize}

%    % ...
%
%  (Con hyperref aggiungere anche \renewcommand{\theHchapter}{z.\arabic{chapter}}
%  nella parte zero, per evitare ancore duplicate.)
%  Tutte le etichette di questo capitolo hanno il prefisso "z1:" per non
%  collidere con quelle dei capitoli del corso.
% =====================================================================

\chapter[La cultura del feedback]{Lavorare in squadra: la cultura del feedback}
\label{cap:z1-feedback}

Questo capitolo apre le \emph{lezioni zero}: incontri che precedono l'inizio
vero e proprio del corso di Game Design e che non riguardano una tecnica di
design in senso stretto, ma \textbf{la figura professionale} di chi lavora
nell'industria dei videogiochi. Prima di imparare a progettare un gioco
conviene infatti capire in che tipo di ambiente lo si progetterà: un ambiente
fatto di team, di gerarchie, di scadenze e, soprattutto, di persone che devono
riuscire a lavorare insieme.

La lezione è la prima giornata di un \textbf{workshop di due giorni} (circa
dieci ore in totale, pause pranzo comprese) tenuto da Stefano
Cozzi\footnote{Il cognome è riportato come compare nella trascrizione
automatica, che in più punti storpia i nomi propri; andrebbe verificato.
Nella stessa trascrizione lo studio per cui lavora appare come
``Stormine'', ``Stormline'' e ``Stormind'': si tratta di Stormind Games,
studio con sede a Catania, coerentemente con quanto detto in lezione.}, un
professionista con circa quindici anni di esperienza nel settore, tra
Blizzard Entertainment, Riot Games e, oggi, Stormind Games. Il workshop è
nato da un confronto con i coordinatori dell'accademia\footnote{Nella
trascrizione l'accademia è resa in modi diversi (``Hive'', ``a dive'',
``D.I.V.''); i coordinatori citati sono Enrico e Andrea, insieme al
direttore Diego Ricchiuti. Nel resto del capitolo si usa semplicemente
``l'accademia'', come negli altri capitoli.}: serviva un momento dedicato a
\emph{come ci si comporta in un ambiente professionale}, cosa che riguarda sia
il lavoro futuro sia i progetti di gruppo svolti durante il percorso di
studi, dove ci si deve già abituare a comportarsi ``come se si fosse in
un'azienda''.

La prima giornata ha un unico grande tema, sezionato in più parti:
\textbf{la cultura del feedback}. La seconda giornata (che sarà oggetto della
prossima lezione zero) sarà più pratica: \emph{time management},
\emph{brainstorming} e il ``guidare da dentro'' un team anche senza esserne
il leader (sezione~\ref{sec:z1-domani}).

Il percorso del capitolo segue quello della lezione:
\begin{enumerate}
  \item il formato del workshop e il percorso professionale del relatore
        (sezione~\ref{sec:z1-relatore});
  \item perché le aziende tengono tanto al lavoro di squadra e al
        \emph{fit} culturale (sezione~\ref{sec:z1-squadra});
  \item che cos'è il feedback e da che cosa dipende il modo di darlo
        (sezione~\ref{sec:z1-cose-feedback});
  \item le cinque fasi di vita di un team e il valore del
        \emph{post-mortem} (sezioni~\ref{sec:z1-fasi-team}
        e~\ref{sec:z1-qa-fasi});
  \item il confronto con i colleghi e alcune domande sul lavoro nei grandi
        studi (sezioni~\ref{sec:z1-rivalita} e~\ref{sec:z1-industria});
  \item il circolo virtuoso tra feedback, miglioramento e fiducia, e i suoi
        effetti sulla \emph{retention} e sulla cultura aziendale
        (sezioni~\ref{sec:z1-triangolo} e~\ref{sec:z1-cultura});
  \item il ``gioco del traditore'' e il principio del
        \textbf{default to trust} (sezioni~\ref{sec:z1-among-us}
        e~\ref{sec:z1-default-trust});
  \item le traiettorie di carriera e la valutazione delle performance
        (sezione~\ref{sec:z1-traiettorie});
  \item come si dà, come si riceve e quando si dà un feedback
        (sezioni~\ref{sec:z1-dare}, \ref{sec:z1-ricevere}
        e~\ref{sec:z1-quando});
  \item il feedback come regalo, le domande finali e le anticipazioni
        della seconda giornata (sezioni~\ref{sec:z1-regalo},
        \ref{sec:z1-qa-finale} e~\ref{sec:z1-domani}).
\end{enumerate}

Una premessa di metodo: il relatore ha alternato continuamente teoria,
domande alla classe e racconti della propria esperienza. I racconti non sono
digressioni: sono gli esempi concreti con cui i principi vengono resi
tangibili, e per questo sono riportati con cura.

% ---------------------------------------------------------------------
\section{Il workshop e il relatore}
\label{sec:z1-relatore}

\subsection{Formato e regole di partecipazione}

Il relatore ha cercato di ``riassumere quindici anni di esperienza in dieci
ore''. Ha materiale per circa il doppio delle ore disponibili: le cose
strettamente necessarie vengono trattate per prime, il resto è tenuto ``in
riserva'' come \emph{nice to have}. Le regole del workshop:
\begin{itemize}
  \item le domande si fanno \textbf{alla fine di ogni slide}, così da chiudere
        prima il discorso; non c'è un'unica sessione finale di domande e
        risposte (anche se ci sarà), perché il relatore preferisce fermarsi
        spesso e non fare ``uno spiegone di dieci ore'';
  \item le domande non devono per forza riguardare la slide appena vista: se
        una slide dà lo spunto per altro, se ne può parlare; al massimo sarà
        il relatore a rispondere in breve o a rimandare a una slide
        successiva;
  \item si può chiedere anche di esperienze passate, opinioni e sensazioni;
        se per NDA o per serietà professionale non può rispondere, lo dirà
        esplicitamente (``non lasciate che sia un \emph{blocker}'');
  \item su Zoom si usa l'alzata di mano o la chat, che il relatore tiene
        aperta sul secondo monitor;
  \item non ci sono pause ufficiali: la lezione è registrata e ognuno è
        libero di assentarsi quando serve, cercando però di non disturbare il
        flusso della lezione;
  \item le slide verranno fornite integralmente tramite la segreteria
        dell'accademia (Vanessa), insieme alle registrazioni: le slide senza
        il parlato servono poco, ma sono un buon promemoria.
\end{itemize}

\subsection{Il percorso professionale}

Il relatore si presenta (scherzosamente) come la ``guida escursionistica''
che accompagna la classe nella ``foresta incantata'' della cultura del
feedback, perché per molti è un territorio sconosciuto. Il suo percorso:
\begin{itemize}
  \item in Italia ha lavorato per aziende che fornivano servizi a terzi, tra
        cui il lavoro su \emph{Warhammer Online}, e per un canale televisivo
        sul digitale terrestre dedicato ai videogiochi\footnote{Nella
        trascrizione: ``Digital Game Network Channel''. Il nome esatto del
        canale andrebbe verificato.};
  \item a \textbf{Blizzard Entertainment} (circa sei anni, a Parigi) ha fatto
        il \emph{community manager} per l'Italia (\emph{StarCraft~II},
        \emph{World of Warcraft} da \emph{Mists of Pandaria} in poi,
        \emph{Hearthstone}, \emph{Diablo~III}, le espansioni di
        \emph{StarCraft}, il lancio di Battle.net), per poi diventare il
        manager dei community manager europei;
  \item a \textbf{Riot Games} (circa cinque anni, a Dublino) è stato
        brevemente community manager italiano e poi \emph{publishing
        manager} di un team di circa trenta persone per l'Europa. Il
        \emph{publishing} non è solo community: è l'unione di
        \textbf{community, marketing e PR};
  \item a \textbf{PG Esports} si è occupato di eventi dal vivo e online
        (ad esempio la Milano Games Week e tornei di \emph{Valorant} e
        \emph{League of Legends}), restando in contatto con Riot e Blizzard;
  \item da circa un anno è tornato nello sviluppo, a \textbf{Stormind
        Games}, dove oggi si occupa di risorse umane e organizzazione del
        lavoro; lo studio verrà presentato più ampiamente al \emph{Career
        Day} dell'accademia insieme alla sua HR manager;
  \item fa inoltre consulenze aziendali anche al di fuori dei videogiochi e
        insegna in master universitari (IULM di Milano).
\end{itemize}

Il suo ruolo ricorrente, sia a Blizzard sia a Riot, è stato quello di
\emph{Mr.~Wolf} (``per chi ha visto \emph{Pulp Fiction}''): quello che
risolve i problemi. Veniva inserito nei team più in difficoltà, ad esempio
dal punto di vista delle performance, per sistemarli e passare al successivo.
Era presente in alcuni dei momenti più critici, come il lancio di
\emph{Diablo~III} o il lancio di espansioni di \emph{World of
Warcraft}\footnote{La trascrizione riporta ``il lancio con i Gary's'' e
``World of Trenor'': si tratta verosimilmente di \emph{Mists of Pandaria} e
di \emph{Warlords of Draenor}.}, a fare \emph{damage control}.

La scoperta fondamentale di questi anni è che i problemi da risolvere
\textbf{quasi mai riguardavano il prodotto} (non era certo lui a sistemare i
server di \emph{Diablo~III}), ma i team che c'erano dietro: persone che
litigavano, non collaboravano, non riuscivano a stare nella stessa stanza
senza urlarsi contro. E nel 90\% dei casi tutto questo si riduce a
\textbf{problemi di comunicazione}: è raro che perfino le incompatibilità
caratteriali non siano risolvibili parlando.

\subsection{Struttura non vuol dire burocrazia}

A Stormind il relatore sta accompagnando un team passato da 30 a 95 persone
in meno di un anno, attraverso tutti i cambiamenti di cui parla il workshop.
Darsi una struttura viene spesso percepito male, come ``più processi, più
procedure'', ma bisogna distinguere:
\begin{itemize}
  \item un processo come ``timbrare il cartellino'' può essere davvero
        invasivo e noioso;
  \item definire \emph{come funziona un documento di game design, chi lo
        approva, chi lo visiona e chi no} è invece una struttura necessaria.
\end{itemize}
Uno studio di 30 persone nella stessa sede può farne a meno: quando c'è un
dubbio ci si alza e si va alla scrivania del collega. Uno studio di 95
persone sparse per l'Italia (e in parte per l'Europa) non può più farlo:
servono una struttura diversa, un modo di lavorare diverso e una
pianificazione del futuro diversa.

% ---------------------------------------------------------------------
\section{Perché il lavoro è sempre lavoro di squadra}
\label{sec:z1-squadra}

\subsection{Non si va in ufficio per fare amicizia}

Il primo messaggio è netto: \textbf{non si va in ufficio per fare amici} o
per trovare compagni di vita. Si va in ufficio perché c'è un lavoro da fare, e
lo scopo è riuscire a farlo \emph{anche con chi sta antipatico}. Con le
persone simpatiche è facile: diventeranno amici, e non sarà mai un problema.
È con chi sta antipatico che le cose si complicano, ed è lì che emergono i
\textbf{veri professionisti}: quelli che riescono a lavorare con chiunque
senza che l'antipatia pesi sul modo in cui lavorano e collaborano, mettendo il
prodotto e il team davanti al proprio ego e alle proprie preferenze personali.

La frase che fa da filo conduttore a tutto il workshop (e che il relatore
promette di riproporre ``fino alla nausea'') è:

\begin{quote}
  \emph{La capacità di un gruppo di persone di fare cose straordinarie
  dipende dal modo in cui queste persone operano insieme.}
\end{quote}

Le aziende raramente pensano che l'innesto di una singola persona possa
cambiare il destino di un progetto. Questo è tanto più vero quanto più grande
è il team, ma vale anche per i piccoli studi, anche per un gruppo di cinque
persone. Come negli sport di squadra: l'individuo conta fino a un certo punto,
e anche se si ha Dončić in squadra, se si gioca uno contro cinque si perde
comunque\footnote{Il riferimento è alle finali NBA in corso nei giorni del
workshop tra Dallas Mavericks e Boston Celtics (giugno 2024), il che
permette di datare approssimativamente la lezione.}.

\subsection{Lo \emph{useless genius} e il costo delle persone tossiche}

Molti dipartimenti HR usano la definizione di \textbf{useless
genius}\footnote{Il concetto è più noto in letteratura manageriale come
\emph{brilliant jerk}; in lezione è usata l'espressione \emph{useless
genius}.}: si può essere bravissimi, a volte i migliori nel proprio lavoro,
ma se non ci si sa comportare, relazionare e operare in un gruppo si è
fondamentalmente inutili ai fini del gruppo. Nel migliore dei casi inutili,
nel peggiore \emph{dannosi}.

In passato le aziende tendevano a ``chiudere un occhio'': la persona è
fondamentale, la teniamo. Oggi la domanda è diventata: \emph{quanto costa
all'azienda, oltre al suo stipendio, tenere una persona che crea un ambiente
tossico e peggiora chi le sta intorno?} Ci sono costi sommersi che le aziende
cominciano a vedere, e la tendenza è liberarsi di queste persone il prima
possibile. Il ragionamento è semplice:
\begin{itemize}
  \item le \textbf{competenze si possono trasmettere} o sostituire: si trova
        qualcuno che fa lo stesso lavoro magari nel doppio del tempo, oppure
        più persone che se lo dividono;
  \item il \textbf{carattere delle persone} invece cambia molto raramente.
\end{itemize}

\subsection{Il fit culturale}

Quando qualcuno \textbf{non è un \emph{fit} culturale} per un'azienda è
meglio separarsi presto. Ma non essere adatti a \emph{quell'}azienda non
significa non essere adatti a lavorare: è perfettamente legittimo che una
persona inadatta a un'azienda sia adattissima a un'altra.

Il relatore l'ha vissuto in prima persona. Dopo sei anni a Blizzard viene
contattato da Riot e affronta un processo di selezione lunghissimo (``quasi
undici colloqui''), al termine del quale viene scartato. Scrive allora una
mail che gli ha ``cambiato la traiettoria di carriera'': chiede
semplicemente perché non è stato ritenuto adatto. La risposta: in nessun
colloquio si era esposto abbastanza da far capire \emph{perché} volesse
lasciare Blizzard, e quindi pensavano che in realtà stesse bene dov'era. Il
motivo era che, per educazione professionale, non gli piace parlare male del
proprio posto di lavoro; portata all'eccesso, però, questa riservatezza era
diventata un ``è tutto bellissimo, ma vorrei provare Riot'', senza un perché.
Lo invitano comunque a un incontro dal vivo (``dal vivo è diverso''), e alla
fine viene assunto.

Una volta arrivato, i sei mesi di prova (\emph{probation}) sono serviti tutti
per smettere di ragionare e lavorare ``come a Blizzard''. La differenza
culturale era enorme, tanto che in Riot esisteva un gruppo di supporto
interno tra ex dipendenti Blizzard: un collega gliel'ha descritto come
``uscire da una relazione tossica e doversi riabituare a una relazione
normale''\footnote{Il relatore precisa che parla della Blizzard di allora,
che non esiste più: oggi è parte di Microsoft e nessuno dei suoi ex colleghi
vi lavora ancora, anche a causa dei licenziamenti di massa.}.
\begin{description}
  \item[Blizzard] era \emph{developer first}: il focus era chiaramente sugli
        sviluppatori, prima di tutti gli altri. Dopo sei anni c'erano ancora
        meeting a cui non poteva partecipare perché non abbastanza
        \emph{senior}, informazioni che non poteva sapere, piani che non
        poteva vedere.
  \item[Riot] era \emph{player first}: prima i giocatori, e i dipendenti come
        mezzo per trasferire il feedback dei giocatori dentro il gioco. Il
        secondo giorno lo portano in una stanza e gli mostrano il piano dei
        cinque anni successivi. Lui, sorpreso, pensa a un errore (``sono qui
        da oggi, non vorrei mettervi nei guai''); la risposta è: ``sei un
        dipendente Riot o no? Allora non ci sono livelli di accesso: devi
        sapere queste cose per fare meglio il tuo lavoro''.
\end{description}

Da questa esperienza ha imparato molto sulla \textbf{trasparenza}: condividere
le informazioni aiuta il lavoro. Non sempre, però: a volte non si condivide
qualcosa non per segretezza, ma perché non è ancora pronta per essere
presentata, esattamente come non si mostra un gioco quando è ancora solo
teoria (``te lo faccio vedere quando ho qualcosa di tangibile'').

La lezione per gli studenti è doppia: le aziende valuteranno il \emph{fit}
culturale sempre di più, ma capire se si è un \emph{fit} richiede anche
\textbf{lavoro da parte del candidato}. Bisogna avere la presenza di spirito
di dire: ``conosco i valori di questa azienda, non mando nemmeno il CV,
perché so che durerei un mese e mezzo''.

% ---------------------------------------------------------------------
\section{Che cos'è il feedback}
\label{sec:z1-cose-feedback}

\subsection{Le prime risposte della classe}

La prima domanda alla classe è: \emph{quando sentite la parola
``feedback'', a cosa pensate?} Le risposte raccolte, a voce e in chat:
\begin{itemize}
  \item un \textbf{riscontro} su ciò che si fa e su come lo si fa, sul lavoro
        o sul modo in cui ci si interfaccia con il gruppo; può essere
        lavorativo oppure comunicativo e comportamentale. Chi ha risposto lo
        vede come qualcosa che, se l'azienda non dà, va attivamente
        richiesto, perché è importante sapere come sta andando per
        migliorarsi;
  \item né positivo né negativo, ma \textbf{costruttivo}: in entrambi i casi
        serve a migliorarsi;
  \item uno \emph{scambio di opinioni}, un \emph{riscontro produttivo};
  \item un feedback ``di reazione'', non verbale: si fa una cosa e si
        osserva come cambia l'ambiente o il comportamento delle persone
        intorno (anche perché a voce la gente spesso si tiene le cose
        negative);
  \item per un designer, anche il \textbf{linguaggio del corpo}: osservando
        un giocatore giocare si può dedurre l'emozione che prova (frustrazione,
        incomprensione).
\end{itemize}

Sulla forma (scritto o orale?) le opinioni divergono: c'è chi lo preferisce
\emph{mille volte} scritto; chi pensa che dipenda dalle dimensioni
dell'azienda (orale nelle piccole, scritto in quelle con team di centinaia di
persone); chi pensa che dipenda dall'importanza del feedback (scritto se
tecnico o pieno di dati, da tenere sotto gli occhi; orale se è un consiglio
veloce).

Il relatore usa due esempi per mostrare che \emph{dipende}:
\begin{itemize}
  \item un feedback sulla programmazione ha senso darlo scritto, magari
        commentando direttamente il codice, più che ``andando a prendere un
        caffè'';
  \item ma se il feedback è ``ieri, durante il meeting, ho percepito male il
        fatto che mi interrompessi di continuo'', ce lo si aspetta
        \textbf{orale}.
\end{itemize}

\subsection{Da che cosa dipende il modo di dare feedback}

Lo \emph{spoiler} della giornata è che \textbf{non esiste una risposta univoca}
su come dare feedback. Il workshop non insegna ``il'' modo, ma fornisce
\emph{tutte le opzioni}, perché il modo giusto cambia in base a molti fattori:
\begin{itemize}
  \item a chi lo si dà: un collega, un superiore, il proprio capo, un proprio
        \emph{report} (un sottoposto);
  \item se è personale o professionale;
  \item se è la prima volta o la terza volta che si dà lo stesso feedback;
  \item se riguarda un progetto importante o secondario.
\end{itemize}

Manca però l'elemento fondamentale, che la classe individua subito:
\textbf{l'altro}. A parità di tutti gli altri fattori, il \emph{carattere
della persona che riceve il feedback} dovrebbe far cambiare il modo in cui lo
si dà. Per questo, prima della lezione, agli studenti era stato chiesto di
fare un test di personalità: non un test di psicologia avanzata, ma uno
strumento per rendersi conto di quanta varietà ci sia anche solo in una
classe di 53 persone (il che non vuol dire trovare 52 modi diversi di dare
feedback, ma lavorarci un po' sì).

Il relatore porta se stesso come esempio. Preferisce:
\begin{itemize}
  \item i feedback \textbf{positivi scritti}, perché è a disagio a ricevere
        complimenti in pubblico;
  \item i feedback sul proprio lavoro \textbf{diretti}: ``vai subito al
        punto e non girarmi intorno''.
\end{itemize}
È molto orgoglioso del suo lavoro e, quando qualcuno lo critica, la sua
primissima reazione istintiva è ``come ti permetti''. Negli anni ha imparato
a farla durare un secondo, ma chi gli dà feedback di persona dice di
vederla ancora negli occhi. Per mitigarla, dice esplicitamente a chi gli dà
feedback: \emph{``ignora il modo in cui reagisco nei primi cinque secondi,
vai avanti, poi mi passa''}. È considerato da tutti molto aperto al feedback
e non ha problemi a dire ``hai ragione, facciamo diversamente'': ma ci sono
reazioni istintive e caratteriali che è difficile controllare al cento per
cento.

\subsection{Con chi si lavora meglio}

Guardando le persone con cui si lavora abitualmente (anche nei progetti
dell'accademia), si noterà che si lavora meglio con due profili:
\begin{enumerate}
  \item chi è \textbf{esattamente come noi}, ma qui è un 50 e 50: alcuni
        caratteri vanno d'accordo con i propri simili, altri non vogliono
        nessuno con lo stesso carattere nel gruppo (e si potrebbe rientrare
        in questa categoria);
  \item chi è \textbf{affine}: ha qualcosa in comune con noi, ma non tutto,
        e proprio per questo può dare punti di vista differenti.
\end{enumerate}
Il lavoro della giornata è capire come valorizzare il secondo caso: ottenere
feedback da persone con caratteri, visioni e background diversi, per
completare una visione del proprio lavoro che è necessariamente limitata. È
come esplorare una mappa con la nebbia di guerra: un'altra persona, al di là
del muro, vede cose che noi non vediamo.

\subsection{Il feedback come somma di piccoli miglioramenti}

Per il relatore il feedback è questo: \textbf{applicare tanti piccoli
miglioramenti periodici}, che si sommano nel tempo. Migliorare dello 0,1\%
ogni giorno porta a fine anno a un miglioramento enorme, e questo vale per i
prodotti, per i giochi, per le persone e per i team.

\paragraph{L'esempio del ciclismo britannico.} La squadra di ciclismo
britannica, fino a una quindicina d'anni fa, non aveva vinto quasi nulla.
Cambiato il team manager, tutti si aspettavano grandi cambiamenti sportivi
(convocazioni, strategie di gara). Arrivarono invece \emph{minuscoli
cambiamenti non sportivi}: un altro chef in trasferta, materassi migliori nei
camper per riposare meglio, un sellino più comodo a fine tappa. Cambiamenti
marginali che da soli non avrebbero cambiato niente; eppure, dopo due o tre
anni di questo lavoro quotidiano, la squadra ha cominciato a vincere a
ripetizione, Olimpiadi comprese, cambiando il modo in cui il ciclismo è
visto in Gran Bretagna\footnote{È la nota strategia delle \emph{marginal
gains} associata a Dave Brailsford (British Cycling e Team Sky); il nome non
è citato in lezione.}. Certo, si è investito anche sulle persone giuste, ma
tutto è partito dai piccoli miglioramenti periodici.

Per gli studenti vale lo stesso: chiarire un conflitto con una persona con
cui si ha ancora poco a che fare, o una situazione che ha lasciato
``l'amaro in bocca'', sembra inutile oggi, ma darà i suoi frutti tra tre, sei,
nove mesi. Sono \emph{investimenti sul futuro}.

\paragraph{Peggiorare ogni giorno.} Esiste anche il percorso inverso:
lasciar correre, ``buttare sotto il tappeto'' le cose perché ``tanto è
inutile parlarne''. Queste cose tornano sempre fuori dopo otto mesi, come il
\emph{libro dei rancori} dei nani di Warhammer: ``ti ricordi quando, un anno
e mezzo fa, all'inizio del progetto, mi hai detto quella cosa? Mi sono
offeso''. A quel punto è quasi impossibile risolvere, perché nessuno ricorda
più con precisione cosa sia successo (``ma quale cena?'', ``forse era in
pizzeria'', ``no, eravamo da un'altra parte''). \textbf{Tenersi le cose nel
cassetto per tirarle fuori quando servono è sempre sconsigliato}: non servono
ed è impossibile usarle in modo costruttivo.

% ---------------------------------------------------------------------
\section{Le cinque fasi di vita di un team}
\label{sec:z1-fasi-team}

\subsection{Perché conoscerle}

Quasi tutti gli studenti lavorano già in gruppo (al primo anno su un gioco da
tavolo, negli anni successivi su un videogioco) e si saranno accorti che un
team attraversa delle fasi: il modo di lavorare insieme cambia nel tempo.
Spesso ci si fa prendere di sorpresa: ``eravamo così affiatati, cosa è
successo?''.

Un esempio da una consulenza fuori dai videogiochi: un'azienda farmaceutica
gli affida un team commerciale (rappresentanti che presentano medicinali a
ospedali, medici e privati). Il problema: ``prima andava tutto bene; ora che
siamo in diciotto dovremmo lavorare meno e meglio, e invece va peggio''. Va
peggio perché le persone nuove sono appena arrivate, mentre il management si
aspettava già di vedere i risultati dell'investimento.

Fino a pochi anni fa le fasi di un team venivano disegnate come una
\textbf{linea dritta verso l'alto}: le fasi erano le stesse, ma viste come un
miglioramento costante. Non si vedevano le parti \emph{fisiologiche} del
lavoro in team, cioè anche i peggioramenti. Accettarle è il primo passo per
accettare i fallimenti: nessun percorso personale o di team è lineare.

Il relatore ammette che il suo modo di lavorare è nato per il 70\% dal
proprio impegno a imparare, fare bene e rispettare le persone, e per il 30\%
dall'aver lavorato con ``incapaci colossali'' che gli hanno insegnato come
\emph{non} fare le cose. A Blizzard un manager gli fece capire che, pur non
sapendo ancora che tipo di \emph{people manager} sarebbe diventato, sapeva
benissimo quale tipo di manager \emph{non} voleva essere.

\subsection{Forming, storming, norming, performing, adjourning}

Le fasi sono cinque, un po' come le fasi del lutto. Molte aziende ne
considerano solo quattro (``e sono quelle che solitamente chiudono''); le
aziende al passo con i tempi ne considerano cinque\footnote{Il modello è
quello di Bruce Tuckman (1965), a cui la quinta fase, \emph{adjourning}, fu
aggiunta nel 1977. Il nome di Tuckman non è citato in lezione; la
trascrizione riporta la quinta fase come ``journing''.}.

\begin{description}
  \item[Forming] Si viene messi insieme e si comincia a lavorare. Non si sa
        nulla degli altri; spesso non si conoscono nemmeno i nomi, ed è
        normale.
  \item[Storming] Si cominciano a ``dare le testate come le capre in
        montagna per determinare il territorio'': questa cosa la faccio io,
        questa la vorrei fare così, chi decide? Si stabilisce chi fa cosa e
        come si divide il lavoro. Nel grafico è una \textbf{fase di
        peggioramento}: che il team litighi o discuta civilmente non cambia
        nulla, il team deve investire parte di tempo ed energie per risolvere
        queste questioni. È normale e \textbf{non è evitabile}.
  \item[Norming] Si comincia a capire come lavorano gli altri, cosa sanno
        fare e come far funzionare il gruppo; iniziano i miglioramenti. È una
        forma di \emph{accettazione} degli altri, che deriva dall'aver capito
        cosa devo fare io e cosa devono fare loro.
  \item[Performing] C'è uno sforzo \textbf{proattivo}: ``questa cosa
        toccherebbe a me, ma so che Alessio la fa molto meglio, quindi gli
        chiedo aiuto''. Si sfruttano i punti di forza degli altri, perché li
        si conosce e perché si sa che passare il testimone non avrà effetti
        negativi (il capo non dirà ``perché non l'hai fatta tu?'').
  \item[Adjourning] Il progetto finisce, si chiude o viene riorganizzato, il
        team viene sciolto. Quello che andrebbe fatto è analizzare
        \textbf{cosa ha funzionato, cosa no e cosa migliorare} per il
        progetto successivo.
\end{description}

\subsection{La fase che tutti saltano}

``Spoiler incredibile'': l'\emph{adjourning} è la fase che molte aziende
saltano, \textbf{perché c'è sempre il progetto dopo}.

A PG Esports, per la Milano Games Week (come per qualsiasi evento dal vivo a
cadenza annuale: Lucca Comics, Etna Comics, Comicon\dots), finita l'edizione
di un anno si cominciava dal giorno dopo a lavorare a quella successiva,
perché un anno sembra tanto ma non basta mai. Nei videogiochi è lo stesso:
finito un gioco, a parte le patch di mantenimento e del \emph{day one}, si
passa subito al successivo.

Il cambiamento introdotto dal relatore (in parte già praticato, ma da lui
strutturato) è stato: finita la Games Week, \textbf{tutti in ferie per una
settimana}; poi ci si ritrova per definire cosa ha funzionato e cosa no, con
un \textbf{post-mortem} (o \emph{debrief}); solo dopo si inizia a lavorare
all'edizione successiva. Senza questa fase manca la parte di
\emph{learning}, cioè di ciò che si è imparato dal progetto.

Il consiglio vale per gli studenti: quando si finisce il gioco da tavolo, e
poi il gioco digitale, bisogna trovare il tempo per il post-mortem,
altrimenti si butta via \textbf{il 90\% dell'utilità del progetto}. Una
struttura semplice per farlo è descritta nella sezione~\ref{sec:z1-qa-finale}.

% ---------------------------------------------------------------------
\section{Domande sulle fasi del team}
\label{sec:z1-qa-fasi}

\subsection{Lo storming che ritorna}

\paragraph{Capita che lo storming si ritagli spazio anche nelle altre fasi?}
Lo \emph{storming} è una fase precisa, che però molti manager, coordinatori e
team cercano di \textbf{saltare}: ``se io sono bravo come manager e loro sono
bravi come team, passeremo da forming a norming senza problemi''. È un errore,
perché senza una fase in cui si discute davvero non si sviluppa una
\emph{skill} fondamentale: \textbf{saper discutere}. Tutti conoscono persone
non abituate a farlo, per cui ogni minimo disaccordo diventa un affronto
personale (``avevamo deciso la pizza, perché andiamo al sushi? Perché mi fate
questo affronto?''): sono persone che non sono abituate ad avere conflitti
costruttivi.

Se si pensa di essere in \emph{norming} e continuano a emergere conflitti, il
punto è che \textbf{non si era davvero in norming}: lo storming non ``si
trascina'' nelle fasi successive, semplicemente non se n'è mai usciti.
Allo stesso modo, a chi chiede se ogni mini-fase di storming durante la vita
del gruppo sia frutto di uno storming fatto male o solo di carattere, la
risposta è già contenuta qui.

Un partecipante osserva che spesso, ancor prima di finire un progetto, si è
già messi su quello successivo: è il modo migliore per essere in costante
\emph{damage control}. Era l'esperienza del relatore a Blizzard: veniva
spostato fisicamente da un team all'altro, e finito il lancio di
\emph{Diablo~III} era già su \emph{World of Warcraft}, senza vivere il
``dopo''. Si è dovuto ritagliare da solo lo spazio per un debrief personale,
ed è lì che ha deciso che tipo di manager non voleva essere. Il consiglio:
\textbf{se non lo fa il team, fatelo voi}.

\subsection{Dove si passa la maggior parte del tempo}

\emph{In quale fase si passa più tempo nella propria carriera?} Gli studenti
rispondono \emph{storming} e \emph{norming}. Il relatore, in quindici anni,
calcola di aver passato in \emph{performing} circa \textbf{dieci mesi}. Non è
pessimismo: arrivare a performing deve essere sempre l'obiettivo, sia da
manager sia da membro del team; ma è difficile arrivarci, ed è normale che non
sempre ci si arrivi.

Lo storming può durare a lungo, ma con team molto esperti dura pochissimo,
perché c'è chi lo sa gestire. In realtà \textbf{la maggior parte del tempo si
passa in forming}, perché ogni equilibrio raggiunto vale per \emph{quelle}
persone, \emph{quel} progetto, \emph{quelle} tempistiche. Basta che cambi un
qualsiasi fattore (una deadline spostata, una persona nuova, qualcuno che si
dimette, un progetto cancellato) e si ricomincia, perché i bilanciamenti
precedenti non valgono più.

Esempio: un team di cinque persone affiatatissimo, in pieno performing. La
dirigenza, entusiasta, ``raddoppia le risorse'' e lo porta a dieci. O si
lascia tutto com'è (e i cinque nuovi non fanno niente), o bisogna ridiscutere
chi fa cosa e come i nuovi possono aiutare: che piaccia o no, si ricomincia
da forming.

Quindi la bravura di un manager, ma anche di un membro del team,
\textbf{non sta nel passare tanto tempo in performing, ma nell'accelerare il
passaggio da una fase all'altra}, e nel fare bene la fase di
\emph{adjourning}.

\subsection{Analizzare anche i successi}

C'è la tendenza a fare post-mortem solo dei progetti andati male: ``questa
cosa non ha funzionato, vediamo perché''. Va benissimo. Ma quando un progetto
va bene si dà per scontato che sia andato bene ``perché era buono''. Bisogna
invece chiedersi \emph{perché} è andato bene: era il prodotto? La
tempistica? Il modo in cui abbiamo fatto una cosa piuttosto che un'altra?
\textbf{Se non si capisce davvero il motivo del proprio successo, si
tenderanno a replicare le cose sbagliate.}

\paragraph{Blockbuster.} Quando arrivò la notizia che Netflix avrebbe
lanciato un servizio di streaming, Blockbuster (la catena di noleggio di
videocassette e DVD) aveva tutto il tempo per preparare un'offerta
concorrente. In una riunione del management raccontata in un libro (di cui il
relatore non ricorda il titolo), la direzione sembrava orientata a investire
nello streaming, finché un dirigente fece notare che il 12\% del fatturato
derivava dalle \emph{multe per riconsegna tardiva}, impossibili con lo
streaming, mentre secondo le stime lo streaming avrebbe prodotto il 10\% del
fatturato. ``Non possiamo rinunciare a questo 2\%'': niente streaming. Oggi
non c'è più un Blockbuster ``manco a cercarlo col lanternino''. Non avevano
capito qual era il vero motivo del loro successo, cioè il loro modello di
business\footnote{Le percentuali sono quelle riferite in lezione e non sono
state verificate.}.

\paragraph{The Founder.} Nel film sulla nascita di McDonald's c'è una scena in
cui il futuro socio del protagonista gli dice: ``tu pensi di essere nella
\emph{food industry}, e invece sei nella \emph{real estate industry}''. Non
dovrebbe investire sugli hamburger, ma sui terreni su cui costruisce i
ristoranti. Anche per un team vale lo stesso: bisogna capire che cosa
l'ha reso davvero di successo, oppure no.

\subsection{L'autoanalisi personale}

\paragraph{E quando si cambia azienda?} Uno studente chiede se il debrief sia
possibile anche quando si lascia un'azienda. Non si può tornare in ufficio a
discutere della propria esperienza, ma lo si può fare \textbf{per se stessi}.
Il relatore l'ha fatto, anche se all'epoca in modo poco consapevole, ed è
così che è passato dal \emph{community management} al \emph{people
management}. Analizzando cosa aveva fatto bene a Blizzard si era reso conto
che:
\begin{itemize}
  \item aveva costruito un ottimo rapporto di fiducia con i giocatori (con
        alcuni proprietari di fansite parla ancora oggi);
  \item pur senza la seniority, era diventato un punto di riferimento per il
        team, visto come imparziale ed equilibrato anche su temi che non lo
        riguardavano direttamente;
  \item sapeva \textbf{risolvere problemi}: durante il lancio di
        \emph{Diablo~III} era il più tranquillo di tutti, non perché
        sminuisse il problema, ma perché sapeva che si sarebbe risolto con i
        suoi tempi; il lavoro non era spegnere l'incendio ma ``far sfollare
        la gente'' e aspettare che tutto tornasse in ordine.
\end{itemize}
E aveva capito cosa sapeva fare meno bene, ad esempio \textbf{gestire il
proprio manager}. Avere un manager non significa che il manager gestisca te:
il manager gestisce il tuo lavoro, ma tu devi saper gestire i tuoi
\emph{stakeholder}. Spesso lo stakeholder è il proprio manager; da manager,
gli stakeholder saranno CEO, CFO e così via. Gestirli significa conoscerli,
sapere come portare loro le informazioni, come far presente che qualcosa non
va. Fa parte del lavoro, anche se a lui non piace, e ci ha lavorato.

Nel mese e mezzo di pausa prima di entrare in Riot si è concentrato su
questo: ``cosa porto come mio punto di forza?''. L'autoanalisi è culminata
nel chiedere espressamente alla persona che l'aveva assunto: ``prima mi avete
detto di no, ora mi avete preso: perché?''.

\paragraph{Come si capisce se un progetto è andato bene?} Dipende
dall'\textbf{obiettivo}. Se l'obiettivo era vendere milioni di copie e il
gioco le ha vendute, bisogna capire perché. Se, più ragionevolmente,
l'obiettivo era un gioco da tavolo che permettesse di mettere in pratica le
conoscenze del corso, va valutato per quello: ha permesso di farlo? Se sì,
perché? Se no, perché no? L'obiettivo va fissato \textbf{all'inizio del
progetto, non alla fine}, proprio per poterlo valutare dopo.

\subsection{Quando il team cambia di continuo}

\paragraph{Tutti questi cambiamenti non sono dannosi per il progetto?}
Assolutamente sì, ma quasi mai sono volontari. I progetti andati male sono
spesso progetti grandi.

Il relatore racconta di \emph{Cyberpunk 2077}. Conosceva persone in CD
Projekt Red (``il mondo è piccolissimo, le cose si sanno''). Alla Gamescom va
con Diego Ricchiuti, direttore dell'accademia, a vedere la presentazione del
primo gameplay, quella ``in cui regalavano il giubbottino da samurai''.
Finito il video, tra gli applausi, si guardano e capiscono entrambi che
\textbf{il gioco non è pronto}, a soli quattro mesi dall'uscita prevista:
Diego aveva notato il combattimento e la guida. Gli amici di CD Projekt Red
(uno aveva lavorato con lui a Blizzard per sei anni) gli confermano tra le
righe che ``ce la faremo, ma ci sono ancora cose da sistemare''. Il gioco fu
poi rimandato più volte\footnote{In lezione il relatore dice ``due, forse
tre volte''; l'uscita di \emph{Cyberpunk 2077} è stata effettivamente
rinviata tre volte prima del dicembre 2020.}.

Per deformazione professionale, chi ha sviluppato l'occhio critico non
riesce più a godersi un trailer: in un trailer di \emph{The Day Before} Diego
gli aveva fatto notare che l'effetto della fiammata di un fucile era stato
chiaramente aggiunto in postproduzione e non faceva parte del gameplay. ``È
il triste epilogo di questa storia'': non riescono più a godersi un trailer.

Il punto è che più il team è grande, più queste variabili vengono imposte, e
più è facile che il fallimento \textbf{non dipenda dal fatto che il prodotto
sia brutto} (\emph{Cyberpunk} è un bellissimo gioco sotto molti punti di
vista, e il relatore lo sta rigiocando ora che è stato corretto ed espanso).
Nei grandi gruppi qualcuno arriva, qualcuno si dimette, cambia il budget,
arriva il rinvio; nel caso di \emph{Anthem}, il team ha fatto così tante volte
forming--storming--forming--storming che ne è uscito ``un prodotto che
dovevamo espellere: l'abbiamo espulso, buttato fuori dal nostro organismo pur
di liberarcene''.

Tra i cambiamenti c'è anche la crescita verticale: se il capo viene promosso
perché lavora benissimo, bisogna scegliere un nuovo capo, ed è sempre un
caos (chi ci resta male, chi pensava di meritarlo, chi voleva il suo amico,
chi non sopporta il prescelto). Per questo il relatore è \textbf{contrario
alle promozioni durante i progetti}: si rischia di rovinare tutto il lavoro di
storming.

\subsection{Obiettivi chiari nei piccoli gruppi}

Nei gruppi piccoli, come quelli degli studenti, il problema deriva più
spesso dal \textbf{non avere chiaro che cosa si vuole ottenere}: un prodotto
puramente educativo e formativo, o un prodotto richiesto dall'accademia che,
se viene bene, si potrebbe anche vendere? Un gioco commerciale o un gioco
artistico? Va chiarito subito, perché ci saranno sempre momenti in cui
scegliere tra A e B, e a guidare la scelta sarà la domanda: \emph{quale delle
due opzioni ci porta più vicino all'obiettivo?} Se l'obiettivo non è chiaro,
non si sa cosa scegliere.

Questo porta molte aziende a fare \textbf{scelte di design che in realtà sono
scelte commerciali}. Le microtransazioni sono scelte commerciali: nessuno ha
mai pensato ``sai cosa piacerebbe tanto ai giocatori? Pagare le bustine''.
Prima le bustine si guadagnavano giocando; ora la tendenza è ``potresti
giocare, ma se paghi è meglio''. Si pensi alle polemiche su Ubisoft per
l'opzione di ``saltare il gioco'' pagando, o ai famosi \emph{boost} di livello
di \emph{World of Warcraft}: ``devo pagarti per saltare il gioco?''.

Una studentessa (Silvia) cita \emph{Final Fantasy XIV}, che a un certo punto fu
``resettato'' perché non funzionava: un esempio positivo. Il relatore aggiunge
quello negativo di \emph{Age of Conan}, che non è mai riuscito a fare quel
reset iniziale (``è una schifezza, ve lo ripresentiamo quando è pronto''). Ci
vuole consapevolezza, ma anche i mezzi: un conto è Square Enix, che può
permetterselo, un altro è un'azienda che non può chiudere tutto per un anno
continuando a pagare gli stipendi.

Riassumendo: nei piccoli gruppi i problemi sono più di
\textbf{comunicazione e chiarezza} (obiettivi non chiari o che si spostano);
nelle grandi aziende di più di \textbf{persone che vanno e vengono, obiettivi
e tempistiche che cambiano}. In accademia le scadenze sono fisse (il gioco da
tavolo si consegna a fine anno scolastico, a giugno, e giugno resta); in
azienda può succedere che una feature prevista per la \emph{vertical slice}
di luglio venga anticipata a maggio per portarla alla Gamescom, e un sistema
perfetto per consegnare a giugno può non esserlo più per consegnare a maggio.

\subsection{L'ambiente determina la qualità del prodotto}

\paragraph{È vero che chi passa per Electronic Arts finisce per fare brutti
giochi?} Uno studente cita BioWare, che faceva ottimi giochi finché era
autonoma. Il problema non è dei developer, che non sono ``improvvisamente
diventati incapaci'': è evidente che il problema è altrove, e che
\textbf{l'ambiente in cui si fa un gioco determina molto spesso la qualità del
prodotto}.

Un esempio locale positivo: Tommaso Bonanni, conosciuto dal relatore insieme a
Diego e già docente dell'accademia, ha fondato a Roma Caracal Games Studio,
applicando fin dall'inizio questi concetti di cultura aziendale. L'azienda è
rimasta piccola, ma ha molto chiaro cosa vuole fare e come. Quando si
predispone fin dall'inizio l'azienda e il team a lavorare in un certo modo è
fattibilissimo; è molto più difficile quando si è EA, si acquisisce
un'azienda e le si dice ``da oggi si fa così''.

Il relatore ha vissuto la transizione da Blizzard ad Activision Blizzard (uno
dei motivi per cui se n'è andato, volontariamente). Ha visto la cultura
interna cambiare per gradi:
\begin{enumerate}
  \item ``facciamo giochi fighissimi, perché i giochi devono essere
        fighissimi'';
  \item ``facciamo giochi fighi così vendono di più'';
  \item ``vendiamo di più facendo giochi fighi'' (sembra uguale, ma è molto
        diverso);
  \item ``facciamo più soldi''.
\end{enumerate}
Lo si capiva dalle scelte di design, osservando quale fosse il peso più
grande sulla bilancia. Lo stesso varrà per gli studenti, nei loro giochi e nei
colloqui di lavoro: capire subito che cosa vuole l'altra parte.

\paragraph{Leader e visione.} Nei gruppi dell'accademia, dal secondo anno, c'è
una persona che assume il ruolo di \emph{lead} o \emph{scrum master}. È
diverso dall'essere un \emph{game director} (``il gioco si fa come dico io,
aiutatemi a farlo con la mia visione''), come avviene nei grandi studi: da lì
nascono i nomi nei titoli, come \emph{Sid Meier's Civilization}, che indicano
la persona che ha fatto quel gioco, oppure un puro branding, come \emph{Tom
Clancy's Rainbow Six}.

\paragraph{Inclusion come jolly.} A margine, uno studente osserva che in molte
aziende ``la risposta giusta è sempre \emph{inclusion}'', usata come un
jolly. È comunissimo, conferma il relatore: \emph{diversity and inclusion}
viene spesso usata per dire ``così includiamo più clienti'', ma non è quello
il suo punto (``altro discorso'').

% ---------------------------------------------------------------------
\section{Confrontarsi con chi è più bravo}
\label{sec:z1-rivalita}

Arrivato a Riot dopo sei anni a Blizzard, il relatore pensava di avere un buon
bagaglio di esperienza. Si è trovato invece ``come un fante buttato in mezzo
ai Navy SEAL'': persone che avevano superato 11, 12, 13 colloqui, con un
livello molto alto.

Individua subito un collega francese che, nella sua testa, diventa ``la
versione migliore di me'': stesso ruolo (community manager), esperienza
simile (arrivava da Ubisoft, come quasi tutti i francesi di Riot), ma molto
più spigliato, sempre con la battuta pronta, ``l'anima della festa''. Il
relatore invece è introverso: un workshop come questo per lui è
\emph{emotivamente drenante}, non ama parlare in pubblico pur amando
condividere le proprie esperienze, e non riesce a sostenere conversazioni
come ``cosa hai fatto questo weekend?''. Lavorando raramente insieme, aveva
sviluppato una rivalità ``da allenatore di Pokémon'' (lui era Ash, il
relatore l'antagonista).

Dopo sei mesi, a un evento di \emph{team building} aziendale (a Tenerife), il
suo capo, molto attento a queste dinamiche, li mette apposta nello stesso
gruppo. A fine serata, complice l'open bar, il relatore gli confessa di averlo
mal sopportato dal primo giorno per invidia della sua estroversione. Il
collega ride e risponde che pochi giorni prima aveva detto al capo la stessa
cosa al contrario: invidiava a Stefano il suo essere sempre organizzato, sul
pezzo, presentarsi nella sala riunioni giusta (mentre lui finiva in quella
sbagliata), arrivare con gli appunti, con i dati a supporto e le analisi di
mercato, mentre lui spesso non aveva nemmeno letto la mail del meeting.
Un'invidia in due direzioni, che li ha portati a lavorare molto insieme e a
migliorarsi a vicenda.

Le lezioni:
\begin{itemize}
  \item non bisogna demoralizzarsi vedendo qualcuno che fa lo stesso lavoro
        ed è più avanti sotto alcuni aspetti: spesso c'è un bilanciamento. I
        ``punti caratteristica'' sono limitati: se li ha messi in carisma,
        non li avrà messi in saggezza;
  \item conviene \textbf{chiedere a queste persone come mai} sono così brave
        in qualcosa, come si fa con un progetto andato bene: magari non lo si
        farà mai bene quanto loro, ma si comincerà ad avere il
        \emph{mindset} giusto; e viceversa, loro saranno spinti a chiedere a
        noi;
  \item spesso \textbf{le proprie doti le scopriamo dagli altri}. Il relatore
        non si era mai accorto di essere ``quello che resta calmo durante le
        emergenze'' (il lancio di \emph{Diablo~III}, quello di
        \emph{Draenor}, i problemi dal vivo alla Games Week) finché non
        gliel'hanno fatto notare tutti: per lui era la normalità. A chi pensa
        ``non sono capace, forse non è il mio percorso'', può darsi che
        semplicemente nessuno gli abbia ancora detto ``questa cosa la sai
        fare davvero bene, concentrati su questo''.
\end{itemize}

% ---------------------------------------------------------------------
\section{Domande sul lavoro nell'industria}
\label{sec:z1-industria}

Diverse domande degli studenti hanno riguardato le condizioni di lavoro nei
grandi studi internazionali. Le risposte sono raccolte qui per argomento.

\subsection{Presenza, remoto e trasferimenti}

\paragraph{Grandi studi: quasi solo in presenza.} Sia a Blizzard sia a Riot si
lavorava \textbf{solo in presenza}, senza smart working: sei anni a Parigi e
cinque a Dublino. Si parlava solo inglese: anche nei meeting tra soli
italiani si tendeva a comunicare e a scrivere in inglese, per correttezza
verso chi si fosse aggiunto e per lasciare una traccia comprensibile a tutti.
Il francese serviva per fare la spesa e sopravvivere, non in ufficio.
L'ambiente era estremamente eterogeneo (francesi, polacchi, spagnoli e
molti altri): durante il colloquio dal vivo a Dublino, mentre si preoccupava
di essersi vestito in modo adeguato, vide passare dietro la parete a vetro un
collega svedese con i dreadlock, a piedi scalzi, con una pianta sotto un
braccio e un aquilone sotto l'altro, e capì di essere ``a posto così''.
Un'esposizione a molte culture e modi di pensare che il relatore considera
bellissima.

Anche oggi, tranne alcuni studi particolari, la maggior parte del lavoro si
svolge di persona: amici a Paradox lavorano a Stoccolma, i pochi rimasti a
Blizzard si sono dovuti trasferire a Irvine. Per posizioni \emph{junior} o
\emph{entry level} i grandi studi raramente concedono il lavoro da casa; per
profili senior può esserci qualche eccezione (ad esempio nove mesi da remoto
e tre in sede).

\paragraph{Relocation.} Spesso tocca spostarsi, ma il relatore non ha mai
pagato un trasloco di tasca propria: queste aziende offrono la
\emph{relocation} (aiuto a trovare casa, ad aprire il conto in banca e
supporto pratico), perché è un investimento sulla persona.

\paragraph{Colloqui mentre si lavora altrove.} Per i colloqui in presenza
prendeva ferie, senza dire dove andava. Il consiglio generale: per fare
colloqui \textbf{prendete ferie o permessi}, non datevi malati, ``che non è un
bel modo di comportarsi''.

\paragraph{Remoto ``doppio''.} Un collega lavora per l'azienda coreana che
gestisce \emph{PUBG}, con ufficio ad Amsterdam, ma lavora da Milano: un
``doppio remoto''. Si trova, ma in genere un ambiente anglofono con sedi in
Europa chiederà di trasferirsi in Europa.

\paragraph{Famiglia e trasferimenti.} Il relatore ha avuto la fortuna di
sposare una persona che fa lo stesso lavoro (conosciuta lavorando su
\emph{Warhammer Online}, poi con lui a Blizzard e a Riot, nella
localizzazione e nel doppiaggio, ad esempio di \emph{World of Warcraft} in
italiano): avendo competenze rare, spesso venivano assunti entrambi, al
massimo a pochi mesi di distanza. Per i developer la soluzione frequente è
che la famiglia segua la persona assunta, perché lo stipendio basta per
tutti. Ma non sempre c'è equilibrio, e a volte uno dei due deve sacrificarsi:
sua moglie ha lasciato Blizzard per seguirlo a Dublino (per poi essere
assunta da Riot); lui ha lasciato Riot per tornare in Italia quando lei ha
trovato lavoro qui. Ha dovuto rinunciare più volte a lavorare in Islanda per
CCP (\emph{EVE Online}), perché non era una scelta solo sua. Si fanno
compromessi, ma si riesce a far quadrare.

\paragraph{Ci si sposta per ogni progetto?} No: si sono spostati
esclusivamente per cambiare azienda. Cambiando progetto nella stessa azienda
è raro, quasi mai (``non l'ho mai visto succedere''), che venga chiesto di
trasferirsi. In Francia stavano benissimo (vivevano fuori Parigi; ``i
problemi con i francesi sono in realtà problemi con i parigini''), parlavano
francese ed erano vicini alla doppia cittadinanza: se ne sono andati perché
non vedevano un futuro a Blizzard, per il modo in cui l'azienda trattava i
colleghi.

\paragraph{Il mercato oggi.} Non esiste un \emph{silver bullet} per trovare
lavoro domani dove si vuole. Con tutti i \emph{layoff} recenti ci sono
migliaia di esperti del settore che cercano lavoro: serve pazienza, ma è
fattibile.

\paragraph{Studi indie e remoto.} Negli studi più piccoli il remoto è più
diffuso: a Stormind circa 60 persone su 95 lavorano da remoto, ma lo studio
si è strutturato fin dall'inizio per farlo. Serve consapevolezza. Bending
Spoons (che sviluppa app per mobile, non giochi, ed è l'azienda dell'app
Immuni), pur essendo grandissima e con sede a Milano, lavora da remoto e
lascia l'ufficio a chi vuole: lo smart working è posizionato quasi come un
\emph{benefit}, ma permette anche di avere un ufficio da 300 persone invece
che da 700. Le sedi fisiche costano decine di milioni all'anno (Blizzard ha
chiuso le sedi di Amsterdam e di Parigi). Tuttavia, per lavori molto
verticali e specifici gli studi preferiscono avere le persone dal vivo; i
ruoli più spesso remoti sono quelli che non fanno sviluppo (customer service,
PR\dots).

Anche gli artisti lavorano spesso in presenza: l'accademia può dare solo una
\emph{simulazione} di quanto stretto sarà il contatto tra programmazione,
design e art. ``Sarà ancora di più.''

\paragraph{La comunicazione da remoto.} È più difficile. Pur essendo
estremamente introverso, il relatore \textbf{obbliga tutti a usare la webcam}
nei meeting in cui conta: non in un aggiornamento con novanta persone
collegate, ma certamente in un colloquio a due sulla performance dell'ultimo
anno. La componente fisica è una parte enorme della
comunicazione\footnote{In lezione si parla di ``quasi l'ottanta per cento'';
il dato va preso come indicativo, poiché le percentuali di questo tipo
circolano in forme diverse e sono spesso citate fuori contesto.}: anche se si
è timidi, accendere la webcam è l'approssimazione più vicina all'essere dal
vivo. Negli ambienti dove si lavorava senza webcam la situazione era
``esponenzialmente peggiore''.

\subsection{Differenze culturali nel team}

\paragraph{Gestire sensibilità diverse.} A Riot i team erano divisi per
\emph{lingua}, non per nazione: il team di lingua spagnola includeva anche un
messicano, quello di lingua russa includeva russi e ucraini, proprio nel
periodo della guerra in Crimea. Tra i colleghi non c'era alcun problema (erano
amici), ma ci si è posti la domanda: \emph{come comunichiamo con una community
che sarà divisa su questo tema?} Nei giochi online la divisione era evidente:
in \emph{Rocket League} c'era chi, con un compagno di squadra con la bandiera
avversaria, si faceva otto autogol pur di farlo perdere. Allo stesso modo,
con molti americani mandati a Dublino a coprire posizioni scoperte, c'erano le
elezioni (Clinton--Trump, Biden--Trump) e varie guerre: serviva sempre una
certa sensibilità.

Imparare a navigare queste situazioni lo ha reso un professionista migliore:
gli ha insegnato a fermarsi e chiedersi ``io la percepisco così; gli altri
dieci, con background, storie e lingue diverse, come la penseranno?'', a
ragionare più \emph{long term} e più \emph{big picture}, con più empatia.

\paragraph{Stili di comunicazione.} Un partecipante che ha lavorato nel Regno
Unito chiede dello stile comunicativo italiano, a volte frainteso. Uno dei più
grandi complimenti ricevuti dal relatore è stato ``non sembri italiano''
(``gesticoli poco per essere italiano'', ``parli bene inglese per essere
italiano''). Ha imparato ad \textbf{adattare la comunicazione alla persona}:
con spagnoli e francesi la battuta sulla nazione funziona, perché c'è un
``odio--amore'' condiviso; con russi e giapponesi no. Doveva adattare anche
l'umorismo. Un collega metà russo e metà messicano diceva di aver preso ``il
lato peggiore di entrambe le nazionalità'', ed è stato un ottimo maestro per
imparare a leggere culture diverse.

Non bisogna generalizzare, ma gli stereotipi possono essere un punto di
partenza ragionevole, \emph{sapendo che troverete sempre le eccezioni}. E
vanno letti correttamente: gli americani tendono a essere più estroversi e ad
avere sempre qualcosa da dire non perché ``sono fatti così'', ma perché
vengono educati così (all'università tra i corsi c'è il \emph{debate}, saper
argomentare in pubblico); un collega giapponese può essere più introverso per
ragioni culturali. Nel dubbio si parte da lì, poi si conosce la persona.

\subsection{Negoziare, discriminazione e stipendi}

\paragraph{La famiglia come leva.} Avere un partner riduce le opzioni, ma le
aziende lo sanno: uno spostamento è per due. Si può usare come leva per
negoziare (un compenso diverso, un supporto per il trasloco): se un'azienda ha
deciso di volere quella persona, 2, 3, 4 o 10.000 euro in più all'anno, per
un'azienda che fattura centinaia di milioni o miliardi, non cambiano nulla.

\paragraph{Essere sposati è un motivo per non essere assunti?} No. In alcuni
paesi è proprio illegale discriminare su questo: in Francia era illegale
chiedere CV con data di nascita, foto, sesso o religione, e chiedere se si
vogliono figli. Purtroppo, dai racconti dei suoi studenti di master (in
settori come \emph{beauty and wellness} e \emph{food and wine}), in piccole
realtà italiane queste domande (``hai 25 anni, hai intenzione di
sposarti?'') vengono ancora fatte, e servono a filtrare i candidati. Nelle
grandi aziende ci sono sistemi per evitarlo; anzi, molte danno \emph{più}
benefit a chi ha famiglia (``porta pure la famiglia, è tutto pagato''),
perché le persone competenti hanno spesso una certa età e quindi,
ragionevolmente, una famiglia.

\paragraph{Gli stipendi.} Dipendono molto da dove si vive. Lo stipendio a Riot
non era paragonabile a quello di Blizzard perché in Francia era tassato al
52\% e in Irlanda al 17\%\footnote{Aliquote come riportate in lezione,
riferite alla situazione personale del relatore.}: nessuno stipendio
competitivo può compensare il 35\% di differenza di tasse. Negli Stati Uniti,
con pochissima previdenza pubblica, gli stipendi sono molto più alti perché
poi ci si paga tutto (``se hai un raffreddore paghi tu''). Bisogna valutare il
potere d'acquisto di quei soldi nel paese in cui si andrà a vivere.

\subsection{Freelance e \emph{work for hire}}

All'inizio della sessione pomeridiana uno studente che lavora nell'audio
chiede quali siano le condizioni per chi lavora come \emph{freelance}.
\begin{itemize}
  \item Dipende molto dal progetto: spesso il contratto è a progetto, per
        portare a termine una parte o un intero progetto.
  \item Non tutte le aziende usano freelance: Riot e Blizzard lo facevano
        pochissimo e preferivano assumere (magari trattando poi la persona
        ``come un freelance'', lasciandola a casa a fine progetto).
  \item Spesso i lavori offerti come freelance sono in realtà
        \textbf{sostitutivi di un full time} (``freelance, ma sei ore al
        giorno'').
  \item Si può fare, ma bisogna essere molto chiari nel definire l'impegno
        (quante ore, come vengono tracciate): non tutti specificano cosa
        significa davvero ``freelance''. Lavorando così per due o tre aziende
        si può arrivare a un tempo pieno.
  \item È più comune il \emph{work for hire}: aziende che forniscono personale
        ad altre aziende; si è assunti da X ma si lavora ``embeddati'' presso
        Y (una sorta di subappalto). Anche lo studio del relatore lo fa.
  \item Il puro freelance di quattro mesi esiste, ma di solito per quei
        quattro mesi diventa un lavoro a tempo pieno; è raro ``quattro mesi,
        due ore al giorno''. Nelle grandi aziende è improbabile lavorare nello
        stesso mese per Ubisoft, Riot e Blizzard tre ore ciascuna; è più
        realistico lavorare per più startup indie italiane, che non possono
        permettersi un'assunzione ma possono pagare due ore al giorno.
  \item Molto passa dai contatti: trovare le aziende giuste e rendere esplicite
        le condizioni.
\end{itemize}

% ---------------------------------------------------------------------
\section{Il circolo virtuoso: feedback, miglioramento, fiducia}
\label{sec:z1-triangolo}

\subsection{Il triangolo}

L'obiettivo a cui il workshop vuole arrivare è un \textbf{triangolo che si
autoalimenta}:
\begin{enumerate}
  \item si dà un \textbf{feedback} con l'obiettivo di migliorare le
        performance (inclusi elementi caratteriali);
  \item la persona \textbf{migliora} e vede di averne beneficiato;
  \item questo genera \textbf{fiducia}, che rende più facile dare e ricevere
        il feedback successivo.
\end{enumerate}
Togliendo uno qualsiasi dei tre elementi, il circolo non è altrettanto forte.

In qualsiasi fase si parta, il rispetto per i colleghi è inizialmente dettato
dall'\emph{educazione}, non dalla fiducia. La fiducia arriva quando i colleghi
cominciano a dire cose giuste e sensate e a \textbf{mantenere la parola}.
Dare e ricevere feedback è il modo in cui, in un certo senso,
``inganniamo il cervello'' per creare una \emph{case history} di fiducia: si
ha motivo di fidarsi di più, di ascoltare di più, di essere più ricettivi. E
si diventa sempre più bravi ad accettare il feedback e ad agire di
conseguenza.

\subsection{Accettare un feedback non vuol dire fare come dice l'altro}

Dieci anni fa il relatore non era ostile ai feedback, ma molto
\emph{difensivo}: ``sì, sì, va bene, ora ci guardo''; oppure, a chi gli diceva
``avresti dovuto farla così'', rispondeva spiegando perché l'aveva fatta in
quel modo. Oggi ha imparato a dire ``ok, grazie, ora ci ragiono'': non deve
difendere la propria posizione da chi viene solo a dire ``hai pensato a
questa cosa?''. Prima lo viveva come un attacco (``mi viene a dire che non so
fare il mio lavoro''), ora come un ``ottimo punto di vista, lo prendo in
considerazione''.

Con una postilla fondamentale: \textbf{accettare il feedback non significa
fare come dice l'altro}. ``Mi hai detto che preferisci questo mob verde. Ho
pensato anche a farlo verde? Sì. Penso comunque che giallo sia meglio? Sì, lo
faccio giallo.'' Il relatore accetta \emph{sempre} il feedback, valido o no:
garantisce di tenerlo a mente quando lavora. Che poi faccia come gli viene
detto è un altro discorso.

\subsection{Fiducia e scadenze}

La fiducia funziona allo stesso modo con gli impegni. Se si chiede a qualcuno
una cosa per venerdì e la consegna venerdì, la volta dopo gliela si chiederà
per venerdì. Se invece arriva il giovedì della settimana dopo, la volta dopo
gli si darà una \textbf{deadline anticipata}, come con chi arriva sempre in
ritardo agli appuntamenti (``se ti dico alle 7 arrivi alle 7 e mezza; la
prossima volta ti dico alle 6 e mezza'').

Al lavoro questo è dannosissimo, perché crea \textbf{deadline fittizie} che
non sono mai quelle vere, e nascono discussioni del tipo ``perché devo farlo
in due giorni se si fa in cinque?'' (``sì, se tu lo facessi in cinque''). La
conversazione giusta è un'altra: \emph{``che problemi hai avuto nel
consegnare venerdì?'' ``Questo, questo e questo.'' ``Cosa posso fare per
evitare che succeda?''}

Il relatore non vuole mai sentirsi dire ``non sono riuscito a farlo''.
Preferisce di gran lunga sentirsi dire, tre giorni prima, \textbf{``potrei non
riuscire a farlo, per questi motivi''}. Chi ha la visione d'insieme può dare
più budget, più persone, spostare la deadline, cambiare lo scopo del progetto:
può fare qualsiasi cosa per aiutare. Ma se lo si scopre il venerdì, cinque
minuti prima della deadline, non può fare più niente. Il messaggio è:
\textbf{``aiutami a poterti aiutare''}. Non penserà mai male di chi dice ``non
ce la faccio per questi motivi'', perché non sta dicendo ``non ho voglia'': sta
indicando i problemi, e lavorare su quei problemi è compito del manager. Non
può invece fidarsi di chi fino al giorno prima dice ``tranquillo, ce la
faccio'' e il giorno dopo ``non ce l'ho fatta''.

Questo è fondamentale per chi lavora in programmazione, design o art, perché
ci saranno sempre deadline: bisogna sempre \textbf{lavorare sulle aspettative},
dire al proprio manager (o al proprio report) cosa è lecito aspettarsi. Non
bisogna fare come i pinguini di \emph{Madagascar} (``sorridete e salutate, va
tutto bene''): se c'è un problema, va detto, altrimenti si peggiora solo la
propria situazione.

Da italiani si tende a non voler mai alzare la mano per chiedere aiuto
(``no, no, ce la facciamo, lo faccio io''). A volte non ce la si può fare, non
perché non si è capaci, ma perché vari fattori lo impediscono. Accettarlo è
importante: \textbf{è molto meglio avvisare e chiedere una mano che dire
``non ce l'ho fatta''}.

\subsection{Domande sul dare e ricevere feedback}

\paragraph{``Ne parliamo a fare, se poi fai come vuoi?''} Uno studente (Cristian)
riceve feedback che, dopo averli analizzati, non ritiene congrui; quando non
li segue gli viene detto ``ne parliamo a fare?''. C'è un problema di
\textbf{aspettative}: l'idea che accettare un feedback voglia dire fare come
dice l'altro. Accettarlo significa prendere in considerazione ciò che è stato
detto. Chi decide diversamente ha però un \textbf{onere}: argomentare perché
si fa comunque a modo proprio (``mi hai detto verde, la facciamo gialla
perché\dots''). Poi bisogna solo aspettare i risultati per vedere chi aveva
ragione.

La decisione spetta sempre \textbf{a chi ha la responsabilità di quella
decisione}. Se il game director deve scegliere tra A e B, anche se il
feedback dice B, la scelta è sua perché la responsabilità è sua. Non sarebbe
corretto essere ritenuti responsabili di una scelta delegata a chi dà
feedback. La risposta può essere semplicemente: ``l'ho tenuto in
considerazione, ma visto che ne sono responsabile io e le conseguenze sono
mie, ho scelto di fare X''.

\paragraph{E se l'altro fa ripicca?} Un partecipante chiede di chi, vedendo
il proprio feedback non accolto, si vendica ignorando a sua volta feedback
buoni. Per il relatore è un problema caratteriale: chi si comporta così ``è
già sulla strada per uscire''. Non c'è spazio per l'ego in team di questo
tipo. Il fenomeno nasce dall'idea che \emph{il feedback negativo faccia fare
brutta figura}, che il fatto che si potesse fare meglio tolga qualcosa alle
proprie capacità.

Non è così: \textbf{il feedback non è un giudizio}. Si può passare un anno a
Riot ricevendo solo feedback negativi e avere una valutazione annuale
stellare, perché ricevere molti feedback negativi significa essere una persona
approcciabile, che agisce sul feedback e invita gli altri a darlo. Il relatore
ha visto valutazioni molto peggiori per chi riceveva solo feedback positivi,
perché la gente aveva paura di dargliene di negativi: ma solo ``pacche sulle
spalle'' non fanno migliorare. Il problema va affrontato parlando con la
persona; a volte le conversazioni iniziano con \textbf{``ti do un feedback sul
feedback''}, cioè sul modo in cui l'altro ha dato il suo.

Esempi dalla sua esperienza:
\begin{itemize}
  \item lui chiede spesso di \textbf{non ricevere feedback senza esempi
        tangibili}, perché fa fatica ad astrarre: ``sei stato maleducato'' non
        significa nulla, ``ieri al meeting mi hai risposto male su questa
        cosa'' sì;
  \item ad altri serve l'opposto: ha ricevuto il feedback ``quando devi dirmi
        che ho sbagliato, puoi girarci un po' intorno e darmi del contesto?'',
        perché il suo stile diretto (``questa cosa non ha funzionato, dovremmo
        fare così'') feriva. Ha adattato il suo modo di dare feedback: con
        alcuni è diretto, con altri gira intorno. L'importante è saperlo.
\end{itemize}

\subsection{Non portare la soluzione}

Un consiglio centrale: \textbf{quando si dà un feedback, non si dà una
soluzione}. Non si dice ``io farei così'', ma ``secondo me questa cosa si
potrebbe fare diversamente, per questo e questo motivo: come potremmo farla
diversamente?''. È come chi smette di fumare perché vuole e chi smette perché
glielo dicono moglie e figli: se non si è convinti, spesso non si smette.

Quando il relatore diceva ``c'è questo problema, io farei così'', il meeting
successivo iniziava inevitabilmente con ``ci ho provato come hai detto tu, con
me non ha funzionato''. Sempre. Perché la soluzione era sua. La
responsabilità della soluzione deve essere invece \textbf{congiunta}: ogni
feedback deve creare un piccolo team che ha come obiettivo risolvere quel
problema. Questo toglie l'antagonismo (non ``ti dico che non sai fare il tuo
lavoro'', ma ``ti voglio aiutare a farlo meglio, lavoriamoci insieme'') e
sposta il focus: il problema non sei tu, è ciò che sta accadendo.

\paragraph{E se sono i piani alti a tenere all'oscuro dei problemi?} È una
questione di cultura aziendale. Il relatore fa sempre presente al proprio capo
il concetto di \textbf{service level agreement} (che verrà ripreso nella
seconda giornata): ``se mi dai questi input, ti garantisco questo output in
questi tempi''. Se il capo chiede delle cose senza dire budget, tempi e
perché, non può aspettarsi lo stesso risultato. È una comunicazione a due
vie: si dice al manager cosa serve per fare il lavoro, e lui deve fornirlo o
accettare che il risultato non sarà quello voluto.

\subsection{Amici, parenti e parole proibite}

\paragraph{``Non si lavora con amici e parenti''?} Per sua definizione
personale il relatore ha pochissimi amici (``stanno sulle dita di una mano'')
e una soglia molto alta per considerare amico un collega. Allo stesso tempo,
in quindici anni e con centinaia di colleghi, ha avuto solo \textbf{due}
esperienze in cui non è riuscito a lavorare con qualcuno: non serve essere
amici per far funzionare le cose. Con gli amici, anzi, gli viene più facile:
se dice a un amico che una cosa ``fa schifo'', l'amico non pensa ``come si
permette'', ma ``se perfino Stefano, che è mio amico, me lo dice, allora è
proprio così''. Diventa un rafforzativo, grazie alla fiducia costruita prima.
Il problema del lavorare con amici e parenti nasce dal \emph{non voler} dire
all'amico che sta sbagliando; se invece la confidenza permette di dirlo, si
può saltare la ``confezione'' del messaggio e passare direttamente al
contenuto.

\paragraph{Le frasi peggiori.} Alla domanda se esistano parole peggiori di
``ottimo lavoro'', il relatore indica la sua classifica: ``ottimo lavoro''
(detto in modo vuoto)\footnote{Il passaggio della trascrizione è
frammentario: si intende verosimilmente il complimento generico e di
circostanza, che non dice nulla di utile.}; poi la frase che lo ``fa uscire di
testa'', \textbf{``abbiamo sempre fatto così''} (``e infatti avete chiuso'',
come Blockbuster); infine \textbf{``siamo una famiglia''}, riferito al posto di
lavoro. Un'azienda non è mai una famiglia, perché non ci si può aspettare lo
stesso impegno da chi la possiede e da chi ci lavora: non c'è un
corrispettivo proporzionale tra impegno e ritorno. Per quanto inclusiva,
flessibile e razionale, non sarà mai una famiglia.

\subsection{Ego, muri e il peggior modo di scegliere un leader}

\paragraph{Chi mette davanti l'ego e alza muri.} Esistono, ma di solito durano
pochissimo in azienda. Non si tratta di avere un ego forte, ma di \textbf{non
saper ricevere feedback}. Spesso sono persone bravissime nel loro lavoro che
poi si chiedono perché tutti intorno diventino più bravi di loro: perché gli
altri accettano feedback e hanno più fonti da cui attingere, mentre loro si
sono rintanati nella loro torre.

Molti manager tendono ad \textbf{assumere persone un filo peggiori di loro},
per non fare brutta figura. Il relatore assume solo persone più brave di lui,
``perché non è una gara'': ha lavorato con molti community manager più bravi
di lui come community manager, ma lui era più bravo a fare il manager. Non
viene valutato su quanto sa fare il community manager, ma su quanto il suo
team sa farlo. Un partecipante chiede se la paura di avere qualcuno di
migliore esista anche in ambito accademico. Assolutamente sì: al relatore è
capitato di non essere preso come docente in un corso di formazione sui
videogiochi perché il responsabile aveva percepito che ne sapesse più di
lui\footnote{La trascrizione dice ``che non sapevo più di lui'', ma il senso
del passaggio richiede il contrario.}. ``Ma perché te ne importa? Tu sei quello che ha portato chi
sa fare il suo lavoro: ben venga.'' O si ha questa cultura, o si ha paura per
il proprio posto.

\paragraph{Promuovere il migliore.} Quando bisogna scegliere un leader per un
team (programmazione, design, art), molto spesso si sceglie \textbf{il migliore
tecnico}: il miglior programmatore, il miglior designer. È spesso la scelta
sbagliata: è come prendere un idraulico per fare l'elettricista. Fare il
\emph{people manager} non è come fare il designer o il programmatore; si può
essere un ottimo programmatore e un pessimo manager, e viceversa. Le aziende
l'hanno capito e lo fanno meno di prima. A Blizzard c'era un grave problema di
\emph{middle management}: molte persone in ruoli importanti erano state
promosse cercando le skill sbagliate, erano frustrate dal non saper fare il
proprio lavoro e schiacciavano chi stava sotto per paura di fare brutta
figura.

Un partecipante cita una teoria simile: una persona viene promossa finché non
raggiunge \emph{il proprio livello di incompetenza}, dove si blocca\footnote{È
il \emph{principio di Peter} (Laurence J.\ Peter, 1969), citato in lezione
come ``la teoria di Peter''.}. Un altro cita Michael Scott di \emph{The
Office}, su cui il relatore tornerà (sezione~\ref{sec:z1-quando}).

% ---------------------------------------------------------------------
\section{Perché la cultura conta: retention, costi e fiducia}
\label{sec:z1-cultura}

\subsection{L'ambiente di lavoro e la retention}

Perché tutto questo è così importante? Perché l'ambiente di lavoro conta
moltissimo per la \textbf{retention}, cioè per trattenere le persone e farle
lavorare in modo efficiente. Se contasse solo lo stipendio, molti non
lascerebbero mai Riot o Blizzard; e più si va avanti con l'età, più si impara
a dare valore anche al resto.

Inoltre, secondo gli studi citati nelle slide, chi vive il proprio lavoro in
modo negativo è \textbf{meno efficiente nel problem solving}: l'energia
mentale va verso ``come faccio a non farmi dare la colpa'' o ``come faccio a
fare bella figura'', non verso la soluzione del problema. In ambienti negativi,
dove ci si deve preoccupare della propria ``sopravvivenza'', si lavora
peggio.

\subsection{Il costo di chi se ne va}

Ogni volta che qualcuno si dimette o viene licenziato, il costo di trovare
una nuova persona, fare i colloqui, formarla e farle attraversare forming,
storming, norming\dots\ è circa \textbf{un terzo del suo stipendio}. Per questo
le aziende tendono a fare più \emph{retention} che \emph{acquisition}, almeno
con le figure considerate utili alla produzione.

Allo stesso tempo si vedono sempre più spesso \emph{layoff} e \emph{hiring
wave}: si licenziano 3.000 persone e sei mesi dopo se ne riassumono 2.500 o
3.500, ma tutte più junior. Il relatore ne ha scritto in un articolo su
Eurogamer.it, in cui spiega anche alcuni ``trucchi del mestiere'': perché i
licenziamenti avvengono sempre in certi periodi e sempre ``a gruppi'' (prima
Bethesda, poi Blizzard, poi un'altra\dots).

\paragraph{L'offerta di Riot.} A Riot esisteva una pratica (di cui il relatore
non ricorda il nome) per cui, firmato il contratto, si aveva \textbf{un anno
di tempo per dire ``non mi piace, non me la sento, me ne vado''} e ricevere
subito come buonuscita un anno di stipendio. Avevano stimato che fosse molto
meglio pagare subito un anno a qualcuno e lasciarlo andare che tenerlo due,
tre, quattro anni come \emph{underperformer}: la responsabilità della scelta
di assumere era dell'azienda\footnote{La pratica, nota come ``The Offer'', è
descritta in fonti pubbliche con condizioni diverse (un periodo più breve e
una somma pari a una frazione dello stipendio annuo). Si riportano le
condizioni come raccontate in lezione; andrebbero verificate.}. L'hanno usata
pochissimi, ma il solo fatto che esistesse faceva lavorare meglio anche lui da
manager, perché sapeva che chi non si trovava bene poteva andarsene con un
anno di tempo per trovare un altro lavoro.

\paragraph{Il ricircolo.} C'è un forte ricircolo nell'industria: molti ex
colleghi passano da Ubisoft a Paradox, da Paradox a Funcom, da Funcom a
Blizzard. Il relatore ha scelto di uscirne lavorando solo per aziende
\textbf{indipendenti}: Stormind non è un indie di cinque persone, ma è
indipendente, e non deve accontentare azionisti. Quando vede un'azienda con
utili miliardari licenziare migliaia di persone si chiede ``dov'è la linea''
(\emph{when enough is enough}): con azionisti che vogliono un ritorno, se
quello è l'obiettivo stanno facendo la cosa giusta, ma lui non condivide
l'obiettivo\footnote{In lezione si parla di ``6,5 miliardi di utile'' e di
``3.000 persone licenziate''; le cifre sono riportate come dette e non sono
state verificate.}.

\paragraph{L'open space vuoto.} A Riot, durante un'ondata di licenziamenti,
del suo team di trenta persone è rimasto solo lui: ``teniamo te come
coordinatore, poi ricostruirai il team''. Ha rifiutato di ricostruire un team
cancellato ``con un colpo di spugna'' e ha deciso che non avrebbe più lavorato
per una grande azienda. Nessuno aveva pensato di spostarlo, e ha lavorato per
tre mesi in un open space di cento metri quadri con la sua scrivania e altre
ventinove vuote. Su una lavagna aveva scritto i nomi dei suoi ragazzi e li
cancellava man mano che, anche grazie ai suoi \emph{referral}, trovavano
lavoro. Cancellato l'ultimo nome, si è dimesso: per lui il lavoro era finito.

\paragraph{Team che si dissolvono.} Uno studente cita un articolo su Larian
Studios, il cui management è contrario a ricostruire ogni volta il team a fine
progetto. È lo stesso problema: si pensa ai ``giochi Blizzard'' o ai ``giochi
BioWare'', ma le persone che li hanno fatti non ci sono più; è rimasto solo il
brand, e non ha senso avere le stesse aspettative da un gioco fatto da persone
diverse. E anche chi sopravvive ai layoff vive pensando che il prossimo
potrebbe essere lui: alcune persone, licenziate nell'ultima ondata a
Blizzard, gli hanno detto ``erano due anni che lo sapevo''. Vivevano a Cork,
``un paesotto irlandese dove o vai per lavorare a Blizzard o non ci vai'',
avevano comprato casa: non ci si può aspettare chissà che morale.

\subsection{La cultura è forte quando si lavora per gli altri}

\begin{quote}
  \emph{La cultura è forte quando le persone lavorano per gli altri e a favore
  degli altri, e non per se stesse.}
\end{quote}

Se si dà feedback per far migliorare l'altro e non per imporre il proprio
modo, se si aiuta per il bene del progetto e non per fare bella figura, la
cultura si rafforza. È dai \textbf{piccoli gesti} che emerge la cultura
aziendale: ad esempio, dal non dover chiedere ``perché mi stai aiutando?''.

\paragraph{L'incognita.} \textbf{L'incognita è spesso dannosa quanto le
cattive intenzioni.} In un esperimento (``mi pare a Stoccolma, cinque o sei
anni fa'') si distribuivano per strada banconote da 500 euro, e la maggior
parte delle persone le rifiutava chiedendosi ``dov'è la fregatura?''
\footnote{L'esperimento è riportato in modo vago in lezione e non è stato
possibile identificarlo con certezza.}. Lo stesso accadeva al relatore,
arrivato a PG e a Stormind da Riot: a ``vedo che sei in difficoltà, posso
aiutarti?'' la risposta era sempre ``perché?''. Alla fine rispondeva: ``perché
quando sarò in difficoltà io vorrei che tu mi aiutassi''.

Quando si comincia a lavorare per gli altri e si smette di pensare ``così avrò
la promozione, così avrò l'aumento'', promozioni, aumenti e rispetto arrivano
da soli. A Blizzard aveva fallito tre volte i colloqui per essere promosso a
manager (inteso come coordinatore del team); quando ha deciso di smettere di
pensarci e di fare solo il bene del proprio lavoro, la volta dopo è stato
promosso. Funziona così: chi aiuta i colleghi e fa andare bene il progetto
viene notato (``questa persona fa migliorare cinque persone, non solo se
stessa: vale di più'').

\subsection{Il lavoro è valutato come opera d'insieme}

Il relatore riconosce che buona parte del suo successo è venuto dal fatto che
il suo team lavorava bene, non dalla sua bravura personale. Anche i progetti
degli studenti, presentati a Diego, a Enrico, ad Andrea, verranno valutati
\textbf{come opera d'insieme}. Nessuno dice ``il \emph{combat system} di questo
gioco l'ha fatto X, che bravo'': si dice ``BioWare ha fatto un bel gioco'',
``Blizzard ha fatto un bel gioco''.

Lo ha replicato anche lato HR. In molte aziende ci sono \textbf{bonus
personali} legati a obiettivi (se il gioco vende X, se la feature è pronta in
tempo). Il relatore li ha dimezzati: metà restano personali, metà diventano
\textbf{bonus di team}, che tutti ricevono se il team raggiunge gli obiettivi.
Un incentivo concreto a collaborare: non solo si dice che il lavoro di gruppo è
importante, lo si dimostra dandogli più valore economico. Aveva proposto il
100\% di bonus di team, ma gli è stato concesso il 50 e 50; ha funzionato
comunque, anche se sarebbe stato ancora più potente. Ovviamente, alla
domanda ``perché mi aiuti?'' si può anche rispondere ``perché voglio il
bonus'': un aiuto dato ``per il motivo sbagliato'' era un compromesso che era
disposto ad accettare.

\subsection{Il coraggio di chiedere aiuto}

Un partecipante osserva che aiutarsi dovrebbe essere la norma. Il modello
diffuso, invece, è ``devi essere forte e indipendente, non chiedere mai una
mano'', e vale in tutte le direzioni (dal maschilismo alla ``donna forte e
indipendente che non ha bisogno di nessuno''): \textbf{non ammettere di avere
bisogno di aiuto è sempre uno dei motivi per cui le aziende vanno male}, che si
tratti di uomini, donne o persone non binarie. Il problema è non avere la
sicurezza di poter alzare la mano.

Il relatore ha cercato di togliere ogni conseguenza dall'aver alzato la mano:
se si \emph{depenalizza} qualcosa, le persone non sono più disincentivate a
farla. Ha detto espressamente che non avrebbe mai valutato nessuno per il
numero di volte in cui chiedeva aiuto o per il numero di feedback negativi
ricevuti, e questo ha funzionato.

\subsection{Perché si dà un feedback}

\emph{Perché avete dato il vostro ultimo feedback? Che cosa volevate ottenere?}
Le risposte della classe:
\begin{itemize}
  \item per \textbf{riconoscere il valore} di un collega, anche verso il
        manager, perché qualcuno che non aveva visibilità vedesse quanto
        quella persona fosse efficace;
  \item per \textbf{spingere una persona a superarsi};
  \item per aiutare un collega a rendere il prodotto all'altezza di ciò che
        voleva rappresentare;
  \item per \textbf{genuinità}: c'era qualcosa che poteva essere fatto meglio.
\end{itemize}
Il relatore aggiunge che la forma cambia col tempo: all'inizio era ``devo
darti un feedback, hai dieci minuti?''; ora è un messaggio informale su
Discord, perché ha costruito un rapporto di fiducia e sa che non verrà letto
come ``l'HR manager mi scrive, sono nei guai?'', ma come ``Stefano è così, mi
dà feedback, resterà tra noi e lo risolveremo insieme''.

Un altro motivo è \textbf{risolvere incompatibilità caratteriali}: le
conversazioni più difficili, del tipo ``io e te non stiamo funzionando, cosa
possiamo fare diversamente?''. Sono importanti perché, se due responsabili non
vanno d'accordo, il team se ne accorge, comincia a raccontarsi storie (``si
odiano, ci saranno problemi'') e, per difendere il proprio manager, comincia a
giudicare male il lavoro dell'altro gruppo. È la dinamica da fermare:
\textbf{valutare gli altri solo per difendere la propria fazione}.

Dare e ricevere feedback sono \textbf{due muscoli da allenare} (il relatore è
sempre stato più bravo a darlo che a riceverlo), e si impara col tempo. È come
fare colloqui: fare colloqui è una skill a parte, che non indica quanto si è
bravi nel proprio lavoro. Essere scartati non significa non saper fare quel
lavoro: spesso significa aver fatto male il colloquio o che è stato preferito
qualcun altro.

\paragraph{Feedback pubblici e privati.} Uno studente racconta di aver dato un
feedback negativo a un superiore, senza un motivo particolare, per dimostrare
a sé e ai colleghi che non si viene licenziati per questo. Il relatore coglie
l'occasione per precisare che \textbf{i feedback negativi si danno
individualmente}: lui dà moltissimi feedback di cui il resto dell'azienda non
sa nulla. Reagirebbe molto male se, in un meeting, qualcuno alzasse la mano per
dire ``comunque questa cosa non l'hai fatta bene'': quello non è feedback, è
discussione, accettabile solo se è un brainstorming e se l'opinione è stata
chiesta. Alcune aziende hanno piattaforme pubbliche di feedback, ma spesso
servono solo ad applaudire; lo stesso accade nelle aziende che
\emph{obbligano} a lasciare un certo numero di feedback a settimana, che
diventano ``lui è bravo, lei è competente, anche questa settimana l'abbiamo
fatto''. Bisognerebbe chiedere a chi le adotta cosa pensa di ottenere (forse
trasparenza), ma difficilmente una persona ``attaccata'' pubblicamente reagisce
bene.

% ---------------------------------------------------------------------
\section{Among Us: la caccia al traditore}
\label{sec:z1-among-us}

Prima della pausa pranzo il relatore mostra tre film di fantascienza di
qualche decennio fa e chiede che cosa li accomuni\footnote{I titoli non sono
leggibili nella trascrizione: viene citato \emph{Event Horizon} da uno
studente, e il film ``in mezzo'' è indicato come ``l'originale'',
verosimilmente \emph{La cosa} (\emph{The Thing}).}. In ognuno c'è un gruppo di
persone con un obiettivo e la presenza, reale o percepita, di qualcuno che
rema contro: \textbf{l'energia del gruppo viene spesa per trovare il
responsabile}.

\paragraph{L'esperimento degli alieni.} Nei master in presenza il relatore
propone un esercizio. Distribuisce segretamente cartellini plastificati con
scritto ``essere umano'' o ``alieno''; le carte alieno contengono istruzioni di
comportamento con cui gli alieni, che cercano di fingersi umani, si
tradirebbero (dare sempre ragione a chi ha parlato prima, finire ogni frase
con una domanda, parlare in un certo modo). Lo scopo: in dieci minuti formare
gruppi, e \textbf{vince il gruppo più numeroso senza alieni}; basta un alieno
per perdere. Si formano sempre gruppi piccoli (due, cinque, sei persone),
perché la tendenza è ``mi fido al massimo di voi due''.

Ma vincono i gruppi più grandi, per un motivo semplice: il relatore
\textbf{non distribuisce nessuna carta alieno}. Mostra a tutti entrambe le
carte, ma consegna solo carte ``umano''. Vuole simulare ciò che accade in
ufficio: la tendenza a percepire sempre qualcuno che rema contro (i manager,
l'altro team, quel collega), quando in realtà non c'è nessuno, ``è tutto nella
nostra testa''. Quando chiede le motivazioni, emergono sospetti
\emph{fabbricati}: ``faceva sempre domande'', ``parlava solo all'infinito''
(``no, parla così e basta''). Ci si fabbrica da soli i motivi per non fidarsi
degli altri, ed è comunissimo.

\paragraph{Among Us.} L'autosabotaggio è, per tornare ad \emph{Among Us},
\textbf{concentrarsi sul cercare cadaveri e traditori invece di fare le
task}: non si vince perché non si sono completate le task, mentre si cercava il
colpevole. Succede tantissimo negli uffici, ed è il motivo per cui l'80\% dei
problemi che il relatore risolveva erano di comunicazione. Entrando nei team
sentiva immancabilmente: ``facciamo il possibile, ma questa lotta continua con
l'altro gruppo è estenuante''. ``Perché c'è questo problema?'' ``Non lo so,
c'era già quando sono arrivato.'' Un paradosso: c'è un problema, nessuno sa
dire quale sia, ma c'è sempre stato e quindi lo si accetta.

\paragraph{L'esperimento delle scimmie.} È l'esempio, ``finto o reale'', delle
scimmie in gabbia\footnote{È il noto racconto delle ``cinque scimmie'', spesso
presentato come esperimento scientifico ma probabilmente apocrifo (lo stesso
relatore lo definisce ``finto barra reale''); nella trascrizione il numero
delle scimmie varia.}. In una gabbia con alcune scimmie pende una banana
elettrificata: chi la tocca fa prendere la scossa a tutte. Le scimmie imparano
e, quando una si avvicina alla banana, le altre la picchiano. Poi si
sostituisce una scimmia alla volta con una nuova: la nuova va verso la banana,
le altre la picchiano, e smette. Si continua finché nella gabbia ci sono solo
scimmie nuove, poi si toglie la corrente. Nessuna scimmia tocca la banana:
nessuna ha mai preso la scossa, ma tutte hanno imparato che chi tocca la
banana prende le botte.

Nei team è lo stesso: il motivo per cui si ``odia'' l'altro team si
trasforma, ma un vero motivo non c'è. Ne derivano la sfiducia, il non alzare
mai la mano per chiedere aiuto, il non dire mai ``bel lavoro'' agli altri
perché non sia mai che qualcuno ti veda fare un complimento ``a quelli''.
Risolvere queste dinamiche è lungo e tedioso, e sempre pieno di ``non mi
ricordo'', ``io non c'ero'' e ``abbiamo sempre fatto così'', la cosa più
frustrante per chi cerca di risolverle.

% ---------------------------------------------------------------------
\section{Il principio del default to trust}
\label{sec:z1-default-trust}

\subsection{La definizione}

Il relatore indica questa come la slide più importante della presentazione.
\textbf{Default to trust} è uno dei valori aziendali di Riot Games, che lui ha
tradotto, adattato e ``esportato'':

\begin{quote}
  \emph{Approccio ogni situazione di conflitto partendo dal presupposto che una
  persona abbia agito in un certo modo perché, con le informazioni a sua
  disposizione in quel momento, ha pensato che fosse la soluzione migliore, e
  non perché sia un'incapace o volesse danneggiare il progetto
  intenzionalmente.}
\end{quote}

In altre parole: \textbf{di default, mi fido}.

\subsection{L'esempio del ritardo al daily stand-up}

A Riot il relatore aveva un \emph{daily stand-up} dal vivo alle 15 con il suo
team di venti persone. Una ragazza arrivava regolarmente alle 15:15, per una,
due, tre settimane. Alla quarta ci ha parlato. \emph{Come iniziereste la
conversazione?} Le proposte della classe:
\begin{itemize}
  \item sottolineare l'importanza di esserci dall'inizio, perché potrebbe
        perdersi qualcosa che le interessa, e per rispetto dei colleghi
        presenti;
  \item ``devi arrivare in orario'', spiegando perché è importante;
  \item \textbf{chiedere perché}.
\end{itemize}

Il relatore osserva, anche dalle sue consulenze fuori dai videogiochi, che
\textbf{più si sale con la seniority}, più la risposta è ``devi smettere di
arrivare in ritardo'', spesso con un \emph{or else} finale (``altrimenti
succede quest'altra cosa''). Più si scende con l'età, e quindi con meno
abitudini consolidate, più la prima domanda diventa ``dimmi perché arrivi in
ritardo''. Già spiegare perché è importante è ``miglia avanti'' rispetto al
semplice ``smetti''.

Partire dal perché gli ha permesso di scoprire la causa: il meeting precedente
della ragazza, in un'altra zona del campus a un quarto d'ora di distanza,
finiva sempre tardi, e lei doveva attraversare il campus di corsa. Non aveva
la confidenza per dire all'altro manager, che insisteva perché si restasse
fino alla fine, che doveva andare; sacrificava quindi il meeting del relatore,
sapendo che era il più comprensivo dei due.

Il relatore si è offerto di \textbf{mediare}: ha detto all'altro manager che
o anticipava l'inizio del suo meeting, così da finire comunque in tempo,
oppure lasciava uscire la ragazza, perché quel meeting gli serviva. C'è anche
una questione di \textbf{economia di scala}, spesso sottovalutata: se venti
persone aspettano qualcuno per dieci minuti, si sono sprecati dieci minuti per
venti. Per lo stesso motivo il workshop è iniziato alle 10:05: iniziare alle
10:10 avrebbe significato far aspettare cinquanta persone cinque minuti in
più. A Riot c'era una scrivania vuota con un pupazzo gonfiabile, \textbf{Mr.\
Late}: sommando tutti i minuti di ritardo a meeting, call e conversazioni, in
un anno veniva fuori praticamente \emph{una persona a tempo pieno}. Serviva a
ricordare che anche cinque minuti contano, e che bisogna rispettare gli altri
arrivando in orario o almeno avvisando.

\paragraph{Gli effetti.} La conversazione era avvenuta nel primissimo periodo a
Riot, quando il team non aveva ancora motivo di fidarsi di lui. Risolvendo il
problema ha eliminato anche tutto lo strato di micro-dinamiche spiacevoli che
un manager impara a vedere (due colleghi che si guardano con la risatina
quando lei entra in ritardo). La ragazza ha cominciato a entrare tranquilla,
con una certa leggerezza: non le faceva piacere arrivare in ritardo (``ci sono
anche persone che lo fanno apposta, e quello è un problema da affrontare
diversamente''). Soprattutto, con i colleghi non parlava di lui come del capo
che minacciava, ma come di chi le aveva risolto un problema: era diventato
un \textbf{alleato}. La settimana dopo altre tre persone sono andate da lui
con un problema da risolvere. Ripetuto cinque, sei, sette, dodici volte, in un
paio di mesi poteva dire quasi qualsiasi cosa nei meeting sapendo che
richieste e critiche sarebbero state percepite come ``me lo chiede perché vuole
darmi una mano'', non per fare bella figura o per presentare i successi del
team ai suoi superiori.

\subsection{Condividere i successi, schermare i fallimenti}

Un altro elemento fondamentale (ripreso nella seconda giornata): il manager
\textbf{condivide i successi e protegge il team dai fallimenti}. ``Se va male
qualcosa è colpa mia; se va bene qualcosa, voglio che ne parliate voi.'' Il
lavoro del manager non dovrebbe essere valutato su ciò che presenta del lavoro
del team: tutti sanno comunque chi è il manager. Tanto vale lasciar presentare
il team, che potrà dire con orgoglio di aver presentato a Mike
Morhaime\footnote{Cofondatore ed ex presidente di Blizzard; nella
trascrizione ``Mike Morayme''.} la propria campagna europea per
\emph{StarCraft~II}. Ci sono manager che dicono ``come mi prendo le colpe, mi
prendo i meriti'', ma i meriti li hanno già: è inutile ostentarli. Con i
fallimenti invece il manager è la \textbf{rete di sicurezza} del team. Il
relatore ne ha fatto una questione di principio: \textbf{non scaricare mai la
colpa su qualcun altro}.

\paragraph{Colpevole o causa.} Un esempio che il relatore attribuisce a Pirate
Software, creatore di contenuti ed ex dipendente Blizzard\footnote{Il relatore
aggiunge che si tratterebbe del figlio dello sviluppatore a cui si sarebbe
ispirato il personaggio del giocatore ``senza vita'' nell'episodio di
\emph{South Park} dedicato a \emph{World of Warcraft}. L'informazione è
riportata come detta in lezione e non è verificata.}. Contattando il servizio
clienti Blizzard per motivi personali, aveva scoperto che i colleghi gli
avevano scritto nei ticket informazioni confidenziali che non avrebbero dovuto
condividere in quel modo, nemmeno con un collega. Non ha segnalato le persone
da punire: ha detto ``c'è un problema, evidentemente le policy non sono
chiare''. Quando gli hanno chiesto i nomi, ha rifiutato: ``se vi dico chi è,
punite quelli, e gli altri cinquecento continueranno a farlo; non avremo
cambiato la causa alla fonte''.

\begin{quote}
  \emph{È inutile trovare il colpevole. Voglio trovare la causa. Perché il
  colpevole cambia.}
\end{quote}

Il colpevole a volte è vittima della situazione: per questo la definizione
dice \emph{``con le informazioni a sua disposizione in quel momento''}. Se il
relatore avesse rimproverato la ragazza in ritardo, non solo non avrebbe
risolto il problema, ma avrebbe creato un'\textbf{ingiustizia}: lei, che
attraversava il campus di corsa per una colpa non sua e di cui non poteva
parlare, sarebbe stata anche rimproverata.

Nessuno fa un brutto lavoro apposta: nessuno è contento quando sbaglia o
perde. Perché al lavoro diamo per scontato che gli altri sbaglino
intenzionalmente? Se qualcuno ha fatto una cosa male, è perché in quel momento
credeva fosse la scelta migliore. La domanda giusta non è ``perché hai fatto
questa cosa che è stata un fallimento?'', ma:

\begin{quote}
  \emph{Che cosa ti ha fatto pensare che questa fosse la scelta giusta?}
\end{quote}

Magari erano limiti di budget, di tempo, di scopo: si agisce su quelli, non
sulla persona. Se si cambia la persona ma le si danno gli stessi fattori, è
molto probabile che arrivi alle stesse conclusioni, e ci si troverà a dover
``disciplinare'' anche lei.

\paragraph{La reputazione costruita sui rumor.} Durante il giro di
presentazione in un ufficio (non in nessuna delle aziende citate), al relatore
è stato detto: ``questo è il dipartimento di design; se ti serve qualcosa
chiedi a lei, non a lui, perché lui consegna sempre in ritardo''. Oltre a
trovarlo poco professionale, una settimana dopo ha chiesto: ``cosa ti ha fatto
questa persona?'' ``A me niente, però ho sentito che\dots''. Si torna alle
scimmie: magari quella persona non consegna in ritardo da mesi, ma ormai lo
dicono tutti. È un circolo vizioso da fermare subito.

\paragraph{Niente escalation inutili.} Lo stesso vale per il feedback: quando
qualcuno va dal relatore dicendo ``ho un problema con questa persona'', lui
chiede: \textbf{``perché ne stai parlando a me e non a lei?''}. ``Perché sei il
mio capo'' non basta. Passare da lui significa dare per scontato che l'altra
persona non abbia interesse a risolvere, mentre magari basta parlarle. È la
strada più facile (``lo dico a Stefano, ci pensa Stefano''), e proprio per
questo bisogna mettere dei paletti: il manager interviene quando serve, ma le
cose che si possono gestire da soli vanno gestite da soli. (Quando invece il
confronto diretto fallisce, l'escalation è legittima: si veda la
sezione~\ref{sec:z1-qa-finale}.)

% ---------------------------------------------------------------------
\section{Traiettorie di carriera e valutazione delle performance}
\label{sec:z1-traiettorie}

\subsection{Quale traiettoria vorreste?}

Il relatore mostra un grafico con quattro curve colorate, contrassegnate
anche da lettere (A, B, C, D) per eventuali persone daltoniche (``a me è
capitato che qualcuno scoprisse di esserlo durante una lezione''). Le curve
rappresentano la performance lavorativa nel tempo. Due domande: \emph{quale
vorreste essere? E quale pensate realisticamente di essere?} Un partecipante
aspira a B e pensa di essere C; quasi tutti scelgono B.

Perché non A, che sale più rapidamente? Un partecipante risponde che una
crescita così repentina può essere più dannosa che utile, come una ``rush'' di
studio molto faticosa che poi non si riesce a sostenere. Il relatore propone
una lettura suggerita da un amico psicologo, che spiega la forma del grafico:
di B, C e D \textbf{vediamo dove finiscono}; di A no. Potrebbe continuare a
salire oltre la scala del grafico, o crollare sotto D. Non saperlo ci spinge a
preferire B. (A chi obietta che anche le altre proseguono oltre il bordo, il
relatore precisa: il bordo destro è ``oggi'', e di A non sappiamo come sia
arrivata fin lì.)

\subsection{La carriera non è una linea retta}

Tutti immaginano la carriera come una linea: si parte da X e, attraverso
stage e posizioni, si arriva a Y. In realtà carriera e performance avranno un
andamento irregolare, con alti e bassi. Come non si accetta che un team possa
temporaneamente \emph{underperformare}, tanto meno lo si accetta da se stessi.
Bisogna accettare che ci saranno momenti bassi.

Per questo il relatore \textbf{detesta le valutazioni annuali delle
performance}: prendono un punto casuale della traiettoria, che se si è
fortunati è il picco più alto e se si è sfortunati il più basso, e lo usano
per valutare un intero anno. Insegna ai manager a fare più punti di controllo
e a guardare altre cose:
\begin{enumerate}
  \item \textbf{quanto tempo passa tra un picco e l'altro}, cioè quanto
        \emph{cooldown} serve a una persona. È il motivo per cui, dopo la
        Games Week, tutti staccavano una o due settimane. Il relatore stesso,
        che si impegna moltissimo e tende al burnout, non potrebbe rifare lo
        stesso workshop il lunedì e il martedì successivi con la stessa
        qualità;
  \item \textbf{la differenza tra la miglior performance e la peggiore}: che
        cosa ha portato a lavorare così bene, e così male? Con chi si
        lavorava, su cosa ci si concentrava, quanto tempo e budget c'erano?
        L'obiettivo è capire come \emph{replicare} le condizioni che fanno
        lavorare bene, invece del solito ``hai lavorato male, è colpa tua''.
        Il relatore ha avuto molti manager per cui quando lavorava male era
        colpa sua e quando lavorava bene ``che bravo'', ma sapeva che raramente
        era l'unico fattore della propria performance. È la fallacia del ``se
        sai fare 100, perché non fai sempre 100?'', come nella scena di
        \emph{Schindler's List} sull'operaio che produce una vite in pochi
        secondi: non sempre un ritmo è sostenibile;
  \item \textbf{il fallimento di oggi confrontato con quello di sei mesi fa}.
        Il relatore non è contento della presentazione di oggi (l'avrebbe
        migliorata per altri due anni), ma rispetto a quelle di cinque anni fa
        vede un grande miglioramento di chiarezza: il suo peggior fallimento di
        oggi equivale al suo primo successo di inizio carriera.
\end{enumerate}
Quello che conta è la \textbf{traiettoria} dall'inizio dell'anno a oggi, non
i singoli punti. E se si riesce a capire insieme quando e perché si ha
successo e quando e perché si fallisce, si può lavorare per replicare solo le
situazioni che fanno lavorare bene.

\paragraph{Lo stesso vale per i giochi.} È il motivo principale per cui i
giochi delle grandi case vanno bene o male. Quando si riescono a replicare le
condizioni che hanno prodotto \emph{Warcraft~II}, esce \emph{Warcraft~III};
quando non ci si riesce, esce ``il nuovo \emph{Dragon Age}'': la gente non è
arrabbiata perché il gioco non è bello, ma perché non è come gli altri
\emph{Dragon Age}; non ne replica il \emph{feeling}, anche se potrebbe
essere un gioco favoloso. L'eccezione positiva è il passaggio da \emph{Dune} a
\emph{Dune~II}: il primo era di fatto un'avventura grafica, il secondo ha
generato il genere degli strategici in tempo reale. Gli sviluppatori avevano
capito che il successo non stava nel tipo di gioco, ma nell'ambientazione che
attirava le persone: un bel gioco in quell'ambientazione sarebbe piaciuto
comunque. Il relatore comprò \emph{Dune~II} aspettandosi un seguito
dell'avventura e, pur sorpreso, lo finì.

\subsection{Lying, hiding, faking}

Quando si è valutati su un singolo punto nel tempo, e da quel punto dipendono
bonus, stipendio e il rischio di licenziamento, si comincia a fare
\textbf{lying, hiding e faking}:
\begin{description}
  \item[Lying] dire che è sempre andato tutto bene, che si finirà sempre in
        orario, che non c'è nessun problema;
  \item[Hiding] nascondere: ``non era un mio task'', ``non c'ero'', ``non ho
        visto niente'';
  \item[Faking] fingere, anche in cose apparentemente innocue come ``sì, sì,
        ho capito tutto'' quando si ha solo paura di dire ``me lo rispieghi?
        Non ho capito niente''.
\end{description}

\subsection{``Non ho capito'' e la conferma delle istruzioni}

Dire ``non ho capito, ripetimelo'' è ancora visto come una figura da sciocco.
Il relatore non si arrabbierà mai se qualcuno gli chiede di rispiegare
un'intera slide: le cose ovvie per lui possono non esserlo per altri. Più si
conosce qualcosa, \textbf{più è difficile spiegarla}, perché si danno per
scontate molte cose. Ne parlava proprio con Diego, che gli diceva che fare
l'insegnante diventa sempre più difficile man mano che si diventa esperti. È
come spiegare \emph{EVE Online}, uno dei giochi più complessi in circolazione:
chi lo conosce bene dà per scontato che ``ovviamente non puoi fare questa
cosa'', mentre per un nuovo giocatore non è affatto ovvio. A scuola chi dice
``non ho capito'' viene visto come lo ``scemo''; nessuno si chiede ``non è che
l'ho spiegato male?''.

Il racconto più ``traumatico'' dei ragazzi che fanno stage nelle grandi
aziende è sempre lo stesso: ``Cinzia, devi fare questa cosa.'' Cinzia la fa.
``Non te l'avevo chiesta così.'' Il problema è che nessuno ha
\textbf{confermato che Cinzia avesse capito}. Il consiglio, valido sempre:
\begin{itemize}
  \item quando si \emph{ricevono} istruzioni, ripeterle: ``ho capito bene?
        Vuoi che faccia questa cosa in questo modo?'';
  \item quando si \emph{danno} istruzioni, chiedere a chi le riceve di
        ripeterle, non perché lo si ritenga incapace, ma per capire che cosa
        ha capito dalla propria spiegazione (``quindi vuoi un quadrato?''
        ``No, ho detto un triangolo\dots'' ``Però mi hai detto una forma con
        quattro angoli.'' ``Hai ragione: voglio un quadrato'').
\end{itemize}
Frasi come ``non te lo ripeto due volte'' o ``stammi dietro'' portano solo a
fare le cose come si pensa che vadano fatte e a sentirsi dire a fine progetto
``non la volevo così'': che si tratti di una feature, di un design o di un
modello 3D. (Un partecipante osserva che lo stesso accade insegnando le lingue:
chi insegna dà per scontato, ad esempio, che si sappia dell'esistenza del
genere neutro in tedesco.)

\paragraph{Come convincere il team a chiedere?} Il relatore mette sempre come
premessa che \textbf{l'obiettivo non è trovare la persona più intelligente
nella stanza, ma fare un buon lavoro}. Che Gianni capisca alla prima e un altro
alla quinta non gli interessa, se alla fine entrambi lavorano bene: ripetere
cinque volte è comunque meno lavoro che far rifare le cose sette volte. È un
investimento a lungo termine, e non dissimula alcun fastidio: davvero non gli
dà fastidio. Ma ripetere le stesse parole non aiuta: ``non ho sentito'' si
risolve ripetendo; ``non ho capito'' si risolve cambiando spiegazione. Meglio
ancora: ``dimmi che cosa hai capito tu, e vediamo cosa non torna''. Se
intravedesse nel team uno sguardo del tipo ``è la quinta volta che lo
chiede'', interverrebbe: spesso chi giudica è proprio chi dice ``ho capito'' e
poi sbaglia. Preferisce di gran lunga chi chiede cinque volte.

\subsection{Il feedback più utile che abbiate ricevuto}

\emph{Qual è il feedback più utile che abbiate mai ricevuto?} Il relatore ne ha
almeno dieci che gli hanno cambiato il modo di lavorare. Le risposte:
\begin{itemize}
  \item un semplice \textbf{``grazie''} per un lavoro fatto, più utile di un
        ``questa cosa fa schifo'';
  \item un feedback sulla propria \textbf{impulsività} nel rispondere: fermarsi
        un attimo e capire di cosa si sta parlando prima di intervenire;
  \item \textbf{essere se stessi}, anche al lavoro, e non scusarsi
        continuamente per ogni cosa.
\end{itemize}
I feedback che toccano di più sono quelli \textbf{veri}. Il ``grazie''
significa ``ho visto quello che hai fatto e l'ho apprezzato''. Quelli negativi
che si ricordano sono quelli che hanno fatto realizzare ``questa cosa che
pensavo di saper fare, non la so fare'', oppure ``la so fare, ma questa volta
non mi è venuta'' (non è che se una cosa una volta non riesce non la si sa più
fare: può capitare di sbagliare una porta). Arrivano sempre da un
\textbf{punto cieco}: qualcosa che non si sapeva fosse percepito dagli altri.
Conta più ciò che le persone ti dicono di ciò che pensi che pensino di te; e
ciò che il relatore pensa che gli altri pensino di lui è sempre terribile,
mentre il loro giudizio si rivela molto più positivo del previsto.

\paragraph{Il denominatore comune siete voi.} Si potrebbe pensare che questi
concetti non siano applicabili da chi arriva in un'azienda senza poterne
cambiare la cultura. Ma c'è sempre un denominatore comune a cui applicarli:
\emph{se stessi}. La cultura aziendale può non spingere a dare feedback, può
vedere male chi chiede aiuto; ma il relatore ha visto che quando lui lo fa, le
persone intorno tendono a farlo, e un gruppo che lo fa lavora meglio. Perfino
i CEO più ostici a un certo punto gli chiedevano: ``perché il tuo team ora
performa meglio degli altri?''. La differenza è il default to trust: una
pratica che col tempo fa operare il team \emph{come un team} e non come un
insieme di individui. Un videogioco è un progetto corale: c'è chi fa giochi da
solo, ma non si fa da soli \emph{Red Dead Redemption~2}; e se il gruppo non è
affiatato e allineato non esce \emph{Red Dead Redemption~2}, esce
\emph{Anthem} (``oggi la mia fissazione è \emph{Anthem}, domani sarà
un'altra'').

% ---------------------------------------------------------------------
\section{Come dare un feedback}
\label{sec:z1-dare}

\subsection{Concentrarsi sull'azione e sui suoi effetti}

Il primo principio: \textbf{il feedback non riguarda mai la persona}. Non si
dice ``tu hai fatto questa cosa che non va bene'', ma ``questa cosa non va
bene: come facciamo a non farla più?''. Che sia un comportamento, un'azione,
una decisione o una strategia, ci si concentra su quello, e si
\textbf{esplicita l'effetto} che ha. Nell'esempio del ritardo l'effetto era
``quindici minuti di attesa moltiplicati per venti persone''. In un progetto
può essere una catena:
\begin{quote}
  se non rispetti le deadline, l'animazione che viene dopo resta ferma; se
  loro sono in ritardo, la \emph{vertical slice} è in ritardo; se la vertical
  slice non è pronta per la Gamescom, non andiamo alla Gamescom e non troviamo
  un publisher; se non troviamo un publisher, chiudiamo.
\end{quote}
Più la conseguenza è grave, più dovrebbe sembrare piccolo il gesto di dire
``datemi una mano'': è il \emph{male minore}. Poi, a cose fatte, ci si dà
feedback: cosa ha funzionato, cosa no, cosa fare di più, di meno, cosa smettere
di fare.

\subsection{Il pericolo degli ambienti solo positivi}

Gli ambienti dove ci si fa sempre grandi complimenti e nessuno parla dei
problemi creano difficoltà vere alle aziende. Il problema non è mai la
mancanza di complimenti, ma un \textbf{eccesso di ottimismo}: si comincia a
pensare che le cose andranno come si spera, perché tutti intorno danno solo
feedback positivi.

A Blizzard i feedback determinavano gli \textbf{aumenti annuali}: ognuno
doveva indicare cinque persone che avrebbero dato feedback su di lui, e
ovviamente tutti indicavano gli amici, o colleghi di progetti andati bene. Un
sistema nato con le migliori intenzioni diventava uno \emph{scambio di voti}
(``se mi metti follow ti metto follow''), senza alcun miglioramento.

Al contrario, \textbf{le conversazioni più brutte sono state le più utili}:
affrontare i licenziamenti a Riot, o la cancellazione di \emph{Project Titan}
a Blizzard, il progetto che doveva essere il nuovo MMO dopo \emph{World of
Warcraft} e da cui poi è nato \emph{Overwatch}\footnote{Nella trascrizione si
legge anche ``doveva essere Warcraft 2'', probabilmente un lapsus o un errore
di trascrizione.}. Andando a vedere perché lo si cancellava, c'erano solo
motivazioni razionali e pragmatiche, di business, chiare a tutti, non un ``Mike
Morhaime ci odia''. Sono anche le conversazioni in cui ha lavorato su se
stesso: è molto \textbf{avverso ai conflitti}, e questo a volte gli ha
procurato conseguenze molto più gravi di quanto sarebbe costato un piccolo
disagio. Dieci anni fa faceva fatica ad accettare questo difetto; ha dovuto
subire le conseguenze di conflitti evitati e poi cresciuti.

\subsection{Che cosa fare}

Le linee guida della slide, che il relatore legge rapidamente perché saranno
disponibili nel materiale:
\begin{description}
  \item[Attenzione al linguaggio del corpo] dare un feedback a braccia
        incrociate o guardando altrove è diverso dal guardare l'altro negli
        occhi.
  \item[Sii specifico] con esempi tangibili.
  \item[Esplicita gli effetti] come visto sopra.
  \item[Discuti il prima possibile] non aspettare un mese per parlare di
        qualcosa successo ieri, altrimenti nessuno ricorderà più bene cosa sia
        accaduto.
  \item[Concentrati sull'azione e sul problema, non sulla persona] il nemico è
        l'azione, non la persona che si ha davanti. Non bisogna far cambiare
        l'altro: bisogna evitare che il problema si ripresenti.
\end{description}

\paragraph{Gli insulti sono rivolti al ruolo.} Un esempio solo apparentemente
scollegato. Il relatore ha lavorato su \emph{League of Legends}, ``non
esattamente noto per avere una community da orsetti del cuore'', in varie
\emph{task force} sulla tossicità. A chi smetteva di giocare veniva inviato un
sondaggio sul perché; il 33\% indicava la tossicità della community. Dopo varie
campagne di sensibilizzazione e cambiamenti si è scesi al 30\%: sembra solo
un 3\% in meno, ma valeva circa \textbf{120 milioni di revenue in più
all'anno}, perché quei giocatori restavano esposti a skin e contenuti e
continuavano a comprare. Anche qui, piccoli guadagni che si sommano.

Lavorando con degli psicologi, aveva notato che gli insulti \textbf{non sono
quasi mai rivolti alla persona, ma al ruolo}: ``il mio jungler fa schifo'',
``il mio DPS non è capace'', ``il nostro portiere è una pippa''. Chiamare
qualcuno per ruolo lo \emph{depersonalizza} e rende più facile insultarlo. Lo
stesso nell'esempio del giro dell'ufficio: non ``non chiedere a Giulio'', ma
``non chiedere a quel grafico, consegna sempre in ritardo''. Concentrarsi
sull'azione, senza attaccare la persona, fa sì che il feedback non sia (e non
venga percepito come) un attacco personale.

\paragraph{Considera il carattere.} Qui tornano utili i test di personalità.
Servono più agli studenti che al relatore (che sa già che usciranno 53 profili
diversi, senza un pattern): bisogna conoscere il proprio profilo e
\textbf{condividerlo con chi lavora con noi}, chiedendo anche agli altri di
fare lo stesso. In ufficio il relatore ha usato una versione più avanzata, a
pagamento, con un report di circa trenta pagine: il \textbf{profilo
DISC}\footnote{Nella trascrizione ``disk''. Si tratta del modello DISC
(Dominance, Influence, Steadiness, Conscientiousness).}. Ha mappato tutta
l'azienda e ha fatto stampare il profilo di ciascuno sulla sua scrivania, così
che tutti sapessero con chi stavano parlando. Capire come ragionano gli altri
(``è introverso, non gliene parlo davanti a tutti'') è una delle cose che gli
sono tornate più utili: \emph{adattare ciò che si vuole dire a chi si ha
davanti}.

\paragraph{Assicurati di non disturbare.} Sotto il profilo di ognuno c'era un
piccolo \textbf{semaforo}: verde, ``parlami pure, interrompimi''; giallo,
``interrompimi solo se è urgente''; rosso, ``mi sto concentrando, non mi
parlare''.

\paragraph{Lascia finire il lavoro.} Durante gli eventi dal vivo il relatore si
era reso conto di quanto fosse grave dare feedback mentre qualcuno sta ancora
lavorando. Se una cosa sta andando male, dirlo in quel momento non aiuta: tra
la pressione di finire e quella del feedback, peggiorano entrambe le cose, e
anche il rapporto. Meglio lasciarla finire, anche male, e parlarne dopo.

\subsection{Che cosa non fare}

\begin{description}
  \item[Non basarti su supposizioni] ``mi hanno detto che ieri in un meeting
        hai fatto\dots'': o c'eri e l'hai visto, o non esiste.
  \item[Non giudicare la persona.]
  \item[Non fornire feedback vago.]
  \item[Non accettare e non fornire feedback di terze parti] Se qualcuno vi
        dicesse ``mi hanno detto che hai fatto un brutto lavoro'', il primo
        pensiero non sarebbe ``cosa ho fatto male?'', ma \emph{``chi te l'ha
        detto? Chi è la vipera che ho in seno? Chi è l'impostore?''}. Siamo
        fatti così. Ogni conversazione che inizia con ``mi hanno detto che'' è
        inutile: l'altro non ascolterà, penserà solo a chi dare la caccia, e
        spesso penserà alla persona sbagliata, creando antipatie tra persone
        che non c'entravano nulla. Quando al relatore dicono ``mi hanno detto
        che lui ha fatto\dots'', risponde ``vai a parlarne con lui''; se gli
        dicono ``mi hanno detto che tu hai fatto\dots'', risponde ``mandami la
        persona che te l'ha detto, altrimenti non ne parliamo''. Non per
        cattiveria: se chiedesse ``che cosa ho detto di preciso a Marco?'' e la
        risposta fosse ``non lo so, non c'ero'', di cosa si starebbe parlando?
        \textbf{Parlate con la fonte}, altrimenti fate come se quel feedback non
        esistesse.
  \item[Non dare consigli non richiesti] come visto, non portare la soluzione:
        la soluzione deve venire da chi riceve il feedback.
  \item[``Non mi piace'' non è un feedback] ``Penso che questo design non
        funzioni perché, visto che il nostro obiettivo è\dots'' ha argomenti;
        ``non mi piace il combat system'' è un'opinione. Anche ``io l'avrei
        fatto diversamente'' non è un feedback.
  \item[Accettare non vuol dire fare come proponi] deve essere chiaro a chi dà
        il feedback: l'aspettativa è essere ascoltati e presi in
        considerazione, ma se la responsabilità è dell'altro, deciderà
        l'altro.
  \item[Non offenderti] se il feedback non viene seguito. Restarci male
        significa pensare di saperne più dell'altro, cioè \emph{venire meno al
        default to trust}. ``Decidi tu'' significa ``mi fido di te, perché sei
        lo specialista, e mi fido che mi ascolterai e farai la cosa migliore
        per tutti'', che potrebbe anche non essere ciò che ho detto io. Se
        voglio che l'altro faccia sempre come dico io, mi sto fidando solo di
        me stesso.
\end{description}

\subsection{Il primo feedback e il debriefing}

La cosa più difficile è dare il \textbf{primo} feedback: è un muscolo da
allenare. Chi non l'ha mai fatto (soprattutto con questa struttura) dovrebbe
provarci, anche su minuzie (``abbiamo concordato di usare WhatsApp nel gruppo:
non potremmo usare Telegram?''), anche solo per vedere come va. Altre
indicazioni:
\begin{itemize}
  \item feedback \textbf{solo uno a uno}, mai davanti ad altri, soprattutto se
        negativo;
  \item \textbf{prepararsi} prima;
  \item fare il \textbf{debriefing del feedback}: dopo, fermarsi a pensare a
        com'è andato. Il relatore, quando media un feedback tra due persone,
        si sente spesso rispondere ``benissimo'' da una e ``malissimo''
        dall'altra: la percezione è molto diversa. Non si fa problemi a tornare
        dalla stessa persona: ``ieri ti ho dato un feedback: com'è andata? Sei
        soddisfatto? Come poteva andare meglio?''. Con le persone con cui si ha
        meno a che fare, questo aiuta molto la volta successiva.
\end{itemize}

% ---------------------------------------------------------------------
\section{Come ricevere un feedback}
\label{sec:z1-ricevere}

Per il relatore è la parte più difficile. Si riassume in quattro passi:
\begin{enumerate}
  \item \textbf{Ascolta}: si parte sempre ascoltando.
  \item \textbf{Chiedi dettagli}, se serve, ma per capire meglio, non per
        difenderti.
  \item \textbf{Accetta, a prescindere}: non si sta decidendo chi ha ragione;
        si sta ascoltando un'opinione. Non c'è nulla da cui difendersi.
  \item \textbf{Agisci}: fatti un piano. Può essere banale (``mi hanno detto
        che interrompo sempre: cerco di non interrompere''), oppure ``devo
        migliorare questa cosa: come?''. Pensa anche a \emph{chi} potrebbe
        aiutarti e a \emph{che supporto} ti serve: il piano non deve essere
        per forza ``devo risolvere tutto da solo''.
\end{enumerate}
Una domanda che molti sottovalutano: \emph{``sono junior in un'azienda: sarò
mica il primo ad avere questo problema?''}. Probabilmente qualcun altro l'ha
già avuto: basta chiedere ``tu come hai fatto?''.

\subsection{Il feedback positivo}

Tutto questo vale per i feedback costruttivi negativi, a cui si pensa di
default; ma esistono anche i \textbf{feedback positivi}, e vanno usati. Quando
si riceve un ``grazie'' o un ``bravo'', la domanda è: \emph{come faccio a
ricreare le stesse condizioni?}

Il relatore ha lavorato anche in aziende dove il feedback era solo negativo.
Il rischio? Si immagini l'impiegato ideale: vuole sempre migliorare, è sempre
pronto a dire ``magari sbaglio io''. Se nessuno gli dice cosa sta facendo
\emph{bene}, spinto dal desiderio di migliorare rischia di cambiare cose che
erano già buone, e di farle peggio. Il rinforzo positivo \textbf{non è solo una
gentilezza}: è un atto pratico, che evita che le persone peggiorino ciò che
sanno già fare bene.

\subsection{Il feedback accidentale}

Esiste infine il \textbf{feedback accidentale}, molto frequente negli uffici:
\begin{itemize}
  \item il ``chiedi a lei, non a quel grafico'' del giro dell'ufficio è un
        feedback negativo su una persona, dato accidentalmente, né richiesto
        né voluto;
  \item a Riot li chiamavano \textbf{drive-by}, come nelle scene dei film in
        cui passa un'auto, si abbassa il finestrino, si spara e si riparte. In
        molte aziende con una ``cultura del feedback'' in cui tutti devono
        dire la loro, si riceve una valanga di commenti del tipo ``comunque il
        sito a me non piace, ciao''. Non si sa perché, non si sa cosa fare di
        diverso, non si sa che rischi ci siano: \emph{a che cosa è servita
        questa conversazione?} A niente.
\end{itemize}

% ---------------------------------------------------------------------
\section{Quando, quanto e con chi}
\label{sec:z1-quando}

\subsection{Frequentemente e recentemente}

\emph{Quando avete ricevuto il vostro ultimo feedback?} Per gli studenti,
qualche giorno fa (nel corso) o qualche mese fa (al lavoro, dopo una
consegna). Come aveva osservato Silvia, il feedback andrebbe dato
spesso.

Il relatore usa l'esempio di una nave da crociera (``la Costa Concordia'') che
deve andare da A a B. La rotta ideale (azzurra) è quasi impossibile: nessun
progetto va esattamente come previsto. C'è la rotta ``panoramica'' (verde), che
gira intorno all'isola d'Elba; ci sono poi rotte (arancione e rosa) che a un
certo punto si allontanano molto e poi correggono bruscamente: qualcuno si è
accorto che si stava sbagliando, ``ma non è quello che avevo chiesto'',
``non è dove dovevamo andare''. Sono i famosi \textbf{sanity check} o
\emph{checkpoint}.

L'ideale è avere \textbf{tanti piccoli check per fare tante micro-correzioni}:
correggere la rotta di un paio di gradi è molto più rapido che farlo dopo
essersi allontanati. Bastano trenta secondi di ``ehi, vado bene?''. Spesso non
lo si fa per paura di sembrare insicuri o poco autonomi, ma per il relatore è
l'opposto: quando qualcuno gli chiede ``vado bene così?'' e lui risponde
``sì'', si alleggerisce anche del proprio carico mentale. Un capo che assegna
un compito e lo vede solo alla fine vive sempre col dubbio (``avranno capito?
Lo staranno facendo giusto? Arriverà in tempo?''); con un check di dieci
secondi al giorno (``tutto ok, Stefano'') è tranquillo, e la persona va avanti
sapendo che quanto fatto finora va bene. Non serve fare ogni giorno il giro del
dipartimento: anche ogni uno o due giorni basta a evitare la rotta verde.

Per il feedback è lo stesso: \textbf{va dato subito}. Dire oggi ``ieri hai
fatto questa cosa'' è molto diverso da dire ``negli ultimi sei mesi hai fatto
questa cosa''. Il relatore si è sentito dire più volte, alla valutazione di
dicembre, ``da maggio hai fatto questa cosa che non andava''. La sua risposta:
``e aspetti dicembre per dirmelo? Mi hai lasciato sbagliare per sei mesi?''. È
una mancanza di rispetto, e la conversazione non può più essere costruttiva.
Per questo non ama le valutazioni annuali: i manager crescono abituandosi a
prendere appunti delle cose negative per tirarle fuori alla valutazione
(``a marzo hai fatto questa cosa'' ``sì, però poi ho smesso'' ``sì, però a
marzo l'hai fatta'').

In sintesi:
\begin{description}
  \item[Frequentemente] riguarda l'intervallo tra un check e l'altro.
  \item[Recentemente] significa non lasciar passare un anno prima di dare il
        feedback, altrimenti non serve più a niente.
\end{description}

\subsection{Quanto dura un feedback}

Un feedback dura \textbf{da trenta secondi a quindici minuti}; raramente di
più. Chi non ne ha mai dato uno dovrebbe sforzarsi di darne almeno uno, magari
positivo per cominciare. Ogni feedback dato o ricevuto è un
\textbf{deposito nel proprio ``fondo fiduciario''} (\emph{trust fund}), che
costruisce la reputazione futura. Il relatore si è costruito la reputazione di
chi dà feedback il più possibile oggettivi e analitici: giusti o sbagliati che
siano, costringono a pensare. Non impone mai il suo punto di vista (``si fa
come dico io''), ma chiede ``tu come faresti? Come posso aiutarti?''. Il
feedback, per lui, serve a \emph{far crescere qualcun altro}.

\subsection{I manager peggiori: ``a che numero sto pensando?''}

I peggiori manager che il relatore abbia avuto erano quelli che volevano che si
facesse come volevano loro, anche quando chiedevano ``tu come la faresti?''.
In realtà la domanda era ``\emph{a che numero sto pensando?}'': se dicevi 5,
``no, dovremmo fare 9''; se dicevi 9, ``no, 5''. Lo stress costante di dover
indovinare cosa vuole il capo nasce dalla mancanza di fiducia (``se mi chiedi
come fare, fammelo fare'').

Lo stesso gli succede con i CEO che lo chiamano per ``sistemare la cultura
aziendale'' e poi dicono ``non sono sicuro, io farei diversamente''. La sua
risposta: ``lo so che faresti diversamente, per questo mi hai chiamato''. Se si
chiama un esperto, bisogna lasciargli fare il suo lavoro: non si va dal
meccanico a cambiare le gomme per poi dirgli ``non dovresti girare un po' di
più quel bullone?''. Chiedere un parere a un esperto significa dargli anche il
potere di scegliere. Il problema è il \emph{micromanagement}, o in generale il
non accettare punti di vista diversi dal proprio (magari perché si è il
fondatore e si pensa di essere l'unico a saper fare le cose).

\subsection{Dare feedback a un superiore egocentrico}

\paragraph{Come si dà un feedback negativo a un superiore con un grande ego?}
Il trucco è \textbf{non puntare sull'ego, ma sull'azione e sui suoi effetti}.
Non ``tu fai così e io mi sento così'': anzi, si alimenta il suo ego
rendendolo il \emph{main character} dell'avventura, l'eroe che risolve il
problema (``tu, grande manager, aiutaci a risolvere questo problema che
affligge noi che facciamo l'operatività''). Non è l'ideale (``io cambierei
manager''), ma pochissimi possono scegliersi i colleghi: il relatore stesso ha
quasi sempre ricevuto team preassemblati (``ecco le chiavi, questa è la tua
macchina''). Bisogna giocare con le carte che si hanno in mano.

\paragraph{Un esempio pratico.} Con un manager che fa micromanagement e mette
sempre ``il becco'' nel proprio lavoro, si può dire: ``se ti fidi di me per
questo lavoro, con le tue capacità potresti dedicarti a questa cosa più
strategica e di lungo termine: pensa quanto sarebbe più utile per l'azienda''.
In sostanza: ``guarda cosa potresti raggiungere se mi lasciassi fare il mio
lavoro''. Non come un ``tu fai il tuo e io faccio il mio'', ma come un
``ti sollevo da questa parte del carico, che so che ritieni importante: lascia
che ci pensi io, così puoi dedicarti ad altro'', che magari gli dà anche più
visibilità.

Si gioca su ciò che quella persona pensa di saper fare, e che probabilmente sa
anche fare. Anche Michael Scott di \emph{The Office}, per quanto preso in giro,
è raccontato come un venditore bravissimo, che ha chiuso i contratti più
importanti per l'azienda, e rispetto ai manager che gli succedono è forse
``uno dei meno peggiori''. Tutti, anche i manager peggiori, hanno punti di
forza: \textbf{giocate su quelli}.

% ---------------------------------------------------------------------
\section{Il feedback come regalo}
\label{sec:z1-regalo}

Per chiudere, il relatore risponde alla domanda d'apertura, che aveva lasciato
apposta in sospeso: che cos'è il feedback per lui?

\begin{quote}
  \emph{Il feedback è un regalo che mi viene fatto da qualcuno che so che vuole
  il mio interesse.}
\end{quote}

Ha costruito un team di persone che si fidano di lui e che non hanno motivo di
sabotarlo o di ferirlo: per questo interpreta anche le cose peggiori con il
default to trust. È un regalo, e cerca di prenderlo in considerazione il più
possibile.

Ci è arrivato partendo da se stesso. Oltre a essere avverso ai conflitti, si
definisce ``avverso alla comunicazione umana'': paradossalmente, una persona
introversa e timida finita a fare comunicazione, marketing e ora HR (``succede
spesso''). Prima di ordinare una pizza al telefono ha dieci minuti di
preparazione mentale, e se l'ordine delle domande cambia va in confusione. Per
lui un feedback di dieci minuti sono spesso \textbf{tre ore di preparazione}:
costruisce un \textbf{albero dei dialoghi}, come in un'avventura grafica
personale, con tutte le opzioni. Poi spesso finisce con ``hai fatto questa
cosa, andrebbe migliorata'' ``hai ragione, grazie'', e basta.

Così ha pensato: chissà quanti si fanno gli stessi problemi con lui, e passano
ore a prepararsi per dirgli ``questa cosa l'hai fatta male'', sapendo che
riceveranno comunque quell'occhiata fugace del ``come ti permetti''. Eppure
vengono a dirglielo. Non può far altro che rispondere: ``grazie, hai ragione,
cercherò di fare meglio''. È così che ha metabolizzato il ricevere feedback:
capendo che non c'è nessuna cattiva intenzione e che, se non proprio felice
(``ci ho messo una pietra sopra''), può almeno essere \emph{grato} che ci sia
ancora qualcuno che ha così a cuore il suo interesse da dargli feedback.

\paragraph{Il silenzio fa più paura.} Non bisogna spaventarsi dei feedback
negativi che emergono dai sondaggi aziendali: bisogna spaventarsi del
\textbf{silenzio}, quando nessuno dice più nulla. Significa che nessuno si
fida più che si sappiano risolvere i problemi.

\paragraph{La sindrome dell'impostore.} Il suo motivatore personale, all'inizio,
era soprattutto ``non voglio fallire'': per questo cercava di migliorare, e
migliorare è diventato accettare feedback, perché più feedback riceve, più
migliora. Il feedback è un \textbf{acceleratore}, e lui cerca di incentivare
gli altri a darglielo. Resta però sempre la vocina del ``e se si accorgono che
non sei capace?''. Con la \emph{impostor syndrome} convive da anni e
continuerà a conviverci; se ce l'avete, conviveteci anche voi. Avrete sempre
la sensazione che ``se ne accorgeranno'': forse sì, ma sapete fare anche
altre cose. Non concentratevi solo su ciò che non sapete fare. Gli errori
capitano: la perfezione non è di questo mondo.

\paragraph{Esercitarsi in accademia.} Per fare pratica con il feedback,
l'ambiente migliore è quello dei colleghi di corso: persone nella stessa
situazione, con esperienza simile e lo stesso obiettivo, in un contesto in cui
non si può essere licenziati e l'unico effetto a lungo termine è ``ho parlato
con questa persona, ora le cose vanno meglio''. E bisogna avere il coraggio di
parlare anche con i docenti (``dire a Diego: secondo me questa lezione potevi
farla in questo modo''; ``so che è una persona che lo accetta''), senza farsi
bloccare dalla gerarchia. Nell'azienda giusta (come quella in cui lavora ora) il
relatore va sempre dal CEO a dare feedback, sapendo che non verrà licenziato:
il CEO gli ha detto ``ti ho preso per darmi feedback''.

La frase di chiusura:

\begin{quote}
  \emph{Per superare le vostre difficoltà, tutto quello che vi viene chiesto è
  il coraggio di chiedere aiuto.}
\end{quote}

Chiedere aiuto non è una debolezza: è un segno di consapevolezza di sé
(\emph{self-awareness}) e di forza personale. Se tutti lo facessero, i
progetti andrebbero molto meglio; spesso invece ci si fa frenare dal ``non
voglio fare brutta figura'', ``non voglio passare per quello scemo'', e sono
proprio queste paure a rallentarci e a impedirci di raggiungere gli obiettivi.

% ---------------------------------------------------------------------
\section{Domande e risposte finali}
\label{sec:z1-qa-finale}

L'ultima mezz'ora è dedicata a domande libere.

\paragraph{E se una persona peggiora per uno o due anni di fila?} Se nessuno
gliel'ha mai detto, c'è un grosso problema, non suo ma dell'azienda. Se invece
c'è consapevolezza, il percorso è: ``cosa vuoi fare per risolvere?''; ``ho
provato questo, questo e questo''; ``visto che non ha funzionato, vuoi un mio
consiglio? Proviamo con quello?''; ``che supporto ti serve?''. Se, con tutto il
supporto, il tempo e le soluzioni concordate per mesi, il miglioramento non
arriva, per il relatore la soluzione è: \textbf{``dimmi un prezzo per andare
via''}. Preferisce pagare una buonuscita (\emph{severance package}) che
avviare la strada dei richiami formali sulla performance fino al
licenziamento per giusta causa. A un certo punto andare avanti non è un favore
per nessuno.

\paragraph{È vero che nei videogiochi si viene sfruttati? Meglio un piccolo
studio italiano?} Per il relatore si è quasi più a rischio a Blizzard che in un
piccolo studio di Roma o Torino. Nel piccolo studio si sa dove ci si va a
cacciare: non ci sarà il team building a Santo Domingo, ma ci si concentra su
un gioco alla volta e, se va bene, si fa il successivo. A Blizzard si può
lavorare su \emph{Overwatch}, che vende milioni, e perdere comunque il posto,
in modo totalmente scollegato dal proprio lavoro. Gli stipendi più alti
all'estero compensano anche questo rischio, che spesso non dipende da come va
il gioco (anzi, ``il gioco non ha fatto come volevamo'' è spesso una scusa per
licenziamenti che non avevano motivo).

\paragraph{Cambiare settore a trent'anni, da web developer full remote?} Lo
stipendio indicato dallo studente non è fuori dalla portata di vari studi
italiani, anzi è superabile. Non va esclusa la possibilità di lavorare da
remoto per aziende piccole ma estere. Alcuni amici del relatore, pagati con
stipendi italiani, si sono trasferiti a Bali o in paesi con un costo della vita
più basso; il collega di \emph{PUBG} percepisce uno stipendio olandese vivendo
a Milano, che per quanto cara non costa come i Paesi Bassi.

\paragraph{Come si diventa manager?} Dipende dall'azienda: alcune promuovono i
migliori (processo tendenzialmente sbagliato), altre propongono un percorso di
\emph{people management}. Diventare manager significa \textbf{abbandonare ciò
che si sa fare per imparare a fare altro}. Il relatore ha dovuto accettare che
non sarebbe più stato ``il miglior community manager italiano del mondo''
(``eravamo in pochi, quindi era facile''): dedicandosi alle persone non poteva
più fare il suo lavoro bene come prima. È stata una negoziazione con se
stesso: fare il manager significa non essere più l'esperto di qualcosa, ma
saper gestire e far crescere persone. Significa anche \textbf{cedere il
controllo della propria carriera}: se il team decidesse di non fare più nulla,
il manager ne pagherebbe le conseguenze. Il suo lavoro diventa scegliere le
persone giuste e farle lavorare bene. Lui lo esplicita sempre: sa di affidare
al team parte della propria carriera, e il team sa che lui combatterà per
loro, anche prendendosi le colpe.

\paragraph{L'inglese.} \textbf{Va assolutamente imparato}, anche lavorando in
Italia. Non saperlo preclude circa il 95\% del mercato (tutto quello fuori
dall'Italia) e soprattutto il \emph{networking} con persone che potrebbero dare
occasioni future: l'inglese è spesso l'unica lingua in comune con chi lavora
all'estero ed è oggi la lingua universale dei videogiochi.
\begin{itemize}
  \item \emph{Che livello serve?} Non importa B2 o C2: bisogna saper sostenere
        una conversazione, guardare almeno Netflix con sottotitoli in inglese,
        comunicare senza dodici pause e senza ``come si dice\dots''. Il
        certificato non verrà mai chiesto: verrà chiesto se si sa parlare
        inglese. Il relatore ha sentito persone con un C2 che non sapevano
        parlarlo.
  \item \emph{Come allenarlo senza occasioni di conversazione?} Netflix, Prime
        Video, Now TV, Udemy, Coursera, YouTube, Duolingo, siti di scambio
        linguistico con madrelingua, Discord: non sono mai esistiti così tanti
        mezzi gratuiti. Il relatore ha imparato l'inglese giocando (con una
        pronuncia sbagliata, perché leggeva senza sentire, su cui ha poi
        dovuto lavorare); guardare un film in inglese con i sottotitoli in
        italiano insegna anche la traduzione.
  \item Oltre all'ascolto passivo è fondamentale \textbf{parlare}: in una
        flotta di \emph{EVE Online} con 150 persone quasi tutte madrelingua
        inglese si è esposti a decine di accenti e modi di dire (mille modi di
        dire ``come stai'' oltre a \emph{how are you}); lavorando all'estero si
        impara a distinguere gli accenti (Manchester, Londra) e ad apprezzarli
        anche nei media.
\end{itemize}

\paragraph{Quanto deve sapere un junior?} Dipende da quanto è tecnico il
ruolo. Un game designer junior che non sa qualcosa può essere affiancato da un
senior; un programmatore che non sa programmare, in quel momento, non è utile
all'azienda.

\paragraph{Annunci ``junior con cinque anni di esperienza''.} ``Cercasi
ventenne con trent'anni di esperienza.'' Non è l'annuncio per un junior: è un
errore dell'azienda, che in realtà cerca un \emph{mid}. Riceverà profili
junior, si stupirà che non abbiano cinque anni di esperienza, rifarà l'annuncio
e avrà perso sei mesi. Il consiglio: tenere d'occhio quella pagina, perché
dopo qualche mese l'annuncio verrà corretto.

\paragraph{Una struttura per il post-mortem.} A chi chiede un formato di base,
il relatore anticipa quanto vedrà nella seconda giornata a proposito di
brainstorming: una lavagna \textbf{divisa in quattro aree}, su cui ognuno
attacca tutti i post-it che vuole:
\begin{center}
  \begin{tabular}{@{}p{0.44\textwidth}p{0.44\textwidth}@{}}
    \toprule
    \textbf{Smettere di fare} & \textbf{Iniziare a fare} \\
    ciò che non ha funzionato & ciò che non abbiamo fatto e andava fatto \\
    \midrule
    \textbf{Fare di meno} & \textbf{Fare di più} \\
    ciò che ha fatto perdere tempo rispetto all'impatto ottenuto &
    ciò che ha funzionato ma, fatto con troppa leggerezza, non ha dato i
    frutti sperati \\
    \bottomrule
  \end{tabular}
\end{center}
Con i post-it ognuno ha voce e tutto risulta catalogato; per il report basta
una foto della lavagna (o uno screenshot, se la lavagna è digitale, come Miro).

\paragraph{Come sbloccare la paralisi del feedback verso il management?} Uno
studente ha visto manager dire ``non abbiate paura, la mia porta è sempre
aperta'' senza che cambiasse nulla. C'è una dissonanza tra ciò che dicono e
ciò che fanno: è ``la mamma che dice vieni qua che non ti faccio niente''.
Qualcosa deve aver fatto pensare alle persone che non è sicuro. Prima va
sistemato quello; poi si può provare con i \textbf{feedback anonimi}.
L'obiettivo del relatore resta il feedback aperto e firmato, ma a Stormind ha
iniziato con feedback anonimi verso il management per verificare se
l'anonimato fosse un fattore. Non lo era (evidentemente non c'erano problemi a
parlare con il CEO), e la volta successiva intende abbandonarli. Ma è un buon
\emph{check}: se senza anonimato non arriva nulla e con l'anonimato arrivano
cento feedback, il problema è chiaro.

\paragraph{Quando è giusto rivolgersi a un superiore?} La strada giusta resta
parlare direttamente con l'interessato. Ma se si prova e non si trova una
soluzione, la cosa più ovvia è andare dal proprio manager insieme: ``ci
aiuti?''. Quando il relatore parlava di mediare i dialoghi, era proprio questo:
fare da moderatore, aiutare a sbloccare i punti fermi o, se non si riesce a
scegliere tra A e B, scegliere lui. L'escalation, quando necessaria, va fatta.

\paragraph{Il tempo libero veniva rispettato?} No, e se l'è dovuto imporre
lui. Il lavoro gli piaceva così tanto che continuare a casa non gli sembrava
lavorare. Gli si è ritorto contro con il lancio di \emph{Diablo~III}, fatto in
orario americano (le due o l'una di notte in Europa), con l'ufficio chiuso
perché a nessuno era stato chiesto di restare. Il lancio va malissimo, i
giocatori protestano, e lui, sveglio, risponde sui forum fino alle tre, quattro,
cinque del mattino: ``non posso risolvere il problema, ma resto qui ad
aggiornarvi con quello che riesco a sapere dal team di sviluppo''. Il giorno
dopo trova sulla scrivania un \textbf{richiamo formale dell'HR}: aver lavorato
fuori orario senza autorizzazione, condotta non professionale. Per lui è stato
un interruttore: ``ora me ne vado''. Due settimane dopo faceva colloqui con
Riot. Ancora oggi è contrario a quel richiamo, ma gli ha fatto capire come
ragiona un'azienda: ``hai un contratto 9--18, non mi interessa se l'hai fatto
con le migliori intenzioni''. Eppure per anni ha goduto di un'ottima
reputazione in quelle community come ``quello che è rimasto con noi fino alle
cinque del mattino'', e anche l'azienda ne ha beneficiato.

Ci sono poi aziende in cui, soprattutto per i developer, viene richiesto il
famigerato \textbf{crunch}: lavoro straordinario continuativo, senza pause.
Dipende molto dall'azienda; sta venendo (``grazie al cielo'') deprecato, perché
molte aziende capiscono che non è sano ed è visto male anche dal pubblico, ma
esiste ancora.

\paragraph{E se il tuo capo demotiva qualcuno del tuo team?} Qui ``diventa
personale'': il team del relatore si fida di lui perché sa che in questi casi
``diventa una bestia''. ``Te la puoi prendere con me quanto vuoi, ma il mio
team non lo tocchi.'' Lo gestisce direttamente con l'altro manager, per quanto
brutta possa diventare la conversazione. Lo sforzo più grande è \textbf{non
andarne poi a parlare con il team} (``ragazzi, ho fatto questa cosa''), perché
minerebbe la fiducia verso l'alto: non vuole un rapporto ottimo fino a lui e
pessimo oltre. L'obiettivo è che l'altro manager ``si faccia perdonare''
cambiando il proprio comportamento, e che il team percepisca da solo il
cambiamento.

\paragraph{Chiusura della giornata.} Il relatore chiede un feedback collettivo
in chat (``se domani ci dai \emph{more of the same} siamo contenti, oppure
ditemi cosa limare''), si scusa per i ripetuti crash della connessione e
ringrazia per la proattività nelle domande, che gli servono anche come
\emph{check}: se nessuno chiede nulla su un argomento, significa che può
andare avanti. Racconta di aver passato una settimana a decidere che cosa
mettere nelle slide e di averle poi preparate in poco tempo (rifacendone
alcune fino alle tre di notte). Ricorda infine di usare il test di
personalità per sé e per il proprio gruppo, facendolo fare a chi non l'ha
ancora fatto: tornerà molto utile per cooperare meglio.

% ---------------------------------------------------------------------
\section{Anticipazioni sulla seconda giornata}
\label{sec:z1-domani}

La seconda giornata del workshop sarà più pratica: ``oggi vi ho spiegato
perché, domani vi insegno come lo fanno queste aziende''. Gli argomenti
previsti:
\begin{description}
  \item[Time management] come organizzarsi una giornata professionale in
        azienda: come strutturare una \emph{task list}, cosa fare e cosa non
        fare, come e cosa prioritizzare.
  \item[Brainstorming] i gruppi dovranno fare moltissimi meeting (come facciamo
        il gioco, come lo chiamiamo, che scopo gli diamo). Se non si sanno fare
        bene, i brainstorming sono \textbf{dannosi ed eterni}: un brainstorming
        dovrebbe durare \textbf{al massimo due ore}. Verranno mostrati metodi
        che accorciano la durata, massimizzano l'effetto e, soprattutto,
        \textbf{fanno partecipare tutti}: non c'è niente di peggio di un
        brainstorming in cui parlano in due per il 90\% del tempo. Il classico
        ``ragazzi, non esistono idee sbagliate, sparate'' è il modo peggiore di
        farlo; non è colpa di nessuno, ma si può fare diversamente.
  \item[Guidare da dentro] come essere leader nel proprio team anche senza
        esserne il leader: sviluppare responsabilità personale e verso gli
        altri. Non essere il manager di qualcuno non significa non esserne
        responsabili: la persona accanto è responsabilità di tutti, perché si
        è un team e non un insieme di individui. ``Se qualcuno cade, tutti lo
        facciamo rialzare: non c'è l'addetto al rialzamento.'' È la versione
        \emph{advanced} della giornata: come comportarsi da manager senza
        averne il titolo.
\end{description}
Torneranno inoltre il \emph{service level agreement} e il tema della
condivisione dei successi e della protezione del team dai fallimenti.

% ---------------------------------------------------------------------
\section*{Punti chiave del capitolo}
\addcontentsline{toc}{section}{Punti chiave del capitolo}

\begin{itemize}
  \item ``La capacità di un gruppo di persone di fare cose straordinarie
        dipende dal modo in cui queste persone operano insieme.'' Il 90\% dei
        problemi dei team sono \textbf{problemi di comunicazione}.
  \item Le aziende si liberano sempre più presto dello \textbf{useless
        genius}: le competenze si trasmettono, il carattere no. Il
        \textbf{fit culturale} va valutato anche dal candidato.
  \item Non esiste un unico modo di dare feedback: dipende dal contesto e
        soprattutto da \textbf{chi lo riceve}. Il feedback è una somma di
        piccoli miglioramenti periodici; tenersi le cose nel cassetto non
        serve mai.
  \item I team attraversano \textbf{forming, storming, norming, performing e
        adjourning}. Lo storming non si salta; si passa la maggior parte del
        tempo in forming; la bravura sta nell'accelerare i passaggi e nel
        fare bene il \textbf{post-mortem}, anche dei successi.
  \item Gli obiettivi di un progetto vanno fissati \textbf{all'inizio}: solo
        così si può valutare se è andato bene e scegliere tra A e B.
  \item Il circolo virtuoso è \textbf{feedback $\rightarrow$ miglioramento
        $\rightarrow$ fiducia}. Accettare un feedback non significa fare come
        dice l'altro: decide chi ha la responsabilità, ma deve argomentare.
  \item ``Aiutami a poterti aiutare'': meglio dire in anticipo ``potrei non
        farcela'' che ``non ce l'ho fatta''. Chiedere aiuto non è una
        debolezza.
  \item Quando si dà feedback \textbf{non si porta la soluzione}: la si
        costruisce insieme a chi lo riceve.
  \item La cultura è forte quando si lavora \textbf{per gli altri}; il lavoro
        è valutato come opera d'insieme (anche con bonus di team).
  \item Come in \emph{Among Us}, i team sprecano energie a cercare il
        traditore invece di fare le task. \textbf{Default to trust}: chi ha
        sbagliato ha agito con le informazioni che aveva. Cercare
        \textbf{la causa, non il colpevole}: ``che cosa ti ha fatto pensare
        che fosse la scelta giusta?''.
  \item Carriere e performance non sono lineari: vanno valutate le
        \textbf{traiettorie}, non un punto dell'anno. Le valutazioni puntuali
        generano \emph{lying, hiding, faking}.
  \item Confermare sempre le istruzioni facendosele ripetere; ``non ho
        capito'' non è una colpa.
  \item Dare feedback: sull'\textbf{azione} e sui suoi \textbf{effetti}, specifico,
        presto, uno a uno, adattato al carattere, senza terze parti, senza
        ``non mi piace'', senza offendersi se non viene seguito; poi fare il
        debriefing.
  \item Ricevere feedback: \textbf{ascolta, chiedi dettagli, accetta, agisci}.
        Il feedback positivo serve a non peggiorare ciò che si fa bene; il
        feedback accidentale e il \emph{drive-by} non servono a nulla.
  \item Feedback \textbf{frequenti e recenti}: tanti piccoli \emph{sanity
        check} costano meno di una grande correzione di rotta.
  \item Il feedback è \textbf{un regalo} di chi ha a cuore il nostro
        interesse; fa paura il silenzio, non le critiche.
\end{itemize}
